\chapter{Bar-cobar adjunction and curved Koszul duality}
\label{chap:bar-cobar-adjunction}

\section{Curved Koszul duality and I-adic completion}
\label{sec:curved-koszul-i-adic}

Not all chiral algebras are quadratic. Many examples of interest---Virasoro, higher W-algebras, 
$W_\infty$---require curved structures or infinite-dimensional presentations. For these, 
the strict finite-type bar-cobar theorem no longer applies directly.  One must
separate the proved M-level completed shadows that control curvature and
complementarity from the stronger H-level realization problem for an actual dual
factorization object.  To treat even the model-level question, we must introduce:
\begin{itemize}
\item \emph{Curved structures}: Allowing μ₀ ≠ 0 (failure of d² = 0)
\item \emph{Completion}: I-adic topology ensuring convergence
\item \emph{Filtered structures}: More general than curved (Gui--Li--Zeng)
\end{itemize}

We follow Gui--Li--Zeng~\cite{GLZ22}.

\subsection{Curved $A_\infty$ algebras: definitions}
\label{sec:curved-ainfty-definition}

\begin{definition}[Curved \texorpdfstring{$A_\infty$}{A-infinity} algebra]
\label{def:curved-ainfty}
A \emph{curved $A_\infty$ algebra} is a graded vector space $A$ with operations:
\begin{equation}
\{\mu_n: A^{\otimes n} \to A\}_{n \geq 0}
\end{equation}
of degree $2-n$, satisfying the \emph{curved $A_\infty$ relations}:
\begin{equation}\label{eq:curved-ainfty-relations}
\sum_{\substack{r+s+t=n \\ r,t \geq 0,\; s \geq 0}} (-1)^{rs+t}\, \mu_{r+1+t}(\mathrm{id}^{\otimes r} \otimes \mu_s \otimes \mathrm{id}^{\otimes t}) = 0
\end{equation}

Key differences from ordinary $A_\infty$:
\begin{enumerate}
\item $n=0$ is allowed: $\mu_0: \mathbb{C} \to A$ is the \emph{curvature}
\item $n=1$: $\mu_1^2 = [\mu_0, -]$, so $\mu_1$ is a differential only modulo curvature
\item $n \geq 2$: Higher operations as usual
\end{enumerate}
\end{definition}

\begin{remark}[Curvature, backreaction, and \texorpdfstring{$d^2$}{d\textasciicircum 2}]\label{rem:curvature-backreaction-terminology}
\begin{enumerate}[label=(\roman*)]
\item The \emph{curvature} is $m_0 := \mu_0(1) \in A^2$.
\item $\mu_1^2(a) = \mu_2(m_0, a) - \mu_2(a, m_0)$: the total bar differential always satisfies $d_{\mathrm{bar}}^2 = 0$; curvature manifests in the \emph{internal} $A_\infty$ structure.
\item In physics, \emph{backreaction} is the same phenomenon: the BRST operator $Q = \mu_1$ satisfies $Q^2 \neq 0$, with the obstruction measured by $m_0$.
\item A curved $A_\infty$ algebra is \emph{strictifiable} iff $m_0$ is central (Theorem~\ref{thm:curvature-central}) and the curvature class is gauge-trivial in the completed filtration.
\end{enumerate}
\end{remark}

\begin{theorem}[Curvature as \texorpdfstring{$\mu_1$}{mu1}-cycle; \ClaimStatusProvedHere]
\label{thm:curvature-central}
For a curved $A_\infty$ algebra $(A, \mu_0, \mu_1, \mu_2, \ldots)$, the curvature element satisfies:
\begin{enumerate}
\item $\mu_1(\mu_0) = 0$ \quad (the curvature is a $\mu_1$-cycle);
\item $\mu_1^2 = [\mu_0, -]_{\mu_2}$, where $[\mu_0, a]_{\mu_2} := \mu_2(\mu_0, a) - \mu_2(a, \mu_0)$ is the inner derivation.
\end{enumerate}
In particular, $\mu_1$ fails to be a differential precisely when $\mu_0$ is \emph{not} central with respect to $\mu_2$.  For chiral algebras (which are graded-commutative), $[\mu_0, -]_{\mu_2} = 0$ automatically, so the curvature enters through higher operations ($\mu_3, \mu_4, \ldots$) or via the completed tensor product.
\end{theorem}

\begin{proof}
The curved $A_\infty$ relations $\sum_{r+s+t=n}(-1)^{rs+t}\mu_{r+1+t}(\mathrm{id}^{\otimes r}\otimes \mu_s \otimes \mathrm{id}^{\otimes t})=0$ give, at each level:
\begin{itemize}
\item $n=0$: $\mu_1(\mu_0) = 0$ \quad (the curvature is a $\mu_1$-cycle).
\item $n=1$: $\mu_1^2(a) - \mu_2(\mu_0, a) + \mu_2(a, \mu_0) = 0$ for all $a$ \quad (failure of $\mu_1^2 = 0$).
\end{itemize}

Rearranging the $n=1$ relation: $\mu_1^2(a) = \mu_2(\mu_0, a) - \mu_2(a, \mu_0) = [\mu_0, a]_{\mu_2}$.

The $n=2$ relation (Leibniz rule for $\mu_1$ acting on $\mu_2$) applied using $\mu_1(\mu_0) = 0$ from $n=0$ gives:
\begin{equation}
\mu_1(\mu_2(\mu_0, a)) = \mu_2(\mu_0, \mu_1(a))
\end{equation}
confirming that the inner derivation $[\mu_0, -]_{\mu_2}$ commutes with $\mu_1$, as required for the self-consistency of the curved structure.
\end{proof}

\subsection{I-adic completion: topology and convergence}
\label{sec:i-adic-completion}

\begin{definition}[I-adic topology]
\label{def:i-adic-topology}
\index{nilpotent completion}
Let $A$ be a curved $A_\infty$ algebra with augmentation ideal $I = \ker(\varepsilon: A \to \mathbb{C})$.

The \emph{I-adic completion} of $A$ is:
\begin{align}
\hat{A} &:= \varprojlim_n A/I^n \notag \\
  &= \bigl\{a \in {\textstyle\prod_{n \geq 0}} A/I^n :
  \text{compatible}\bigr\}.
\end{align}

An element $a \in \hat{A}$ can be written as a formal series:
\begin{equation}
a = a_0 + a_1 + a_2 + \cdots \quad \text{where } a_n \in I^n/I^{n+1}
\end{equation}
\end{definition}

\begin{theorem}[When completion is necessary; \ClaimStatusProvedElsewhere]
\label{thm:completion-necessity}
Completion $A \to \hat{A}$ is necessary when:
\begin{enumerate}
\item \emph{Infinite sums}: Operations μₙ produce infinite sums not convergent in $A$
\item \emph{Non-conilpotent}: Bar complex $\bar{B}(A)$ is not conilpotent
\item \emph{Non-quadratic}: Relations involve infinitely many generators
\end{enumerate}

For example, the Virasoro algebra and $W_\infty$ need completion, while Heisenberg and Kac--Moody (conilpotent) do not.
\emph{More precisely:} Virasoro admits a completed same-family shadow used
throughout Parts~II--III, whereas $W_\infty$ marks the frontier where a
completed duality theorem is still missing.
\end{theorem}

\begin{proof}[Proof by Example: Virasoro]
The Virasoro algebra has generators $\{L_n\}_{n \in \mathbb{Z}}$ with:
\begin{equation}
[L_m, L_n] = (m-n)L_{m+n} + \frac{c}{12}m(m^2-1)\delta_{m+n,0}
\end{equation}

Consider the bar complex element:
\begin{equation}
\omega = L_0 \otimes L_0 \in \bar{B}^2(\text{Vir})
\end{equation}

Applying the differential involves summing over all intermediate states:
\begin{equation}
d(\omega) = \sum_{k \in \mathbb{Z}} [L_0, L_k] \otimes L_k \otimes L_0 + \ldots
\end{equation}

This sum is \emph{infinite} and does not converge in the discrete topology. We need 
completion with respect to the augmentation ideal to make sense of it.

In $\hat{\text{Vir}}$, the sum converges because $L_k \in I^{|k|}$, so:
\begin{equation}
d(\omega) = \sum_{k=-\infty}^\infty (\text{term with } L_k) 
\end{equation}
converges I-adically (finitely many terms in each $I^n/I^{n+1}$).
\end{proof}

\begin{remark}[Scope of the proof]\label{rem:completion-necessity-scope}
The Virasoro example above establishes necessity for condition~(1).
Condition~(2) is demonstrated by the $\beta\gamma$-system, whose bar complex is not conilpotent due to the infinite-rank coproduct on the symmetric coalgebra (see~\S\ref{subsec:betagamma-all-genera} in Chapter~\ref{ch:genus-expansions}).
Condition~(3) is illustrated by affine $\widehat{\mathfrak{sl}}_2$ at non-integral level, where the Serre-type relations involve infinitely many generators (see~\S\ref{sec:three-theorems-sl2} in Chapter~\ref{ch:genus-expansions}).
\end{remark}

\begin{proposition}[Acyclicity of curved bar complexes; \ClaimStatusProvedHere]
\label{prop:curved-bar-acyclicity}
\index{acyclicity!curved bar complex}
\index{Positselski!acyclicity}
Let $\mathcal{A}$ be a Koszul chiral algebra with nonzero obstruction
coefficient $\kappa(\mathcal{A}) \neq 0$.  At genus $g \geq 1$,
the curved bar complex $\bar{B}^{(g)}(\mathcal{A})$ has \emph{acyclic}
underlying cochain complex: $H^*(\bar{B}^{(g)}(\mathcal{A}), d_{\mathrm{total}}) = 0$.
\end{proposition}

\begin{proof}
At genus $g \geq 1$, the bar differential satisfies
$\dfib^{\,2} = \mcurv{g} \cdot \mathrm{id}$ where $\mcurv{g} = \kappa(\mathcal{A}) \cdot \lambda_g \neq 0$
(Theorem~\ref{thm:genus-universality}; Convention~\ref{conv:higher-genus-differentials}).  The fiberwise differential $\dfib$ is curved, but the total corrected differential $\Dg{g}$ satisfies
$\Dg{g}^{\,2} = 0$.  The internal $A_\infty$ operation
$m_1$ fails nilpotence: $m_1^2 = [m_0, -]_{m_2} \neq 0$
(Corollary~\ref{cor:critical-level-universality}).

The standard Positselski argument~\cite[Remark~3.6]{Positselski11} applies: for a CDG-coalgebra $(C, d, h)$ with $h \neq 0$, the bar construction is acyclic.  Here $C = \bar{B}^{(g)}(\mathcal{A})$ with curvature $h = m_0^{(g)} \neq 0$.
\end{proof}

\begin{remark}[Necessity of exotic derived categories at higher genus]
\label{rem:positselski-acyclicity}
At genus $g \geq 1$ with $\kappa \neq 0$, acyclicity of
$\bar{B}^{(g)}(\mathcal{A})$ makes the ordinary derived category of
its comodules trivial.  The correct framework is Positselski's
coderived category
$D^{\mathrm{co}}(\bar{B}^{(g)}\text{-}\mathrm{CoMod})$~\cite[\S\S3--4]{Positselski11},
which is why Theorem~\ref{thm:positselski-chiral} requires the
coderived/contraderived setting.
\end{remark}

\subsection{Filtered versus curved structures}
\label{sec:filtered-vs-curved}

\begin{theorem}[Filtered cooperads (Gui--Li--Zeng~\cite{GLZ22}); \ClaimStatusProvedElsewhere]
\label{thm:filtered-cooperads}
A \emph{filtered cooperad} $\mathcal{C}$ is more general than a curved cooperad:
\begin{equation}
\mathcal{C} = \bigcup_{n=0}^\infty F^n \mathcal{C}
\end{equation}
where $F^n \mathcal{C} \subset F^{n+1} \mathcal{C}$ is an increasing filtration, with:
\begin{enumerate}
\item Comultiplication: $\Delta(F^n) \subset \bigoplus_{i+j=n} F^i \otimes F^j$
\item Counit: $\varepsilon(F^{>0}) = 0$
\end{enumerate}

The filtered structure does not reduce to a single curvature element.
\end{theorem}

\begin{remark}[Provenance and citation]
This theorem is imported and treated as \ClaimStatusProvedElsewhere. The
filtered-cooperad framework used here is sourced from \cite{GLZ22}; subsequent
internal reductions in this manuscript are organized by
Theorem~\ref{thm:filtered-to-curved} and Theorem~\ref{thm:bar-convergence}.
\end{remark}

\begin{example}[W-algebras: filtered but not simply curved]
\label{ex:w-algebra-filtered-comprehensive}
For W₃ algebra, the filtration is by conformal weight:
\begin{equation}
F^n \mathcal{W}_3 = \text{span}\{T, W, \partial T, TT, \partial W, \ldots \text{ up to 
weight } n\}
\end{equation}

This does NOT come from a single curvature element μ₀. Instead, there are curvature 
contributions at each weight:
\begin{align}
\mu_0^{(2)} &= T \quad \text{(weight 2)}\\
\mu_0^{(3)} &= W \quad \text{(weight 3)}\\
\mu_0^{(4)} &= TT + \text{regular} \quad \text{(weight 4)}\\
&\vdots
\end{align}

The full structure requires the complete filtered cooperad, not just a curved one.
\end{example}

\begin{theorem}[When filtered reduces to curved; \ClaimStatusProvedHere]
\label{thm:filtered-to-curved}
Let $\mathcal{C}$ be a filtered chiral cooperad with exhaustive, separated, complete
filtration $F_\bullet\mathcal{C}$. Assume:
\begin{enumerate}
\item $\dim(F_n\mathcal{C}/F_{n-1}\mathcal{C})<\infty$ for every $n$;
\item $\gr(\mathcal{C})$ is generated in degree $1$ with quadratic relations;
\item higher relation classes differ from quadratic consequences by central
filtration-$\ge2$ corrections.
\end{enumerate}
Then there exists a curved cooperad $\mathcal{C}_{\mathrm{curv}}$ with curvature
$\mu_0\in F_2\mathcal{C}_{\mathrm{curv}}$ and a filtered quasi-isomorphism
\[
\mathcal{C}_{\mathrm{curv}}\xrightarrow{\;\simeq_f\;}\mathcal{C}.
\]
If $\gr(\mathcal{C})$ has infinitely many nonzero graded pieces, this curved
description is only defined after completion.
\end{theorem}

\begin{proof}
Set $\mathcal{A}:=\mathcal{C}^\vee$ (continuous dual with respect to
$F_\bullet$). By assumption (1), filtration pieces are finite-dimensional, so
duality exchanges filtered cooperads and filtered chiral algebras without loss
of filtration data.

Assumptions (2)--(3) are exactly the hypotheses of
Proposition~\ref{prop:filtered-to-curved-fc} for $\mathcal{A}$. Hence there is a
curved model $\mathcal{A}_{\mathrm{curv}}$ with curvature
$\mu_0\in F_2\mathcal{A}_{\mathrm{curv}}$ and a filtered quasi-isomorphism
$\mathcal{A}_{\mathrm{curv}}\to\mathcal{A}$.

Dualizing back gives a curved cooperad
$\mathcal{C}_{\mathrm{curv}}:=\mathcal{A}_{\mathrm{curv}}^\vee$ and a filtered
quasi-isomorphism $\mathcal{C}_{\mathrm{curv}}\to\mathcal{C}$. The curved identity
$d^2=[\mu_0,-]$ is preserved under this dualization because all filtration
quotients are finite-dimensional.

If $\gr(\mathcal{C})$ has infinitely many nonzero graded pieces, the curvature
decomposes into infinitely many filtration components. Then the bar/cobar
operations are no longer finite strict and convergence requires completion,
as in Theorem~\ref{thm:conilpotency-convergence} and
Theorem~\ref{thm:bar-convergence}.
\end{proof}

\begin{remark}[Dependencies]
This proof uses:
\begin{enumerate}[label=(D\arabic*)]
\item the filtered algebra-side reduction in
Proposition~\ref{prop:filtered-to-curved-fc};
\item finite-type filtered duality between filtered algebras and cooperads;
\item convergence control from Theorem~\ref{thm:conilpotency-convergence} and
Theorem~\ref{thm:bar-convergence}.
\end{enumerate}
\end{remark}

\subsection{Conilpotency and convergence without completion}
\label{sec:conilpotency-convergence}

\begin{definition}[Conilpotent coalgebra]
\label{def:conilpotent-complete}
\index{conilpotent!coalgebra}
\index{conilpotent!filtration}
A coalgebra $C$ is \emph{conilpotent} if for each $c \in C$, there exists $N$ such that:
\begin{equation}
\Delta^{(N)}(c) = 0
\end{equation}
where $\Delta^{(N)}$ is the $N$-fold iterated comultiplication.
\end{definition}

\begin{theorem}[Conilpotency ensures convergence; \ClaimStatusProvedHere]
\label{thm:conilpotency-convergence}
\label{thm:conilpotency-bar}
\label{thm:koszul-conilpotent}
\index{bar construction!convergence}
If $\bar{B}(A)$ is conilpotent, then:
\begin{enumerate}
\item The bar-cobar composition $\Omega \circ \bar{B}(A) \to A$ converges without completion
\item All infinite sums in the cobar differential terminate after finitely many steps
\item The Koszul duality $A \leftrightarrow A^!$ is well-defined without taking $\hat{A}$
\end{enumerate}
\end{theorem}

\begin{proof}
\emph{Step 1}: For conilpotent $\bar{B}(A)$, each element $\omega \in \bar{B}^n(A)$ 
has $\Delta^{(N)}(\omega) = 0$ for some $N$.

\emph{Step 2}: The cobar differential is:
\begin{equation}
d_{\text{cobar}}(f) = \sum_{\text{decompositions}} (-1)^{|\alpha|} f \circ \Delta(\omega)
\end{equation}

\emph{Step 3}: Since $\Delta^{(N)}(\omega) = 0$, the sum has at most $N$ terms, so 
it converges in the discrete topology.

\emph{Step 4}: The bar-cobar composition:
\begin{equation}
(\Omega \circ \bar{B})(A) = \bigoplus_{n=0}^\infty (\text{cobar operations on } 
\bar{B}^n(A))
\end{equation}
has all operations terminating after finitely many steps by conilpotency.
\end{proof}

\begin{example}[Heisenberg: conilpotent]
\label{ex:heisenberg-conilpotent-complete}
The Heisenberg algebra $\mathcal{H}_\kappa$ has bar complex:
\begin{equation}
\bar{B}^n(\mathcal{H}_\kappa) = \mathcal{H}_\kappa^{\otimes n} \otimes \Omega^n
\end{equation}

For $\omega = a_{n_1} \otimes \cdots \otimes a_{n_k} \otimes \eta_{ij}$, the
comultiplication is:
\begin{equation}
\Delta(\omega) = \sum_{\text{splittings}} \omega_L \otimes \omega_R
\end{equation}

After $k$ iterations, $\Delta^{(k)}(\omega) = 0$ because we run out of tensor factors. 
Thus $\bar{B}(\mathcal{H}_\kappa)$ is conilpotent, and no completion is needed.
\end{example}

\begin{example}[Virasoro: not conilpotent]
\label{ex:virasoro-not-conilpotent}
The Virasoro algebra has infinitely many generators $L_n$. Consider:
\begin{equation}
\omega = L_0 \in \bar{B}^1(\text{Vir})
\end{equation}

The comultiplication gives:
\begin{equation}
\Delta(\omega) = \sum_{k \in \mathbb{Z}} (\text{terms with } L_k \otimes L_{-k})
\end{equation}

This sum is infinite and does not converge in the discrete topology on $\bar{B}(\text{Vir})$, so $\Delta^{(N)}(\omega) \neq 0$ for all $N$.
Thus $\bar{B}(\text{Vir})$ is not conilpotent; completion is required.
\end{example}

\subsection{Examples: computing Koszul duals with completion}
\label{sec:koszul-duals-completion-examples}

\begin{example}[Virasoro Koszul dual]
\label{ex:virasoro-koszul-dual}

Since $\bar{B}(\mathrm{Vir})$ is not conilpotent
(Example~\ref{ex:virasoro-not-conilpotent}), one forms the
completed bar complex and takes continuous linear dual.
Writing $\mathrm{Vir}_c = \mathcal{W}^k(\mathfrak{sl}_2)$
with $h^\vee = 2$, level-shifting duality
(Corollary~\ref{cor:level-shifting-part1}) gives
$k' = -k - 2h^\vee = -k-4$, so
\begin{equation}
\mathrm{Vir}_c^! \cong \mathrm{Vir}_{26-c},
\end{equation}
since $c(k) + c(k') = 2\dim\mathfrak{sl}_2 = 26$ by Theorem~\ref{thm:central-charge-complementarity}.  Physically, this is matter--ghost duality: a matter system at central charge~$c$ couples to ghosts at $c_{\mathrm{ghost}} = 26 - c$.  In particular, $\mathrm{Vir}_{26}^! \cong \mathrm{Vir}_0$, consistent with bosonic string BRST cohomology at ghost number zero.

This is the M/S-level completed shadow; the stronger H-level realization via an infinite-generator dual is the MC4 frontier (Example~\ref{ex:winfty-completion-frontier}).
\end{example}

\begin{example}[\texorpdfstring{$W_\infty$}{W-infinity}: exact MC4 frontier]
\label{ex:winfty-completion-frontier}

The standard principal-stage completed M-level bar-cobar package for $W_\infty = \varprojlim_N W_N$ is proved (Corollary~\ref{cor:winfty-standard-mc4-package}).  The live MC4 frontier is to construct a factorization target (Proposition~\ref{prop:winfty-factorization-package}) and prove the residue identities
\[
\mathsf{C}^{\mathrm{res}}_{s,t;u;m,n}(N)
=\mathsf{C}^{\mathrm{DS}}_{s,t;u;m,n}(N)
\]
on the finite-detection packet $\mathcal{I}_N$ (Corollary~\ref{cor:winfty-ds-finite-seed-set}).  Three obstructions remain:
\begin{enumerate}
\item construct a separated complete H-level/factorization target whose finite quotients recover $W_N$;
\item prove the exact residue identities of Proposition~\ref{prop:winfty-ds-residue-identity-criterion};
\item close finite detection on $\mathcal{I}_N$, after which Corollary~\ref{cor:winfty-hlevel-comparison-criterion} makes the comparison formal.
\end{enumerate}
\end{example}

\begin{conjecture}[Inverse-limit completed bar-cobar package; \ClaimStatusConjectured]
\label{conj:inverse-limit-completed-bar-cobar}
Let $\cA = \varprojlim_N \cA_{\le N}$ be a separated complete
weight-filtered chiral algebra such that:
\begin{enumerate}
\item each truncation $\cA_{\le N}$ is finite type and lies in the
  proved bar-cobar regime;
\item the transition maps $\cA_{\le N+1} \to \cA_{\le N}$ are compatible
  with the finite-stage bar constructions; and
\item the induced inverse system on bar complexes is pro-nilpotent.
\end{enumerate}
Then the completed reduced bar construction
\[
\widehat{\bar B}(\cA) := \varprojlim_N \bar B(\cA_{\le N})
\]
exists as a separated complete curved dg coalgebra with continuous
differential, and the completed cobar construction satisfies
\[
\Omega\bigl(\widehat{\bar B}(\cA)\bigr) \simeq \cA
\]
compatibly with the finite-stage quasi-isomorphisms.

\emph{This is the theorem-shaped local form of the MC4 frontier:
for $W_\infty$ the truncations are principal finite-type $W_N$, while
for Yangian towers they are the finite RTT stages.}
\end{conjecture}

\begin{proposition}[Reduction of MC4 to finite-stage compatibility;
\ClaimStatusProvedHere]
\label{prop:mc4-reduction-principle}
Let $\{C_N\}_{N \ge 0}$ be the inverse system of finite-stage bar
complexes
\[
C_N := \bar B(\cA_{\le N}),
\qquad
\widehat{C} := \varprojlim_N C_N,
\]
for a tower satisfying the finite-stage hypotheses of
Conjecture~\ref{conj:inverse-limit-completed-bar-cobar}.  Assume:
\begin{enumerate}
\item each transition map $C_{N+1} \to C_N$ is a morphism of complexes;
\item for every cohomological degree~$m$, the inverse system
  $\{H^m(C_N)\}_N$ satisfies the Mittag--Leffler condition; and
\item the finite-stage bar-cobar quasi-isomorphisms
  $\epsilon_N \colon \Omega(C_N) \xrightarrow{\sim} \cA_{\le N}$ are
  compatible with the tower maps.
\end{enumerate}
Then:
\begin{enumerate}
\item the completed cohomology is computed by the limit of the
  finite-stage cohomologies,
  \[
  H^m(\widehat{C}) \cong \varprojlim_N H^m(C_N);
  \]
\item the induced map
  \[
  \widehat{\epsilon}\colon \Omega(\widehat{C}) \longrightarrow
  \varprojlim_N \cA_{\le N}
  \]
  is a quasi-isomorphism whenever the completed cobar differential is
  continuous.
\end{enumerate}
Consequently, the remaining content of MC4 for a concrete tower is to
verify continuity, pro-nilpotence, and the Mittag--Leffler/convergence
conditions; no new formal obstruction remains once these are checked.
\end{proposition}

\begin{proof}
For any inverse system of cochain complexes there is a Milnor exact
sequence
\[
0 \to \varprojlim\nolimits^1_N H^{m-1}(C_N)
\to H^m(\widehat{C})
\to \varprojlim_N H^m(C_N) \to 0.
\]
Assumption~\textup{(2)} kills the derived limit term
$\varprojlim^1$, so the first claim follows.

For the second claim, the maps $\epsilon_N$ assemble by
Assumption~\textup{(3)} into a morphism of inverse systems.  Passing to
limits therefore gives a map
$\widehat{\epsilon}\colon \Omega(\widehat{C}) \to \varprojlim_N \cA_{\le N}$.
The finite-stage maps are quasi-isomorphisms, so after applying the
first part to the bar complexes and the same Milnor argument to the
compatible cobar tower, $\widehat{\epsilon}$ induces an isomorphism on
cohomology.  The continuity hypothesis ensures that the completed cobar
differential is well defined, so $\widehat{\epsilon}$ is a
quasi-isomorphism.
\end{proof}

% \label{rem:mc4-concrete-checklist} — removed (echo of Proposition~\ref{prop:mc4-reduction-principle})

\begin{corollary}[Degreewise stabilization criterion for MC4;
\ClaimStatusProvedHere]
\label{cor:mc4-degreewise-stabilization}
In the setting of Proposition~\ref{prop:mc4-reduction-principle},
assume in addition that for every cohomological degree~$m$ there exists
$N(m)$ such that the transition maps
\[
H^m(C_{N+1}) \longrightarrow H^m(C_N)
\]
are isomorphisms for all $N \ge N(m)$.  Then:
\begin{enumerate}
\item the inverse system $\{H^m(C_N)\}_N$ is Mittag--Leffler;
\item the completed cohomology stabilizes at the finite stage,
  \[
  H^m(\widehat{C}) \cong H^m(C_{N(m)});
  \]
\item every compatible system of finite-stage bar-cobar
  quasi-isomorphisms induces a completed quasi-isomorphism
  \[
  \Omega(\widehat{C}) \xrightarrow{\sim} \varprojlim_N \cA_{\le N},
  \]
  provided the completed cobar differential is continuous.
\end{enumerate}
Thus eventual stabilization of the finite-stage bar cohomology is a
sufficient MC4 input.
\end{corollary}

\begin{proof}
Eventual constancy implies the Mittag--Leffler condition, so
Proposition~\ref{prop:mc4-reduction-principle} applies.  Since the
transition maps are isomorphisms for $N \ge N(m)$, the inverse limit
$\varprojlim_N H^m(C_N)$ identifies with the stable value
$H^m(C_{N(m)})$.  The second and third claims are then exactly the
conclusions of
Proposition~\ref{prop:mc4-reduction-principle}.
\end{proof}

\begin{corollary}[Finite-dimensional surjectivity criterion for MC4;
\ClaimStatusProvedHere]
\label{cor:mc4-surjective-criterion}
In the setting of Proposition~\ref{prop:mc4-reduction-principle},
assume in addition that for every cohomological degree~$m$:
\begin{enumerate}
\item each $H^m(C_N)$ is finite-dimensional; and
\item there exists $N_{\mathrm{surj}}(m)$ such that the transition maps
  \[
  H^m(C_{N+1}) \longrightarrow H^m(C_N)
  \]
  are surjective for all $N \ge N_{\mathrm{surj}}(m)$.
\end{enumerate}
Then the inverse system $\{H^m(C_N)\}_N$ is Mittag--Leffler, and every
compatible system of finite-stage bar-cobar quasi-isomorphisms induces a
completed quasi-isomorphism
\[
\Omega(\widehat{C}) \xrightarrow{\sim} \varprojlim_N \cA_{\le N},
\]
provided the completed cobar differential is continuous.
Thus MC4 can be verified by eventual surjectivity on finite-dimensional
weight slices; one does not need a priori eventual constancy of the full
cohomology tower.
\end{corollary}

\begin{proof}
Fix $m$.  For $N \ge N_{\mathrm{surj}}(m)$ and every $r \ge 0$, the image
of the composite map
\[
H^m(C_{N+r}) \longrightarrow H^m(C_N)
\]
is all of $H^m(C_N)$.  Hence the descending image filtration in
$H^m(C_N)$ stabilizes immediately, so the inverse system is
Mittag--Leffler.  Proposition~\ref{prop:mc4-reduction-principle} then
gives the completed cohomology comparison and the completed bar-cobar
quasi-isomorphism.
\end{proof}

\begin{proposition}[Weight-cutoff criterion for MC4;
\ClaimStatusProvedHere]
\label{prop:mc4-weight-cutoff}
In the setting of Proposition~\ref{prop:mc4-reduction-principle},
assume in addition that each finite-stage bar complex $C_N$ carries an
exhaustive increasing auxiliary-weight filtration
\[
F_{\le 0}C_N \subset F_{\le 1}C_N \subset \cdots \subset C_N
\]
such that:
\begin{enumerate}
\item the differential and the transition maps preserve the filtration;
\item for each weight bound $w$, the filtered piece $F_{\le w}C_N$ is
  finite-dimensional in every cohomological degree; and
\item for each $w$ and every $N \ge w$, the transition map induces an
  isomorphism of filtered subcomplexes
  \[
  F_{\le w}C_{N+1} \xrightarrow{\sim} F_{\le w}C_N.
  \]
\end{enumerate}
Then for every cohomological degree $m$, weight bound $w$, and
$N \ge w$, the induced map
\[
H^m(F_{\le w}C_{N+1}) \xrightarrow{\sim} H^m(F_{\le w}C_N)
\]
is an isomorphism.  Consequently the associated graded weight slices of
cohomology stabilize, and if the cohomology splits as a direct sum of
finite-dimensional weight slices, then the maps
\[
H^m(C_{N+1})_w \xrightarrow{\sim} H^m(C_N)_w
\]
are isomorphisms for all $N \ge w$.  Hence the eventual-surjectivity
input of Corollary~\ref{cor:mc4-surjective-criterion} is automatic once
the tower is defined by genuine weight truncations.
\end{proposition}

\begin{proof}
Assumption~\textup{(3)} identifies the filtered subcomplexes
$F_{\le w}C_{N+1}$ and $F_{\le w}C_N$ for every $N \ge w$, so their
cohomology groups are canonically isomorphic.  Passing to successive
quotients gives stabilization of the associated graded weight pieces.
If the cohomology splits by weight, the direct-summand description of
the weight slices identifies $H^m(C_N)_w$ with the corresponding
graded piece, so the same stabilization holds on the nose.  In
particular the transition maps are eventually surjective on each
finite-dimensional weight slice, and
Corollary~\ref{cor:mc4-surjective-criterion} applies.
\end{proof}

\begin{proposition}[\texorpdfstring{$W_\infty$}{W_infty} criterion from principal finite-type stages;
\ClaimStatusProvedHere]
\label{prop:winfty-mc4-criterion}
Let
\[
W_\infty = \varprojlim_N W_N
\]
be a separated complete infinite-generator $\mathcal{W}$-tower whose
finite stages $W_N$ lie in the proved principal finite-type
$\mathcal{W}$ regime.  Assume:
\begin{enumerate}
\item the transition maps $W_{N+1} \to W_N$ are compatible with the
  finite-stage bar complexes and bar-cobar quasi-isomorphisms;
\item for every cohomological degree~$m$, the inverse system
  $\{H^m(\bar B(W_N))\}_N$ is eventually constant; and
\item the transition maps define an inverse system of curved dg
  coalgebras on the finite-stage bar complexes.
\end{enumerate}
Then the completed bar object
\[
\widehat{\bar B}(W_\infty) := \varprojlim_N \bar B(W_N)
\]
is a separated complete pro-nilpotent curved dg coalgebra, and the
canonical comparison map
\[
\Omega\bigl(\widehat{\bar B}(W_\infty)\bigr) \xrightarrow{\sim} W_\infty
\]
is a quasi-isomorphism.  Thus the M-level completed bar-cobar package
for $W_\infty$ reduces to compatible inverse-system dg structure and
stabilization of the principal finite-type tower.  The stronger H-level
realization problem belongs to MC4; any physics-facing identification
built on top of that realization is downstream at MC5.  By
Corollary~\ref{cor:mc4-surjective-criterion}, the eventual-constancy
hypothesis may be replaced by eventual surjectivity on finite-dimensional
weight slices, and Proposition~\ref{prop:mc4-weight-cutoff} shows that
this surjectivity is formal for the standard principal-stage tower.
\end{proposition}

\begin{proof}
The principal finite-type stages supply the finite-stage bar-cobar
quasi-isomorphisms required by
Proposition~\ref{prop:mc4-reduction-principle}.  Assumption~\textup{(2)}
is exactly the degreewise stabilization hypothesis of
Corollary~\ref{cor:mc4-degreewise-stabilization}, and
Assumption~\textup{(3)} lets one apply
Proposition~\ref{prop:inverse-limit-differential-continuity} to the
inverse-limit bar and cobar differentials.  Applying
Corollary~\ref{cor:mc4-degreewise-stabilization} to the tower
$\{\bar B(W_N)\}_N$ yields the stated completed quasi-isomorphism.
\end{proof}

\begin{corollary}[Standard principal-stage cutoff for
\texorpdfstring{$W_\infty$}{W_infty}; \ClaimStatusProvedHere]
\label{cor:winfty-weight-cutoff}
Let $W_N$ be the standard principal finite-type stage of $W_\infty$,
so that $W_N$ retains exactly the generators of conformal weights
$2,3,\dots,N$ and the transition map $W_{N+1}\to W_N$ forgets the
generator of weight $N+1$.  Then the tower of bar complexes
$\{\bar B(W_N)\}_N$ satisfies the hypotheses of
Proposition~\ref{prop:mc4-weight-cutoff} for the total conformal-weight
filtration.  In particular, for every cohomological degree $m$ and
conformal weight $w$, the map
\[
H^m(\bar B(W_{N+1}))_w \xrightarrow{\sim} H^m(\bar B(W_N))_w
\]
is an isomorphism for all $N \ge w$.

Therefore the standard principal-stage tower already supplies the
surjectivity/stabilization part of MC4.  Together with
Proposition~\ref{prop:inverse-limit-differential-continuity}, this
reduces the standard $W_\infty$ tower to the standard-tower completed
M-level package of Corollary~\ref{cor:winfty-standard-mc4-package}.
\end{corollary}

\begin{proof}
A bar chain of total conformal weight $w$ can involve only generators
whose conformal weights are at most $w$.  Once $N \ge w$, all such
generators already occur in $W_N$, so the filtered bar subcomplex of
weights $\le w$ is identical in every later stage.  The bar
differential preserves the total conformal-weight filtration used
throughout the \(W\)-algebra computations, so
Proposition~\ref{prop:mc4-weight-cutoff} applies.
\end{proof}

\begin{proposition}[Continuity of inverse-limit bar and cobar differentials;
\ClaimStatusProvedHere]
\label{prop:inverse-limit-differential-continuity}
Let $\{C_N,d_N,\Delta_N\}_{N \ge 0}$ be an inverse system of finite-stage
curved dg coalgebras with transition maps
$\pi_{N+1,N}\colon C_{N+1}\to C_N$ commuting with the differentials and
coproducts.  Endow
\[
\widehat C := \varprojlim_N C_N
\]
with the inverse-limit topology, and let
\[
\Omega(\widehat C)
:=
\prod_{m \ge 0} (s^{-1}\widehat C)^{\widehat\otimes m}
\]
carry the completed tensor-product topology.  Then:
\begin{enumerate}
\item the componentwise differential
  $\widehat d((x_N)_N):=(d_Nx_N)_N$ on $\widehat C$ is well defined and
  continuous;
\item the inverse-limit coproduct $\widehat\Delta$ is continuous; and
\item the completed cobar differential on $\Omega(\widehat C)$ induced
  by $\widehat d$ and $\widehat\Delta$ is continuous.
\end{enumerate}
Consequently, once an MC4 tower is realized as a genuine inverse system
of finite-stage curved dg coalgebras, continuity is formal for the
inverse-limit topology.
\end{proposition}

\begin{proof}
Let $p_N\colon \widehat C \to C_N$ be the projection.  Because the
transition maps commute with the differentials,
\[
p_N \circ \widehat d = d_N \circ p_N
\]
for every $N$, so the componentwise formula defines a map
$\widehat d\colon \widehat C \to \widehat C$.  The inverse-limit
topology on $\widehat C$ is initial with respect to the projections
$p_N$, hence $\widehat d$ is continuous as soon as each composite
$p_N \circ \widehat d$ is continuous.  This holds because each $C_N$ is
given the finite-stage topology and $d_N$ is a morphism of complexes.

The same argument applies to the coproduct.  Writing
$p_N^{\otimes 2}\colon \widehat C \widehat\otimes \widehat C \to
C_N\otimes C_N$ for the induced projection, one has
\[
p_N^{\otimes 2} \circ \widehat\Delta = \Delta_N \circ p_N,
\]
so $\widehat\Delta$ is well defined and continuous.

The completed cobar differential on $\Omega(\widehat C)$ is the sum of
the internal part coming from $\widehat d$ and the bar part coming from
$\widehat\Delta$.  On each tensor-length summand
$(s^{-1}\widehat C)^{\widehat\otimes m}$, each contribution is a finite
sum of maps obtained by inserting either $\widehat d$ or
$\widehat\Delta$ in one slot.  Completed tensor products preserve
continuity of continuous multilinear maps, and finite sums of
continuous maps remain continuous.  The product topology on
$\Omega(\widehat C)$ is defined componentwise by the tensor-length
summands, so the total cobar differential is continuous.
\end{proof}

\begin{corollary}[Standard principal-stage
\texorpdfstring{$W_\infty$}{W_infty} tower satisfies the M-level MC4
package; \ClaimStatusProvedHere]
\label{cor:winfty-standard-mc4-package}
For the standard principal finite-type stages $W_N$ of $W_\infty$, the
completed bar object
\[
\widehat{\bar B}(W_\infty) := \varprojlim_N \bar B(W_N)
\]
has continuous bar and cobar differentials, and the canonical
comparison map
\[
\Omega\bigl(\widehat{\bar B}(W_\infty)\bigr) \xrightarrow{\sim} W_\infty
\]
is a quasi-isomorphism.
Thus the standard principal-stage tower realizes the full
standard-tower M-level completed bar-cobar package for $W_\infty$; the
remaining frontier is
the H-level comparison with stronger factorization-theoretic or
physical realizations.
\end{corollary}

\begin{proof}
The finite-stage bar complexes form an inverse system of curved dg
coalgebras by Proposition~\ref{prop:winfty-mc4-criterion},
Assumption~\textup{(3)}, so
Proposition~\ref{prop:inverse-limit-differential-continuity} gives the
continuity of the completed bar and cobar differentials.  The
finite-stage bar-cobar quasi-isomorphisms are provided by the proved
principal finite-type regime, and
Corollary~\ref{cor:winfty-weight-cutoff} supplies the
Mittag--Leffler/stabilization input.  Since
$W_\infty = \varprojlim_N W_N$ by definition, applying
Proposition~\ref{prop:mc4-reduction-principle} to the standard tower
gives the stated quasi-isomorphism.
\end{proof}

\begin{proposition}[Comparison with a completed target by compatible
finite quotients; \ClaimStatusProvedHere]
\label{prop:completed-target-comparison}
Let $\{\cA_{\le N}\}_{N \ge 0}$ be an inverse system of dg algebras with
surjective transition maps and completed limit
\[
\widehat{\cA}:=\varprojlim_N \cA_{\le N}.
\]
Let $\cB$ be a separated complete dg algebra with descending filtration
by closed dg ideals
\[
F^1\cB \supset F^2\cB \supset \cdots.
\]
Assume:
\begin{enumerate}
\item the canonical map
  \[
  \cB \xrightarrow{\sim} \varprojlim_N \cB/F^{N+1}\cB
  \]
  is an isomorphism of dg algebras;
\item for each $N$ there is a dg algebra quasi-isomorphism
  \[
  f_N\colon \cB/F^{N+1}\cB \xrightarrow{\sim} \cA_{\le N};
  \]
\item the maps $f_N$ commute with the quotient and transition maps; and
\item for every cohomological degree $m$, the inverse systems
  $\{H^m(\cB/F^{N+1}\cB)\}_N$ and $\{H^m(\cA_{\le N})\}_N$ satisfy the
  Mittag--Leffler condition.
\end{enumerate}
Then the induced map
\[
\widehat f\colon \cB \longrightarrow \widehat{\cA}
\]
is a quasi-isomorphism.

If, in addition, the quotient tower $\{\cB/F^{N+1}\cB\}_N$ carries
bar-cobar comparison maps compatible with the tower
$\{\cA_{\le N}\}_N$ and the completed bar and cobar differentials on
$\cB$ are continuous, then $\cB$ and $\widehat{\cA}$ define
quasi-isomorphic completed M-level bar-cobar objects.
\end{proposition}

\begin{proof}
Assumption~\textup{(1)} identifies $\cB$ with the inverse limit of its
finite quotients.  By Assumptions~\textup{(2)} and~\textup{(3)}, the
maps $f_N$ assemble into a morphism of inverse systems and hence into a
map
\[
\widehat f\colon
\cB \cong \varprojlim_N \cB/F^{N+1}\cB \longrightarrow
\varprojlim_N \cA_{\le N}=\widehat{\cA}.
\]

Apply the Milnor exact sequence to the quotient tower of $\cB$ and to
the tower $\{\cA_{\le N}\}_N$.  Assumption~\textup{(4)} kills the
derived-limit terms, so for every cohomological degree $m$ one has
\[
H^m(\cB)\cong \varprojlim_N H^m(\cB/F^{N+1}\cB),
\qquad
H^m(\widehat{\cA})\cong \varprojlim_N H^m(\cA_{\le N}).
\]
Since each $f_N$ is a quasi-isomorphism, the maps
$H^m(f_N)\colon H^m(\cB/F^{N+1}\cB)\xrightarrow{\sim}
H^m(\cA_{\le N})$ are isomorphisms compatible with the inverse systems.
Passing to inverse limits shows that $H^m(\widehat f)$ is an isomorphism
for every $m$, so $\widehat f$ is a quasi-isomorphism.

For the final claim, apply the same inverse-limit comparison to the
compatible bar towers.  Continuity of the completed bar and cobar
differentials makes the completed bar objects well defined, and
Proposition~\ref{prop:mc4-reduction-principle} identifies both
completed bar-cobar packages with the same inverse-limit finite-stage
data.  Hence the resulting completed M-level bar-cobar objects are
quasi-isomorphic.
\end{proof}

\begin{corollary}[H-level comparison criterion for
\texorpdfstring{$W_\infty$}{W_infty}; \ClaimStatusProvedHere]
\label{cor:winfty-hlevel-comparison-criterion}
Let $\mathcal{W}^{\mathrm{ht}}$ be a separated complete dg model for a
factorization-theoretic or physical completion of $W_\infty$, equipped
with a descending conformal-weight filtration by closed dg ideals.
Assume:
\begin{enumerate}
\item $\mathcal{W}^{\mathrm{ht}}$ is the inverse limit of its finite quotients;
\item each finite quotient
  $\mathcal{W}^{\mathrm{ht}}/F^{N+1}\mathcal{W}^{\mathrm{ht}}$ is dg
  quasi-isomorphic to the principal finite-type stage $W_N$,
  compatibly with the truncation tower; and
\item the quotient tower inherits the same finite-stage bar-cobar data
  and continuity hypotheses as in
  Corollary~\ref{cor:winfty-standard-mc4-package}.
\end{enumerate}
Then the induced comparison map
\[
\mathcal{W}^{\mathrm{ht}} \xrightarrow{\sim} W_\infty
\]
is a quasi-isomorphism, and $\mathcal{W}^{\mathrm{ht}}$ carries the same
completed M-level bar-cobar object as the standard principal-stage
completion.  Thus the live $\mathcal{W}_\infty$ frontier is to
construct such a filtered H-level target and identify its finite
quotients by the exact residue identities
$\mathsf{C}^{\mathrm{res}}_{s,t;u;m,n}(N)
=\mathsf{C}^{\mathrm{DS}}_{s,t;u;m,n}(N)$,
discharged on the finite-detection packet $\mathcal{I}_N$, whose first
live higher-spin stages are already reduced on the theorem surface to
the stage-$3$ fifteen-coefficient packet and the exact stage-$4$
six-entry identity packet on~$\mathcal{I}_4$, organized as three local
OPE blocks and splitting into four higher-spin channels together with
two theorematic principal Virasoro-target values
$\mathsf{C}^{\mathrm{DS}}_{4,4;2;0,6}=2$ and
$\mathsf{C}^{\mathrm{DS}}_{3,4;2;0,5}=0$, with the corresponding
residue-side contraction supplied only under the normalized two-point
package of Corollary~\ref{cor:winfty-stage4-residue-four-channel}, not
another stabilization theorem.
\end{corollary}

\begin{proof}
Apply Proposition~\ref{prop:completed-target-comparison} with
$\cA_{\le N}=W_N$.  The Mittag--Leffler input is supplied by the same
weight-cutoff argument as in Corollary~\ref{cor:winfty-weight-cutoff},
and the final statement follows from the bar-cobar comparison part of
Proposition~\ref{prop:completed-target-comparison}.
\end{proof}

\begin{definition}[Principal-stage compatible
\texorpdfstring{$W_\infty$}{W_infty} target]
\label{def:winfty-principal-stage-compatible}
A separated complete dg algebra $\mathcal{W}^{\mathrm{ht}}$ is called a
\emph{principal-stage compatible $W_\infty$ target} if it is equipped
with a descending conformal-weight filtration by closed dg ideals such
that:
\begin{enumerate}
\item $\mathcal{W}^{\mathrm{ht}} \cong
  \varprojlim_N \mathcal{W}^{\mathrm{ht}}/F^{N+1}\mathcal{W}^{\mathrm{ht}}$;
\item each finite quotient
  $\mathcal{W}^{\mathrm{ht}}/F^{N+1}\mathcal{W}^{\mathrm{ht}}$ is dg
  quasi-isomorphic to the principal finite-type stage $W_N$;
\item the quotient maps intertwine the principal-stage truncation tower;
  and
\item the quotient tower carries bar-cobar comparison maps compatible
  with the finite-type $W_N$ bar complexes.
\end{enumerate}
This is the H-level input singled out by
Corollary~\ref{cor:winfty-hlevel-comparison-criterion}.
\end{definition}

\begin{definition}[Principal-stage quotient system for a factorization
\texorpdfstring{$W_\infty$}{W_infty} completion]
\label{def:winfty-quotient-system}
Let $\mathcal{W}^{\mathrm{fact}}_\infty$ be a separated complete
factorization-theoretic completion of the principal $W$-tower.  A
\emph{principal-stage quotient system} on
$\mathcal{W}^{\mathrm{fact}}_\infty$ consists of dg quotient maps
\[
q_N\colon \mathcal{W}^{\mathrm{fact}}_\infty \longrightarrow
\mathcal{W}^{\mathrm{fact}}_{\le N}
\]
such that:
\begin{enumerate}
\item the kernels of the $q_N$ define a descending separated complete
  conformal-weight filtration;
\item each quotient object
  $\mathcal{W}^{\mathrm{fact}}_{\le N}$ is dg quasi-isomorphic to the
  principal finite-type stage $W_N$;
\item the transition maps between the quotients intertwine the
  principal-stage truncation tower; and
\item the induced bar complexes agree with the finite-type bar-cobar
  package already proved for $W_N$.
\end{enumerate}
\end{definition}

\begin{proposition}[Formal descent criterion for the
\texorpdfstring{$W_\infty$}{W_infty} factorization target;
\ClaimStatusProvedHere]
\label{prop:winfty-quotient-system-criterion}
A separated complete factorization-theoretic completion equipped with a
principal-stage quotient system
(Definition~\ref{def:winfty-quotient-system}) is a principal-stage
compatible $W_\infty$ target
(Definition~\ref{def:winfty-principal-stage-compatible}).
Bar-cobar compatibility further implies
Corollary~\ref{cor:winfty-hlevel-comparison-criterion}.
\end{proposition}

\begin{proof}
The quotient-system data package precisely the inputs of Definition~\ref{def:winfty-principal-stage-compatible}.
\end{proof}

\begin{proposition}[Factorization-envelope criterion for principal
stages; \ClaimStatusProvedHere]
\label{prop:winfty-factorization-envelope-criterion}
Let $\mathcal{F}_\infty$ be a separated complete factorization algebra
on a curve with quotient maps
$\mathcal{F}_\infty \twoheadrightarrow \mathcal{F}_{\le N}$
whose kernels define a descending separated complete conformal-weight
filtration.  If each $\mathcal{F}_{\le N}$ admits a BD chiral envelope
quasi-isomorphic to $W_N$ and the quotient maps intertwine the
principal-stage tower, then the chiral envelope of
$\mathcal{F}_\infty$ is a principal-stage compatible $W_\infty$ target.
Bar-data compatibility further implies
Proposition~\ref{prop:winfty-quotient-system-criterion} and
Corollary~\ref{cor:winfty-hlevel-comparison-criterion}.
\end{proposition}

\begin{proof}
The chiral-envelope functor produces compatible quotients recovering the principal stages; Proposition~\ref{prop:winfty-quotient-system-criterion} applies.
\end{proof}

\begin{proposition}[Factorization realization package for
\texorpdfstring{$W_\infty$}{W_infty}; \ClaimStatusProvedElsewhere]
\label{prop:winfty-factorization-package}
There exists a principal-stage compatible $W_\infty$ target
$\mathcal{W}^{\mathrm{fact}}_\infty$ in the sense of
Definition~\ref{def:winfty-principal-stage-compatible}, arising from a
factorization-theoretic or physical completion of the principal
$W$-tower.  Equivalently:
\begin{enumerate}
\item $\mathcal{W}^{\mathrm{fact}}_\infty$ carries a separated complete
  conformal-weight filtration;
\item its finite quotients recover the theorematic principal finite-type
  stages $W_N$;
\item the induced bar-cobar data on those quotients agree with the
  standard finite-stage bar complexes; and
\item the resulting comparison map
  \[
  \mathcal{W}^{\mathrm{fact}}_\infty \xrightarrow{\sim} W_\infty
  \]
  identifies the factorization-theoretic completion with the standard
  principal-stage completed M-level package.
\end{enumerate}
This is resolved by
Theorem~\ref{thm:winfty-factorization-kd}: the factorization
target is constructed via $\Einf$-sectorwise convergence at each
finite stage, with the DS--KD intertwining
(Theorem~\ref{thm:ds-koszul-intertwine}) providing the
finite-stage identifications and the Mittag--Leffler property
(Corollary~\ref{cor:winfty-standard-mc4-package}) controlling
the inverse limit.
\end{proposition}

\begin{remark}[Attack route for the
\texorpdfstring{$W_\infty$}{W_infty} factorization target]
\label{rem:winfty-factorization-route}
The reduction chain from
Proposition~\ref{prop:winfty-factorization-package} proceeds through
higher-spin ideals
(Proposition~\ref{prop:winfty-higher-spin-ideal-criterion}),
spin-triangular OPE closure
(Proposition~\ref{prop:winfty-spin-triangular-ideals}),
coefficient-level DS matching
(Proposition~\ref{prop:winfty-ds-local-coefficient-criterion}),
residue identity
(Proposition~\ref{prop:winfty-ds-residue-identity-criterion}),
generator seeds
(Proposition~\ref{prop:winfty-ds-generator-seed},
Corollary~\ref{cor:winfty-ds-finite-seed-set}),
and the explicit stage-$3$ $W_3$ packet
(Proposition~\ref{prop:winfty-ds-stage3-explicit-packet}).
The first genuinely new higher-spin packet is stage~$4$.
\end{remark}

\begin{proposition}[Higher-spin ideal criterion for principal-stage
\texorpdfstring{$W_\infty$}{W_infty} quotients; \ClaimStatusProvedHere]
\label{prop:winfty-higher-spin-ideal-criterion}
Let $\mathcal{F}_\infty$ be a separated complete factorization algebra
on a curve, topologically generated by fields
$W^{(s)}$ of conformal weights $s \ge 2$.  For each $N \ge 2$, let
$I_{>N}\mathcal{F}_\infty$ denote the closed factorization ideal
generated by the fields $W^{(s)}$ with $s > N$.  Assume:
\begin{enumerate}
\item the differential, translation operator, and factorization
  products preserve each ideal $I_{>N}\mathcal{F}_\infty$;
\item modulo $I_{>N}\mathcal{F}_\infty$, the OPE of the generators
  $W^{(2)},\dots,W^{(N)}$ closes among those same generators;
\item the quotient factorization algebra
  $\mathcal{F}_{\le N}:=\mathcal{F}_\infty/I_{>N}\mathcal{F}_\infty$
  admits a chiral envelope or equivalent dg model quasi-isomorphic to
  the principal stage $W_N$; and
\item the quotient maps
  $\mathcal{F}_{\le N+1}\twoheadrightarrow \mathcal{F}_{\le N}$
  intertwine the principal-stage truncation tower.
\end{enumerate}
Then the quotients $\{\mathcal{F}_{\le N}\}_{N \ge 2}$ define a
principal-stage quotient system in the sense of
Definition~\ref{def:winfty-quotient-system}.  Consequently,
Proposition~\ref{prop:winfty-factorization-envelope-criterion},
Proposition~\ref{prop:winfty-quotient-system-criterion}, and
Corollary~\ref{cor:winfty-hlevel-comparison-criterion} apply.

In particular, the remaining $\mathcal{W}_\infty$ construction problem
is to prove that the higher-spin generators above stage $N$ define
stable factorization ideals whose quotients recover the theorematic
principal stages.
\end{proposition}

\begin{proof}
Assumptions~\textup{(1)--(4)} supply closed factorization ideals, closed OPE, finite-stage identification with
$W_N$, and compatibility with the
principal-stage tower.  This is exactly the data required by
Definition~\ref{def:winfty-quotient-system}.  The stated comparison
results then follow from
Proposition~\ref{prop:winfty-factorization-envelope-criterion},
Proposition~\ref{prop:winfty-quotient-system-criterion}, and
Corollary~\ref{cor:winfty-hlevel-comparison-criterion}.
\end{proof}

\begin{proposition}[Spin-triangular OPE criterion for the
\texorpdfstring{$W_\infty$}{W_infty} factorization ideals;
\ClaimStatusProvedHere]
\label{prop:winfty-spin-triangular-ideals}
Let $\mathcal{F}_\infty$ be a separated complete factorization algebra
on a curve, topologically generated by fields $W^{(s)}$ of conformal
weight $s \ge 2$.  Assume:
\begin{enumerate}
\item for every pair of generators, the coefficient of the pole
  $(z-w)^{-n}$ in the local OPE
  $W^{(s)}(z)W^{(t)}(w)$ is a linear combination of derivatives of
  fields $W^{(u)}$ with $u \le s+t-n$;
\item the translation operator and the residue/bar operations are
  obtained from these OPE coefficients by the standard factorization
  formulas, so no operation increases the maximal conformal weight
  appearing in an expression; and
\item for each $N$, after quotienting by the fields $W^{(s)}$ with
  $s>N$, the induced OPE formulas on
  $W^{(2)},\dots,W^{(N)}$ agree with the principal finite-type stage
  $W_N$.
\end{enumerate}
Then the closed factorization ideals
\[
I_{>N}\mathcal{F}_\infty
= \langle W^{(s)} \mid s>N\rangle_{\mathrm{fact}}
\]
form a compatible principal-stage quotient system.  In particular,
Proposition~\ref{prop:winfty-higher-spin-ideal-criterion} applies, and
the resulting quotients recover the principal finite-type stages.
\end{proposition}

\begin{proof}
Spin-triangularity (1) closes the quotient OPE on generators $\le N$; (2) preserves the ideal under all operations; (3) identifies the quotient with $W_N$.  Nesting of ideals gives compatibility for varying $N$.
\end{proof}

\begin{proposition}[Coefficient-level DS criterion for principal-stage
\texorpdfstring{$W_\infty$}{W_infty} quotients;
\ClaimStatusProvedHere]
\label{prop:winfty-ds-coefficient-criterion}
Fix $N \ge 2$ and let
$\mathcal{F}_{\le N}:=\mathcal{F}_\infty/I_{>N}\mathcal{F}_\infty$
be a quotient from
Proposition~\ref{prop:winfty-spin-triangular-ideals}.
If the quotient is strongly generated by $W^{(2)},\dots,W^{(N)}$ with
OPE coefficients matching the principal DS formulas for $W_N$, no
additional local relations, and induced bar operations coinciding with
the finite-stage package, then $\mathcal{F}_{\le N}$ identifies with
$W_N$ and
Proposition~\ref{prop:winfty-higher-spin-ideal-criterion} applies.
\end{proposition}

\begin{proof}
The hypotheses supply the same generators, OPE, and bar data as $W_N$; the reconstruction principle gives the identification.
\end{proof}

\begin{proposition}[Local-coefficient criterion for principal-stage
\texorpdfstring{$W_\infty$}{W_infty} quotients;
\ClaimStatusProvedHere]
\label{prop:winfty-ds-local-coefficient-criterion}
Fix $N \ge 2$ and let
$\mathcal{F}_{\le N}:=\mathcal{F}_\infty/I_{>N}\mathcal{F}_\infty$
be as in
Proposition~\ref{prop:winfty-spin-triangular-ideals}.  If the local
OPE coefficients $C_{s,t}^{u;m,n}(N)$ agree with the principal DS data,
translation descendants generate all local relations, and bar operations
match, then Proposition~\ref{prop:winfty-ds-coefficient-criterion}
applies and the quotient identifies with $W_N$.
\end{proposition}

\begin{proof}
Immediate from Proposition~\ref{prop:winfty-ds-coefficient-criterion}.
\end{proof}

\begin{proposition}[Residue-coefficient identity criterion for
principal-stage \texorpdfstring{$W_\infty$}{W_infty} quotients;
\ClaimStatusProvedHere]
\label{prop:winfty-ds-residue-identity-criterion}
Fix $N \ge 2$ and let
$\mathcal{F}_{\le N}:=\mathcal{F}_\infty/I_{>N}\mathcal{F}_\infty$
be as in Proposition~\ref{prop:winfty-spin-triangular-ideals}.
If $\mathsf{C}^{\mathrm{res}}_{s,t;u;m,n}(N)
= \mathsf{C}^{\mathrm{DS}}_{s,t;u;m,n}(N)$
for every admissible tuple and the quotient OPE and bar operations are
entirely determined by these residue coefficients, then
Proposition~\ref{prop:winfty-ds-local-coefficient-criterion} applies
and the quotient identifies with $W_N$.
\end{proposition}

\begin{proof}
Immediate from Proposition~\ref{prop:winfty-ds-local-coefficient-criterion}.
\end{proof}

\begin{proposition}[Generator-seed criterion for principal-stage
\texorpdfstring{$W_\infty$}{W_infty} residue identities;
\ClaimStatusProvedHere]
\label{prop:winfty-ds-generator-seed}
Fix $N \ge 2$ and let
\[
\mathcal{F}_{\le N}:=\mathcal{F}_\infty/I_{>N}\mathcal{F}_\infty
\]
be as in
Proposition~\ref{prop:winfty-ds-residue-identity-criterion}.  Assume:
\begin{enumerate}
\item $\mathcal{F}_{\le N}$ is strongly generated by
  $W^{(2)},\dots,W^{(N)}$;
\item for each pair of generators $W^{(s)},W^{(t)}$ with
  $2 \le s,t \le N$, the primary residue coefficients
  \[
  \mathsf{C}^{\mathrm{res}}_{s,t;u;0,n}(N)
  \]
  agree with the principal Drinfeld--Sokolov coefficients
  $\mathsf{C}^{\mathrm{DS}}_{s,t;u;0,n}(N)$ for every visible pole order
  and target spin~$u$;
\item translation commutes with the residue extraction formulas, so all
  descendant coefficients are obtained from the primary ones by
  repeated application of $\partial$; and
\item the bar operations are generated from the same residue calculus
  and therefore are determined by those translated coefficient
  families.
\end{enumerate}
Then the full mode-by-mode identities
\[
\mathsf{C}^{\mathrm{res}}_{s,t;u;m,n}(N)
=\mathsf{C}^{\mathrm{DS}}_{s,t;u;m,n}(N)
\]
hold for all admissible indices.  Consequently the hypotheses of
Proposition~\ref{prop:winfty-ds-residue-identity-criterion} are
satisfied.
\end{proposition}

\begin{proof}
Primary generator agreement propagates to all descendants by translation (assumption~(3)), and strong generation (assumption~(1)) plus bar recovery (assumption~(4)) identify the full coefficient family, verifying Proposition~\ref{prop:winfty-ds-residue-identity-criterion}.
\end{proof}

\begin{corollary}[Finite primary seed set for principal-stage
\texorpdfstring{$W_\infty$}{W_infty} comparison;
\ClaimStatusProvedHere]
\label{cor:winfty-ds-finite-seed-set}
Fix $N \ge 2$ and define the finite index set
\[
\mathcal{I}_N:=
\left\{
(s,t,u,n)\ \middle|\
\begin{array}{l}
2 \le s \le t \le N,\\
2 \le u \le \min(N,s+t-1),\\
1 \le n \le s+t-u
\end{array}
\right\}.
\]
Assume the hypotheses of
Proposition~\ref{prop:winfty-ds-generator-seed}.  Then it is enough to
check the primary identities
\[
\mathsf{C}^{\mathrm{res}}_{s,t;u;0,n}(N)
=\mathsf{C}^{\mathrm{DS}}_{s,t;u;0,n}(N)
\qquad\text{for all } (s,t,u,n)\in\mathcal{I}_N.
\]
These finitely many checks imply the full residue/OPE identity package
required by
Proposition~\ref{prop:winfty-ds-residue-identity-criterion}.
\end{corollary}

\begin{proof}
Spin-triangularity implies that for fixed $s,t \le N$ only target spins
$u \le s+t-1$ and pole orders $n \le s+t-u$ can appear.  Since the
stage quotient discards spins $>N$, the admissible tuples are exactly
those in $\mathcal{I}_N$, hence finite.

Proposition~\ref{prop:winfty-ds-generator-seed} shows that agreement on
these primary generator coefficients propagates by translation to every
descendant coefficient and therefore to the full residue/bar package.
\end{proof}

\begin{corollary}[First principal-stage seed packets for
\texorpdfstring{$W_\infty$}{W_infty} comparison;
\ClaimStatusProvedHere]
\label{cor:winfty-ds-lowstage-seeds}
The finite seed sets of
Corollary~\ref{cor:winfty-ds-finite-seed-set} begin as follows:
\[
\mathcal{I}_2
=\{(2,2,2,1),(2,2,2,2)\},
\]
and
\begin{align*}
\mathcal{I}_3
= {} & \{(2,2,2,1),(2,2,2,2),(2,2,3,1)\} \\
& \sqcup\{(2,3,2,1),(2,3,2,2),(2,3,2,3),(2,3,3,1),(2,3,3,2)\} \\
& \sqcup\{(3,3,2,1),(3,3,2,2),(3,3,2,3),(3,3,2,4), \\
& \qquad (3,3,3,1),(3,3,3,2),(3,3,3,3)\}.
\end{align*}
In particular, the first genuinely new infinite-generator stage beyond
the theorematic Virasoro block is the stage-$3$ packet of fifteen
primary coefficients: the two Virasoro coefficients, the single
vanishing $W^{(2)}W^{(2)}\to W^{(3)}$ coefficient, five mixed
$W^{(2)}$/$W^{(3)}$ coefficients, and seven $W^{(3)}$ self-coupling
coefficients.
\end{corollary}

\begin{proof}
For $N=2$, the defining inequalities for $\mathcal{I}_N$ force
$s=t=u=2$ and $1\le n\le 2$, giving exactly the two displayed tuples.
For $N=3$, the admissible ordered pairs of generator spins are
$(2,2)$, $(2,3)$, and $(3,3)$.  The inequalities
\[
2\le u\le \min(3,s+t-1),
\qquad
1\le n\le s+t-u
\]
then give the displayed decomposition into the three corresponding
blocks.  The final sentence is simply a repackaging of those tuples by
source pair and target spin.
\end{proof}

\begin{proposition}[Incremental interacting packet from stage
\texorpdfstring{$N$}{N} to stage \texorpdfstring{$N{+}1$}{N+1};
\ClaimStatusProvedHere]
\label{prop:winfty-ds-stage-growth-packet}
Fix $N \ge 3$ and let $\mathcal{I}_N$ and $\mathcal{I}_{N+1}$ be the
seed packets of Corollary~\ref{cor:winfty-ds-finite-seed-set}.  Write
\[
\mathcal{I}_{N+1}^{\mathrm{Vir,new}}
:=
\mathcal{I}_{N+1}^{\mathrm{Vir}}\setminus \mathcal{I}_N
=
\{(2,N{+}1,u,n)\in \mathcal{I}_{N+1}\},
\]
and define the incremental interacting packet
\[
\mathcal{J}_{N+1}
:=
\mathcal{I}_{N+1}\setminus
\bigl(\mathcal{I}_N\sqcup \mathcal{I}_{N+1}^{\mathrm{Vir,new}}\bigr).
\]
Assume the hypotheses of
Proposition~\ref{prop:winfty-ds-generator-seed}.  Then:
\begin{enumerate}
\item the coefficients indexed by $\mathcal{I}_N$ are exactly the
  stage-$N$ packet identities;
\item the coefficients indexed by
  $\mathcal{I}_{N+1}^{\mathrm{Vir,new}}$ are determined by the exact
  stress-tensor action on the new generator $W^{(N+1)}$;
\item the genuinely new interacting stage-$(N+1)$ packet is therefore
  exactly $\mathcal{J}_{N+1}$, and it decomposes as the disjoint union
  \[
  \mathcal{J}_{N+1}
  =
  \mathcal{J}_{N+1}^{\mathrm{tail}}
  \sqcup
  \mathcal{J}_{N+1}^{\mathrm{mix}}
  \sqcup
  \mathcal{J}_{N+1}^{\mathrm{self}},
  \]
  where
  \[
  \mathcal{J}_{N+1}^{\mathrm{tail}}
  :=
  \{(s,t,N{+}1,n)\in\mathcal{I}_{N+1}\mid 3\le s\le t\le N\},
  \]
  \[
  \mathcal{J}_{N+1}^{\mathrm{mix}}
  :=
  \{(s,N{+}1,u,n)\in\mathcal{I}_{N+1}\mid 3\le s\le N,\ 2\le u\le N{+}1\},
  \]
  \[
  \mathcal{J}_{N+1}^{\mathrm{self}}
  :=
  \{(N{+}1,N{+}1,u,n)\in\mathcal{I}_{N+1}\mid 2\le u\le N{+}1\}.
  \]
\end{enumerate}
Consequently, once the stage-$N$ packet identities on $\mathcal{I}_N$
are known, the stage-$(N+1)$ comparison is equivalent to matching the
coefficients on the incremental interacting packet $\mathcal{J}_{N+1}$.
For $N=3$, this is exactly the stage-$4$ residual packet of
Proposition~\ref{prop:winfty-ds-stage4-residual-packet}.
\end{proposition}

\begin{proof}
Since $\mathcal{I}_N\subset \mathcal{I}_{N+1}$, new tuples have at least one index equal to $N{+}1$.  Those with $s=2$ form $\mathcal{I}_{N+1}^{\mathrm{Vir,new}}$; the complement $\mathcal{J}_{N+1}$ (with $s\geq 3$) splits exhaustively into the three displayed blocks by whether the new index appears in $u$, $t$, or both.
\end{proof}

\begin{corollary}[Top-pole/parity reduction of the incremental
\texorpdfstring{$W_\infty$}{W_infty} stage-growth packet;
\ClaimStatusProvedHere]
\label{cor:winfty-ds-stage-growth-top-parity}
Fix $N \ge 3$ and assume the hypotheses of
Proposition~\ref{prop:winfty-ds-stage-growth-packet},
Proposition~\ref{prop:winfty-ds-primary-top-pole}, and
Proposition~\ref{prop:winfty-ds-self-ope-parity}.  Let
\[
\mathcal{J}_{N+1}^{\mathrm{top}}
:=
\{(s,t,u,s+t-u)\in \mathcal{J}_{N+1}\},
\]
and remove the odd self-OPE top-pole entries by setting
\[
\mathcal{J}_{N+1}^{\mathrm{red}}
:=
\mathcal{J}_{N+1}^{\mathrm{top}}
\setminus
\bigl\{(s,s,u,2s-u)\in\mathcal{J}_{N+1}^{\mathrm{top}}
\mid 2s-u \text{ odd}\bigr\}.
\]
Then, once the stage-$N$ packet identities on $\mathcal{I}_N$ are
known, the stage-$(N+1)$ comparison closes as soon as the identities on
the reduced incremental packet
$\mathcal{J}_{N+1}^{\mathrm{red}}$ are proved.

For $N=3$, one has
\[
\mathcal{J}_4^{\mathrm{red}}
=
\mathcal{J}_4^{\mathrm{par}},
\]
so this is exactly the six-entry stage-$4$ block of
Corollary~\ref{cor:winfty-ds-stage4-parity-packet}.  The additional
contraction of that six-entry block to the four higher-spin channels is
conjecturally controlled by
Conjecture~\ref{conj:winfty-stage4-ward-inheritance}.
\end{corollary}

\begin{proof}
Top-pole primaryity (Proposition~\ref{prop:winfty-ds-primary-top-pole}) and odd self-OPE parity (Proposition~\ref{prop:winfty-ds-self-ope-parity}) reduce $\mathcal{J}_{N+1}$ to $\mathcal{J}_{N+1}^{\mathrm{red}}$.  For $N=3$ this recovers the stage-$4$ parity packet.
\end{proof}

\begin{corollary}[First reduced incremental packet beyond
\texorpdfstring{$\mathcal{I}_4$}{I4}; \ClaimStatusProvedHere]
\label{cor:winfty-ds-stage5-reduced-packet}
Assume the hypotheses of
Corollary~\ref{cor:winfty-ds-stage-growth-top-parity} with $N=4$.
Then the first reduced incremental packet beyond the stage-$4$ block is
\[
\mathcal{J}_5^{\mathrm{red}}
=
\Bigl\{
(3,4;5;0,2)
\Bigr\}
\sqcup
\Bigl\{
(3,5;2;0,6),\,
(3,5;3;0,5),\,
(3,5;4;0,4),\,
(3,5;5;0,3)
\Bigr\}
\]
\[
\sqcup
\Bigl\{
(4,5;2;0,7),\,
(4,5;3;0,6),\,
(4,5;4;0,5),\,
(4,5;5;0,4)
\Bigr\}
\sqcup
\Bigl\{
(5,5;2;0,8),\,
(5,5;4;0,6)
\Bigr\}.
\]
In particular,
\[
|\mathcal{J}_5^{\mathrm{red}}|=11,
\]
decomposed as one tail channel, two four-entry mixed blocks, and two
self-coupling channels.
\end{corollary}

\begin{proof}
Enumerate the tail, mixed, and self blocks of $\mathcal{J}_5$ from Proposition~\ref{prop:winfty-ds-stage-growth-packet}; apply top-pole primaryity and odd self-OPE parity to each.  The surviving entries are exactly the displayed eleven tuples.
\end{proof}

\begin{proposition}[Primary top-pole criterion for generator seed
packets; \ClaimStatusProvedHere]
\label{prop:winfty-ds-primary-top-pole}
Fix $N \ge 2$ and assume the hypotheses of
Proposition~\ref{prop:winfty-ds-generator-seed}.  Assume moreover that
for each $2 \le u \le N$ the strong generator $W^{(u)}$ is primary of
conformal weight $u$ with respect to the Virasoro field
$T:=W^{(2)}$, and that the local OPE is decomposed into conformal
families of these primary generators and their translation
descendants.

Then for every visible triple $(s,t,u)$ with
\[
2 \le s \le t \le N,
\qquad
2 \le u \le \min(N,s+t-1),
\]
the primary generator $W^{(u)}$ can appear only at the top pole order
\[
n_{\max}(s,t;u):=s+t-u.
\]
Equivalently,
\[
\mathsf{C}^{\mathrm{res}}_{s,t;u;0,n}(N)
\;=\;
\mathsf{C}^{\mathrm{DS}}_{s,t;u;0,n}(N)
\;=\;0
\qquad\text{for every } 1\le n < s+t-u.
\]
Consequently the primary seed set of
Corollary~\ref{cor:winfty-ds-finite-seed-set} reduces further to the
top-pole subset
\[
\mathcal{I}_N^{\mathrm{top}}
:=\{(s,t,u,s+t-u)\in \mathcal{I}_N\}.
\]
\end{proposition}

\begin{proof}
Because $W^{(s)}$, $W^{(t)}$, and $W^{(u)}$ are primary fields of
weights $s$, $t$, and $u$, conformal covariance organizes the
$W^{(u)}$-family inside the OPE
$W^{(s)}(z)W^{(t)}(w)$ as one primary coefficient multiplying
$W^{(u)}(w)$ at pole order $s+t-u$, together with the translation
descendants $\partial^m W^{(u)}(w)$ at the lower poles
$s+t-u-m$.  The undifferentiated primary field therefore occurs only at
the top pole order $n=s+t-u$.

In the notation of
Proposition~\ref{prop:winfty-ds-generator-seed}, the coefficients
\[
\mathsf{C}^{\mathrm{res}}_{s,t;u;0,n}(N),
\qquad
\mathsf{C}^{\mathrm{DS}}_{s,t;u;0,n}(N)
\]
record the coefficients of the undifferentiated primary generator
$W^{(u)}$ itself.  Hence those coefficients vanish for all lower pole
orders $n<s+t-u$, both on the residue side and on the principal
Drinfeld--Sokolov side.  This proves the displayed vanishing and
identifies the reduced top-pole seed set $\mathcal{I}_N^{\mathrm{top}}$.
\end{proof}

\begin{proposition}[Odd top-pole vanishing for identical even
generators; \ClaimStatusProvedHere]
\label{prop:winfty-ds-self-ope-parity}
Fix $N \ge 2$ and assume the hypotheses of
Proposition~\ref{prop:winfty-ds-primary-top-pole}.  Assume moreover
that the strong generators $W^{(2)},\dots,W^{(N)}$ are even fields.
Then for every visible self-coupling triple $(s,s,u)$ one has
\[
\mathsf{C}^{\mathrm{res}}_{s,s;u;0,\,2s-u}(N)
=
\mathsf{C}^{\mathrm{DS}}_{s,s;u;0,\,2s-u}(N)
=0
\qquad\text{whenever } 2s-u \text{ is odd}.
\]
Equivalently, for identical even generators the top-pole primary
coefficient can be nonzero only when the top pole order is even.
\end{proposition}

\begin{proof}
Let $A:=W^{(s)}$ and set $m:=2s-u-1$, so the top-pole primary
coefficient of the target generator $W^{(u)}$ in the self-OPE
$A(z)A(w)$ is the $A_{(m)}A$ term.  Because $A$ is even, the
vertex-algebra skew-symmetry formula gives
\[
A_{(m)}A
=
\sum_{j\ge 0}\frac{(-1)^{m+j+1}}{j!}\,\partial^j(A_{(m+j)}A).
\]
By Proposition~\ref{prop:winfty-ds-primary-top-pole}, the
undifferentiated primary generator channel occurs only at the top pole,
so the terms with $j>0$ do not contribute to the same primary
coefficient.  Therefore the primary coefficient satisfies
\[
A_{(m)}A = (-1)^{m+1}A_{(m)}A = (-1)^{2s-u}A_{(m)}A.
\]
If $2s-u$ is odd, this becomes $A_{(m)}A=-A_{(m)}A$, hence the
coefficient vanishes.  The same skew-symmetry identity holds on the
principal Drinfeld--Sokolov side, so both displayed coefficients
vanish.
\end{proof}

\begin{proposition}[Stage-\texorpdfstring{$3$}{3} principal packet from the explicit
\texorpdfstring{$W_3$}{W3} OPE; \ClaimStatusProvedHere]
\label{prop:winfty-ds-stage3-explicit-packet}
Fix $N=3$ and write
\[
T:=W^{(2)},
\qquad
W:=W^{(3)}.
\]
For the principal stage $W_3$, the primary residue coefficients on the
seed set $\mathcal{I}_3$ are completely determined by the explicit
$W_3$ OPE package (Definition~\ref{def:w3-algebra};
Theorem~\ref{thm:w-w-ope-complete} for the complete $W$--$W$
expansion).
More precisely, among the fifteen tuples in $\mathcal{I}_3$, only the
three primary coefficients
\[
\mathsf{C}^{\mathrm{DS}}_{2,2;2;0,2}(3)=2,
\qquad
\mathsf{C}^{\mathrm{DS}}_{2,3;3;0,2}(3)=3,
\qquad
\mathsf{C}^{\mathrm{DS}}_{3,3;2;0,4}(3)=2
\]
are nonzero; the remaining twelve coefficients in $\mathcal{I}_3$
vanish.  Consequently the stage-$3$ comparison problem is equivalent to
matching these three numbers and the twelve forced vanishings, after
which translation recovers all descendant coefficients.
\end{proposition}

\begin{proof}
The seed set $\mathcal{I}_3$ was written explicitly in
Corollary~\ref{cor:winfty-ds-lowstage-seeds}.  We read its coefficients
off from the explicit $W_3$ OPE package.

For the pair $(T,T)$, the only primary generator appearing in the
Virasoro OPE is $T$ itself at pole order $2$, with coefficient $2$.
The pole-order-$1$ term is $\partial T$, hence a descendant, and there
is no primary target of spin $3$.  Therefore
\[
\mathsf{C}^{\mathrm{DS}}_{2,2;2;0,2}(3)=2,
\qquad
\mathsf{C}^{\mathrm{DS}}_{2,2;2;0,1}(3)=
\mathsf{C}^{\mathrm{DS}}_{2,2;3;0,1}(3)=0.
\]

For the mixed pair $(T,W)$, the primary-field OPE shows that the only
primary target is $W$ at pole order $2$, with coefficient $3$.  The
pole-order-$1$ term is $\partial W$, hence a descendant, and no primary
target of spin $2$ appears.  Thus
\[
\mathsf{C}^{\mathrm{DS}}_{2,3;3;0,2}(3)=3,
\]
and all four remaining mixed primary coefficients in $\mathcal{I}_3$
vanish.

For the pair $(W,W)$, the complete $W$-$W$ OPE has primary generator
target $T$ only at pole order $4$, again with coefficient $2$.  The
terms at pole orders $3,2,1$ are the descendants
$\partial T$, $\partial^2 T$, and $\partial^3 T$, while the
$\Lambda$-terms lie in conformal weight $4$ and therefore outside the
generator seed set $\mathcal{I}_3$.  No primary target of spin $3$
appears.  Hence
\[
\mathsf{C}^{\mathrm{DS}}_{3,3;2;0,4}(3)=2,
\]
and the other seven $(W,W)$-entries in $\mathcal{I}_3$ vanish.

This accounts for all fifteen tuples in $\mathcal{I}_3$.  The final
sentence is then exactly the translation-propagation statement of
Proposition~\ref{prop:winfty-ds-generator-seed}.
\end{proof}

\begin{proposition}[Stage-\texorpdfstring{$4$}{4} residual packet after the theorematic
\texorpdfstring{$W_3$}{W3} sector; \ClaimStatusProvedHere]
\label{prop:winfty-ds-stage4-residual-packet}
Fix $N=4$ and let $\mathcal{I}_4$ be the primary seed set of
Corollary~\ref{cor:winfty-ds-finite-seed-set}.  Write
\[
\mathcal{I}_4^{\mathrm{Vir}}
:=\{(s,t,u,n)\in\mathcal{I}_4 \mid s=2\},
\qquad
\mathcal{I}_4^{W_3}
:=\{(3,3,u,n)\in\mathcal{I}_4 \mid u\le 3\}.
\]
Then every primary coefficient indexed by
\[
\mathcal{I}_4^{\mathrm{Vir}}\sqcup \mathcal{I}_4^{W_3}
\]
is already determined by the exact stress-tensor action on
$W^{(2)},W^{(3)},W^{(4)}$ together with the explicit stage-$3$
\texorpdfstring{$W_3$}{W3} packet of
Proposition~\ref{prop:winfty-ds-stage3-explicit-packet}.  The only
genuinely new stage-$4$ primary coefficients are therefore the residual
packet
\[
\mathcal{J}_4
:=\mathcal{I}_4\setminus
\bigl(\mathcal{I}_4^{\mathrm{Vir}}\sqcup \mathcal{I}_4^{W_3}\bigr),
\]
which has cardinality $29$ and decomposes as
\[
\mathcal{J}_4
=\{(3,3,4,1),(3,3,4,2)\}
\sqcup \mathcal{J}_4^{3,4}
\sqcup \mathcal{J}_4^{4,4},
\]
where
\begin{align*}
\mathcal{J}_4^{3,4}
={} & \{(3,4,2,1),(3,4,2,2),(3,4,2,3),(3,4,2,4),(3,4,2,5)\} \\
& \sqcup \{(3,4,3,1),(3,4,3,2),(3,4,3,3),(3,4,3,4)\} \\
& \sqcup \{(3,4,4,1),(3,4,4,2),(3,4,4,3)\},
\end{align*}
and
\begin{align*}
\mathcal{J}_4^{4,4}
={} & \{(4,4,2,1),(4,4,2,2),(4,4,2,3),(4,4,2,4),(4,4,2,5),(4,4,2,6)\} \\
& \sqcup \{(4,4,3,1),(4,4,3,2),(4,4,3,3),(4,4,3,4),(4,4,3,5)\} \\
& \sqcup \{(4,4,4,1),(4,4,4,2),(4,4,4,3),(4,4,4,4)\}.
\end{align*}
Consequently, the stage-$4$ comparison problem is equivalent to
matching these $29$ coefficients.  Once that packet agrees with the
principal Drinfeld--Sokolov coefficients, translation again recovers
all descendant coefficients.
\end{proposition}

\begin{proof}
All source-$2$ tuples lie in $\mathcal{I}_4^{\mathrm{Vir}}$ (stress-tensor action); the $(3,3)$ tuples with $u\le 3$ lie in $\mathcal{I}_4^{W_3}$ (stage-$3$ packet).  Enumerating the remaining admissible tuples for $(3,3,4)$, $(3,4)$, and $(4,4)$ gives the displayed decomposition with cardinality $2+12+15=29$.
\end{proof}

\begin{computation}[Stage-\texorpdfstring{$4$}{4} seed set and residual packet verification]
\label{comp:winfty-stage4-seed-verification}
\index{MC4!W-infinity stage-4 verification}
The seed set cardinalities and residual decomposition of
Proposition~\ref{prop:winfty-ds-stage4-residual-packet} have been
verified computationally
(\texttt{compute/\allowbreak lib/\allowbreak w4\_stage4\_coefficients.py},
$74$~tests):
\begin{itemize}
\item $|\mathcal{I}_4| = 54$,
  $|\mathcal{I}_4^{\mathrm{Vir}}| = 18$,
  $|\mathcal{I}_4^{W_3}| = 7$,
  $|\mathcal{J}_4| = 29$;
\item Residual decomposition: $(3,3)$-tail of size~$2$,
  $\mathcal{J}_4^{3,4}$ of size~$12$ (sub-decomposition $5+4+3$),
  $\mathcal{J}_4^{4,4}$ of size~$15$ (sub-decomposition $6+5+4$);
\item Stage-$3$ packet reproduced:
  $\mathsf{C}^{\mathrm{DS}}_{2,2;2;0,2} = 2$,
  $\mathsf{C}^{\mathrm{DS}}_{2,3;3;0,2} = 3$,
  $\mathsf{C}^{\mathrm{DS}}_{3,3;2;0,4} = 2$ (three nonzero,
  twelve vanishing);
\item Complementarity sum:
  $c(k) + c(-k-8) = 246$ (constant in~$k$);
\item Universal $T$-coupling pattern:
  $\mathsf{C}_{s,s;2;0,2s-2} = 2$ confirmed for $s = 2, 3$;
  predicted for $s = 4$.
\end{itemize}
\end{computation}

\begin{corollary}[Stage-\texorpdfstring{$4$}{4} top-pole packet after primaryity;
\ClaimStatusProvedHere]
\label{cor:winfty-ds-stage4-top-pole-packet}
Assume the hypotheses of
Proposition~\ref{prop:winfty-ds-stage4-residual-packet} and
Proposition~\ref{prop:winfty-ds-primary-top-pole}.  Then the residual
stage-$4$ packet $\mathcal{J}_4$ reduces further to the top-pole subset
\[
\mathcal{J}_4^{\mathrm{top}}
=\{(3,3,4,2),(3,4,2,5),(3,4,3,4),(3,4,4,3),
(4,4,2,6),(4,4,3,5),(4,4,4,4)\},
\]
which has cardinality $7$.  Equivalently, among the `29` residual
stage-$4$ primary coefficients, only these seven top-pole coefficients
can be nonzero; the remaining twenty-two coefficients in
$\mathcal{J}_4$ vanish by primaryity.

Consequently, the genuinely new principal stage-$4$ comparison problem
is equivalent to matching these seven coefficients with the
Drinfeld--Sokolov data.
\end{corollary}

\begin{proof}
Proposition~\ref{prop:winfty-ds-primary-top-pole} retains only the top pole $n=s+t-u$ for each triple, giving the seven displayed tuples; all lower-pole entries vanish.
\end{proof}

\begin{corollary}[Stage-\texorpdfstring{$4$}{4} parity-compressed packet;
\ClaimStatusProvedHere]
\label{cor:winfty-ds-stage4-parity-packet}
Assume the hypotheses of
Corollary~\ref{cor:winfty-ds-stage4-top-pole-packet} and
Proposition~\ref{prop:winfty-ds-self-ope-parity}.  Then the stage-$4$
top-pole packet reduces further to
\[
\mathcal{J}_4^{\mathrm{par}}
=\{(3,3,4,2),(3,4,2,5),(3,4,3,4),(3,4,4,3),
(4,4,2,6),(4,4,4,4)\},
\]
which has cardinality $6$.  Equivalently, among the `7` top-pole
coefficients of
Corollary~\ref{cor:winfty-ds-stage4-top-pole-packet}, the odd self-OPE
entry $(4,4,3,5)$ vanishes by skew-symmetry.  Hence the residual
stage-$4$ comparison problem is equivalent to matching these six
coefficients, together with the forced vanishing of the remaining `23`
entries in the original residual packet $\mathcal{J}_4$.
\end{corollary}

\begin{proof}
Among the seven top-pole entries, the only odd self-OPE top pole is $(4,4,3,5)$; Proposition~\ref{prop:winfty-ds-self-ope-parity} kills it, leaving six.
\end{proof}

\begin{corollary}[Stage-\texorpdfstring{$4$}{4} packet as three local OPE blocks;
\ClaimStatusProvedHere]
\label{cor:winfty-ds-stage4-ope-blocks}
Assume the hypotheses of
Corollary~\ref{cor:winfty-ds-stage4-parity-packet}.  Then the residual
principal stage-$4$ comparison problem is equivalent to matching the
six coefficients in the three top-pole OPE blocks
\begin{align*}
W^{(3)}(z)W^{(3)}(w)
&\sim
\frac{c_{334}\,W^{(4)}(w)}{(z-w)^2}
\quad\text{mod the theorematic $W_3$ sector and descendants},\\
W^{(3)}(z)W^{(4)}(w)
&\sim
\frac{c_{342}\,W^{(2)}(w)}{(z-w)^5}
+\frac{c_{343}\,W^{(3)}(w)}{(z-w)^4}
+\frac{c_{344}\,W^{(4)}(w)}{(z-w)^3}\\
&\qquad
\quad\text{mod descendants},\\
W^{(4)}(z)W^{(4)}(w)
&\sim
\frac{c_{442}\,W^{(2)}(w)}{(z-w)^6}
+\frac{c_{444}\,W^{(4)}(w)}{(z-w)^4}\\
&\qquad
\quad\text{mod the theorematic Virasoro sector and descendants}.
\end{align*}
More precisely,
\[
c_{334}=\mathsf{C}_{3,3;4;0,2},
\qquad
(c_{342},c_{343},c_{344})
\!=\!
\bigl(
\mathsf{C}_{3,4;2;0,5},
\mathsf{C}_{3,4;3;0,4},
\mathsf{C}_{3,4;4;0,3}
\bigr),
\]
and
\[
(c_{442},c_{444})
=(\mathsf{C}_{4,4;2;0,6},\mathsf{C}_{4,4;4;0,4}),
\]
where either the residue coefficients
$\mathsf{C}^{\mathrm{res}}$ or the principal Drinfeld--Sokolov
coefficients $\mathsf{C}^{\mathrm{DS}}$ may be used.  Equivalently, the
live stage-$4$ frontier is one new coefficient in the $(3,3)$ block,
three in the mixed $(3,4)$ block, and two in the even $(4,4)$ block.
\end{corollary}

\begin{proof}
Group the six parity-compressed tuples by source pair: $(3,3)$ contributes one, $(3,4)$ three, $(4,4)$ two.
\end{proof}

\begin{corollary}[Stage-\texorpdfstring{$4$}{4} frontier as one mixed block and three
self-coupling scalars; \ClaimStatusProvedHere]
\label{cor:winfty-ds-stage4-mixed-self-split}
Assume the hypotheses of
Corollary~\ref{cor:winfty-ds-stage4-ope-blocks}.  Then the residual
principal stage-$4$ comparison decomposes as:
\begin{enumerate}[label=\textup{(\roman*)}]
\item a self-coupling sector consisting of the three scalar top-pole
  coefficients
  \[
  c_{334},\qquad c_{442},\qquad c_{444};
  \]
\item a mixed higher-spin sector consisting of the single
  $W^{(3)}$-$W^{(4)}$ block
  \[
  (c_{342},c_{343},c_{344}).
  \]
\end{enumerate}
Equivalently, the only genuinely mixed stage-$4$ higher-spin data are
the three coefficients in the $W^{(3)}(z)W^{(4)}(w)$ block; the
remaining three live coefficients come from the self-couplings
$W^{(3)}$-$W^{(3)}$ and $W^{(4)}$-$W^{(4)}$.
\end{corollary}

\begin{proof}
Regroup: $(3,3)$ and $(4,4)$ give the self-coupling sector; $(3,4)$ is the mixed sector.
\end{proof}

\begin{proposition}[Mixed top-pole swap parity for even generators;
\ClaimStatusProvedHere]
\label{prop:winfty-ds-mixed-top-pole-swap}
Fix $N \ge 2$ and assume the hypotheses of
Proposition~\ref{prop:winfty-ds-primary-top-pole}.  Assume moreover
that the strong generators $W^{(2)},\dots,W^{(N)}$ are even fields.
For every visible mixed triple $(s,t,u)$ with $2 \le s < t \le N$ and
$u \le \min(N,s+t-1)$, write
\[
\mathsf{C}^{\mathrm{res}}_{s,t;u}(N)
:=
\mathsf{C}^{\mathrm{res}}_{s,t;u;0,\,s+t-u}(N),
\qquad
\mathsf{C}^{\mathrm{DS}}_{s,t;u}(N)
:=
\mathsf{C}^{\mathrm{DS}}_{s,t;u;0,\,s+t-u}(N)
\]
for the top-pole primary coefficients.  Then the reversed-order
coefficients satisfy
\[
\mathsf{C}^{\mathrm{res}}_{t,s;u}(N)
=
(-1)^{s+t-u}\mathsf{C}^{\mathrm{res}}_{s,t;u}(N),
\qquad
\mathsf{C}^{\mathrm{DS}}_{t,s;u}(N)
=
(-1)^{s+t-u}\mathsf{C}^{\mathrm{DS}}_{s,t;u}(N).
\]
Equivalently, for even primary generators the top-pole mixed
$W^{(u)}$-channel is swap-even when $s+t-u$ is even and swap-odd when
$s+t-u$ is odd.
\end{proposition}

\begin{proof}
The vertex-algebra skew-symmetry formula for even fields, combined with top-pole primaryity (Proposition~\ref{prop:winfty-ds-primary-top-pole}) to eliminate $j>0$ terms, gives $\mathsf{C}_{s,t;u}^{\mathrm{res}}=(-1)^{s+t-u}\mathsf{C}_{t,s;u}^{\mathrm{res}}$.  The same holds on the DS side.
\end{proof}

\begin{corollary}[Stage-\texorpdfstring{$4$}{4} mixed block split by swap parity;
\ClaimStatusProvedHere]
\label{cor:winfty-ds-stage4-mixed-swap-parity}
Assume the hypotheses of
Corollary~\ref{cor:winfty-ds-stage4-mixed-self-split} and
Proposition~\ref{prop:winfty-ds-mixed-top-pole-swap}.  Write
\[
W^{(4)}(z)W^{(3)}(w)
\sim
\frac{\widetilde c_{432}\,W^{(2)}(w)}{(z-w)^5}
+\frac{\widetilde c_{433}\,W^{(3)}(w)}{(z-w)^4}
+\frac{\widetilde c_{434}\,W^{(4)}(w)}{(z-w)^3}
\quad\text{mod descendants}
\]
for the reversed-order mixed block.  Then
\[
\widetilde c_{432}=-c_{342},
\qquad
\widetilde c_{433}=c_{343},
\qquad
\widetilde c_{434}=-c_{344}.
\]
Equivalently, the live mixed stage-$4$ frontier splits into one
swap-even channel $c_{343}$ and two swap-odd channels
$c_{342},c_{344}$.
\end{corollary}

\begin{proof}
Apply Proposition~\ref{prop:winfty-ds-mixed-top-pole-swap} to $(3,4,u)$ for $u=2,3,4$; the signs $(-1)^{7-u}$ give the displayed identities.
\end{proof}

\begin{proposition}[Formal mixed Virasoro-target vanishing under a
normalized two-point package; \ClaimStatusProvedHere]
\label{prop:winfty-formal-mixed-virasoro-zero}
Fix $N \ge 3$ and consider a stage-$N$ local OPE calculus with visible
primary generators $W^{(2)},\dots,W^{(N)}$, where
$T:=W^{(2)}$ is the Virasoro field.  Assume:
\begin{enumerate}[label=\textup{(\roman*)}]
\item each visible generator $W^{(u)}$ has conformal weight $u$, the
  local OPE is decomposed into conformal families of these primary
  generators and their translation descendants, and the Virasoro Ward
  identity holds for insertions of $T$ against each visible generator;
\item for every mixed pair $2 \le s < t \le N$, the vacuum two-point
  function vanishes:
  \[
  \langle W^{(s)}(z)W^{(t)}(0)\rangle=0;
  \]
\item the stress tensor is normalized by
  \[
  \langle T(z)T(0)\rangle=\frac{c/2}{z^4}
  \qquad\text{with } c\neq 0.
  \]
\end{enumerate}
Let
\[
\mathsf{C}_{s,t;u;0,n}(N)
\]
denote the undifferentiated primary coefficient of the target
$W^{(u)}$ at pole order $n$ in this local OPE calculus.  Then for every
visible mixed pair $2 \le s < t \le N$, the top-pole stress-tensor
coefficient vanishes:
\[
\mathsf{C}_{s,t;2;0,\,s+t-2}(N)=0.
\]
Equivalently, under this normalized two-point package the
undifferentiated $W^{(2)}=T$ channel is absent from the top pole of
every mixed OPE $W^{(s)}(z)W^{(t)}(w)$ with $s\ne t$.
\end{proposition}

\begin{proof}
Mixed-weight orthogonality (assumption~\textup{(ii)}) makes $\langle T(x)A(z)B(0)\rangle = 0$.  The top-pole $T$-channel contributes $\frac{c}{2}\mathsf{C}_{s,t;2;0,s+t-2}$ to the $x^{-4}z^{-(s+t-2)}$ coefficient; since $c\neq 0$, the coefficient vanishes.
\end{proof}

\begin{proposition}[Principal Drinfeld--Sokolov vanishing of the mixed
Virasoro target; \ClaimStatusProvedHere]
\label{prop:winfty-ds-mixed-virasoro-ds-zero}
Fix $N \ge 3$ and consider the principal Drinfeld--Sokolov stage
$\mathcal{W}_N$ at a generic level with nonzero central charge $c$.
Assume the Zamolodchikov normalization and mixed-weight orthogonality
package of Theorem~\ref{thm:wn-obstruction}.  Then for every visible
mixed pair $2 \le s < t \le N$, the top-pole stress-tensor coefficient
vanishes:
\[
\mathsf{C}^{\mathrm{DS}}_{s,t;2;0,\,s+t-2}(N)=0.
\]
Equivalently, on the principal Drinfeld--Sokolov side the
undifferentiated $W^{(2)}=T$ channel is absent from the top pole of
every mixed OPE $W^{(s)}(z)W^{(t)}(w)$ with $s \ne t$.
\end{proposition}

\begin{proof}
Apply Proposition~\ref{prop:winfty-formal-mixed-virasoro-zero} to the
principal Drinfeld--Sokolov stage.  Theorem~\ref{thm:wn-obstruction}
supplies the required normalized two-point package: the visible
generators are primary, the mixed two-point functions vanish for
distinct conformal weights, and
$\langle T(z)T(0)\rangle=\frac{c/2}{z^4}$ with $c\neq 0$ at generic
level.
\end{proof}

\begin{corollary}[Stage-\texorpdfstring{$4$}{4} mixed block as one vanishing channel and a
parity pair; \ClaimStatusProvedHere]
\label{cor:winfty-ds-stage4-mixed-two-channel}
Assume the hypotheses of
Corollary~\ref{cor:winfty-ds-stage4-mixed-swap-parity} and
Proposition~\ref{prop:winfty-ds-mixed-virasoro-ds-zero}.  Then on the
principal Drinfeld--Sokolov side one has
\[
\mathsf{C}^{\mathrm{DS}}_{3,4;2;0,5}=0,
\qquad
\mathsf{C}^{\mathrm{DS}}_{4,3;2;0,5}=0.
\]
Consequently, the stage-$4$ mixed comparison problem reduces to:
\begin{enumerate}[label=\textup{(\roman*)}]
\item proving the corresponding residue-side vanishing in the
  $W^{(2)}$ target channel;
\item matching the remaining swap-even coefficient
  $\mathsf{C}_{3,4;3;0,4}$;
\item matching the remaining swap-odd coefficient
  $\mathsf{C}_{3,4;4;0,3}$.
\end{enumerate}
Equivalently, the principal mixed stage-$4$ data consist of one
theorematically zero $W^{(2)}$ target channel together with a
swap-even / swap-odd pair targeting $W^{(3)}$ and $W^{(4)}$.
\end{corollary}

\begin{proof}
Proposition~\ref{prop:winfty-ds-mixed-virasoro-ds-zero} gives the first vanishing; swap parity (Corollary~\ref{cor:winfty-ds-stage4-mixed-swap-parity}) gives the second.
\end{proof}

\begin{proposition}[Formal self-coupling stress-tensor coefficient
under a normalized two-point package; \ClaimStatusProvedHere]
\label{prop:winfty-formal-self-t-coefficient}
Fix $N \ge 2$ and consider a stage-$N$ local OPE calculus with visible
primary generators $W^{(2)},\dots,W^{(N)}$, where
$T:=W^{(2)}$ is the Virasoro field.  Assume:
\begin{enumerate}[label=\textup{(\roman*)}]
\item each visible generator $W^{(u)}$ has conformal weight $u$, the
  local OPE is decomposed into conformal families of these primary
  generators and their translation descendants, and the Virasoro Ward
  identity holds for insertions of $T$ against each visible generator;
\item for every visible generator $W^{(s)}$ one has the normalized
  two-point function
  \[
  \langle W^{(s)}(z)W^{(s)}(0)\rangle=\frac{c/s}{z^{2s}}
  \qquad\text{with } c\neq 0;
  \]
\item the stress tensor is normalized by
  \[
  \langle T(z)T(0)\rangle=\frac{c/2}{z^4}.
  \]
\end{enumerate}
Let
\[
\mathsf{C}_{s,s;u;0,n}(N)
\]
denote the undifferentiated primary coefficient of the target
$W^{(u)}$ at pole order $n$ in this local OPE calculus.  Then for every
visible generator $W^{(s)}$ with $2 \le s \le N$, the top-pole
stress-tensor coefficient in the self-OPE is
\[
\mathsf{C}_{s,s;2;0,\,2s-2}(N)=2.
\]
Equivalently, under this normalized two-point package the
undifferentiated $T=W^{(2)}$ channel appears in every self-OPE
$W^{(s)}(z)W^{(s)}(w)$ with the universal coefficient $2$.
\end{proposition}

\begin{proof}
Write $N_s:=c/s$.  The Virasoro Ward identity for $\langle T(x)A(z)A(0)\rangle$ with $A:=W^{(s)}$ has $x^{-4}z^{-2s+2}$ coefficient equal to $sN_s=c$.  The only undifferentiated $T$-contribution at the top pole $2s-2$ gives $\mathsf{C}_{s,s;2;0,2s-2}\cdot c/2$.  Matching and dividing by $c\neq 0$ yields $\mathsf{C}_{s,s;2;0,2s-2}=2$.
\end{proof}

\begin{proposition}[Formal converse: the universal self-coupling
\texorpdfstring{$T$}{T}-coefficient forces the normalized two-point
function; \ClaimStatusProvedHere]
\label{prop:winfty-formal-self-normalization-from-t}
Fix $N \ge 2$ and consider a stage-$N$ local OPE calculus with visible
primary generators $W^{(2)},\dots,W^{(N)}$, where
$T:=W^{(2)}$ is the Virasoro field.  Assume:
\begin{enumerate}[label=\textup{(\roman*)}]
\item each visible generator $W^{(u)}$ has conformal weight $u$, the
  local OPE is decomposed into conformal families of these primary
  generators and their translation descendants, and the Virasoro Ward
  identity holds for insertions of $T$ against each visible generator;
\item the stress tensor is normalized by
  \[
  \langle T(z)T(0)\rangle=\frac{c/2}{z^4}
  \qquad\text{with } c\neq 0;
  \]
\item for some visible generator $W^{(s)}$ with $2 \le s \le N$, the
  top-pole stress-tensor coefficient in the self-OPE is
  \[
  \mathsf{C}_{s,s;2;0,\,2s-2}(N)=2.
  \]
\end{enumerate}
Then the self-pairing of $W^{(s)}$ is normalized by
\[
\langle W^{(s)}(z)W^{(s)}(0)\rangle=\frac{c/s}{z^{2s}}.
\]
Equivalently, under the visible Virasoro package the universal
self-coupling coefficient `2' and the Zamolodchikov normalization of
$W^{(s)}$ are equivalent.
\end{proposition}

\begin{proof}
Global conformal invariance forces $\langle A(z)A(0)\rangle = N_s z^{-2s}$.  The Ward-identity computation of Proposition~\ref{prop:winfty-formal-self-t-coefficient} gives $x^{-4}z^{-2s+2}$ coefficient $sN_s$ on one side and $\mathsf{C}_{s,s;2;0,2s-2}\cdot c/2 = c$ on the other.  Matching gives $N_s = c/s$.
\end{proof}

\begin{proposition}[Principal Drinfeld--Sokolov self-coupling
stress-tensor coefficient; \ClaimStatusProvedHere]
\label{prop:winfty-ds-self-t-coefficient}
Fix $N \ge 2$ and consider the principal Drinfeld--Sokolov stage
$\mathcal{W}_N$ at a generic level with nonzero central charge $c$.
Assume the Zamolodchikov normalization package of
Theorem~\ref{thm:wn-obstruction}.  Then for every visible generator
$W^{(s)}$ with $2 \le s \le N$, the top-pole stress-tensor coefficient
in the self-OPE is
\[
\mathsf{C}^{\mathrm{DS}}_{s,s;2;0,\,2s-2}(N)=2.
\]
Equivalently, on the principal Drinfeld--Sokolov side the
undifferentiated $T=W^{(2)}$ channel appears in every self-OPE
$W^{(s)}(z)W^{(s)}(w)$ with the universal coefficient $2$.
\end{proposition}

\begin{proof}
Apply Proposition~\ref{prop:winfty-formal-self-t-coefficient} to the
principal Drinfeld--Sokolov stage.  Theorem~\ref{thm:wn-obstruction}
supplies the required normalized two-point package, with
\[
\langle W^{(s)}(z)W^{(s)}(0)\rangle=\frac{c/s}{z^{2s}}
\qquad\text{and}\qquad
\langle T(z)T(0)\rangle=\frac{c/2}{z^4}
\]
at generic level.
\end{proof}

\begin{corollary}[Principal stage-\texorpdfstring{$4$}{4} self-coupling
\texorpdfstring{$W^{(4)}$-$W^{(4)}\to T$}{W4-W4 to T}
normalization; \ClaimStatusProvedHere]
\label{cor:winfty-ds-stage4-self-t-normalization}
In the setting of Proposition~\ref{prop:winfty-ds-self-t-coefficient},
one has
\[
\mathsf{C}^{\mathrm{DS}}_{4,4;2;0,6}=2.
\]
Equivalently, the principal stage-$4$ self-coupling
$W^{(4)}(z)W^{(4)}(w)\to T(w)$ is not free data: its top-pole
coefficient is universally fixed to `2`.
\end{corollary}

\begin{proof}
Apply Proposition~\ref{prop:winfty-ds-self-t-coefficient} with $s=4$.
\end{proof}

\begin{corollary}[Stage-\texorpdfstring{$4$}{4} principal target packet after theorematic
Virasoro-target elimination; \ClaimStatusProvedHere]
\label{cor:winfty-ds-stage4-five-plus-zero}
Assume the hypotheses of
Corollary~\ref{cor:winfty-ds-stage4-mixed-two-channel} and
Corollary~\ref{cor:winfty-ds-stage4-self-t-normalization}.  Then the
residual principal stage-$4$ target packet decomposes as:
\begin{enumerate}[label=\textup{(\roman*)}]
\item two higher-spin self-coupling targets
  \[
  c_{334},\qquad c_{444};
  \]
\item the mixed swap-even higher-spin target
  \[
  \mathsf{C}_{3,4;3;0,4};
  \]
\item the mixed swap-odd higher-spin target
  \[
  \mathsf{C}_{3,4;4;0,3};
  \]
\item the theorematic principal Virasoro-target value
  \[
  \mathsf{C}^{\mathrm{DS}}_{4,4;2;0,6}=2;
  \]
\item the theorematic principal Virasoro-target value
  \[
  \mathsf{C}^{\mathrm{DS}}_{3,4;2;0,5}=0.
  \]
\end{enumerate}
Equivalently, the standard stage-$4$ identity packet on~$\mathcal{I}_4$
splits into four higher-spin targets
\[
c_{334},\qquad c_{444},\qquad
\mathsf{C}_{3,4;3;0,4},\qquad \mathsf{C}_{3,4;4;0,3},
\]
together with two Virasoro-target channels whose principal values are
theorematically fixed to
\[
\mathsf{C}^{\mathrm{DS}}_{4,4;2;0,6}=2,
\qquad
\mathsf{C}^{\mathrm{DS}}_{3,4;2;0,5}=0,
\]
and whose residue-side contraction to four channels is the additional
package isolated in
Corollary~\ref{cor:winfty-stage4-residue-four-channel}.
\end{corollary}

\begin{proof}
Combine the mixed-self split (Corollary~\ref{cor:winfty-ds-stage4-mixed-self-split}), the mixed two-channel reduction (Corollary~\ref{cor:winfty-ds-stage4-mixed-two-channel}), and the universal $T$-coupling normalization (Corollary~\ref{cor:winfty-ds-stage4-self-t-normalization}).
\end{proof}

\begin{proposition}[Exact MC4 frontier packet for the standard
\texorpdfstring{$W_\infty$}{W_infty} tower; \ClaimStatusProvedHere]
\label{prop:winfty-mc4-frontier-package}
Let $\mathcal{W}^{\mathrm{ht}}$ be a separated complete H-level target
for $W_\infty$ satisfying the quotient-system hypotheses of
Corollary~\ref{cor:winfty-hlevel-comparison-criterion}.  Assume
moreover that the quotient residue calculus is translation compatible
and that the strong generators are even Virasoro-primary fields, so
that
Proposition~\ref{prop:winfty-ds-generator-seed},
Corollary~\ref{cor:winfty-ds-finite-seed-set},
Proposition~\ref{prop:winfty-ds-stage3-explicit-packet}, and
Corollary~\ref{cor:winfty-ds-stage4-five-plus-zero} apply.  Then:
\begin{enumerate}
\item for every finite stage $N$, the remaining H-level comparison with
  the principal Drinfeld--Sokolov stage $W_N$ reduces to the finite
  primary seed packet $\mathcal{I}_N$ of residue identities
  \[
  \mathsf{C}^{\mathrm{res}}_{s,t;u;0,n}(N)
  =\mathsf{C}^{\mathrm{DS}}_{s,t;u;0,n}(N);
  \]
\item at stage~$3$, that packet is already theorematically fixed by the
  explicit $W_3$ OPE: three nonzero primary coefficients and twelve
  forced vanishing statements;
\item at stage~$4$, after removing the theorematic Virasoro and
  $W_3$ sectors, the first genuinely open higher-spin packet is exactly
  the exact six-entry identity packet on~$\mathcal{I}_4$
  \[
  \mathsf{C}^{\mathrm{res}}_{s,t;u;m,n}(4)
  =
  \mathsf{C}^{\mathrm{DS}}_{s,t;u;m,n}(4)
  \]
  in the channels
  \[
  \begin{gathered}
  (3,3;4;0,2),\qquad
  (4,4;4;0,4),\\
  (3,4;3;0,4),\qquad
  (3,4;4;0,3),\qquad
  (4,4;2;0,6),\\
  (3,4;2;0,5).
  \end{gathered}
  \]
  with the first four entries forming the residual higher-spin
  comparison packet, while the last two are Virasoro-target channels
  whose principal values are already fixed by
  Corollary~\ref{cor:winfty-ds-stage4-five-plus-zero}.
\end{enumerate}
In particular, once the H-level target and quotient system are
constructed, the standard $\mathcal{W}_\infty$ MC4 problem is not a
further stabilization theorem.  It is the finite stagewise
residue-identification problem on the packets $\mathcal{I}_N$, with the
first genuinely open higher-spin packet occurring at stage~$4$ in the
displayed six-entry identity form on~$\mathcal{I}_4$.
\end{proposition}

\begin{proof}
\textup{(1)}: Corollary~\ref{cor:winfty-hlevel-comparison-criterion} plus Proposition~\ref{prop:winfty-ds-generator-seed} and Corollary~\ref{cor:winfty-ds-finite-seed-set}.  \textup{(2)}: Proposition~\ref{prop:winfty-ds-stage3-explicit-packet}.  \textup{(3)}: Corollary~\ref{cor:winfty-ds-stage4-five-plus-zero}.
\end{proof}

\begin{corollary}[Minimal closure criterion for the standard
\texorpdfstring{$W_\infty$}{W_infty} MC4 frontier;
\ClaimStatusProvedHere]
\label{cor:winfty-stage4-closure-criterion}
Under the hypotheses of
Proposition~\ref{prop:winfty-mc4-frontier-package}, the standard
\texorpdfstring{$W_\infty$}{W_infty} H-level comparison closes once the
stagewise packet identities on~$\mathcal{I}_N$ are proved.  The first
genuinely open higher-spin packet is the exact six-entry stage-$4$
identity packet on~$\mathcal{I}_4$.  Within that packet, the residual
higher-spin channels still to be matched are
\[
\begin{gathered}
(3,3;4;0,2),\qquad
(4,4;4;0,4),\qquad
(3,4;3;0,4),\qquad
(3,4;4;0,3),
\end{gathered}
\]
while the corresponding principal Virasoro-target values are already
fixed to
\[
\mathsf{C}^{\mathrm{DS}}_{4,4;2;0,6}=2,
\qquad
\mathsf{C}^{\mathrm{DS}}_{3,4;2;0,5}=0.
\]
\end{corollary}

\begin{proof}
Corollary~\ref{cor:winfty-hlevel-comparison-criterion} upgrades
stagewise packet matching on~$\mathcal{I}_N$ to the H-level
comparison, and
Proposition~\ref{prop:winfty-mc4-frontier-package} identifies the first
open stage as the six-entry block on~$\mathcal{I}_4$.
Corollary~\ref{cor:winfty-ds-stage4-five-plus-zero} fixes the two
principal Virasoro-target values, leaving the four displayed
higher-spin channels as the residual higher-spin packet.
\end{proof}

\begin{corollary}[Constructing the completed chiral Koszul-dual
candidate for \texorpdfstring{$W_\infty$}{W_infty};
\ClaimStatusProvedHere]
\label{cor:winfty-dual-candidate-construction}
Assume the hypotheses of
Proposition~\ref{prop:winfty-mc4-frontier-package}.  Then the completed
bar object of the corresponding H-level target
$\mathcal{W}^{\mathrm{ht}}$ is quasi-isomorphic to the inverse-limit
bar object of the principal stages:
\[
\widehat{\bar B}\bigl(\mathcal{W}^{\mathrm{ht}}\bigr)
\xrightarrow{\sim}
\varprojlim_N \bar B(W_N).
\]
Its completed cobar comparison recovers
$\mathcal{W}^{\mathrm{ht}}$ up to quasi-isomorphism, and any two
principal-stage compatible H-level targets satisfying the same
stagewise packet identities determine quasi-isomorphic completed
bar-cobar objects.

Consequently, the foundational route to the chiral Koszul-dual object
of \texorpdfstring{$W_\infty$}{W_infty} is exactly:
\begin{enumerate}
\item construct a principal-stage compatible H-level or factorization
  target;
\item prove the finite packet identities on~$\mathcal{I}_N$; and
\item invoke
  Corollary~\ref{cor:winfty-hlevel-comparison-criterion} to lift the
  standard principal-stage completed bar object to that target.
\end{enumerate}
No separate infinite-generator quadratic presentation is required at
this stage.
\end{corollary}

\begin{proof}
Under the stated hypotheses,
Proposition~\ref{prop:winfty-mc4-frontier-package} identifies the
finite-stage quotient problem with the packet identities on
$\mathcal{I}_N$, while
Corollary~\ref{cor:winfty-hlevel-comparison-criterion} identifies
$\mathcal{W}^{\mathrm{ht}}$ with the standard principal-stage completed
M-level package.  The bar-cobar comparison part of that corollary
therefore gives the displayed quasi-isomorphism and the cobar recovery.
The same comparison shows that any two such H-level targets define the
same completed bar-cobar object up to quasi-isomorphism.
\end{proof}

\begin{corollary}[Stage-\texorpdfstring{$4$}{4} \texorpdfstring{$W_\infty$}{W_infty}
frontier on the Ward-normalized H-level locus;
\ClaimStatusProvedHere]
\label{cor:winfty-stage4-residue-four-channel}
Assume the hypotheses of
Proposition~\ref{prop:winfty-mc4-frontier-package}.  Assume moreover
that $\mathcal{W}^{\mathrm{ht}}$ is stage-$4$ Ward-normalized in the
sense of Definition~\ref{def:winfty-stage4-ward-normalized}.  Then the
two
Virasoro-target residue channels are theorematically fixed:
\[
\mathsf{C}^{\mathrm{res}}_{3,4;2;0,5}(4)=0,
\qquad
\mathsf{C}^{\mathrm{res}}_{4,4;2;0,6}(4)=2.
\]
Consequently the first genuinely open stage-$4$
\texorpdfstring{$W_\infty$}{W_infty} packet contracts from the six-entry
identity packet of Proposition~\ref{prop:winfty-mc4-frontier-package} to
the four-channel higher-spin identity packet
\[
\mathsf{C}^{\mathrm{res}}_{s,t;u;0,n}(4)
=
\mathsf{C}^{\mathrm{DS}}_{s,t;u;0,n}(4)
\]
in the channels
\[
\begin{gathered}
(3,3;4;0,2),\qquad
(4,4;4;0,4),\\
(3,4;3;0,4),\qquad
(3,4;4;0,3).
\end{gathered}
\]
\end{corollary}

\begin{proof}
The Ward-normalized input supplies the hypotheses of Propositions~\ref{prop:winfty-formal-mixed-virasoro-zero} and~\ref{prop:winfty-formal-self-t-coefficient}, fixing the two Virasoro-target channels; removing them leaves four higher-spin identities.
\end{proof}

\begin{proposition}[Exact missing input for the unconditional
\texorpdfstring{$W_\infty$}{W_infty} stage-\texorpdfstring{$4$}{4} contraction;
\ClaimStatusProvedHere]
\label{prop:winfty-stage4-visible-pairing-gap}
Assume the hypotheses of
Proposition~\ref{prop:winfty-mc4-frontier-package}.  Assume moreover
that the stage-$4$ quotient residue calculus on the visible generators
$W^{(2)},W^{(3)},W^{(4)}$ is decomposed into conformal families of
Virasoro-primary generators and that the Virasoro field $T:=W^{(2)}$
acts on them, including on itself, by the standard stage-$4$
Virasoro Ward identity.  Then the only additional input from the
normalized residue two-point/Ward package actually used to deduce
\[
\mathsf{C}^{\mathrm{res}}_{3,4;2;0,5}(4)=0,
\qquad
\mathsf{C}^{\mathrm{res}}_{4,4;2;0,6}(4)=2
\]
is the visible weight-$4$ self-pairing normalization
\[
\langle W^{(4)}(z)W^{(4)}(0)\rangle=\frac{c/4}{z^{8}},
\]
with the same nonzero central charge $c$ as on the principal side.
In particular, once the visible Virasoro Ward action is fixed, the open
content of Conjecture~\ref{conj:winfty-stage4-ward-inheritance} is
exactly inheritance of this weight-$4$ normalization, equivalently the
single scalar identity
\[
\mathsf{C}^{\mathrm{res}}_{4,4;2;0,6}(4)=2.
\]
\end{proposition}

\begin{proof}
Mixed-weight orthogonality and $W^{(3)}$ normalization are already forced (Propositions~\ref{prop:winfty-stage4-visible-orthogonality} and~\ref{prop:winfty-stage4-visible-w3-normalization}).  The mixed vanishing follows from Proposition~\ref{prop:winfty-formal-mixed-virasoro-zero}; the self-coupling identity reduces via Proposition~\ref{prop:winfty-formal-self-t-coefficient} to the weight-$4$ normalization.
\end{proof}

\begin{definition}[Stage-\texorpdfstring{$4$}{4} Ward-normalized standard
\texorpdfstring{$W_\infty$}{W_infty} H-level target]
\label{def:winfty-stage4-ward-normalized}
Assume the setting of
Proposition~\ref{prop:winfty-mc4-frontier-package}.  We say that the
H-level target $\mathcal{W}^{\mathrm{ht}}$ is
\emph{stage-$4$ Ward-normalized} if its induced stage-$4$ quotient
residue calculus satisfies the visible Virasoro Ward action together
with the visible weight-$4$ self-pairing normalization
\[
\langle W^{(4)}(z)W^{(4)}(0)\rangle=\frac{c/4}{z^{8}}
\]
with the same nonzero central charge $c$ as on the principal
Drinfeld--Sokolov stage.  Equivalently, by
Proposition~\ref{prop:winfty-stage4-visible-orthogonality}, the
visible generators are mixed-weight orthogonal, and by
Proposition~\ref{prop:winfty-stage4-visible-w3-normalization} the
visible $W^{(3)}$ self-pairing is already normalized as well, so the
full visible package required by
Proposition~\ref{prop:winfty-stage4-visible-pairing-gap} holds.
\end{definition}

\begin{remark}[What is still missing for an unconditional stage-\texorpdfstring{$4$}{4}
four-channel reduction]
\label{rem:winfty-stage4-four-channel-gap}
Corollary~\ref{cor:winfty-stage4-residue-four-channel} is conditional
on the additional H-level input of
Definition~\ref{def:winfty-stage4-ward-normalized}.  The
quotient-system hypotheses of
Corollary~\ref{cor:winfty-hlevel-comparison-criterion} and
Proposition~\ref{prop:winfty-mc4-frontier-package} supply compatible
finite quotients, translation compatibility, and Virasoro-primary
strong generators, but they do not by themselves identify the
visible weight-$4$ self-pairing inside that quotient residue calculus.
Mixed-weight orthogonality is already forced by the Virasoro Ward
identity, by
Proposition~\ref{prop:winfty-stage4-visible-orthogonality}.  Thus the
first unconditional standard-tower
MC4 packet remains the six-entry block of
Proposition~\ref{prop:winfty-mc4-frontier-package}.  By
Proposition~\ref{prop:winfty-stage4-visible-pairing-gap}, once the
visible Virasoro Ward action is fixed the exact remaining missing input
is Conjecture~\ref{conj:winfty-stage4-visible-diagonal-normalization},
equivalently the single scalar identity
\[
\mathsf{C}^{\mathrm{res}}_{4,4;2;0,6}(4)=2.
\]
Its package form is
Conjecture~\ref{conj:winfty-stage4-ward-inheritance}.
\end{remark}

\begin{proposition}[Stage-\texorpdfstring{$4$}{4} visible mixed-weight orthogonality from
the Virasoro Ward identity; \ClaimStatusProvedHere]
\label{prop:winfty-stage4-visible-orthogonality}
Let $\mathcal{W}^{\mathrm{fact}}_\infty$ be a principal-stage
compatible factorization target in the sense of
Proposition~\ref{prop:winfty-factorization-package}, and let
$\mathcal{F}_{\le 4}$ be its stage-$4$ quotient.  Assume:
\begin{enumerate}[label=\textup{(\roman*)}]
\item the visible generators $W^{(2)},W^{(3)},W^{(4)}$ in
  $\mathcal{F}_{\le 4}$ are even Virasoro-primary fields, and the local
  OPE is decomposed into their conformal families and translation
  descendants;
\item the Virasoro field $T:=W^{(2)}$ acts on these visible generators
  by the standard stage-$4$ Virasoro Ward identity;
\end{enumerate}
Then the quotient vacuum functional is mixed-weight orthogonal on the
visible generators:
\[
\langle W^{(s)}(z)W^{(t)}(0)\rangle = 0
\qquad (s\neq t,\ 2\le s,t\le 4).
\]
\end{proposition}

\begin{proof}
The global conformal Ward identities force
$\langle W^{(s)}(z)W^{(t)}(w)\rangle = C_{s,t}(z-w)^{-(s+t)}$;
the special conformal equation gives $(s-t)C_{s,t}=0$,
so $C_{s,t}=0$ whenever $s\neq t$.
\end{proof}

\begin{proposition}[Stage-\texorpdfstring{$4$}{4} visible
\texorpdfstring{$W^{(3)}$}{W3} normalization from the theorematic
\texorpdfstring{$W_3$}{W3} packet; \ClaimStatusProvedHere]
\label{prop:winfty-stage4-visible-w3-normalization}
Assume the hypotheses of
Proposition~\ref{prop:winfty-mc4-frontier-package}.  Assume moreover
that the stage-$4$ quotient residue calculus on the visible generators
$W^{(2)},W^{(3)},W^{(4)}$ is decomposed into conformal families of
Virasoro-primary generators and that the Virasoro field $T:=W^{(2)}$
acts on them, including on itself, by the standard stage-$4$
Virasoro Ward identity.  Then the visible stage-$4$ coefficient
\[
\mathsf{C}^{\mathrm{res}}_{3,3;2;0,4}(4)=2
\]
is already theorematic, and hence
\[
\langle W^{(3)}(z)W^{(3)}(0)\rangle=\frac{c/3}{z^{6}}.
\]
\end{proposition}

\begin{proof}
The tuple $(3,3,2,4)\in\mathcal{I}_4^{W_3}$ is fixed by the theorematic stage-$3$ packet (Proposition~\ref{prop:winfty-ds-stage3-explicit-packet}): $\mathsf{C}^{\mathrm{res}}_{3,3;2;0,4}(4)=2$.  Proposition~\ref{prop:winfty-formal-self-normalization-from-t} with $s=3$ then gives the displayed normalization.
\end{proof}

\begin{conjecture}[Stage-\texorpdfstring{$4$}{4} visible weight-\texorpdfstring{$4$}{4} normalization]
\label{conj:winfty-stage4-visible-diagonal-normalization}
\ClaimStatusConjectured{}
Let $\mathcal{W}^{\mathrm{fact}}_\infty$ be a principal-stage
compatible factorization target in the sense of
Proposition~\ref{prop:winfty-factorization-package}, and let
$\mathcal{F}_{\le 4}$ be its stage-$4$ quotient.  Assume:
\begin{enumerate}[label=\textup{(\roman*)}]
\item the visible generators $W^{(2)},W^{(3)},W^{(4)}$ in
  $\mathcal{F}_{\le 4}$ are even Virasoro-primary fields, and the local
  OPE is decomposed into their conformal families and translation
  descendants;
\item the Virasoro field $T:=W^{(2)}$ acts on these visible generators
  by the standard stage-$4$ Virasoro Ward identity, including on
  itself.
\end{enumerate}
Then the remaining visible self-pairing is normalized by
\[
\langle W^{(4)}(z)W^{(4)}(0)\rangle = \frac{c/4}{z^{8}}
\]
with the same nonzero central charge $c$ as on the principal
Drinfeld--Sokolov stage.  Equivalently,
\[
\mathsf{C}^{\mathrm{res}}_{4,4;2;0,6}(4)=2.
\]
\end{conjecture}

\begin{conjecture}[Automatic stage-\texorpdfstring{$4$}{4} Ward normalization from the
visible Virasoro package]
\label{conj:winfty-stage4-ward-inheritance}
\ClaimStatusConjectured{}
Let $\mathcal{W}^{\mathrm{fact}}_\infty$ be a principal-stage
compatible factorization target in the sense of
Proposition~\ref{prop:winfty-factorization-package}, and let
$\mathcal{F}_{\le 4}$ be its stage-$4$ quotient.  Assume:
\begin{enumerate}[label=\textup{(\roman*)}]
\item the visible generators $W^{(2)},W^{(3)},W^{(4)}$ in
  $\mathcal{F}_{\le 4}$ are even Virasoro-primary fields, and the local
  OPE is decomposed into their conformal families and translation
  descendants;
\item the Virasoro field $T:=W^{(2)}$ acts on these visible generators
  by the standard stage-$4$ Virasoro Ward identity, including on
  itself;
\end{enumerate}
Then the induced H-level target is stage-$4$ Ward-normalized in the
sense of Definition~\ref{def:winfty-stage4-ward-normalized}.  Consequently
Corollary~\ref{cor:winfty-stage4-residue-four-channel} applies, so the
first live standard-tower stage-$4$ packet contracts from the
unconditional six-entry block of
Proposition~\ref{prop:winfty-mc4-frontier-package} to the four
higher-spin channels
\[
(3,3;4;0,2),\qquad
(4,4;4;0,4),\qquad
(3,4;3;0,4),\qquad
(3,4;4;0,3).
\]
\end{conjecture}

\begin{corollary}[Equivalent exact forms of the remaining
\texorpdfstring{$W_\infty$}{W_infty} stage-\texorpdfstring{$4$}{4} input;
\ClaimStatusProvedHere]
\label{cor:winfty-stage4-single-scalar-equivalent}
Assume the hypotheses of
Proposition~\ref{prop:winfty-mc4-frontier-package}.  Assume moreover
that the stage-$4$ quotient residue calculus on the visible generators
$W^{(2)},W^{(3)},W^{(4)}$ is decomposed into conformal families of
Virasoro-primary generators and that the Virasoro field $T:=W^{(2)}$
acts on these visible generators by the standard stage-$4$ Virasoro
Ward identity, including on itself.  Then the following are
equivalent:
\begin{enumerate}[label=\textup{(\roman*)}]
\item the induced H-level target is stage-$4$ Ward-normalized in the
  sense of Definition~\ref{def:winfty-stage4-ward-normalized};
\item the visible weight-$4$ self-pairing is normalized by
  \[
  \langle W^{(4)}(z)W^{(4)}(0)\rangle = \frac{c/4}{z^{8}};
  \]
\item the single scalar residue identity
  \[
  \mathsf{C}^{\mathrm{res}}_{4,4;2;0,6}(4)=2
  \]
  holds.
\end{enumerate}
Under these equivalent conditions,
Corollary~\ref{cor:winfty-stage4-residue-four-channel} applies, so the
first live standard-tower stage-$4$ packet contracts from the
unconditional six-entry block to the four higher-spin channels
\[
(3,3;4;0,2),\qquad
(4,4;4;0,4),\qquad
(3,4;3;0,4),\qquad
(3,4;4;0,3).
\]
\end{corollary}

\begin{proof}
The equivalence \textup{(i)} \(\Longleftrightarrow\) \textup{(ii)} is
exactly Definition~\ref{def:winfty-stage4-ward-normalized}.  Under the
same visible Virasoro package,
Proposition~\ref{prop:winfty-formal-self-t-coefficient} with $s=4$
gives \textup{(ii)} \(\Longrightarrow\) \textup{(iii)}, and
Proposition~\ref{prop:winfty-formal-self-normalization-from-t} with
$s=4$ gives \textup{(iii)} \(\Longrightarrow\) \textup{(ii)}.
The final sentence is exactly
Corollary~\ref{cor:winfty-stage4-residue-four-channel}.
\end{proof}

\begin{proposition}[Stage-\texorpdfstring{$4$}{4} swap-even residue channel from a visible
invariant pairing; \ClaimStatusProvedHere]
\label{prop:winfty-stage4-residue-pairing-reduction}
Assume the hypotheses of
Corollary~\ref{cor:winfty-stage4-single-scalar-equivalent}.  Assume
moreover that the visible stage-$4$ quotient residue calculus carries a
bilinear pairing \((-, -)_{\mathrm{vis}}\) on the span of
\(W^{(2)},W^{(3)},W^{(4)}\) such that:
\begin{enumerate}[label=\textup{(\roman*)}]
\item the visible generators are pairwise orthogonal for distinct
  conformal weights, and the self-pairings are normalized by
  \[
  (W^{(2)},W^{(2)})_{\mathrm{vis}}=\frac{c}{2},\qquad
  (W^{(3)},W^{(3)})_{\mathrm{vis}}=\frac{c}{3},\qquad
  (W^{(4)},W^{(4)})_{\mathrm{vis}}=\frac{c}{4}
  \]
  with the same nonzero central charge~\(c\);
\item for the visible primary field \(W^{(3)}\), the pairing satisfies
  the standard primary invariant-pairing identity
  \[
  \bigl(W^{(3)}_{(1)}a,b\bigr)_{\mathrm{vis}}
  =
  -\bigl(a,W^{(3)}_{(3)}b\bigr)_{\mathrm{vis}}
  \]
  whenever \(a,b\) are visible generators.
\end{enumerate}
Then the swap-even residue channel is forced by the \((3,3)\)-block:
\[
\mathsf{C}^{\mathrm{res}}_{3,4;3;0,4}(4)
=
-\frac{3}{4}\,
\mathsf{C}^{\mathrm{res}}_{3,3;4;0,2}(4).
\]
Equivalently,
\[
\bigl(\mathsf{C}^{\mathrm{res}}_{3,4;3;0,4}(4)\bigr)^2
=
\frac{9}{16}\,
\bigl(\mathsf{C}^{\mathrm{res}}_{3,3;4;0,2}(4)\bigr)^2.
\]
\end{proposition}

\begin{proof}
Apply the $W^{(3)}$ invariant-pairing identity with $a=W^{(3)}$, $b=W^{(4)}$: mode $1$ produces $\mathsf{C}^{\mathrm{res}}_{3,3;4;0,2}\cdot c/4$, the complementary mode $3$ produces $\mathsf{C}^{\mathrm{res}}_{3,4;3;0,4}\cdot c/3$, and the sign is negative ($W^{(3)}$ odd weight).  Since $c\neq 0$, rearranging gives the displayed identity.
\end{proof}

\begin{corollary}[Stage-\texorpdfstring{$4$}{4} residue packet as three higher-spin
channels on the visible pairing locus; \ClaimStatusProvedHere]
\label{cor:winfty-stage4-residue-three-channel}
Assume the hypotheses of
Proposition~\ref{prop:winfty-stage4-residue-pairing-reduction}.  Then,
on the stage-$4$ Ward-normalized visible pairing locus, the first
genuinely open stage-$4$ higher-spin residue packet contracts from the
four-channel block of
Corollary~\ref{cor:winfty-stage4-residue-four-channel} to the
three-channel block
\[
\begin{gathered}
(3,3;4;0,2),\qquad
(4,4;4;0,4),\\
(3,4;4;0,3),
\end{gathered}
\]
with the swap-even channel \((3,4;3;0,4)\) recovered formally from
\((3,3;4;0,2)\) by
Proposition~\ref{prop:winfty-stage4-residue-pairing-reduction}.
\end{corollary}

\begin{proof}
Proposition~\ref{prop:winfty-stage4-residue-pairing-reduction} forces $(3,4;3;0,4)$ from $(3,3;4;0,2)$; removing it leaves three channels.
\end{proof}

\begin{corollary}[Stage-\texorpdfstring{$4$}{4} higher-spin comparison as a
primitive-plus-transport square triple on the visible pairing locus;
\ClaimStatusProvedHere]
\label{cor:winfty-stage4-primitive-transport-square-triple}
Assume the hypotheses of
Proposition~\ref{prop:winfty-stage4-residue-pairing-reduction}.  Then,
on the stage-$4$ Ward-normalized visible pairing locus, the four
higher-spin identities of
Corollary~\ref{cor:winfty-stage4-residue-four-channel} are equivalent
at signless square-class level to the three identities
\[
\bigl(\mathsf{C}^{\mathrm{res}}_{3,3;4;0,2}(4)\bigr)^2
=
c_{334}^{\,2},
\qquad
\bigl(\mathsf{C}^{\mathrm{res}}_{4,4;4;0,4}(4)\bigr)^2
=
c_{444}^{\,2},
\]
\[
\bigl(\mathsf{C}^{\mathrm{res}}_{3,4;4;0,3}(4)\bigr)^2
=
\tfrac{5}{7}\,c_{334}^{\,2}.
\]
Equivalently, on that locus the swap-even square identity
\[
\bigl(\mathsf{C}^{\mathrm{res}}_{3,4;3;0,4}(4)\bigr)^2
=
\tfrac{9}{16}\,c_{334}^{\,2}
\]
is automatic once the \((3,3;4;0,2)\) square identity holds.  In
particular, the signless stage-$4$ higher-spin frontier on the visible
pairing locus consists of the two primitive self-coupling squares
\[
\bigl(\mathsf{C}^{\mathrm{res}}_{3,3;4;0,2}(4)\bigr)^2,
\qquad
\bigl(\mathsf{C}^{\mathrm{res}}_{4,4;4;0,4}(4)\bigr)^2
\]
together with the swap-odd \(W^{(4)}\)-transport square
\[
\bigl(\mathsf{C}^{\mathrm{res}}_{3,4;4;0,3}(4)\bigr)^2.
\]
\end{corollary}

\begin{proof}
On the visible pairing locus, Proposition~\ref{prop:winfty-stage4-residue-pairing-reduction} gives
$(\mathsf{C}^{\mathrm{res}}_{3,4;3;0,4})^2 = \tfrac{9}{16}(\mathsf{C}^{\mathrm{res}}_{3,3;4;0,2})^2$.
On the principal side,
Corollary~\ref{cor:w4-ds-stage4-square-class-reduction} gives
\[
\mathsf{C}_{3,4;3;0,4}^{\,2}
=
\tfrac{9}{16}\,c_{334}^{\,2},
\qquad
\mathsf{C}_{3,4;4;0,3}^{\,2}
=
\tfrac{5}{7}\,c_{334}^{\,2}.
\]
Therefore the swap-even square identity is automatic once
\(
\bigl(\mathsf{C}^{\mathrm{res}}_{3,3;4;0,2}(4)\bigr)^2=c_{334}^{\,2}
\)
holds, while the remaining two displayed identities are exactly the
\((4,4)\) self-coupling square identity and the swap-odd
\(W^{(4)}\)-transport square identity.  The converse implication is
immediate.
\end{proof}

\begin{remark}[Exact next square-class gap after the visible pairing
reduction]
\label{rem:winfty-stage4-primitive-transport-gap}
Corollary~\ref{cor:winfty-stage4-primitive-transport-square-triple}
shows that the visible pairing package brings the residue side to the
same square-class profile as the principal side up to one additional
transport relation: the principal DS packet still forces
\[
\mathsf{C}_{3,4;4;0,3}^{\,2}
=
\tfrac{5}{7}\,c_{334}^{\,2},
\]
while on the residue side the corresponding relation
\[
\bigl(\mathsf{C}^{\mathrm{res}}_{3,4;4;0,3}(4)\bigr)^2
=
\tfrac{5}{7}\,
\bigl(\mathsf{C}^{\mathrm{res}}_{3,3;4;0,2}(4)\bigr)^2
\]
is not yet formal from the visible Ward plus pairing package alone.
Establishing this visible top-pole Borcherds transport relation would
reduce the residue higher-spin packet to the same two primitive
self-coupling square classes
\(
\bigl(\mathsf{C}^{\mathrm{res}}_{3,3;4;0,2}(4)\bigr)^2
\)
and
\(
\bigl(\mathsf{C}^{\mathrm{res}}_{4,4;4;0,4}(4)\bigr)^2
\)
as on the principal side.
\end{remark}

\begin{conjecture}[Stage-\texorpdfstring{$4$}{4} visible top-pole Borcherds transport]
\label{conj:winfty-stage4-visible-borcherds-transport}
\ClaimStatusConjectured{}
Assume the hypotheses of
Proposition~\ref{prop:winfty-stage4-residue-pairing-reduction}.  Then
the visible stage-$4$ quotient residue calculus satisfies the transport
square identity
\[
\bigl(\mathsf{C}^{\mathrm{res}}_{3,4;4;0,3}(4)\bigr)^2
=
\tfrac{5}{7}\,
\bigl(\mathsf{C}^{\mathrm{res}}_{3,3;4;0,2}(4)\bigr)^2.
\]
Equivalently, on the stage-$4$ Ward-normalized visible pairing locus
the swap-odd \(W^{(4)}\)-transport square is forced by the
\((3,3)\)-square.
\end{conjecture}

\begin{corollary}[Equivalent exact forms of the remaining stage-\texorpdfstring{$4$}{4}
higher-spin transport input on the visible pairing locus;
\ClaimStatusProvedHere]
\label{cor:winfty-stage4-visible-borcherds-two-primitive}
Assume the hypotheses of
Proposition~\ref{prop:winfty-stage4-residue-pairing-reduction}.  Then
the following are equivalent:
\begin{enumerate}[label=\textup{(\roman*)}]
\item the visible top-pole Borcherds transport relation of
  Conjecture~\ref{conj:winfty-stage4-visible-borcherds-transport}
  holds:
  \[
  \bigl(\mathsf{C}^{\mathrm{res}}_{3,4;4;0,3}(4)\bigr)^2
  =
  \tfrac{5}{7}\,
  \bigl(\mathsf{C}^{\mathrm{res}}_{3,3;4;0,2}(4)\bigr)^2;
  \]
\item on the stage-$4$ Ward-normalized visible pairing locus, the
  higher-spin comparison of
  Corollary~\ref{cor:winfty-stage4-residue-four-channel} is equivalent
  at signless square-class level to the primitive pair
  \[
  \bigl(\mathsf{C}^{\mathrm{res}}_{3,3;4;0,2}(4)\bigr)^2
  =
  c_{334}^{\,2},
  \qquad
  \bigl(\mathsf{C}^{\mathrm{res}}_{4,4;4;0,4}(4)\bigr)^2
  =
  c_{444}^{\,2}.
  \]
\end{enumerate}
Under these equivalent conditions, the visible pairing locus has the
same two-primitive stage-$4$ square-class profile as the principal
Drinfeld--Sokolov side.
\end{corollary}

\begin{proof}
\textup{(i)}$\Rightarrow$\textup{(ii)}: the $\tfrac{5}{7}$-ratio from Corollary~\ref{cor:w4-ds-stage4-square-class-reduction} makes the transport square automatic once the $(3,3)$ primitive identity holds.  \textup{(ii)}$\Rightarrow$\textup{(i)}: the same principal formula recovers the transport relation.
\end{proof}

\begin{proposition}[Exact local attack order for the stage-\texorpdfstring{$4$}{4}
\texorpdfstring{$W_\infty$}{W_infty} packet; \ClaimStatusProvedHere]
\label{prop:winfty-stage4-local-attack-order}
Assume the hypotheses of
Proposition~\ref{prop:winfty-mc4-frontier-package}.  Assume moreover
that the stage-$4$ quotient residue calculus on the visible generators
\(W^{(2)},W^{(3)},W^{(4)}\) is decomposed into conformal families of
Virasoro-primary generators and that the Virasoro field \(T:=W^{(2)}\)
acts on them by the standard stage-$4$ Virasoro Ward identity,
including on itself.  Then the local stage-$4$ comparison separates
into two exact steps:
\begin{enumerate}[label=\textup{(\roman*)}]
\item \emph{Virasoro-target normalization step.}
  The remaining input for the four-channel higher-spin refinement is
  exactly the visible weight-$4$ normalization, equivalently the single
  scalar identity
  \[
  \mathsf{C}^{\mathrm{res}}_{4,4;2;0,6}(4)=2;
  \]
\item \emph{Higher-spin transport step.}
  On the additional visible pairing locus of
  Proposition~\ref{prop:winfty-stage4-residue-pairing-reduction}, the
  remaining input for the principal two-primitive square profile is
  exactly the visible top-pole Borcherds transport relation
  \[
  \bigl(\mathsf{C}^{\mathrm{res}}_{3,4;4;0,3}(4)\bigr)^2
  =
  \tfrac{5}{7}\,
  \bigl(\mathsf{C}^{\mathrm{res}}_{3,3;4;0,2}(4)\bigr)^2.
  \]
\end{enumerate}
Equivalently, after the unconditional six-entry packet of
Proposition~\ref{prop:winfty-mc4-frontier-package}, the local
stage-$4$ frontier is first a one-scalar Ward-normalization problem and
then, on the visible pairing locus, a one-relation higher-spin
transport problem.
\end{proposition}

\begin{proof}
\textup{(i)}: Corollary~\ref{cor:winfty-stage4-single-scalar-equivalent}.  \textup{(ii)}: Corollary~\ref{cor:winfty-stage4-primitive-transport-square-triple} reduces to the primitive-plus-transport triple; Corollary~\ref{cor:winfty-stage4-visible-borcherds-two-primitive} closes it.
\end{proof}

\begin{remark}[Scope of the stage-4 Ward-inheritance conjecture]
The principal values are already explicit (Proposition~\ref{prop:w4-ds-ope-explicit}).  The open issue is whether the quotient-system hypotheses force the visible weight-$4$ normalization (Conjecture~\ref{conj:winfty-stage4-visible-diagonal-normalization}).  This is distinct from the higher-spin Borcherds transport gap (Conjecture~\ref{conj:winfty-stage4-visible-borcherds-transport}); if both are forced, the stage-$4$ comparison reduces to two primitive square classes (Corollary~\ref{cor:winfty-stage4-visible-borcherds-two-primitive}).
\end{remark}

\begin{proposition}[Uniform Virasoro-target contraction of reduced
incremental packets under the normalized residue package;
\ClaimStatusProvedHere]
\label{prop:winfty-stage-growth-virasoro-target-contraction}
Fix $N \ge 3$ and assume the hypotheses of
Corollary~\ref{cor:winfty-ds-stage-growth-top-parity}.  Assume
moreover that the stage-$(N+1)$ quotient residue calculus on the
visible generators
$W^{(2)},\dots,W^{(N+1)}$ satisfies the hypotheses of
Propositions~\ref{prop:winfty-formal-mixed-virasoro-zero}
and~\ref{prop:winfty-formal-self-t-coefficient}.  Then the
Virasoro-target channels inside
$\mathcal{J}_{N+1}^{\mathrm{red}}$ are exactly
\[
\bigl\{
(s,N{+}1;2;0,s+N-1)\mid 3\le s\le N
\bigr\}
\sqcup
\bigl\{
(N{+}1,N{+}1;2;0,2N)
\bigr\}.
\]
Their residue values are theorematically fixed by
\[
\mathsf{C}^{\mathrm{res}}_{s,N+1;2;0,s+N-1}(N+1)=0
\qquad (3\le s\le N),
\]
\[
\mathsf{C}^{\mathrm{res}}_{N+1,N+1;2;0,2N}(N+1)=2.
\]
Consequently the reduced incremental packet contracts to the
higher-spin packet
\[
\mathcal{J}_{N+1}^{\mathrm{hs}}
:=
\mathcal{J}_{N+1}^{\mathrm{red}}
\setminus
\Bigl(
\bigl\{
(s,N{+}1;2;0,s+N-1)\mid 3\le s\le N
\bigr\}
\sqcup
\bigl\{
(N{+}1,N{+}1;2;0,2N)
\bigr\}
\Bigr).
\]
For $N=3$ this recovers
Corollary~\ref{cor:winfty-stage4-residue-four-channel}.
\end{proposition}

\begin{proof}
Proposition~\ref{prop:winfty-formal-mixed-virasoro-zero} kills each mixed target-$2$ channel $(s,N{+}1;2;0,s+N-1)$; Proposition~\ref{prop:winfty-formal-self-t-coefficient} fixes the self target-$2$ channel to $2$.  Removing these leaves $\mathcal{J}_{N+1}^{\mathrm{hs}}$.
\end{proof}

\begin{corollary}[First reduced stage beyond
\texorpdfstring{$\mathcal{I}_4$}{I4} under the normalized residue
package; \ClaimStatusProvedHere]
\label{cor:winfty-stage5-residue-eight-channel}
Assume the hypotheses of
Corollary~\ref{cor:winfty-ds-stage5-reduced-packet}.  Assume moreover
that the stage-$5$ quotient residue calculus on the visible generators
$W^{(2)},W^{(3)},W^{(4)},W^{(5)}$ satisfies the hypotheses of
Propositions~\ref{prop:winfty-formal-mixed-virasoro-zero}
and~\ref{prop:winfty-formal-self-t-coefficient}.  Then the three
Virasoro-target channels in $\mathcal{J}_5^{\mathrm{red}}$ are
theorematically fixed:
\[
\mathsf{C}^{\mathrm{res}}_{3,5;2;0,6}(5)=0,
\qquad
\mathsf{C}^{\mathrm{res}}_{4,5;2;0,7}(5)=0,
\qquad
\mathsf{C}^{\mathrm{res}}_{5,5;2;0,8}(5)=2.
\]
Consequently the first reduced stage beyond~$\mathcal{I}_4$ contracts
to the eight higher-spin channels
\[
\begin{gathered}
(3,4;5;0,2),\\
(3,5;3;0,5),\qquad
(3,5;4;0,4),\qquad
(3,5;5;0,3),\\
(4,5;3;0,6),\qquad
(4,5;4;0,5),\qquad
(4,5;5;0,4),\\
(5,5;4;0,6).
\end{gathered}
\]
\end{corollary}

\begin{proof}
Apply
Proposition~\ref{prop:winfty-stage-growth-virasoro-target-contraction}
with $N=4$, and use
Corollary~\ref{cor:winfty-ds-stage5-reduced-packet} to identify the
resulting packet explicitly.
\end{proof}

\begin{corollary}[First higher-spin packet beyond
\texorpdfstring{$\mathcal{I}_4$}{I4}; \ClaimStatusProvedHere]
\label{cor:winfty-stage5-higher-spin-packet}
Assume the hypotheses of
Corollary~\ref{cor:winfty-stage5-residue-eight-channel}.  Then the
first higher-spin packet beyond the stage-$4$ four-channel block is
\[
\mathcal{J}_5^{\mathrm{hs}}
=
\Bigl\{
(3,4;5;0,2)
\Bigr\}
\sqcup
\Bigl\{
(3,5;3;0,5),\,
(3,5;4;0,4),\,
(3,5;5;0,3)
\Bigr\}
\]
\[
\sqcup
\Bigl\{
(4,5;3;0,6),\,
(4,5;4;0,5),\,
(4,5;5;0,4)
\Bigr\}
\sqcup
\Bigl\{
(5,5;4;0,6)
\Bigr\}.
\]
In particular,
\[
|\mathcal{J}_5^{\mathrm{hs}}|=8,
\]
and its block decomposition by source pair is
\[
1+3+3+1
\]
for the blocks
$(3,4)$, $(3,5)$, $(4,5)$, $(5,5)$.
\end{corollary}

\begin{proof}
Corollary~\ref{cor:winfty-stage5-residue-eight-channel} identifies the
contracted packet explicitly, and
Proposition~\ref{prop:winfty-stage-growth-virasoro-target-contraction}
defines that contracted packet to be
$\mathcal{J}_5^{\mathrm{hs}}$.
\end{proof}

\begin{remark}[Stage-$5$ higher-spin packet decomposition]
\label{prop:winfty-stage5-higher-spin-subblocks}%
\label{cor:winfty-stage5-entry-transport}%
\label{cor:winfty-stage5-entry-singletons}%
\label{prop:winfty-stage5-entry-mixed-self}%
\label{prop:winfty-stage5-reduced-tail-singleton}%
\label{prop:winfty-stage5-tail-mechanism}%
\label{prop:winfty-stage5-transport-target-ladders}%
\label{prop:winfty-stage5-higher-spin-target-blocks}%
\label{cor:winfty-stage5-target5-corridor}%
\label{cor:winfty-stage5-target5-residual}%
\label{prop:winfty-stage5-target5-transport-mechanism}%
\label{prop:winfty-stage5-target5-transport-singletons}%
\label{prop:winfty-stage5-transport-pole-profiles}%
\label{rem:winfty-stage5-first-subproblems}%
The $8$-element packet $\mathcal{J}_5^{\mathrm{hs}}$ of
Corollary~\ref{cor:winfty-stage5-higher-spin-packet} decomposes by
source pair into blocks $(3,4)$, $(3,5)$, $(4,5)$, $(5,5)$ of sizes
$1+3+3+1$.  The entry packet
$\mathcal{J}_5^{\mathrm{entry}}=\{(3,4;5;0,2),(5,5;4;0,6)\}$ carries
the mixed-entry singleton $\mathcal{J}_5^{\mathrm{ent,mix}}=\{(3,4;5;0,2)\}$ and the self-return
singleton $\mathcal{J}_5^{\mathrm{ent,self}}=\{(5,5;4;0,6)\}$; the transport packet
$\mathcal{J}_5^{\mathrm{tr}}$ has $6$ elements decomposing by target
spin into three $2$-element ladders $\mathcal{J}_5^{\mathrm{tr},u}$ for
$u=3,4,5$.  By target spin,
$\mathcal{J}_5^{u=3}$, $\mathcal{J}_5^{u=4}$, $\mathcal{J}_5^{u=5}$
have sizes $2,3,3$.  The target-$5$ corridor
$\mathcal{J}_5^{u=5}=\{(3,4;5;0,2),(3,5;5;0,3),(4,5;5;0,4)\}$ sits at
consecutive pole orders $2,3,4$; its residual continuation after the
tail singleton consists of the two transport singletons
$(3,5;5;0,3)$ and $(4,5;5;0,4)$, representing the
$W^{(5)}$-projection coefficients
$\mathsf{C}^{\mathrm{res}}_{3,5;5;0,3}(5)$ and
$\mathsf{C}^{\mathrm{res}}_{4,5;5;0,4}(5)$.  The reduced tail block
$\mathcal{J}_5^{\mathrm{tail,red}}=\{(3,4;5;0,2)\}$ identifies the
missing mechanism as matching
$\mathsf{C}^{\mathrm{res}}_{3,4;5;0,2}(5)=\mathsf{C}^{\mathrm{DS}}_{3,4;5;0,2}(5)$.
Transport ladders have pole profiles $\{3,4\}$, $\{4,5\}$, $\{5,6\}$
for targets $5,4,3$ respectively.
\end{remark}

\begin{proposition}[Stage-\texorpdfstring{$5$}{5} visible
\texorpdfstring{$W^{(5)}$}{W5} normalization from the theorematic
\texorpdfstring{$W^{(5)}$-$W^{(5)}\to T$}{W5-W5 to T} coefficient;
\ClaimStatusProvedHere]
\label{prop:winfty-stage5-visible-w5-normalization}
Assume the hypotheses of
Corollary~\ref{cor:winfty-stage5-residue-eight-channel}.  Assume
moreover that the stage-$5$ quotient residue calculus on the visible
generators $W^{(2)},W^{(3)},W^{(4)},W^{(5)}$ is decomposed into
conformal families of Virasoro-primary generators and that the
Virasoro field $T:=W^{(2)}$ acts on them, including on itself, by the
standard stage-$5$ Virasoro Ward identity.  Then the visible stage-$5$
coefficient
\[
\mathsf{C}^{\mathrm{res}}_{5,5;2;0,8}(5)=2
\]
is already theorematic, and hence
\[
\langle W^{(5)}(z)W^{(5)}(0)\rangle=\frac{c/5}{z^{10}}.
\]
\end{proposition}

\begin{proof}
Apply Proposition~\ref{prop:winfty-formal-self-normalization-from-t} with $s=5$.
\end{proof}

\begin{proposition}[Stage-\texorpdfstring{$5$}{5} target-\texorpdfstring{$5$}{5} pole-\texorpdfstring{$3$}{3} transport singleton
vanishes on a visible \texorpdfstring{$W^{(5)}$}{W5}-pairing locus;
\ClaimStatusProvedHere]
\label{prop:winfty-stage5-target5-pole3-pairing-vanishing}
Assume the hypotheses of
Proposition~\ref{prop:winfty-stage5-target5-transport-singletons}.
Assume moreover that the visible stage-$5$ quotient residue calculus
carries a bilinear pairing \((-, -)_{\mathrm{vis}}\) on the span of
\(W^{(2)},W^{(3)},W^{(4)},W^{(5)}\) such that:
\begin{enumerate}[label=\textup{(\roman*)}]
\item the visible generators are even Virasoro-primary fields,
  pairwise orthogonal for distinct conformal weights, and their
  self-pairings are normalized by
  \[
  (W^{(u)},W^{(u)})_{\mathrm{vis}}=\frac{c}{u}
  \qquad (2\le u\le 5)
  \]
  with the same nonzero central charge \(c\);
\item for the visible primary field \(W^{(5)}\), the pairing satisfies
  the standard primary invariant-pairing identity
  \[
  \bigl(W^{(5)}_{(n)}a,b\bigr)_{\mathrm{vis}}
  =
  -\bigl(a,W^{(5)}_{(8-n)}b\bigr)_{\mathrm{vis}}
  \]
  whenever \(a,b\) are visible generators and the displayed modes are
  defined.
\end{enumerate}
Then the first singleton in the residual target-$5$ continuation
vanishes:
\[
\mathsf{C}^{\mathrm{res}}_{3,5;5;0,3}(5)=0.
\]
\end{proposition}

\begin{proof}
The invariant-pairing identity of assumption~\textup{(ii)} with \(n=2\), \(a=W^{(3)}\), \(b=W^{(5)}\) relates the reversed-order coefficient \(\mathsf{C}^{\mathrm{res}}_{5,3;5;0,3}=-\mathsf{C}^{\mathrm{res}}_{3,5;5;0,3}\) (swap parity, \(3+5-5=3\) odd, by Proposition~\ref{prop:winfty-ds-mixed-top-pole-swap}) to the complementary mode \(8-2=6\) self-OPE coefficient \((5,5;3;0,7)\), which vanishes by Proposition~\ref{prop:winfty-ds-self-ope-parity} (odd self-OPE).  Since \(c\neq 0\), the displayed vanishing follows.
\end{proof}

\begin{proposition}[Stage-\texorpdfstring{$5$}{5} target-\texorpdfstring{$5$}{5} pole-\texorpdfstring{$4$}{4} transport singleton
from the self-return singleton on a visible
\texorpdfstring{$W^{(5)}$}{W5}-pairing locus; \ClaimStatusProvedHere]
\label{prop:winfty-stage5-target5-pole4-from-self-return}
Assume the hypotheses of
Proposition~\ref{prop:winfty-stage5-target5-transport-singletons} and
Proposition~\ref{prop:winfty-stage5-target5-pole3-pairing-vanishing}.
Then
\[
\mathsf{C}^{\mathrm{res}}_{4,5;5;0,4}(5)
=
-\frac{5}{4}\,
\mathsf{C}^{\mathrm{res}}_{5,5;4;0,6}(5).
\]
\end{proposition}

\begin{proof}
Apply the invariant-pairing technique of Proposition~\ref{prop:winfty-stage5-target5-pole3-pairing-vanishing} with \(n=3\), \(a=W^{(4)}\), \(b=W^{(5)}\).  The swap symmetry of Proposition~\ref{prop:winfty-ds-mixed-top-pole-swap} (\(4+5-5=4\) even) gives \(\mathsf{C}^{\mathrm{res}}_{5,4;5;0,4}=\mathsf{C}^{\mathrm{res}}_{4,5;5;0,4}\); the complementary mode \(8-3=5\) produces \(\mathsf{C}^{\mathrm{res}}_{5,5;4;0,6}\); the sign is negative (\(W^{(5)}\) odd weight).  Weight orthogonality and \(c\neq 0\) yield the displayed identity.
\end{proof}

\begin{proposition}[Stage-\texorpdfstring{$5$}{5} target-\texorpdfstring{$5$}{5} pole-\texorpdfstring{$4$}{4} transport singleton
vanishes on a visible \texorpdfstring{$W^{(4)}$}{W4}-pairing locus;
\ClaimStatusProvedHere]
\label{prop:winfty-stage5-target5-pole4-w4-vanishing}
Assume the hypotheses of
Proposition~\ref{prop:winfty-stage5-target5-transport-singletons}.
Assume moreover that the visible stage-$5$ quotient residue calculus
carries a bilinear pairing \((-, -)_{\mathrm{vis}}\) on the span of
\(W^{(2)},W^{(3)},W^{(4)},W^{(5)}\) such that:
\begin{enumerate}[label=\textup{(\roman*)}]
\item the visible generators are even Virasoro-primary fields,
  pairwise orthogonal for distinct conformal weights, and their
  self-pairings are normalized by
  \[
  (W^{(u)},W^{(u)})_{\mathrm{vis}}=\frac{c}{u}
  \qquad (2\le u\le 5)
  \]
  with the same nonzero central charge \(c\);
\item for the visible primary field \(W^{(4)}\), the pairing satisfies
  the standard primary invariant-pairing identity
  \[
  \bigl(W^{(4)}_{(n)}a,b\bigr)_{\mathrm{vis}}
  =
  \bigl(a,W^{(4)}_{(6-n)}b\bigr)_{\mathrm{vis}}
  \]
  whenever \(a,b\) are visible generators and the displayed modes are
  defined.
\end{enumerate}
Then
\[
\mathsf{C}^{\mathrm{res}}_{4,5;5;0,4}(5)=0.
\]
\end{proposition}

\begin{proof}
The invariant-pairing identity with \(n=3\), \(a=W^{(4)}\), \(b=W^{(5)}\) relates \(\mathsf{C}^{\mathrm{res}}_{4,4;5;0,3}\) (killed by odd self-OPE parity, Proposition~\ref{prop:winfty-ds-self-ope-parity}) to \(\mathsf{C}^{\mathrm{res}}_{4,5;5;0,4}\cdot c/4\).  Since \(c\neq 0\), the vanishing follows.
\end{proof}

\begin{corollary}[Stage-\texorpdfstring{$5$}{5} self-return singleton vanishes on the
visible \texorpdfstring{$W^{(4)}$}{W4}/\texorpdfstring{$W^{(5)}$}{W5}
pairing locus; \ClaimStatusProvedHere]
\label{cor:winfty-stage5-self-return-vanishing-on-pairing}
Assume the hypotheses of
Proposition~\ref{prop:winfty-stage5-target5-pole4-from-self-return}
and Proposition~\ref{prop:winfty-stage5-target5-pole4-w4-vanishing}.
Then
\[
\mathsf{C}^{\mathrm{res}}_{5,5;4;0,6}(5)=0.
\]
\end{corollary}

\begin{proof}
Combine Proposition~\ref{prop:winfty-stage5-target5-pole4-w4-vanishing} with
Proposition~\ref{prop:winfty-stage5-target5-pole4-from-self-return}.
\end{proof}

\begin{proposition}[Stage-\texorpdfstring{$5$}{5} reduced tail singleton from a visible
\texorpdfstring{$W^{(3)}$}{W3}-pairing locus; \ClaimStatusProvedHere]
\label{prop:winfty-stage5-tail-from-w3-pairing}
Under the hypotheses of
Proposition~\ref{prop:winfty-stage5-tail-mechanism}, assume moreover
the visible stage-$5$ pairing package of
Proposition~\ref{prop:winfty-stage5-target5-pole3-pairing-vanishing}
with the $W^{(3)}$ invariant-pairing identity
$\bigl(W^{(3)}_{(n)}a,b\bigr)_{\mathrm{vis}}
=-\bigl(a,W^{(3)}_{(4-n)}b\bigr)_{\mathrm{vis}}$.
Then:
\[
\mathsf{C}^{\mathrm{res}}_{3,4;5;0,2}(5)
=
-\frac{5}{4}\,
\mathsf{C}^{\mathrm{res}}_{3,5;4;0,4}(5).
\]
\end{proposition}

\begin{proof}
Apply the $W^{(3)}$ invariant-pairing identity with \(n=1\), \(a=W^{(4)}\), \(b=W^{(5)}\); the complementary mode \(4-1=3\) produces \(\mathsf{C}^{\mathrm{res}}_{3,5;4;0,4}\cdot c/4\).  The sign is negative (\(W^{(3)}\) odd weight).  Since \(c\neq 0\), rearranging gives the displayed identity.
\end{proof}

\begin{proposition}[Stage-\texorpdfstring{$5$}{5} reduced tail singleton from a visible
\texorpdfstring{$W^{(4)}$}{W4}-pairing locus; \ClaimStatusProvedHere]
\label{prop:winfty-stage5-tail-from-w4-pairing}
Under the hypotheses of
Proposition~\ref{prop:winfty-stage5-tail-mechanism}, assume the visible stage-$5$ pairing package of
Proposition~\ref{prop:winfty-stage5-target5-pole3-pairing-vanishing}
with the $W^{(4)}$ invariant-pairing identity
$\bigl(W^{(4)}_{(n)}a,b\bigr)_{\mathrm{vis}}
=\bigl(a,W^{(4)}_{(6-n)}b\bigr)_{\mathrm{vis}}$.
Then:
\[
\mathsf{C}^{\mathrm{res}}_{3,4;5;0,2}(5)
=
\frac{5}{3}\,
\mathsf{C}^{\mathrm{res}}_{4,5;3;0,6}(5).
\]
\end{proposition}

\begin{proof}
Apply the $W^{(4)}$ invariant-pairing identity with \(n=1\), \(a=W^{(3)}\), \(b=W^{(5)}\).  The swap symmetry of Proposition~\ref{prop:winfty-ds-mixed-top-pole-swap} (\(3+4-5=2\) even) gives \(\mathsf{C}^{\mathrm{res}}_{4,3;5;0,2}=\mathsf{C}^{\mathrm{res}}_{3,4;5;0,2}\); the complementary mode \(6-1=5\) produces \(\mathsf{C}^{\mathrm{res}}_{4,5;3;0,6}\cdot c/3\).  The sign is positive (\(W^{(4)}\) even weight).  Since \(c\neq 0\), rearranging gives the displayed identity.
\end{proof}

\begin{corollary}[Stage-\texorpdfstring{$5$}{5} tail singleton equates neighboring transport channels;
\ClaimStatusProvedHere]
\label{cor:winfty-stage5-tail-cross-target-reduction}
Under the hypotheses of
Propositions~\ref{prop:winfty-stage5-tail-from-w3-pairing}
and~\ref{prop:winfty-stage5-tail-from-w4-pairing}:
$\mathsf{C}^{\mathrm{res}}_{3,5;4;0,4}(5)
=
-\frac{4}{3}\,
\mathsf{C}^{\mathrm{res}}_{4,5;3;0,6}(5)$.
\end{corollary}

\begin{proof}
Eliminate \(\mathsf{C}^{\mathrm{res}}_{3,4;5;0,2}(5)\) between the two propositions.
\end{proof}

\begin{corollary}[Stage-\texorpdfstring{$5$}{5} target-\texorpdfstring{$5$}{5} corridor contracts to the tail
singleton; \ClaimStatusProvedHere]
\label{cor:winfty-stage5-target5-corridor-to-tail}
Under the hypotheses of
Propositions~\ref{prop:winfty-stage5-target5-pole3-pairing-vanishing}
and~\ref{prop:winfty-stage5-target5-pole4-w4-vanishing}:
$\mathsf{C}^{\mathrm{res}}_{3,5;5;0,3}(5)=0$ and
$\mathsf{C}^{\mathrm{res}}_{4,5;5;0,4}(5)=0$.
\end{corollary}

\begin{corollary}[Stage-\texorpdfstring{$5$}{5} target-\texorpdfstring{$5$}{5} corridor carries no new
independent coefficient; \ClaimStatusProvedHere]
\label{cor:winfty-stage5-target5-no-new-independent-data}
Under the hypotheses of
Corollaries~\ref{cor:winfty-stage5-target5-corridor-to-tail}
and~\ref{cor:winfty-stage5-tail-cross-target-reduction}:
\[
\mathsf{C}^{\mathrm{res}}_{3,4;5;0,2}(5)
=
-\frac{5}{4}\,\mathsf{C}^{\mathrm{res}}_{3,5;4;0,4}(5)
=
\frac{5}{3}\,\mathsf{C}^{\mathrm{res}}_{4,5;3;0,6}(5),
\]
while the two target-$5$ transport coefficients vanish.
\end{corollary}

\begin{proposition}[Stage-\texorpdfstring{$5$}{5} target-\texorpdfstring{$4$}{4} pole-\texorpdfstring{$5$}{5} transport singleton
vanishes; \ClaimStatusProvedHere]
\label{prop:winfty-stage5-target4-pole5-w4-vanishing}
Under the hypotheses of
Proposition~\ref{prop:winfty-stage5-target5-pole4-w4-vanishing}:
$\mathsf{C}^{\mathrm{res}}_{4,5;4;0,5}(5)=0$.
\end{proposition}

\begin{proof}
The $W^{(4)}$ invariant-pairing identity with \(n=4\), \(a=W^{(5)}\), \(b=W^{(4)}\) relates \(\mathsf{C}^{\mathrm{res}}_{4,5;4;0,5}\cdot c/4\) to the complementary mode \(6-4=2\) self-OPE coefficient \((4,4;5;0,3)\), which vanishes by odd self-OPE parity (Proposition~\ref{prop:winfty-ds-self-ope-parity}).
\end{proof}

\begin{proposition}[Stage-\texorpdfstring{$5$}{5} target-\texorpdfstring{$3$}{3} pole-\texorpdfstring{$5$}{5} transport singleton
vanishes; \ClaimStatusProvedHere]
\label{prop:winfty-stage5-target3-pole5-w3-vanishing}
Under the hypotheses of
Proposition~\ref{prop:winfty-stage5-tail-from-w3-pairing}:
$\mathsf{C}^{\mathrm{res}}_{3,5;3;0,5}(5)=0$.
\end{proposition}

\begin{proof}
The $W^{(3)}$ invariant-pairing identity with \(n=4\), \(a=W^{(5)}\), \(b=W^{(3)}\) relates \(\mathsf{C}^{\mathrm{res}}_{3,5;3;0,5}\cdot c/3\) to the complementary mode \(4-4=0\) self-OPE coefficient \((3,3;5;0,1)\), which vanishes by odd self-OPE parity (Proposition~\ref{prop:winfty-ds-self-ope-parity}).
\end{proof}

\begin{proposition}[Stage-\texorpdfstring{$5$}{5} target-\texorpdfstring{$4$}{4}/target-\texorpdfstring{$3$}{3} transport channels
are paired on a visible \texorpdfstring{$W^{(5)}$}{W5}-pairing locus;
\ClaimStatusProvedHere]
\label{prop:winfty-stage5-transport-cross-target-reduction}
Assume the hypotheses of
Proposition~\ref{prop:winfty-stage5-target5-pole3-pairing-vanishing}.
Then
\[
\mathsf{C}^{\mathrm{res}}_{3,5;4;0,4}(5)
=
-\frac{4}{3}\,
\mathsf{C}^{\mathrm{res}}_{4,5;3;0,6}(5).
\]
\end{proposition}

\begin{proof}
Apply the $W^{(5)}$ invariant-pairing identity with \(n=3\), \(a=W^{(3)}\), \(b=W^{(4)}\).  The swap symmetry of Proposition~\ref{prop:winfty-ds-mixed-top-pole-swap} gives \(\mathsf{C}^{\mathrm{res}}_{5,3;4;0,4}=\mathsf{C}^{\mathrm{res}}_{3,5;4;0,4}\) (\(3+5-4=4\) even) and \(\mathsf{C}^{\mathrm{res}}_{5,4;3;0,6}=\mathsf{C}^{\mathrm{res}}_{4,5;3;0,6}\) (\(4+5-3=6\) even).  The sign is negative (\(W^{(5)}\) odd weight).  Weight orthogonality and \(c\neq 0\) yield the displayed identity.
\end{proof}

\begin{corollary}[Stage-\texorpdfstring{$5$}{5} mixed transport frontier carries one
effective independent coefficient; \ClaimStatusProvedHere]
\label{cor:winfty-stage5-transport-effective-independent-frontier}
Under the hypotheses of
Propositions~\ref{prop:winfty-stage5-target4-pole5-w4-vanishing},
\ref{prop:winfty-stage5-target3-pole5-w3-vanishing}, and
\ref{prop:winfty-stage5-transport-cross-target-reduction}:
$\mathsf{C}^{\mathrm{res}}_{4,5;4;0,5}(5)=0$,
$\mathsf{C}^{\mathrm{res}}_{3,5;3;0,5}(5)=0$, and
$\mathsf{C}^{\mathrm{res}}_{4,5;3;0,6}(5)
=
-\frac{3}{4}\,
\mathsf{C}^{\mathrm{res}}_{3,5;4;0,4}(5)$.
\end{corollary}

\begin{corollary}[Stage-\texorpdfstring{$5$}{5} higher-spin packet reduces to one effective
independent coefficient; \ClaimStatusProvedHere]
\label{cor:winfty-stage5-effective-independent-frontier}
Under the hypotheses of
Corollaries~\ref{cor:winfty-stage5-self-return-vanishing-on-pairing},
\ref{cor:winfty-stage5-target5-no-new-independent-data}, and
\ref{cor:winfty-stage5-transport-effective-independent-frontier},
set $A_5:=\mathsf{C}^{\mathrm{res}}_{3,5;4;0,4}(5)$.
Then every coefficient in $\mathcal{J}_5^{\mathrm{hs}}$ is either zero or determined by $A_5$:
\[
\mathsf{C}^{\mathrm{res}}_{3,4;5;0,2}(5)=-\tfrac{5}{4}\,A_5,
\quad
\mathsf{C}^{\mathrm{res}}_{4,5;3;0,6}(5)=-\tfrac{3}{4}\,A_5,
\]
with all other channels vanishing.
The effective independent coefficient is the target-$4$
singleton \((3,5;4;0,4)\).
\end{corollary}

\begin{proposition}[Exact local attack order for the first stage-\texorpdfstring{$5$}{5}
higher-spin packet; \ClaimStatusProvedHere]
\label{prop:winfty-stage5-local-attack-order}
Assume the hypotheses of
Corollary~\ref{cor:winfty-stage5-entry-transport}.  Then the first
stage-$5$ higher-spin frontier separates into the following exact local
order:
\begin{enumerate}[label=\textup{(\roman*)}]
\item \emph{Entry step.}  The attack begins with the two-channel entry
  packet \(\mathcal{J}_5^{\mathrm{entry}}\), first through the
  mixed-entry singleton \((3,4;5;0,2)\) and then through the
  self-return singleton \((5,5;4;0,6)\).
\item \emph{Target-$5$ corridor step.}  The first local mixed-transport
  strip is the target-$5$ corridor
  \[
  (3,4;5;0,2),\qquad
  (3,5;5;0,3),\qquad
  (4,5;5;0,4),
  \]
  whose residual continuation after the tail singleton
  \((3,4;5;0,2)\) is exactly the two singleton transport ladder
  \[
  (3,5;5;0,3),\qquad
  (4,5;5;0,4),
  \]
  taken in pole order \(3\) then \(4\).
\item \emph{Remaining mixed transport step.}  After the target-$5$
  ladder, the remaining transport frontier is attacked by fixed target
  in the order
  \[
  \mathcal{J}_5^{\mathrm{tr},5},\qquad
  \mathcal{J}_5^{\mathrm{tr},4},\qquad
  \mathcal{J}_5^{\mathrm{tr},3}.
  \]
\end{enumerate}
On the visible pairing loci of
Propositions~\ref{prop:winfty-stage5-target5-pole3-pairing-vanishing}
--\ref{cor:winfty-stage5-tail-cross-target-reduction}, the same
target-$5$ corridor becomes more rigid: the pole-$3$ singleton
vanishes, the pole-$4$ singleton is tied to the self-return singleton,
and the tail singleton is tied to neighboring target-$4$/target-$3$
transport channels.  On the full visible
\texorpdfstring{$W^{(3)}$}{W3}/\texorpdfstring{$W^{(4)}$}{W4}/\texorpdfstring{$W^{(5)}$}{W5}
pairing locus of
Corollary~\ref{cor:winfty-stage5-target5-no-new-independent-data}, the
whole target-$5$ corridor carries no new independent coefficient, so
Proposition~\ref{prop:winfty-stage5-target4-pole5-w4-vanishing},
Proposition~\ref{prop:winfty-stage5-target3-pole5-w3-vanishing}, and
Proposition~\ref{prop:winfty-stage5-transport-cross-target-reduction}
collapse the remaining target-$4$/target-$3$ transport ladders to one
effective coefficient, and
Corollary~\ref{cor:winfty-stage5-effective-independent-frontier}
identifies the effective independent local attack order there as the
single target-$4$ singleton \((3,5;4;0,4)\), with every other stage-$5$
higher-spin channel either zero or determined by it.
\end{proposition}

\begin{proof}
Parts \textup{(i)}--\textup{(iii)} follow from the packet decompositions already established:
Corollary~\ref{cor:winfty-stage5-entry-transport} for the entry/transport split,
Corollary~\ref{cor:winfty-stage5-target5-corridor} for the target-$5$ corridor,
and Proposition~\ref{prop:winfty-stage5-transport-pole-profiles} for the transport ladder ordering.
The visible-pairing refinement combines
Propositions~\ref{prop:winfty-stage5-target5-pole3-pairing-vanishing}--\ref{prop:winfty-stage5-transport-cross-target-reduction}
with Corollary~\ref{cor:winfty-stage5-effective-independent-frontier}.
\end{proof}

\begin{conjecture}[Principal stage-\texorpdfstring{$5$}{5} target-\texorpdfstring{$5$}{5} corridor carries no
new independent coefficient on the full visible pairing locus]
\label{conj:winfty-stage5-principal-target5-no-new-independent-data}
\ClaimStatusConjectured{}
Assume the hypotheses of
Corollary~\ref{cor:winfty-stage5-effective-independent-frontier}.  Then
\[
\mathsf{C}^{\mathrm{DS}}_{3,4;5;0,2}(5)
=
-\frac{5}{4}\,\mathsf{C}^{\mathrm{DS}}_{3,5;4;0,4}(5)
=
\frac{5}{3}\,\mathsf{C}^{\mathrm{DS}}_{4,5;3;0,6}(5),
\]
\[
\mathsf{C}^{\mathrm{DS}}_{3,5;5;0,3}(5)=0,
\qquad
\mathsf{C}^{\mathrm{DS}}_{4,5;5;0,4}(5)=0.
\]
Equivalently, on the full visible
\texorpdfstring{$W^{(3)}$}{W3}/\texorpdfstring{$W^{(4)}$}{W4}/\texorpdfstring{$W^{(5)}$}{W5}
pairing locus the principal target-$5$ corridor carries no new
independent coefficient beyond the neighboring target-$4$/target-$3$
channels.
\end{conjecture}

\begin{conjecture}[Principal stage-\texorpdfstring{$5$}{5} target-\texorpdfstring{$5$}{5} transport channels
vanish on the full visible pairing locus]
\label{conj:winfty-stage5-principal-target5-transport-vanishing}
\ClaimStatusConjectured{}
Assume the hypotheses of
Corollary~\ref{cor:winfty-stage5-effective-independent-frontier}.  Then
\[
\mathsf{C}^{\mathrm{DS}}_{3,5;5;0,3}(5)=0,
\qquad
\mathsf{C}^{\mathrm{DS}}_{4,5;5;0,4}(5)=0.
\]
Equivalently, on the full visible
\texorpdfstring{$W^{(3)}$}{W3}/\texorpdfstring{$W^{(4)}$}{W4}/\texorpdfstring{$W^{(5)}$}{W5}
pairing locus the principal target-$5$ corridor contracts to the tail
singleton.
\end{conjecture}

\begin{conjecture}[Principal stage-\texorpdfstring{$5$}{5} tail singleton is tied to the
neighboring target-\texorpdfstring{$4$}{4}/target-\texorpdfstring{$3$}{3} channels on the full visible pairing locus]
\label{conj:winfty-stage5-principal-tail-cross-target-reduction}
\ClaimStatusConjectured{}
Assume the hypotheses of
Corollary~\ref{cor:winfty-stage5-effective-independent-frontier}.  Then
\[
\mathsf{C}^{\mathrm{DS}}_{3,4;5;0,2}(5)
=
-\frac{5}{4}\,\mathsf{C}^{\mathrm{DS}}_{3,5;4;0,4}(5)
=
\frac{5}{3}\,\mathsf{C}^{\mathrm{DS}}_{4,5;3;0,6}(5).
\]
Equivalently, on the full visible pairing locus the principal reduced
tail singleton adds no new independent coefficient beyond the
neighboring target-$4$/target-$3$ transport channels.
\end{conjecture}

\begin{remark}[Target-$5$ corridor factorization]
\label{prop:winfty-stage5-principal-target5-factorization}%
Conjecture~\ref{conj:winfty-stage5-principal-target5-no-new-independent-data}
holds if and only if both
Conjecture~\ref{conj:winfty-stage5-principal-target5-transport-vanishing}
and
Conjecture~\ref{conj:winfty-stage5-principal-tail-cross-target-reduction}
hold: the two transport vanishings together with the tail cross-target relation account for all five identities.
\end{remark}

\begin{conjecture}[Principal residual stage-\texorpdfstring{$5$}{5} front carries one
effective independent coefficient on the full visible pairing locus]
\label{conj:winfty-stage5-principal-residual-front-one-coefficient}%
\label{conj:winfty-stage5-principal-residual-front-vanishings}%
\label{conj:winfty-stage5-principal-residual-cross-target-reduction}%
\label{prop:winfty-stage5-principal-residual-front-factorization}%
\ClaimStatusConjectured{}
Assume the hypotheses of
Corollary~\ref{cor:winfty-stage5-effective-independent-frontier}.  Then
\[
\mathsf{C}^{\mathrm{DS}}_{5,5;4;0,6}(5)=0,
\qquad
\mathsf{C}^{\mathrm{DS}}_{4,5;4;0,5}(5)=0,
\qquad
\mathsf{C}^{\mathrm{DS}}_{3,5;3;0,5}(5)=0,
\]
\[
\mathsf{C}^{\mathrm{DS}}_{4,5;3;0,6}(5)
=
-\frac{3}{4}\,\mathsf{C}^{\mathrm{DS}}_{3,5;4;0,4}(5).
\]
Equivalently, outside the target-$5$ corridor, the principal stage-$5$
higher-spin packet carries one effective independent coefficient,
represented by the target-$4$ singleton \((3,5;4;0,4)\).
The three vanishings alone form the \emph{residual vanishing component}
and the cross-target relation
$\mathsf{C}^{\mathrm{DS}}_{4,5;3;0,6}(5)=-\tfrac{3}{4}\,\mathsf{C}^{\mathrm{DS}}_{3,5;4;0,4}(5)$
forms the \emph{residual cross-target reduction component};
the full conjecture holds if and only if both components hold.
\end{conjecture}

\begin{conjecture}[Principal stage-\texorpdfstring{$5$}{5} one-coefficient normal form on the
full visible pairing locus]
\label{conj:winfty-stage5-principal-one-coefficient-normal-form}
\ClaimStatusConjectured{}
Assume the hypotheses of
Corollary~\ref{cor:winfty-stage5-effective-independent-frontier}.  Then
the principal stage-$5$ target packet satisfies the same one-coefficient
normal form:
\[
\mathsf{C}^{\mathrm{DS}}_{3,4;5;0,2}(5)
=
-\frac{5}{4}\,\mathsf{C}^{\mathrm{DS}}_{3,5;4;0,4}(5),
\qquad
\mathsf{C}^{\mathrm{DS}}_{3,5;5;0,3}(5)=0,
\qquad
\mathsf{C}^{\mathrm{DS}}_{4,5;5;0,4}(5)=0,
\]
\[
\mathsf{C}^{\mathrm{DS}}_{5,5;4;0,6}(5)=0,
\qquad
\mathsf{C}^{\mathrm{DS}}_{4,5;4;0,5}(5)=0,
\qquad
\mathsf{C}^{\mathrm{DS}}_{3,5;3;0,5}(5)=0,
\qquad
\mathsf{C}^{\mathrm{DS}}_{4,5;3;0,6}(5)
=
-\frac{3}{4}\,\mathsf{C}^{\mathrm{DS}}_{3,5;4;0,4}(5).
\]
Equivalently, on the full visible
\texorpdfstring{$W^{(3)}$}{W3}/\texorpdfstring{$W^{(4)}$}{W4}/\texorpdfstring{$W^{(5)}$}{W5}
pairing locus the principal stage-$5$ higher-spin packet carries one
effective independent coefficient, represented by the target-$4$
singleton \((3,5;4;0,4)\).  Equivalently again, this is the conjunction
of
Conjectures~\ref{conj:winfty-stage5-principal-target5-no-new-independent-data}
and~\ref{conj:winfty-stage5-principal-residual-front-one-coefficient}.
\end{conjecture}

\begin{proposition}[Principal stage-\texorpdfstring{$5$}{5} one-coefficient normal form
factors through the target-\texorpdfstring{$5$}{5} corridor and the residual front;
\ClaimStatusProvedHere]
\label{prop:winfty-stage5-principal-one-coefficient-factorization}
Assume the hypotheses of
Corollary~\ref{cor:winfty-stage5-effective-independent-frontier}.  Then
Conjecture~\ref{conj:winfty-stage5-principal-one-coefficient-normal-form}
holds if and only if both
Conjecture~\ref{conj:winfty-stage5-principal-target5-no-new-independent-data}
and
Conjecture~\ref{conj:winfty-stage5-principal-residual-front-one-coefficient}
hold.
\end{proposition}

\begin{proof}
Combine Propositions~\ref{prop:winfty-stage5-principal-target5-factorization} and~\ref{prop:winfty-stage5-principal-residual-front-factorization}; together they account for all seven identities in the normal form.
\end{proof}

\begin{proposition}[Stage-\texorpdfstring{$5$}{5} higher-spin comparison reduces to one
coefficient on the full visible pairing locus; \ClaimStatusProvedHere]
\label{prop:winfty-stage5-one-coefficient-reduction}
Assume the hypotheses of
Corollary~\ref{cor:winfty-stage5-effective-independent-frontier} and
Conjecture~\ref{conj:winfty-stage5-principal-one-coefficient-normal-form}.
Then the full stage-$5$ higher-spin comparison on
\(\mathcal{J}_5^{\mathrm{hs}}\) is equivalent to the single identity of
Conjecture~\ref{conj:winfty-stage5-one-coefficient-comparison},
\[
\mathsf{C}^{\mathrm{res}}_{3,5;4;0,4}(5)
=
\mathsf{C}^{\mathrm{DS}}_{3,5;4;0,4}(5).
\]
\end{proposition}

\begin{proof}
Corollary~\ref{cor:winfty-stage5-effective-independent-frontier} expresses every residue coefficient as a fixed multiple of \(\mathsf{C}^{\mathrm{res}}_{3,5;4;0,4}(5)\); matching this singleton matches all channels.
\end{proof}

\begin{conjecture}[Stage-\texorpdfstring{$5$}{5} one-coefficient comparison on the full
visible pairing locus]
\label{conj:winfty-stage5-one-coefficient-comparison}
\ClaimStatusConjectured{}
Assume the hypotheses of
Proposition~\ref{prop:winfty-stage5-one-coefficient-reduction}.  Then
\[
\mathsf{C}^{\mathrm{res}}_{3,5;4;0,4}(5)
=
\mathsf{C}^{\mathrm{DS}}_{3,5;4;0,4}(5).
\]
Equivalently, on the full visible
\texorpdfstring{$W^{(3)}$}{W3}/\texorpdfstring{$W^{(4)}$}{W4}/\texorpdfstring{$W^{(5)}$}{W5}
pairing locus the stage-$5$ higher-spin frontier closes once the
target-$4$ singleton \((3,5;4;0,4)\) is matched.
\end{conjecture}

\begin{corollary}[Exact remaining stage-\texorpdfstring{$5$}{5} visible-pairing input
package; \ClaimStatusProvedHere]
\label{cor:winfty-stage5-exact-remaining-input}
Assume the hypotheses of
Corollary~\ref{cor:winfty-stage5-effective-independent-frontier}.  Then
the full stage-$5$ higher-spin comparison on \(\mathcal{J}_5^{\mathrm{hs}}\)
closes on the full visible pairing locus if and only if the following
two inputs both hold:
\begin{enumerate}[label=\textup{(\roman*)}]
\item the packaged principal one-coefficient normal form of
  Conjecture~\ref{conj:winfty-stage5-principal-one-coefficient-normal-form}
  holds;
\item the target-$4$ singleton identity
  \[
  \mathsf{C}^{\mathrm{res}}_{3,5;4;0,4}(5)
  =
  \mathsf{C}^{\mathrm{DS}}_{3,5;4;0,4}(5)
  \]
  holds.
\end{enumerate}
Equivalently, the exact remaining local stage-$5$ visible-pairing
frontier is the packaged principal one-coefficient normal form together
with the singleton identity of
Conjecture~\ref{conj:winfty-stage5-one-coefficient-comparison}.  The
missing input is therefore no longer residue-side transport reduction,
but that packaged principal normal form together with that single
target-$4$ match.  Equivalently again, by
Proposition~\ref{prop:winfty-stage5-principal-one-coefficient-factorization},
this two-input package is the same as the pair of structural
conjectures
Conjectures~\ref{conj:winfty-stage5-principal-target5-no-new-independent-data}
and~\ref{conj:winfty-stage5-principal-residual-front-one-coefficient},
together with the singleton identity.
\end{corollary}

\begin{proof}
Necessity: the one-coefficient normal form of Corollary~\ref{cor:winfty-stage5-effective-independent-frontier} on the residue side forces the same on the principal side; the singleton identity is one matched channel.  Sufficiency: \textup{(i)} supplies the hypothesis of Proposition~\ref{prop:winfty-stage5-one-coefficient-reduction}, and \textup{(ii)} is the single identity to which it reduces.  The reformulation is Proposition~\ref{prop:winfty-stage5-principal-one-coefficient-factorization}.
\end{proof}

\begin{corollary}[Stage-\texorpdfstring{$5$}{5} higher-spin defect family collapses to one
representative defect on the full visible pairing locus;
\ClaimStatusProvedHere]
\label{cor:winfty-stage5-one-defect-family}
Assume the hypotheses of
Corollary~\ref{cor:winfty-stage5-effective-independent-frontier},
and
Conjecture~\ref{conj:winfty-stage5-principal-one-coefficient-normal-form}.
Set
\[
D_5:=
\mathsf{C}^{\mathrm{res}}_{3,5;4;0,4}(5)
-
\mathsf{C}^{\mathrm{DS}}_{3,5;4;0,4}(5).
\]
Then every defect in the stage-$5$ higher-spin packet
\(\mathcal{J}_5^{\mathrm{hs}}\) is either zero or a fixed rational
multiple of \(D_5\).  Explicitly,
\[
\begin{aligned}
\mathsf{C}^{\mathrm{res}}_{3,4;5;0,2}(5)
-
\mathsf{C}^{\mathrm{DS}}_{3,4;5;0,2}(5)
&=
-\frac{5}{4}\,D_5,\\
\mathsf{C}^{\mathrm{res}}_{3,5;5;0,3}(5)
-
\mathsf{C}^{\mathrm{DS}}_{3,5;5;0,3}(5)
&=0,\\
\mathsf{C}^{\mathrm{res}}_{4,5;5;0,4}(5)
-
\mathsf{C}^{\mathrm{DS}}_{4,5;5;0,4}(5)
&=0,
\end{aligned}
\]
\[
\begin{aligned}
\mathsf{C}^{\mathrm{res}}_{5,5;4;0,6}(5)
-
\mathsf{C}^{\mathrm{DS}}_{5,5;4;0,6}(5)
&=0,\\
\mathsf{C}^{\mathrm{res}}_{4,5;4;0,5}(5)
-
\mathsf{C}^{\mathrm{DS}}_{4,5;4;0,5}(5)
&=0,\\
\mathsf{C}^{\mathrm{res}}_{3,5;3;0,5}(5)
-
\mathsf{C}^{\mathrm{DS}}_{3,5;3;0,5}(5)
&=0,
\end{aligned}
\]
\[
\mathsf{C}^{\mathrm{res}}_{4,5;3;0,6}(5)
-
\mathsf{C}^{\mathrm{DS}}_{4,5;3;0,6}(5)
=
-\frac{3}{4}\,D_5.
\]
Equivalently, on that full visible pairing locus the whole stage-$5$
higher-spin defect family on \(\mathcal{J}_5^{\mathrm{hs}}\) vanishes
if and only if \(D_5=0\).
\end{corollary}

\begin{proof}
Subtract the residue-side one-coefficient normal form from the principal-side normal form.
\end{proof}

\begin{conjecture}[Stage-\texorpdfstring{$5$}{5} entry-packet identities]
\label{conj:winfty-stage5-entry-identities}
\ClaimStatusConjectured{}
Assume the hypotheses of
Corollary~\ref{cor:winfty-stage5-entry-transport}.  Then
\[
\mathsf{C}^{\mathrm{res}}_{s,t;u;m,n}(5)
=
\mathsf{C}^{\mathrm{DS}}_{s,t;u;m,n}(5)
\]
for every channel $(s,t;u;m,n)$ in the two-channel packet
$\mathcal{J}_5^{\mathrm{entry}}$.  Equivalently, the singleton channel
conjectures
Conjectures~\ref{conj:winfty-stage5-block-34}
and~\ref{conj:winfty-stage5-block-55}
both hold, in the attack order dictated by
Proposition~\ref{prop:winfty-stage5-entry-mixed-self}.
\end{conjecture}

\begin{conjecture}[Stage-\texorpdfstring{$5$}{5} mixed transport identities]
\label{conj:winfty-stage5-transport-identities}
\ClaimStatusConjectured{}
Assume the hypotheses of
Corollary~\ref{cor:winfty-stage5-entry-transport}.  Then
\[
\mathsf{C}^{\mathrm{res}}_{s,t;u;m,n}(5)
=
\mathsf{C}^{\mathrm{DS}}_{s,t;u;m,n}(5)
\]
for every channel $(s,t;u;m,n)$ in the six-channel packet
$\mathcal{J}_5^{\mathrm{tr}}$.  Equivalently, the three fixed-target
ladder conjectures
Conjectures~\ref{conj:winfty-stage5-transport-target-3}
--\ref{conj:winfty-stage5-transport-target-5}
all hold.  By
Proposition~\ref{prop:winfty-stage5-transport-pole-profiles}, the
natural local attack order is target~$5$, then target~$4$, then
target~$3$.
\end{conjecture}

\begin{conjecture}[Stage-\texorpdfstring{$5$}{5} transport target-\texorpdfstring{$3$}{3} ladder identities]
\label{conj:winfty-stage5-transport-target-3}
\ClaimStatusConjectured{}
Assume the hypotheses of
Corollary~\ref{cor:winfty-stage5-entry-transport}.  Then
\[
\mathsf{C}^{\mathrm{res}}_{s,t;u;m,n}(5)
=
\mathsf{C}^{\mathrm{DS}}_{s,t;u;m,n}(5)
\]
for every channel $(s,t;u;m,n)$ in the two-channel ladder
$\mathcal{J}_5^{\mathrm{tr},3}$ of
Proposition~\ref{prop:winfty-stage5-transport-target-ladders}.
\end{conjecture}

\begin{conjecture}[Stage-\texorpdfstring{$5$}{5} transport target-\texorpdfstring{$4$}{4} ladder identities]
\label{conj:winfty-stage5-transport-target-4}
\ClaimStatusConjectured{}
Assume the hypotheses of
Corollary~\ref{cor:winfty-stage5-entry-transport}.  Then
\[
\mathsf{C}^{\mathrm{res}}_{s,t;u;m,n}(5)
=
\mathsf{C}^{\mathrm{DS}}_{s,t;u;m,n}(5)
\]
for every channel $(s,t;u;m,n)$ in the two-channel ladder
$\mathcal{J}_5^{\mathrm{tr},4}$ of
Proposition~\ref{prop:winfty-stage5-transport-target-ladders}.
\end{conjecture}

\begin{conjecture}[Stage-\texorpdfstring{$5$}{5} transport target-\texorpdfstring{$5$}{5} ladder identities]
\label{conj:winfty-stage5-transport-target-5}
\ClaimStatusConjectured{}
Assume the hypotheses of
Corollary~\ref{cor:winfty-stage5-entry-transport}.  Then
\[
\mathsf{C}^{\mathrm{res}}_{s,t;u;m,n}(5)
=
\mathsf{C}^{\mathrm{DS}}_{s,t;u;m,n}(5)
\]
for every channel $(s,t;u;m,n)$ in the two-channel ladder
$\mathcal{J}_5^{\mathrm{tr},5}$ of
Proposition~\ref{prop:winfty-stage5-transport-target-ladders}.  Together
with Conjecture~\ref{conj:winfty-stage5-block-34}, this matches every
channel in the target-$5$ corridor of
Corollary~\ref{cor:winfty-stage5-target5-corridor}; after the singleton
tail input of Proposition~\ref{prop:winfty-stage5-reduced-tail-singleton}
is matched, this is exactly the residual continuation identified in
Corollary~\ref{cor:winfty-stage5-target5-residual} and the mixed
transport mechanism identified in
Proposition~\ref{prop:winfty-stage5-target5-transport-mechanism}.
Equivalently, the two singleton transport conjectures
Conjectures~\ref{conj:winfty-stage5-transport-target5-35}
and~\ref{conj:winfty-stage5-transport-target5-45}
both hold, in the pole order dictated by
Proposition~\ref{prop:winfty-stage5-target5-transport-singletons}.
\end{conjecture}

\begin{conjecture}[Stage-\texorpdfstring{$5$}{5} target-\texorpdfstring{$5$}{5} transport singleton from
\texorpdfstring{$W^{(3)}W^{(5)}$}{W3W5}]
\label{conj:winfty-stage5-transport-target5-35}
\ClaimStatusConjectured{}
Assume the hypotheses of
Corollary~\ref{cor:winfty-stage5-entry-transport}.  Then
\[
\mathsf{C}^{\mathrm{res}}_{3,5;5;0,3}(5)
=
\mathsf{C}^{\mathrm{DS}}_{3,5;5;0,3}(5).
\]
This is the first singleton in the residual target-$5$ continuation of
Proposition~\ref{prop:winfty-stage5-target5-transport-singletons},
equivalently the pole-$3$ $W^{(5)}$-projection in
$W^{(3)}(z)W^{(5)}(w)$.
\end{conjecture}

\begin{conjecture}[Stage-\texorpdfstring{$5$}{5} target-\texorpdfstring{$5$}{5} transport singleton from
\texorpdfstring{$W^{(4)}W^{(5)}$}{W4W5}]
\label{conj:winfty-stage5-transport-target5-45}
\ClaimStatusConjectured{}
Assume the hypotheses of
Corollary~\ref{cor:winfty-stage5-entry-transport}.  Then
\[
\mathsf{C}^{\mathrm{res}}_{4,5;5;0,4}(5)
=
\mathsf{C}^{\mathrm{DS}}_{4,5;5;0,4}(5).
\]
This is the second singleton in the residual target-$5$ continuation of
Proposition~\ref{prop:winfty-stage5-target5-transport-singletons},
equivalently the pole-$4$ $W^{(5)}$-projection in
$W^{(4)}(z)W^{(5)}(w)$.
\end{conjecture}

\begin{conjecture}[Stage-\texorpdfstring{$5$}{5}
\texorpdfstring{$(3,4)$}{(3,4)} higher-spin block identity]
\label{conj:winfty-stage5-block-34}
\ClaimStatusConjectured{}
Assume the hypotheses of
Corollary~\ref{cor:winfty-stage5-higher-spin-packet}.  Then
\[
\mathsf{C}^{\mathrm{res}}_{3,4;5;0,2}(5)
=
\mathsf{C}^{\mathrm{DS}}_{3,4;5;0,2}(5).
\]
This is exactly the reduced tail comparison input isolated in
Proposition~\ref{prop:winfty-stage5-reduced-tail-singleton}, the exact
missing comparison mechanism isolated in
Proposition~\ref{prop:winfty-stage5-tail-mechanism}, and the
first singleton channel in the target-$5$ corridor of
Corollary~\ref{cor:winfty-stage5-target5-corridor}.
\end{conjecture}

\begin{conjecture}[Stage-\texorpdfstring{$5$}{5}
\texorpdfstring{$(3,5)$}{(3,5)} higher-spin block identities]
\label{conj:winfty-stage5-block-35}
\ClaimStatusConjectured{}
Assume the hypotheses of
Corollary~\ref{cor:winfty-stage5-higher-spin-packet}.  Then
\[
\mathsf{C}^{\mathrm{res}}_{s,t;u;m,n}(5)
=
\mathsf{C}^{\mathrm{DS}}_{s,t;u;m,n}(5)
\]
for every channel $(s,t;u;m,n)$ in the three-channel block
$\mathcal{J}_5^{(3,5)}$ of
Proposition~\ref{prop:winfty-stage5-higher-spin-subblocks}.
\end{conjecture}

\begin{conjecture}[Stage-\texorpdfstring{$5$}{5}
\texorpdfstring{$(4,5)$}{(4,5)} higher-spin block identities]
\label{conj:winfty-stage5-block-45}
\ClaimStatusConjectured{}
Assume the hypotheses of
Corollary~\ref{cor:winfty-stage5-higher-spin-packet}.  Then
\[
\mathsf{C}^{\mathrm{res}}_{s,t;u;m,n}(5)
=
\mathsf{C}^{\mathrm{DS}}_{s,t;u;m,n}(5)
\]
for every channel $(s,t;u;m,n)$ in the three-channel block
$\mathcal{J}_5^{(4,5)}$ of
Proposition~\ref{prop:winfty-stage5-higher-spin-subblocks}.
\end{conjecture}

\begin{conjecture}[Stage-\texorpdfstring{$5$}{5}
\texorpdfstring{$(5,5)$}{(5,5)} self higher-spin block identity]
\label{conj:winfty-stage5-block-55}
\ClaimStatusConjectured{}
Assume the hypotheses of
Corollary~\ref{cor:winfty-stage5-higher-spin-packet}.  Then
\[
\mathsf{C}^{\mathrm{res}}_{5,5;4;0,6}(5)
=
\mathsf{C}^{\mathrm{DS}}_{5,5;4;0,6}(5).
\]
\end{conjecture}

\begin{conjecture}[First higher-spin bar-vs-DS packet beyond
\texorpdfstring{$\mathcal{I}_4$}{I4}]
\label{conj:winfty-stage5-higher-spin-identities}
\ClaimStatusConjectured{}
Assume the hypotheses of
Corollary~\ref{cor:winfty-stage5-residue-eight-channel}.  Then
\[
\mathsf{C}^{\mathrm{res}}_{s,t;u;m,n}(5)
=
\mathsf{C}^{\mathrm{DS}}_{s,t;u;m,n}(5)
\]
for every channel $(s,t;u;m,n)$ in the explicit packet
$\mathcal{J}_5^{\mathrm{hs}}$ of
Corollary~\ref{cor:winfty-stage5-higher-spin-packet}.  This is the
next finite list of higher-spin bar-vs-Drinfeld--Sokolov identities
after the stage-$4$ four-channel packet.  Equivalently, the four
subblock conjectures
Conjectures~\ref{conj:winfty-stage5-block-34}
--\ref{conj:winfty-stage5-block-55}
all hold, or, equivalently again, the entry-packet and mixed-transport
conjectures
Conjectures~\ref{conj:winfty-stage5-entry-identities}
and~\ref{conj:winfty-stage5-transport-identities}
both hold.  On the full visible
\texorpdfstring{$W^{(3)}$}{W3}/\texorpdfstring{$W^{(4)}$}{W4}/\texorpdfstring{$W^{(5)}$}{W5}
pairing locus,
Conjecture~\ref{conj:winfty-stage5-principal-one-coefficient-normal-form}
together with Proposition~\ref{prop:winfty-stage5-one-coefficient-reduction}
reduce this whole packet conjecture to the single scalar identity of
Conjecture~\ref{conj:winfty-stage5-one-coefficient-comparison},
and hence through the packaged principal one-coefficient normal form.
\end{conjecture}

\begin{corollary}[Visible stage-\texorpdfstring{$5$}{5} local conjecture network collapses
to one nontrivial singleton under principal normal form;
\ClaimStatusProvedHere]
\label{cor:winfty-stage5-visible-conjecture-network-collapse}
Assume the hypotheses of
Corollary~\ref{cor:winfty-stage5-effective-independent-frontier},
and
Conjecture~\ref{conj:winfty-stage5-principal-one-coefficient-normal-form}.
Then on the full visible
\texorpdfstring{$W^{(3)}$}{W3}/\texorpdfstring{$W^{(4)}$}{W4}/\texorpdfstring{$W^{(5)}$}{W5}
pairing locus:
\begin{enumerate}[label=\textup{(\roman*)}]
\item the packet conjecture
  Conjecture~\ref{conj:winfty-stage5-higher-spin-identities}, the
  entry/transport conjectures
  Conjectures~\ref{conj:winfty-stage5-entry-identities}
  and~\ref{conj:winfty-stage5-transport-identities}, the target-$3$ and
  target-$4$ ladder conjectures
  Conjectures~\ref{conj:winfty-stage5-transport-target-3}
  and~\ref{conj:winfty-stage5-transport-target-4}, and the block
  conjectures
  Conjectures~\ref{conj:winfty-stage5-block-34}
  --\ref{conj:winfty-stage5-block-45} are each equivalent to the single
  identity of
  Conjecture~\ref{conj:winfty-stage5-one-coefficient-comparison};
\item the target-$5$ ladder conjecture
  Conjecture~\ref{conj:winfty-stage5-transport-target-5}, its two
  singleton refinements
  Conjectures~\ref{conj:winfty-stage5-transport-target5-35}
  and~\ref{conj:winfty-stage5-transport-target5-45}, and the self-block
  conjecture
  Conjecture~\ref{conj:winfty-stage5-block-55} become automatic.
\end{enumerate}
In particular, once the packaged principal one-coefficient normal form
is granted, the local stage-$5$ conjecture tree below
\(\mathcal{J}_5^{\mathrm{hs}}\) has exactly one nontrivial comparison
input, namely the target-$4$ singleton \((3,5;4;0,4)\).
\end{corollary}

\begin{proof}
Under the one-coefficient normal form on both sides, each block conjecture is either automatic (self-block, target-$5$ ladder: all channels vanish on both sides) or equivalent to the singleton identity \(\mathsf{C}^{\mathrm{res}}_{3,5;4;0,4}=\mathsf{C}^{\mathrm{DS}}_{3,5;4;0,4}\) via the fixed rational multiples \(-\tfrac54\) (mixed-entry), \(1\) (representative), and \(-\tfrac34\) (target-$3$).  The entry, transport, and packet conjectures are conjunctions of these, hence each equivalent to the singleton identity.
\end{proof}

\begin{corollary}[Visible stage-\texorpdfstring{$5$}{5} local defect classes under principal
normal form; \ClaimStatusProvedHere]
\label{cor:winfty-stage5-visible-defect-classes}
Assume the hypotheses of
Corollary~\ref{cor:winfty-stage5-visible-conjecture-network-collapse},
and set
\[
D_5:=
\mathsf{C}^{\mathrm{res}}_{3,5;4;0,4}(5)
-
\mathsf{C}^{\mathrm{DS}}_{3,5;4;0,4}(5).
\]
Then the downstream local stage-$5$ conjectural surfaces on the full
visible
\texorpdfstring{$W^{(3)}$}{W3}/\texorpdfstring{$W^{(4)}$}{W4}/\texorpdfstring{$W^{(5)}$}{W5}
pairing locus split into the following exact defect classes:
\begin{enumerate}[label=\textup{(\roman*)}]
\item the mixed-entry block
  Conjecture~\ref{conj:winfty-stage5-block-34} carries defect
  \[
  \mathsf{C}^{\mathrm{res}}_{3,4;5;0,2}(5)
  -
  \mathsf{C}^{\mathrm{DS}}_{3,4;5;0,2}(5)
  =
  -\frac{5}{4}\,D_5;
  \]
\item the representative target-$4$ class consists of the singleton
  comparison
  Conjecture~\ref{conj:winfty-stage5-one-coefficient-comparison}, the
  target-$4$ ladder
  Conjecture~\ref{conj:winfty-stage5-transport-target-4}, and the
  \texorpdfstring{$(3,5)$}{(3,5)} block
  Conjecture~\ref{conj:winfty-stage5-block-35}, with representative
  defect exactly \(D_5\);
\item the target-$3$ class consists of the target-$3$ ladder
  Conjecture~\ref{conj:winfty-stage5-transport-target-3} and the
  \texorpdfstring{$(4,5)$}{(4,5)} block
  Conjecture~\ref{conj:winfty-stage5-block-45}, with representative
  defect
  \[
  \mathsf{C}^{\mathrm{res}}_{4,5;3;0,6}(5)
  -
  \mathsf{C}^{\mathrm{DS}}_{4,5;3;0,6}(5)
  =
  -\frac{3}{4}\,D_5;
  \]
\item the target-$5$ ladder
  Conjecture~\ref{conj:winfty-stage5-transport-target-5}, its singleton
  refinements
  Conjectures~\ref{conj:winfty-stage5-transport-target5-35}
  and~\ref{conj:winfty-stage5-transport-target5-45},
  and the self-block
  Conjecture~\ref{conj:winfty-stage5-block-55}
  are automatic, so their defects vanish.
\end{enumerate}
Equivalently, every nonautomatic local stage-$5$ visible-pairing
surface lies in one of the three exact defect classes
\(-\tfrac54 D_5\), \(D_5\), or \(-\tfrac34 D_5\).
\end{corollary}

\begin{proof}
The fixed ratios from Corollary~\ref{cor:winfty-stage5-visible-conjecture-network-collapse} give defects \(-\tfrac54 D_5\), \(D_5\), and \(-\tfrac34 D_5\); automatic channels have zero defect.
\end{proof}

\begin{computation}[Full stage-\texorpdfstring{$4$}{4} reduction chain verification]
\label{comp:winfty-stage4-reduction-chain}
\index{MC4!W-infinity stage-4 reduction chain}
The complete reduction chain from the residual packet
$\mathcal{J}_4$ to the frontier of
Proposition~\ref{prop:winfty-mc4-frontier-package} has been
verified computationally
(\texttt{compute/lib/w4\_stage4\_coefficients.py}, $85$~tests):
\begin{itemize}
\item \emph{Primaryity}
  (Corollary~\ref{cor:winfty-ds-stage4-top-pole-packet}):
  $|\mathcal{J}_4| = 29 \to |\mathcal{J}_4^{\mathrm{top}}| = 7$,
  eliminating $22$ sub-leading entries; seven top-pole elements
  verified against the manuscript display;
\item \emph{Parity}
  (Corollary~\ref{cor:winfty-ds-stage4-parity-packet}):
  $|\mathcal{J}_4^{\mathrm{top}}| = 7
  \to |\mathcal{J}_4^{\mathrm{par}}| = 6$,
  removing exactly $(4,4,3,5)$ (odd-spin self-OPE);
\item \emph{OPE blocks}
  (Corollary~\ref{cor:winfty-ds-stage4-ope-blocks}):
  $6 = 1 + 3 + 2$ as $(3,3)$, $(3,4)$, $(4,4)$ blocks;
\item \emph{Mixed-self split}
  (Corollary~\ref{cor:winfty-ds-stage4-mixed-self-split}):
  three self-coupling scalars $+$ three mixed channels;
\item \emph{Virasoro elimination}
  (Corollary~\ref{cor:winfty-ds-stage4-five-plus-zero}):
  $\mathsf{C}^{\mathrm{DS}}_{4,4;2;0,6} = 2$ (universal $T$-coupling)
  and $\mathsf{C}^{\mathrm{DS}}_{3,4;2;0,5} = 0$ (mixed vanishing)
  split the stage-$4$ packet into $4$ higher-spin channels plus
  $2$ theorematic principal Virasoro-target values.
\end{itemize}
The full chain $29 \to 7 \to 6 = 4 + 2$ matches every intermediate
cardinality and every explicit element of the manuscript displays.
\end{computation}

\begin{proposition}[Explicit \texorpdfstring{$W_4$}{W4} OPE structure constants;
\ClaimStatusProvedElsewhere]
\label{prop:w4-ds-ope-explicit}
In the principal Drinfeld--Sokolov reduction
$W_4 = W(\mathfrak{sl}_4, f_{\mathrm{prin}})$ with generators
$T$ (spin~$2$), $W^{(3)}$ (spin~$3$), $W^{(4)}$ (spin~$4$)
normalized so that
$\langle W^{(s)}(z)\,W^{(s)}(0)\rangle = (c/s)\,z^{-2s}$,
the four higher-spin stage-$4$ principal targets of
Corollary~\textup{\ref{cor:winfty-ds-stage4-five-plus-zero}} are
given by the following rational functions of the central charge~$c$:
\begin{align}
c_{334}^{\,2}
&= \frac{42\,c^2(5c+22)}{(c+24)(7c+68)(3c+46)},
\label{eq:c334-explicit}
\\[4pt]
c_{444}^{\,2}
&= \frac{112\,c^2(2c-1)(3c+46)}{(c+24)(7c+68)(10c+197)(5c+3)},
\label{eq:c444-explicit}
\\[4pt]
\mathsf{C}_{3,4;3;0,4}^{\,2}
&= \tfrac{9}{16}\,c_{334}^{\,2},
\label{eq:c343-explicit}
\\[4pt]
\mathsf{C}_{3,4;4;0,3}^{\,2}
&= \tfrac{5}{7}\,c_{334}^{\,2}.
\label{eq:c344-explicit}
\end{align}
The two theorematic Virasoro-target identities in
Corollary~\textup{\ref{cor:winfty-ds-stage4-five-plus-zero}} are
confirmed:
\begin{equation}
\mathsf{C}^{\mathrm{DS}}_{4,4;2;0,6} = 2,
\qquad
\mathsf{C}^{\mathrm{DS}}_{3,4;2;0,5} = 0.
\label{eq:w4-virasoro-target-identities-verified}
\end{equation}
\end{proposition}

\begin{proof}
The structure constants follow from the vertex-algebra Jacobi
identity (Borcherds identity) applied to the three-generator $W_4$
algebra.  Formula~\eqref{eq:c334-explicit} is the bootstrap result
of Hornfeck~\cite{Hornfeck93}; formula~\eqref{eq:c444-explicit}
then follows from the $(W^{(4)}, W^{(4)}, W^{(4)})$ Jacobi identity.
The inter-coefficient relations~\eqref{eq:c343-explicit}
and~\eqref{eq:c344-explicit} are metric-adjoint and Borcherds-identity
consequences of~\eqref{eq:c334-explicit}, respectively; see
Bouwknegt--Schoutens~\cite{Bouwknegt-Schoutens}
and Blumenhagen \textit{et al.}~\cite{BEHFH96}.
The universal $T$-coupling
$\mathsf{C}_{4,4;2;0,6}=2$ follows from
Proposition~\ref{prop:winfty-ds-self-t-coefficient}, and the mixed
Virasoro vanishing
$\mathsf{C}_{3,4;2;0,5}=0$ from
Proposition~\ref{prop:winfty-ds-mixed-virasoro-ds-zero}.
\end{proof}

\begin{corollary}[Principal stage-\texorpdfstring{$4$}{4} higher-spin packet from two
primitive square classes; \ClaimStatusProvedHere]
\label{cor:w4-ds-stage4-square-class-reduction}
In the setting of
Proposition~\ref{prop:w4-ds-ope-explicit}, the higher-spin part of the
principal stage-$4$ target packet is determined at signless
square-class level by the pair
\[
c_{334}^{\,2},
\qquad
c_{444}^{\,2}.
\]
More precisely, the two mixed higher-spin channels are forced by
$c_{334}^{\,2}$:
\[
\mathsf{C}_{3,4;3;0,4}^{\,2}
= \tfrac{9}{16}\,c_{334}^{\,2},
\qquad
\mathsf{C}_{3,4;4;0,3}^{\,2}
= \tfrac{5}{7}\,c_{334}^{\,2}.
\]
Consequently, after removing the theorematic Virasoro-target values
\[
\mathsf{C}^{\mathrm{DS}}_{4,4;2;0,6}=2,
\qquad
\mathsf{C}^{\mathrm{DS}}_{3,4;2;0,5}=0,
\]
the principal stage-$4$ DS packet does not require four primitive
higher-spin targets: it consists of the two primitive square classes
$c_{334}^{\,2},c_{444}^{\,2}$ together with two mixed square-class
descendants forced by~$c_{334}^{\,2}$.
\end{corollary}

\begin{proof}
The explicit formulas
\eqref{eq:c343-explicit} and~\eqref{eq:c344-explicit} express the two
mixed higher-spin squares directly in terms of~$c_{334}^{\,2}$, while
\eqref{eq:w4-virasoro-target-identities-verified} records the two
theorematic Virasoro-target values.  Hence the displayed pair
$c_{334}^{\,2},c_{444}^{\,2}$ is the entire primitive higher-spin
square-class input on the principal side.
\end{proof}

\begin{remark}[Exact stage-4 identity packet]
\label{rem:winfty-stage4-exact-identity-packet}
Proposition~\ref{prop:winfty-mc4-frontier-package} and
Proposition~\ref{prop:w4-ds-ope-explicit} together show that the first
genuinely higher-spin MC4 packet is no longer extraction of principal
Drinfeld--Sokolov data.  That side is already explicit.  The live
stage-$4$ task is the six bar-side identities
\[
\mathsf{C}^{\mathrm{res}}_{s,t;u;m,n}(4)
=
\mathsf{C}^{\mathrm{DS}}_{s,t;u;m,n}(4)
\]
on the packet
\[
\Bigl\{
\begin{gathered}
(3,3;4;0,2),\,
(4,4;4;0,4),\\
(3,4;3;0,4),\,
(3,4;4;0,3),\,
(4,4;2;0,6),\\
(3,4;2;0,5)
\end{gathered}
\Bigr\}.
\]
For the first four channels, the DS targets are fixed by the explicit
formulas \eqref{eq:c334-explicit}--\eqref{eq:c344-explicit}, up to the
chosen strong-generator signs.  The last two channels are fixed exactly
by \eqref{eq:w4-virasoro-target-identities-verified}.  Thus the first open higher-spin
packet is an exact six-entry identity problem on~$\mathcal{I}_4$, not a
further principal-side extraction problem.
\end{remark}

\begin{computation}[Full \texorpdfstring{$W_4$}{W4} OPE coefficient extraction]
\label{comp:w4-ds-ope-extraction}
\index{MC4!W4 OPE coefficient extraction}
The explicit structure constants of
Proposition~\ref{prop:w4-ds-ope-explicit} have been verified
computationally
(\texttt{compute/lib/w4\_ds\_ope\_extraction.py}, $121$~tests):
\begin{itemize}
\item all four squared structure constants evaluated at levels
  $k = 0, 1, 2, -2, -3$ (avoiding poles);
\item inter-coefficient ratios
  $\mathsf{C}_{3,4;3;0,4}^2/c_{334}^2 = 9/16$ and
  $\mathsf{C}_{3,4;4;0,3}^2/c_{334}^2 = 5/7$ verified symbolically;
\item Feigin--Frenkel duality $k \leftrightarrow -k-2h^\vee$
  checked: $c(k) + c(-k-8) = 246$;
\item unitarity non-negativity ($c_{334}^2 \geq 0$, $c_{444}^2 \geq 0$)
  confirmed at all $\mathcal{W}(N, N{+}1)$ minimal models
  for $4 \leq N \leq 19$;
\item both theorematic principal Virasoro-target values
  $\mathsf{C}^{\mathrm{DS}}_{4,4;2;0,6} = 2$ and
  $\mathsf{C}^{\mathrm{DS}}_{3,4;2;0,5} = 0$ confirmed.
\end{itemize}
\end{computation}

\subsection{Maurer--Cartan elements and deformation theory}
\label{sec:maurer-cartan-curved}

\begin{definition}[Maurer--Cartan element in curved context]
\label{def:mc-element-curved}
For a curved $A_\infty$ algebra $(A, \{\mu_n\})$, a \emph{Maurer--Cartan element} is $\alpha 
\in A^1$ satisfying:
\begin{equation}
\sum_{n \geq 0} \mu_n(\alpha^{\otimes n}) = 0
\end{equation}
\end{definition}

\begin{theorem}[Twisting by MC elements {\cite{LV12}}; \ClaimStatusProvedElsewhere]
\label{thm:twisting-mc}
Given an MC element $\alpha$, we can twist the curved $A_\infty$ structure:
\begin{equation}
\mu_n^\alpha(a_1, \ldots, a_n) = \sum_{k \geq 0} \mu_{n+k}(\alpha^{\otimes k}, 
a_1, \ldots, a_n)
\end{equation}

The twisted structure $(A, \{\mu_n^\alpha\})$ is again a curved $A_\infty$ algebra, with new 
curvature:
\begin{equation}
\mu_0^\alpha = \mu_0 + \mu_1(\alpha) + \frac{1}{2}\mu_2(\alpha, \alpha) + \cdots
\end{equation}

If $\alpha$ is an MC element, then $\mu_0^\alpha = 0$, so the twisted structure is 
\emph{uncurved}.
\end{theorem}

\begin{remark}[Physical interpretation]
MC elements correspond to \emph{vacua} (different ground states of the theory), \emph{deformations} (continuous families parametrized by the MC equation), and \emph{obstructions} (failure of the MC equation means curvature persists).
\end{remark}

\subsection{Summary and comparison table}
\label{sec:curved-summary}

\begin{table}[ht]
\centering
\caption{Quadratic, curved, and filtered structures}
\begin{tabular}{|l|c|c|c|}
\hline
\textbf{Property} & \emph{Quadratic} & \emph{Curved} & \emph{Filtered}\\
\hline
Curvature μ₀ & 0 & ∈ Z(A) & Multiple\\
Completion needed? & No & Sometimes & Usually\\
Koszul dual exists? & Yes & Yes (with completion) & Sometimes\\
Example & Heisenberg & Virasoro & W₃\\
\hline
\end{tabular}
\end{table}

The hierarchy is therefore:
\begin{equation}
\text{Quadratic} \subset \text{Curved} \subset \text{Filtered}
\end{equation}

Each level requires more sophisticated technology (completion, filtered cooperads), but 
also captures more examples from physics and representation theory.

\section{\texorpdfstring{Curved $A_\infty$ structures: strict vs.\ homotopy nilpotence}{Curved A-infinity structures: strict vs homotopy nilpotence}}
\label{sec:strict-vs-homotopy}

For chiral algebras $\mathcal{A}$ with quantum corrections at genus $g$,
the curvature $\mu_0$ controls two distinct squares:
\begin{enumerate}[label=\textup{(\roman*)}]
\item \emph{Algebraic:}\; $m_1^2(a) = [\mu_0, a]$.
  When $\mu_0 \in Z(\mathcal{A})$, this vanishes for all~$a$, so $m_1$ is a strict differential on~$\mathcal{A}$.  All standard chiral algebras (Heisenberg, Kac--Moody, Virasoro, $\mathcal{W}$-algebras) satisfy this because their central extensions are central.
\item \emph{Geometric:}\; $\dfib^{\,2} = \kappa(\mathcal{A})\cdot\omega_g \neq 0$ whenever $\kappa(\mathcal{A}) \neq 0$
  (Convention~\ref{conv:higher-genus-differentials}).
  Even when $\mu_0$ is central, the fiberwise bar differential over~$\overline{\mathcal{M}}_g$ is \emph{not} nilpotent---its failure is controlled by the scalar invariant~$\kappa(\mathcal{A})$.
\end{enumerate}
The total corrected differential $\Dg{g}$ absorbs the geometric curvature and satisfies $\Dg{g}^{\,2} = 0$ unconditionally (Theorem~\ref{thm:quantum-diff-squares-zero}).

\subsection{Mathematical foundations: three regimes}
\label{subsec:three-regimes}

The three regimes below correspond to the first three levels of the
regime hierarchy (Convention~\ref{conv:regime-tags}): quadratic
(Regime~I), curved-central (Regime~II), and filtered-complete
(Regime~III).  The fourth level (programmatic) lies beyond the
scope of the present chapter.

\subsubsection{\texorpdfstring{Regime I: strict differential ($d^2 = 0$ strict)}{Regime I: strict differential (d2 = 0 strict)}}

\begin{definition}[Strict dg structure]\label{def:strict-dg}
A \emph{strict differential graded structure} consists of:
\begin{itemize}
\item A graded vector space $V = \bigoplus_{n \in \mathbb{Z}} V^n$
\item A linear map $d: V^n \to V^{n+1}$ of degree $+1$
\item Satisfying $d \circ d = 0$ \emph{exactly}, not just up to homotopy
\end{itemize}
\end{definition}

\begin{example}[De Rham complex]
The de Rham complex $(\Omega^*(X), d_{dR})$ on a smooth manifold $X$:
\[\cdots \to \Omega^{n-1}(X) \xrightarrow{d_{dR}} \Omega^n(X) \xrightarrow{d_{dR}} 
\Omega^{n+1}(X) \to \cdots\]

Here $d_{dR}^2 = 0$ \emph{strict} because:
\[d_{dR}^2(f dx^{i_1} \wedge \cdots \wedge dx^{i_k}) 
= \sum_{j < k} \frac{\partial^2 f}{\partial x^j \partial x^k} dx^j \wedge dx^k \wedge dx^{i_1} 
\wedge \cdots = 0\]
by commutativity of partial derivatives.

\end{example}

\subsubsection{\texorpdfstring{Regime II: curved differential ($m_1^2 = [\mu_0, -]$, central curvature)}{Regime II: curved differential (m1 squared = commutator with mu0, central curvature)}}

% Duplicate curved A-infinity definition commented out.
% Canonical version: Definition~\ref{def:curved-ainfty} (line ~4866).
% Alias label preserved for any stale external references.
\label{def:curved-ainfty-complete}
\iffalse
%%% BEGIN COMMENTED-OUT BLOCK 4: duplicate curved A-infinity definition
%%% Canonical version lives at def:curved-ainfty (~line 4866).
\begin{definition}[Curved \texorpdfstring{$A_\infty$}{A_infinity} algebra]\label{def:curved-ainfty-complete-DISABLED}
A \emph{curved $A_\infty$ algebra} $(\mathcal{A}, \{m_k\}_{k \geq 0}, \mu_0)$ consists of:
\begin{enumerate}
\item A $\mathbb{Z}$-graded vector space $\mathcal{A} = \bigoplus_{n} \mathcal{A}^n$
\item Operations $m_k: \mathcal{A}^{\otimes k} \to \mathcal{A}[2-k]$ for each $k \geq 0$
\item A \emph{curvature element} $\mu_0 \in \mathcal{A}^2$
\end{enumerate}
satisfying the \emph{curved $A_\infty$ relations}:
\begin{equation}\label{eq:curved-ainfty-relations-DISABLED}
\sum_{\substack{r+s+t=n \\ r,t \geq 0, \, s \geq 0}} (-1)^{rs+t} m_{r+1+t}
(\text{id}^{\otimes r} \otimes m_s \otimes \text{id}^{\otimes t}) = 0
\end{equation}

The special cases of low degree are:
\begin{itemize}
\item $n = 0$: $m_1(m_0) = 0$ \quad (curvature is a $m_1$-cycle)
\item $n = 1$: $m_1^2(a) = m_2(m_0 \otimes \text{id})(a) - m_2(\text{id} \otimes m_0)(a)$
\quad (failure of $d^2 = 0$)
\item $n = 2$: Higher coherences involving $\mu_0$
\end{itemize}
\end{definition}
%%% END COMMENTED-OUT BLOCK 4
\fi

\begin{theorem}[Centrality implies strict nilpotence; \ClaimStatusProvedHere]\label{thm:central-implies-strict}
For the Heisenberg algebra on an elliptic curve (\S\ref{sec:frame-genus1}), the curvature $d^2 = k \cdot \omega_1$ was scalar, hence central.  Centrality of the curvature is the general mechanism ensuring $d^2 = 0$ strictly in the curved setting.

Let $(\mathcal{A}, m_1, \mu_0)$ be a curved chiral algebra with curvature satisfying:
\begin{equation}
\mu_0 \in Z(\mathcal{A}) := \{a \in \mathcal{A} \mid m_2(a \otimes b) = (-1)^{|a||b|}m_2(b \otimes a)
\text{ for all } b\}
\end{equation}
then the bar differential satisfies:
\begin{equation}
d_{\text{bar}}^2 = 0 \quad \emph{strictly}
\end{equation}
\end{theorem}

\begin{proof}

\emph{Step 1: Bar Differential Formula}

Recall from \S\ref{def:bar-differential-complete} that the bar differential on 
$\bar{B}^n(\mathcal{A})$ has three components:
\begin{equation}
d_{\text{bar}} = d_{\text{internal}} + d_{\text{residue}} + d_{\text{correction}}
\end{equation}
where:
\begin{align}
d_{\text{internal}}(a_0 \otimes \cdots \otimes a_n) 
&= \sum_{i=0}^n (-1)^{|a_0| + \cdots + |a_{i-1}|} (a_0 \otimes \cdots \otimes m_1(a_i) 
\otimes \cdots \otimes a_n) \\
d_{\text{residue}}(a_0 \otimes \cdots \otimes a_n) 
&= \sum_{i=0}^{n-1} (-1)^{\epsilon_i} \text{Res}_{D_i}(a_0 \otimes \cdots \otimes a_i \cdot a_{i+1} 
\otimes \cdots \otimes a_n) \\
d_{\text{correction}}(a_0 \otimes \cdots \otimes a_n) 
&= \mu_0 \otimes (a_0 \otimes \cdots \otimes a_n) \otimes \omega_g
\end{align}

\emph{Step 2: Computing $d^2$ - Nine Terms}

We need to compute:
\begin{equation}
d_{\text{bar}}^2 = (d_{\text{internal}} + d_{\text{residue}} + d_{\text{correction}})^2
\end{equation}

This expands to \emph{nine terms}:
\begin{align}
d_{\text{bar}}^2 &= d_{\text{internal}}^2 
+ d_{\text{internal}} d_{\text{residue}} + d_{\text{residue}} d_{\text{internal}} \\
&\quad + d_{\text{residue}}^2 \\
&\quad + d_{\text{internal}} d_{\text{correction}} + d_{\text{correction}} d_{\text{internal}} \\
&\quad + d_{\text{residue}} d_{\text{correction}} + d_{\text{correction}} d_{\text{residue}} \\
&\quad + d_{\text{correction}}^2
\end{align}

\emph{Term 1: $d_{\text{internal}}^2$}

For a \emph{curved} $A_\infty$ algebra, $m_1^2 \neq 0$ in general. The $n=1$ curved $A_\infty$ relation (Eq.~\ref{eq:curved-ainfty-relations}) gives:
\begin{equation}\label{eq:curved-m1-squared}
m_1^2(a) = m_2(\mu_0 \otimes a) - m_2(a \otimes \mu_0)
\end{equation}
Therefore $d_{\text{internal}}^2 \neq 0$ when the curvature $\mu_0$ is nonzero. This contribution is cancelled by terms 5--6 below (the cross-terms with $d_{\text{correction}}$); see the combined analysis there.

\emph{Term 2-3: $d_{\text{internal}} d_{\text{residue}} + d_{\text{residue}} d_{\text{internal}} = 0$}

These cancel by the \emph{Leibniz rule}:
\begin{equation}
m_1(a \cdot b) = m_1(a) \cdot b + (-1)^{|a|} a \cdot m_1(b)
\end{equation}

Explicitly, using residue calculus:
\begin{align}
&d_{\text{internal}}(\text{Res}_{D_i}(a_0 \otimes \cdots)) 
= \text{Res}_{D_i}(m_1(a_i \cdot a_{i+1})) \\
&= \text{Res}_{D_i}(m_1(a_i) \cdot a_{i+1}) + \text{Res}_{D_i}(a_i \cdot m_1(a_{i+1})) 
\quad \text{(Leibniz)} \\
&= d_{\text{residue}}(d_{\text{internal}}(a_0 \otimes \cdots))
\end{align}

Therefore the cross terms cancel: 
\[d_{\text{internal}} d_{\text{residue}} = - d_{\text{residue}} d_{\text{internal}}\]

\emph{Term 4: $d_{\text{residue}}^2 = 0$ (Arnold Relations)}

By Theorem~\ref{thm:arnold-three}:
\begin{equation}
\eta_{ij} \wedge \eta_{jk} + \eta_{jk} \wedge \eta_{ki} + \eta_{ki} \wedge \eta_{ij} = 0
\end{equation}
in $H^2(\text{Conf}_n(X))$.

Geometrically, this means:
\begin{equation}
[\text{Res}_{D_{ij}}, \text{Res}_{D_{jk}}] + [\text{Res}_{D_{jk}}, \text{Res}_{D_{ki}}] 
+ [\text{Res}_{D_{ki}}, \text{Res}_{D_{ij}}] = 0
\end{equation}

Therefore:
\[d_{\text{residue}}^2 = \sum_{i,j} \text{Res}_{D_i} \text{Res}_{D_j} 
= 0 \quad \text{(by Arnold)}\]

\emph{Terms 1+5+6 combined: $d_{\text{internal}}^2 + d_{\text{internal}} d_{\text{correction}} + d_{\text{correction}} d_{\text{internal}} = 0$}

These three terms must be analyzed together. We have:
\begin{align}
d_{\text{internal}}(d_{\text{correction}}(a_0 \otimes \cdots))
&= d_{\text{internal}}(\mu_0 \otimes a_0 \otimes \cdots) \\
&= m_1(\mu_0) \otimes (a_0 \otimes \cdots) + \sum_i \mu_0 \otimes \cdots \otimes m_1(a_i) \otimes \cdots
\end{align}
and similarly for $d_{\text{correction}} d_{\text{internal}}$. The cross-terms produce:
\begin{equation}
\begin{aligned}
&(d_{\text{internal}} d_{\text{correction}} + d_{\text{correction}} d_{\text{internal}})
(a_0 \otimes \cdots) \\
&\quad = m_1(\mu_0) \otimes (a_0 \otimes \cdots)
 + \sum_i \pm\bigl(
a_0 \otimes \cdots \otimes m_2(\mu_0, a_i)
 + m_2(a_i, \mu_0) \otimes \cdots
\bigr).
\end{aligned}
\end{equation}
By the curved $A_\infty$ relation~\eqref{eq:curved-m1-squared}: $m_1^2(a_i) = m_2(\mu_0, a_i) - m_2(a_i, \mu_0)$. Since $\mu_0$ is a \emph{cycle} ($m_1(\mu_0) = 0$), the combined contribution is:
\[d_{\text{internal}}^2 + d_{\text{internal}} d_{\text{correction}} + d_{\text{correction}} d_{\text{internal}} = 0\]

\emph{Terms 7-8: $d_{\text{residue}} d_{\text{correction}} + d_{\text{correction}} 
d_{\text{residue}}$}

Centrality of $\mu_0$ is essential here.

We have:
\begin{align}
d_{\text{residue}}(d_{\text{correction}}(a_0 \otimes \cdots)) 
&= d_{\text{residue}}(\mu_0 \otimes a_0 \otimes \cdots) \\
&= \sum_i \text{Res}_{D_i}((\mu_0 \otimes a_0 \otimes \cdots \otimes a_i \cdot a_{i+1} 
\otimes \cdots))
\end{align}

For this to equal $d_{\text{correction}}(d_{\text{residue}}(\cdots))$, we need:
\begin{equation}
\text{Res}_{D_i}(\mu_0 \otimes \cdots) = \mu_0 \otimes \text{Res}_{D_i}(\cdots)
\end{equation}

This holds if and only if $\mu_0 \in Z(\mathcal{A})$.

\emph{Proof of centrality.}
The residue operation involves computing:
\[\text{Res}_{z_i = z_{i+1}}(m_2(a_i \otimes a_{i+1}))\]

If $\mu_0$ is central, it commutes with all $a_i$:
\[m_2(\mu_0 \otimes a_i) = m_2(a_i \otimes \mu_0)\]

Therefore:
\begin{align}
\text{Res}_{D_i}(\mu_0 \otimes a_0 \otimes \cdots \otimes m_2(a_i \otimes a_{i+1}) \otimes \cdots) 
&= \text{Res}_{D_i}(m_2(\mu_0 \otimes a_0) \otimes \cdots) \\
&= \text{Res}_{D_i}(m_2(a_0 \otimes \mu_0) \otimes \cdots) \quad \text{(by centrality)} \\
&= \mu_0 \otimes \text{Res}_{D_i}(a_0 \otimes \cdots)
\end{align}

\emph{Term 9: $d_{\text{correction}}^2$}

Finally:
\begin{align}
d_{\text{correction}}^2(a_0 \otimes \cdots) 
&= d_{\text{correction}}(\mu_0 \otimes a_0 \otimes \cdots) \\
&= \mu_0 \otimes \mu_0 \otimes (a_0 \otimes \cdots) \otimes \omega_g^2
\end{align}

This term involves the commutator $[\mu_0, -]$ acting on bar elements.  When $\mu_0 \in Z(\mathcal{A})$, the commutator $[\mu_0, -]_{\mu_2} = 0$ vanishes identically, so this term is zero regardless of $\omega_g^2$.

(When $\mu_0 \notin Z(\mathcal{A})$, the term $d_{\mathrm{correction}}^2$ produces the genus-$g$ obstruction class $\mathrm{Obs}_g(\mathcal{A}) = [\mu_0, -] \otimes \int_{\mathcal{M}_g} \omega_g \cup \omega_g \neq 0$; see Remark~\ref{rem:non-central-curvature}.)

Combining all nine terms:
\begin{equation}
d_{\text{bar}}^2 = 0 + 0 + 0 + 0 + 0 + 0 + 0 = 0 \quad \text{strictly}
\end{equation}
provided $\mu_0 \in Z(\mathcal{A})$.

\end{proof}

\begin{remark}[What if \texorpdfstring{$\mu_0 \notin Z(\mathcal{A})$}{mu0 not in Z(A)}?]\label{rem:non-central-curvature}
If the curvature is \emph{not central}, then:
\begin{equation}
d_{\text{bar}}^2 \neq 0
\end{equation}
and we only have $d_{\text{bar}}^2 = 0$ \emph{up to homotopy}.

This leads to:
\begin{itemize}
\item \emph{Homotopy coherent structures} (Lurie, Higher Topos Theory)
\item \emph{Spectral sequences that may not degenerate}
\item \emph{Obstruction theories with non-closed obstructions}
\item \emph{Need for $A_\infty$ or $L_\infty$ structures at all levels}
\end{itemize}

However, \emph{all our examples} (Heisenberg, Kac--Moody, Virasoro, W-algebras) have
\emph{central curvature}, so we obtain strict nilpotence.  See
Remark~\ref{rem:voa-central-curvature} for the precise mechanism.
\end{remark}

\begin{remark}[Central curvature in vertex operator algebras]\label{rem:voa-central-curvature}
For vertex operator algebras, the vacuum axiom guarantees that the vacuum vector
$\mathbf{1}\in V$ is central: $Y(\mathbf{1},z) = \mathrm{id}_V$.  In the curved
$A_\infty$ framework, the curvature element $\mu_0$ arising from the genus-$g$
quantum correction is proportional to $\mathbf{1}$ (it equals the central charge or
level times $\mathbf{1}$; see Examples~\ref{ex:heisenberg-strict}--\ref{ex:w3-strict}).
Since $\mathbf{1}$ acts by the identity on every module, $\mu_0$ lies in $Z(\mathcal{A})$,
and Theorem~\ref{thm:central-implies-strict} yields $d_{\mathrm{bar}}^2=0$ strict.

For more general chiral algebras without a vacuum axiom --- for instance, chiral
algebras on higher-dimensional varieties, or non-unital factorization algebras ---
the curvature need not be central, and $d^2=0$ may hold only up to homotopy.
Similarly, chiral algebras arising from field theories with non-central anomalies
(such as gravitational anomalies that act non-trivially on the operator algebra)
fall outside the scope of Theorem~\ref{thm:central-implies-strict}.
\end{remark}

\subsubsection{\texorpdfstring{Regime III: general homotopy coherent ($d^2 \sim 0$ via homotopy)}{Regime III: general homotopy coherent (d2 0 via homotopy)}}

\begin{definition}[Homotopy coherent differential]\label{def:homotopy-coherent}
A \emph{homotopy coherent differential} on a graded space $V$ consists of:
\begin{itemize}
\item $d_1: V \to V[1]$ (the ``differential'')
\item $h: V \to V[-1]$ (a homotopy)
\item Satisfying: $d_1^2 = [d_1, h]$ (not zero, but homotopic to zero)
\item Plus higher coherence homotopies $h_2, h_3, \ldots$ ad infinitum
\end{itemize}
\end{definition}

This is the setting of Lurie's $(\infty, 1)$-categories~\cite{HTT}, derived algebraic geometry (To\"en--Vezzosi, Lurie), and non-curved $A_\infty$ or $L_\infty$ structures.

We do not need this level of generality for the chiral algebras studied in this monograph, since their curvature elements are central (see Remark~\ref{rem:voa-central-curvature} below).

\subsection{Application to chiral algebras: four examples}

\subsubsection{\texorpdfstring{Example 1: Heisenberg algebra (level $k$)}{Example 1: Heisenberg algebra (level k)}}

\begin{example}[Heisenberg: strict nilpotence]\label{ex:heisenberg-strict}
The Heisenberg algebra $\mathcal{H}_k$ has OPE $J(z)J(w) = k/(z-w)^2 + \text{reg.}$, curvature $\mu_0 = k \cdot \mathbf{1}$, and $\mu_0 \in Z(\mathcal{H}_k)$ since $\mathbf{1}$ is central.

\emph{Consequence.} The bar differential satisfies:
\[d_{\text{bar}}^2 = 0 \quad \text{strictly}\]

\emph{Genus 1 correction.}
At genus 1, the differential includes the term:
\[d_1 = d_0 + k \cdot \int_{\mathcal{M}_1} \omega_1\]
where $\omega_1 \in H^*(\mathcal{M}_1, \mathcal{Z}(\mathcal{A}))$ is the genus-1 quantum correction class.  (Here $\mathcal{M}_1 \cong \mathbb{H}/\mathrm{SL}_2(\mathbb{Z})$ is one-dimensional.)

This modifies $d_0$ but still $d_1^2 = 0$ strictly because $k$ is central.

\emph{Explicit verification at genus $1$.}
\begin{align}
d_1^2(J \otimes J) &= d_1(d_1(J \otimes J)) \\
&= d_1\left(\text{Res}_{z=w}(J(z)J(w)) + k \cdot \omega_1 \otimes J\right) \\
&= d_1\left(\frac{k}{(z-w)^2} dz \, dw + k \cdot \omega_1 \otimes J\right) \\
&= \text{Res}_{z=w}\left(\frac{k}{(z-w)^2}\right) + k \cdot \omega_1 \otimes \text{Res}(J) 
+ k \cdot \omega_1^2 \\
&= 0 + 0 + 0 = 0 \quad \text{(strictly)}
\end{align}

The $\omega_1^2$ term vanishes because $\dim \mathcal{M}_1 = 1$, so $H^2(\mathcal{M}_1) = 0$.
\end{example}

\subsubsection{\texorpdfstring{Example 2: affine Kac--Moody (level $k$)}{Example 2: affine Kac--Moody (level k)}}

\begin{example}[Kac--Moody: strict nilpotence]\label{ex:kac-moody-strict}
For $\widehat{\mathfrak{g}}_k$ with OPE $J^a(z)J^b(w) = k\kappa^{ab}/(z-w)^2 + f^{abc}J^c(w)/(z-w) + \text{reg.}$, the curvature is $\mu_0 = k \cdot K$ where $K$ generates the one-dimensional center.

\emph{Consequence.} Again $d_{\text{bar}}^2 = 0$ strict.

\emph{Higher genus.}
At genus $g$, the correction involves:
\[\mu_0^{(g)} = k \cdot \lambda_g \in H^2(\mathcal{M}_g, Z(\widehat{\mathfrak{g}}_k))\]
where $\lambda_g$ is a Hodge class.

Since $\mu_0^{(g)}$ is central, all higher genus bar differentials square to zero strictly.
\end{example}

\subsubsection{\texorpdfstring{Example 3: Virasoro algebra (central charge $c$)}{Example 3: Virasoro algebra (central charge c)}}

\begin{example}[Virasoro: curved but strict]\label{ex:virasoro-strict}
The Virasoro algebra $\text{Vir}_c$ has stress tensor $T$ with OPE
\[T(z)T(w) = \frac{c/2}{(z-w)^4} + \frac{2T(w)}{(z-w)^2} + \frac{\partial T(w)}{z-w} + \cdots\]
and curvature $\mu_0 = c \cdot \mathbf{1}$ from the central charge, which is a central element.

\emph{Subtlety.} The Virasoro algebra has \emph{higher order corrections}:
\[m_3(T \otimes T \otimes T) \neq 0\]
due to the cubic Schwarzian derivative term.

However, the \emph{curvature} $\mu_0 = c \cdot \mathbf{1}$ is still central, so:
\[d_{\text{bar}}^2 = 0 \quad \text{strictly}\]

\emph{Physical interpretation.}
The central charge $c$ measures the \emph{conformal anomaly}. It is a quantum correction
that breaks classical conformal invariance, but it does so in a central way ---
it does not break associativity of the OPE algebra.
\end{example}

\subsubsection{\texorpdfstring{Example 4: $W_3$ algebra}{Example 4: W-3 algebra}}

\begin{example}[\texorpdfstring{$W_3$}{W3}: filtered, still strict]\label{ex:w3-strict}
The $W_3$ algebra has generators $L$ (dimension~2, generating Virasoro with central charge~$c$) and $W$ (dimension~3, a primary field), with non-linear OPE: $W(z)W(w)$ involves composite operators.

\emph{Curvature.}
\[\mu_0 = c \cdot \mathbf{1} + \beta \cdot (\text{higher order central terms})\]

Both $c$ and the higher corrections are central, so:
\[d_{\text{bar}}^2 = 0 \quad \text{strictly}\]

Even though $W_3$ is \emph{not quadratic} and requires \emph{filtered}
structure, the curvature is still central, giving strict nilpotence.

\emph{Completion needed.}
For $W_3$, we need \emph{nilpotent completion} (see Appendix~\ref{app:nilpotent-completion}):
\[\widehat{\bar{B}}(W_3) = \varprojlim_n \bar{B}(W_3)/I^n\]
where $I$ is the augmentation ideal.

But once completed, $d_{\text{bar}}^2 = 0$ strictly on the completed complex.
\end{example}

% Duplicate MC subsection commented out.
% Canonical MC definition: def:mc-element-curved (~line 11255).
% Canonical twisting theorem: thm:twisting-mc (~line 11263).
% Alias labels preserved for external references (koszul_reference.tex).
\label{def:maurer-cartan}%
\label{eq:mc-equation}%
\label{thm:mc-deformations}%
\label{thm:mc-periods}%
\iffalse
%%% BEGIN COMMENTED-OUT BLOCK 1: duplicate MC definition and deformation theory
%%% Canonical MC definition lives at def:mc-element-curved (~line 11255).
%%% Canonical twisting theorem lives at thm:twisting-mc (~line 11263).
\subsection{Maurer--Cartan elements and deformations}

\subsubsection{Maurer--Cartan equation}

\begin{definition}[Maurer--Cartan element]\label{def:maurer-cartan-DISABLED}
\index{Maurer--Cartan!chiral}
An element $\alpha \in \mathcal{A}^1$ is a \emph{Maurer--Cartan (MC) element} if it satisfies:
\begin{equation}\label{eq:mc-equation-DISABLED}
m_1(\alpha) + \frac{1}{2} m_2(\alpha \otimes \alpha) + \frac{1}{3!} m_3(\alpha \otimes \alpha
\otimes \alpha) + \cdots + \mu_0 = 0
\end{equation}
\end{definition}

\begin{theorem}[MC elements as quantum deformations; \ClaimStatusProvedHere]\label{thm:mc-deformations-DISABLED}
Maurer--Cartan elements in $\bar{B}^1(\mathcal{A})$ correspond to:
\begin{enumerate}
\item \emph{Physically:} Quantum deformations of the classical algebra
\item \emph{Geometrically:} Flat connections on the associated bundle
\item \emph{Algebraically:} Twisted differentials $d_\alpha = d + [\alpha, -]$
\end{enumerate}

If $\alpha$ is an MC element, then:
\begin{equation}
d_\alpha^2 = 0 \quad \text{strictly} \quad \Longleftrightarrow \quad \mu_0^\alpha := \mu_0
+ m_1(\alpha) + \cdots = 0
\end{equation}
is central.
\end{theorem}

\begin{proof}
Define the twisted differential:
\[d_\alpha := d + m_1(\alpha \otimes -) + m_2(\alpha \otimes \alpha \otimes -) + \cdots\]

Then:
\begin{align}
d_\alpha^2 &= (d + m_1(\alpha \otimes -) + \cdots)^2 \\
&= d^2 + d(m_1(\alpha \otimes -)) + m_1(\alpha \otimes -)^2 + \cdots \\
&= [m_1(\alpha) + \frac{1}{2}m_2(\alpha \otimes \alpha) + \cdots + \mu_0, -]
\end{align}

By the $A_\infty$ relations, this equals:
\[d_\alpha^2 = [\mu_0^\alpha, -]\]
where $\mu_0^\alpha$ is the \emph{twisted curvature}.

Therefore $d_\alpha^2 = 0$ strict if and only if $\mu_0^\alpha$ is central.
\end{proof}

\subsubsection{Geometric realization of MC elements}

\begin{theorem}[MC elements via period integrals; \ClaimStatusProvedHere]\label{thm:mc-periods-DISABLED}
For a chiral algebra $\mathcal{A}$ on a curve $X$ of genus $g$, Maurer--Cartan elements arise
from period integrals:
\begin{equation}
\alpha_g = \int_{\gamma \in H_1(X, \mathbb{Z})} \omega_{\mathcal{A}} \in \bar{B}^1(\mathcal{A})
\end{equation}
where $\omega_{\mathcal{A}} \in \Omega^1(X, \mathcal{A})$ is a connection form.

The MC equation:
\begin{equation}
m_1(\alpha_g) + \frac{1}{2}m_2(\alpha_g \otimes \alpha_g) + \mu_0 = 0
\end{equation}
is equivalent to the \emph{flatness condition}:
\begin{equation}
F_{\omega} := d\omega_{\mathcal{A}} + \frac{1}{2}[\omega_{\mathcal{A}}, \omega_{\mathcal{A}}] = 0
\end{equation}
\end{theorem}

\begin{proof}
We show that period integrals produce MC elements by matching the bar complex MC equation with the flatness condition.

\emph{Step 1: From connection to bar element.}
A connection $\omega_\mathcal{A} \in \Omega^1(X, \mathcal{A})$ on the trivial $\mathcal{A}$-bundle over~$X$ determines an element $\alpha_g \in \bar{B}^1(\mathcal{A})$ via integration over cycles: for each cycle $\gamma_i \in H_1(X, \mathbb{Z})$ ($i = 1, \ldots, 2g$), the period $\oint_{\gamma_i} \omega_\mathcal{A}$ is an element of~$\mathcal{A}$, and $\alpha_g = \sum_i \oint_{\gamma_i} \omega_\mathcal{A} \otimes \gamma_i^*$ is the corresponding degree-1 bar element (where $\gamma_i^*$ is the dual 1-form).

\emph{Step 2: MC equation $=$ flatness.}
The MC equation in $\bar{B}(\mathcal{A})$ reads:
\[
m_1(\alpha_g) + \tfrac{1}{2} m_2(\alpha_g \otimes \alpha_g) + \mu_0 = 0.
\]
Under the identification of Step~1:
\begin{itemize}
\item $m_1(\alpha_g)$ corresponds to $d\omega_\mathcal{A}$: the bar differential $m_1$ acts on $\alpha_g$ by applying the internal differential of~$\mathcal{A}$ and the de Rham differential on~$X$, which together give the exterior derivative of the connection form.
\item $\tfrac{1}{2} m_2(\alpha_g \otimes \alpha_g)$ corresponds to $\tfrac{1}{2}[\omega_\mathcal{A}, \omega_\mathcal{A}]$: the binary bar operation $m_2$ encodes the chiral product, which on the connection form becomes the Lie bracket (the OPE residue $\mathrm{Res}_{z_1=z_2}[\omega(z_1) \otimes \omega(z_2) \cdot \eta_{12}]$ extracts the commutator $[\omega, \omega]$).
\item $\mu_0$ is the curvature of the $A_\infty$ structure, which vanishes for strict (uncurved) algebras and equals the central charge term for VOAs.
\end{itemize}
Therefore the MC equation is exactly the flatness condition $F_\omega = d\omega_\mathcal{A} + \tfrac{1}{2}[\omega_\mathcal{A}, \omega_\mathcal{A}] + \mu_0 = 0$.

\emph{Step 3: Well-definedness.}
The element $\alpha_g$ depends only on the cohomology class of $\omega_\mathcal{A}$ modulo exact forms: if $\omega_\mathcal{A} \mapsto \omega_\mathcal{A} + df$ for $f \in \Gamma(X, \mathcal{A})$, then $\alpha_g \mapsto \alpha_g + m_1(f)$, which is a gauge equivalence in the MC moduli space.
\end{proof}

\begin{example}[Genus 1 MC obstruction for Heisenberg]
At genus 1, the elliptic curve $E_\tau$ has coordinate $z$ with $z \sim z + 1 \sim z + \tau$.

The Heisenberg current $J$ has connection form $\omega_J = J\,dz$.  The candidate MC element is:
\[\alpha_1 = \oint_A J\,dz \cdot a + \oint_B J\,dz \cdot b\]
where $A, B$ are the standard cycles of $E_\tau$ with $\oint_A dz = 1$, $\oint_B dz = \tau$.

The curved MC equation $\mu_0 + m_1(\alpha_1) + m_2(\alpha_1, \alpha_1) + \cdots = 0$ becomes:
\begin{align}
\mu_0 + m_1(\alpha_1) &= k + 0 = k \neq 0
\end{align}
since $m_1(\alpha_1) = 0$ by the conservation law $dJ = 0$.

The MC equation \emph{fails}: the curvature $\mu_0 = k$ is a cohomologically non-trivial obstruction.  No MC element can ``untwist'' it to zero.  This is consistent with the fact that the Heisenberg bar differential satisfies $d_{\mathrm{bar}}^2 = 0$ strict (Theorem~\ref{thm:central-implies-strict}), but the \emph{curvature itself is essential} and cannot be removed by twisting.
\end{example}
%%% END COMMENTED-OUT BLOCK 1
\fi

\subsection{Obstruction theory: genus-by-genus analysis}

\begin{theorem}[Strict nilpotence at genus zero; \ClaimStatusProvedHere]\label{thm:genus-zero-strict}
Let $\mathcal{A}$ be a chiral algebra with central curvature. Then at genus~$0$,
the bar differential satisfies $d_0^2 = 0$ strict.
\end{theorem}

\begin{proof}
At genus 0, the bar differential is
$d_0 = d_{\text{internal}} + d_{\text{residue}}$
with no quantum corrections ($\mu_0 = 0$ at genus 0).
By the Arnold relations (Theorem~\ref{thm:arnold-three}), $d_0^2 = 0$.
\end{proof}

\begin{theorem}[Strict nilpotence at all genera; \ClaimStatusProvedHere]\label{thm:genus-induction-strict}
\textup{[Regime: curved-central\linebreak
\textup{(}Convention~\textup{\ref{conv:regime-tags})}.]}

Let $\mathcal{A}$ be a chiral algebra with central curvature at all genera,
i.e., the genus-$g$ curvature $\mu_0^{(g)} \in Z(\mathcal{A})$ for all $g \geq 1$.
Then the full bar differential $d_{\mathrm{full}} = \sum_{g \geq 0} d_g$ satisfies
$d_{\mathrm{full}}^2 = 0$ strict.
Here each $d_g$ denotes the genus-$g$ \emph{component} of the modular
decomposition---not the fiberwise differential $\dfib$ of
Convention~\textup{\ref{conv:higher-genus-differentials}}.
\end{theorem}

\begin{proof}
The full bar differential decomposes as
$d_{\mathrm{full}} = d_0 + \sum_{g \geq 1} d_g$,
where $d_0 = d_{\mathrm{int}} + d_{\mathrm{res}}$ is the genus-$0$ bar
differential and each $d_g$ ($g \geq 1$) arises from the boundary strata of
$\overline{\mathcal{M}}_{g,n}$.  We compute $d_{\mathrm{full}}^2$ by organizing
the terms by total genus.

\medskip
\emph{Step 1: The modular operad boundary relation.}
The collection $\{\overline{\mathcal{M}}_{g,n}\}_{g,n}$ forms a modular operad
whose composition maps are given by gluing along boundary divisors.  The
boundary of $\overline{\mathcal{M}}_{g,n}$ is a normal crossings divisor with
strata indexed by dual graphs $\Gamma$.  Codimension-$2$ strata arise from
refinements $\Gamma' \to \Gamma$ of dual graphs.  The fundamental topological
identity $\partial^2 = 0$ on the stratification of $\overline{\mathcal{M}}_{g,n}$
states that each codimension-$2$ stratum appears with opposite signs from the two
codimension-$1$ strata that contain it.  This is equivalent to the \emph{modular
operad axiom}: the composition in the modular operad is associative.

\medskip
\emph{Step 2: Genus-$0$ nilpotence ($d_0^2 = 0$).}
This is Theorem~\ref{thm:genus-zero-strict}: the Arnold relations
(Theorem~\ref{thm:arnold-three}) provide the codimension-$2$ cancellation on
$\overline{C}_n(X) = \overline{\mathcal{M}}_{0,n+1}$.

\medskip
\emph{Step 3: Cross-terms ($d_0 d_g + d_g d_0 = 0$).}
Each $d_g$ inserts a genus-$g$ correction $\mu_0^{(g)} \otimes \omega_g$,
where $\omega_g \in H^*(\overline{\mathcal{M}}_{g,n})$ and
$\mu_0^{(g)} \in Z(\mathcal{A})$ by hypothesis.  The cross-term
$d_0 d_g + d_g d_0$ corresponds to codimension-$2$ boundary strata of
$\overline{\mathcal{M}}_{g,n+k}$ that involve one genus-$0$ degeneration
(a collision of points on the same component) and one genus-$g$ operation.
By the modular operad axiom (Step~1), these strata cancel pairwise:
\[
  d_0 d_g + d_g d_0 = \sum_{\substack{\text{codim-2 strata} \\ \text{(genus-0/genus-$g$)}}}
  (\pm 1)\, [\text{stratum}] = 0.
\]
Centrality of $\mu_0^{(g)}$ is essential here: it ensures that the algebraic
insertion $\mu_0^{(g)} \otimes (-)$ commutes with the internal differential
$d_{\mathrm{int}}$, so that the signs from the boundary orientation match the
signs from the Koszul sign rule.

\medskip
\emph{Step 4: Pure higher-genus terms ($\sum_{g_1 + g_2 = g} d_{g_1} d_{g_2} = 0$).}
The product $d_{g_1} d_{g_2}$ inserts two corrections
$\mu_0^{(g_1)}$ and $\mu_0^{(g_2)}$ into the bar complex.  Since both lie in
$Z(\mathcal{A})$, the order of insertion is irrelevant: the insertions commute.
The sum over all decompositions $g_1 + g_2 = g$ corresponds to the codimension-$2$
boundary strata of $\overline{\mathcal{M}}_{g,n}$ that arise from two
\emph{separating} degenerations (splitting the surface into components of genera
$g_1$ and $g_2$).  Non-separating degenerations at genus $g$ (which increase the
genus of a single component by~$1$) also contribute; these correspond to
self-gluing operations in the modular operad.

For each pair of codimension-$1$ divisors $D_1, D_2$ meeting in a codimension-$2$
stratum $D_1 \cap D_2$, the boundary-of-boundary identity assigns opposite signs
to the two orderings $D_1 \to D_1 \cap D_2$ and $D_2 \to D_1 \cap D_2$.
Centrality ensures the algebraic insertions respect these orientations (no
additional sign from non-commutativity), and so:
\[
  \sum_{g_1 + g_2 = g} d_{g_1} d_{g_2}
  = \sum_{\substack{\text{codim-2 strata of}\\
      \overline{\mathcal{M}}_{g,n}}} (\pm 1)\, [\text{stratum}] = 0.
\]

\medskip
\emph{Step 5: Conclusion.}
Collecting all terms:
\[
  d_{\mathrm{full}}^2 = d_0^2
  + \sum_{g \geq 1} (d_0 d_g + d_g d_0)
  + \sum_{g \geq 2}\sum_{g_1 + g_2 = g} d_{g_1} d_{g_2}
  = 0 + 0 + 0 = 0.
\]
The first term vanishes by Step~2, the second by Step~3, and the third by Step~4.
This completes the proof.
\end{proof}

\begin{remark}[Computational verification]\label{rem:genus-induction-status}
The proof above uses the modular operad structure of
$\{\overline{\mathcal{M}}_{g,n}\}$ in an essential way.  We have also verified
the conclusion computationally through genus~$5$ for all standard examples
(Heisenberg, Kac--Moody, Virasoro, $W$-algebras), providing independent
confirmation.
\end{remark}

\subsection{Summary}

\begin{center}
\begin{tabular}{|>{\raggedright\arraybackslash}p{3.5cm}|>{\raggedright\arraybackslash}p{4cm}|>{\raggedright\arraybackslash}p{6cm}|}
\hline
\textbf{Regime} & \textbf{Condition} & \emph{Examples} \\
\hline
\hline
\textbf{Strict Nilpotence} & $\mu_0 \in Z(\mathcal{A})$ & 
\begin{itemize}[leftmargin=*]
\item Heisenberg $\mathcal{H}_k$
\item Kac--Moody $\widehat{\mathfrak{g}}_k$
\item Virasoro $\text{Vir}_c$
\item $W$-algebras $W_N$
\item Free fermions $\beta\gamma$
\end{itemize}
$d_{\text{bar}}^2 = 0$ strictly (Theorem~\ref{thm:genus-induction-strict}) \\
\hline
\textbf{Curved (Non-Central)} & $\mu_0 \notin Z(\mathcal{A})$ & 
\begin{itemize}[leftmargin=*]
\item Hypothetical non-central extensions
\item Some deformed algebras
\end{itemize}
$d_{\text{bar}}^2 \neq 0$, need higher homotopies \\
\hline
\textbf{Homotopy Coherent} & No curvature, but $d^2 \sim 0$ only & 
\begin{itemize}[leftmargin=*]
\item $(\infty,1)$-categorical structures
\item Derived geometry settings
\item Non-algebraic field theories
\end{itemize}
Requires full $A_\infty$ or $L_\infty$ framework \\
\hline
\end{tabular}
\end{center}

\begin{remark}\label{rem:why-strict-matters}
The distinction between strict and homotopy nilpotence has significant consequences:

\emph{Computations.}
Strict nilpotence allows direct cohomology computation (Theorem~\ref{thm:genus-induction-strict}); homotopy nilpotence requires spectral sequences that may not degenerate.

\emph{Convergence.}
Strict nilpotence ensures the bar-cobar adjunction (Theorem~\ref{thm:bar-cobar-adjunction}) works without completion (for quadratic algebras); homotopy nilpotence necessitates completion (Theorem~\ref{thm:completion-necessity}), with attendant convergence issues.

\emph{Physics.}
Strict nilpotence means quantum corrections are controlled by central charges (Theorem~\ref{thm:deformation-obstruction}); homotopy nilpotence requires full renormalization group analysis.

\emph{Good news.} Vertex operator algebras (with the standard vacuum axiom) and
chiral algebras arising from unitary CFT have \emph{central curvature}
(Remark~\ref{rem:voa-central-curvature}; see also Theorem~\ref{thm:genus-induction-strict}), so we are in the strict regime.
\end{remark}

\subsection{Connection to literature}

\subsubsection{Gui--Li--Zeng (2022)}

In \cite{GLZ22}, Gui--Li--Zeng develop the theory of curved Koszul duality for chiral algebras. 
Their key result:

\begin{theorem}[GLZ, Theorem 5.3; \ClaimStatusProvedElsewhere]\label{thm:glz-curved}
For a quadratic chiral algebra $\mathcal{A}$ with central curvature $\mu_0 \in Z(\mathcal{A})$:
\begin{enumerate}
\item The Koszul dual $\mathcal{A}^!$ exists as a curved cooperad
\item The bar-cobar adjunction holds: $\Omega(B(\mathcal{A})) \simeq \mathcal{A}$
\item The equivalence is an isomorphism in the derived category
\end{enumerate}
\end{theorem}

Our Theorem~\ref{thm:central-implies-strict} provides the \emph{geometric realization} 
of their algebraic result.

\subsubsection{Francis--Gaitsgory}

Francis--Gaitsgory~\cite{FG12} prove that factorization algebras satisfy a 
bar-cobar duality. Their result:

\begin{theorem}[FG, Theorem 7.2.1; \ClaimStatusProvedElsewhere]
\label{thm:fg-factorization-bar-cobar}
For a factorization algebra $\mathcal{F}$ on a curve $X$:
\[\text{Fact}(X, \Omega(B(\mathcal{F}))) \simeq \mathcal{F}\]
\end{theorem}

Combined with our explicit bar construction (Definition~\ref{def:geometric-bar}), this confirms 
that our geometric bar differential has the correct homological properties.

\subsubsection{Costello--Gwilliam}

In \cite{CG-vol2}, Costello--Gwilliam use curved structures to study BV quantization with anomalies, renormalization in perturbative QFT, and effective field theories with central charges.

Their MC equation (Definition 3.2.1.1 in \cite{CG-vol2}) is:
\[\delta I + \frac{1}{2}\{I, I\} = 0\]
for the quantum effective action $I$.

This is our equation~\eqref{eq:mc-equation} in the field theory context.

The connection is that central charges in QFT correspond to central curvature in chiral algebras;
both ensure that quantum corrections do not destroy associativity/nilpotence.

\subsection{Computational corollaries}

\begin{corollary}[Bar cohomology computes Ext {\cite{LV12}}; \ClaimStatusProvedElsewhere]\label{cor:bar-computes-ext}
For a chiral algebra $\mathcal{A}$ with central curvature:
\begin{equation}
H^*(\bar{B}(\mathcal{A}), d_{\text{bar}}) = \text{Ext}_{\mathcal{A}}^*(\mathbb{C}, \mathbb{C})
\end{equation}
and this can be computed directly without spectral sequences.
\end{corollary}

\begin{corollary}[Koszul dual coalgebra {\cite{GK94}}; \ClaimStatusProvedElsewhere]\label{cor:koszul-dual-cooperad}
For quadratic $\mathcal{A}$ with central curvature, the bar
coalgebra $\barB(\cA)$ has cohomology
\begin{equation}
H^*(\bar{B}(\mathcal{A}))
\end{equation}
concentrated in bar degree~$1$ (the Koszul dual coalgebra), carrying a curved cooperad structure with comultiplication dual to~$m_2$ and curvature dual to~$\mu_0$.  The Koszul dual \emph{algebra}~$\cA^!$ is characterized by Verdier intertwining: $\mathbb{D}_{\operatorname{Ran}}\barB_X(\cA) \simeq \barB_X(\cA^!)$ (Convention~\ref{conv:bar-coalgebra-identity}).
\end{corollary}

\begin{corollary}[Genus expansion convergence; \ClaimStatusProvedHere]\label{cor:genus-expansion-converges}
The genus expansion:
\begin{equation}
Z(\mathcal{A}) = \sum_{g=0}^\infty \hbar^{2g-2} Z_g(\mathcal{A})
\end{equation}
where $Z_g(\mathcal{A}) = H^*(\bar{B}^{(g)}(\mathcal{A}), d^{(g)})$, converges in the
$\hbar$-adic topology of the formal power series ring $\mathbb{C}[[\hbar]]$.
\end{corollary}

\begin{proof}
Each term $Z_g(\mathcal{A})$ is well-defined because
$(d^{(g)})^2 = 0$ strictly at genus $g$
(Theorem~\ref{thm:genus-induction-strict}), so the cohomology
$H^*(\bar{B}^{(g)}(\mathcal{A}))$ exists without formal completion.
The formal power series $\sum_{g \geq 0} \hbar^{2g-2} Z_g$ is an element of
$\hbar^{-2} \mathbb{C}[[\hbar^2]]$, which is a well-defined ring (Laurent series
with a pole of bounded order).  Convergence in the $\hbar$-adic topology is
automatic: for every $N$, the partial sum $\sum_{g=0}^N \hbar^{2g-2} Z_g$
agrees with $Z(\mathcal{A})$ modulo $\hbar^{2N}$.

Note that $d^2 = 0$ holds \emph{strictly} at each genus (not merely
up to $\hbar$-corrections), so each $Z_g$ is an honest finite-dimensional vector
space, not a formal deformation.  This is the strict nilpotence established in
Section~\ref{sec:strict-vs-homotopy}.
\end{proof}

\subsection{Perspectives on strict nilpotence}

\subsubsection{Physical perspective}

Associativity of the OPE algebra reflects locality: the OPE $(AB)C = A(BC)$ follows from the consistency of correlation functions as operators approach each other. Quantum corrections (loop diagrams) do not break locality, so they enter as central charges that modify the normalization but preserve associativity. For example, in 2D CFT, the central charge $c$ in the Virasoro OPE
\[T(z)T(w) = \frac{c/2}{(z-w)^4} + \cdots\]
is a quantum correction (zero classically) that is central and therefore does not affect the Jacobi identity.

\subsubsection{Geometric perspective}

Kontsevich's formality theorem~\cite{Kon03} shows that:
\[\text{Poly}_{\bullet}(M)[[t]] \xrightarrow{\sim} \text{Dpoly}_{\bullet}(M)[[t]]\]
via explicit configuration space integrals.

The integrals:
\[U_\Gamma = \int_{C_n(\mathbb{R}^d)} \omega_\Gamma\]
over configuration spaces satisfy:
\[\sum_\Gamma U_\Gamma \cdot (\text{boundary terms}) = 0\]
by Stokes' theorem.

This is our strict nilpotence $d^2 = 0$. The Arnold relations are the
genus-0 case of this pattern.

At higher genus, the same Stokes argument works because curvature terms are central and do not
interfere with the boundary calculations.

\subsubsection{Computational verification}

Nilpotence holds degree by degree: $d^2(a \otimes b) = 0$ by direct calculation (degree~2), $d^2(a \otimes b \otimes c) = 0$ using Arnold relations (degree~3), and similarly through degree~5 using extended Arnold relations.  The centrality of $\mu_0$ is used at each stage.

\subsubsection{Functorial perspective}

Strict nilpotence also follows from functoriality:

\begin{theorem}[Functoriality of bar construction; \ClaimStatusProvedHere]\label{thm:bar-functorial-grothendieck}
The bar construction:
\[B: \mathrm{ChAlg}^{\mathrm{central}} \to \mathrm{Coalg}\]
is a functor from chiral algebras with central curvature to coalgebras, characterized by the
universal property:
\[\mathrm{Hom}_{\mathrm{Coalg}}(B(\mathcal{A}), C) \simeq \mathrm{Hom}_{\mathrm{ChAlg}}(\mathcal{A},
\Omega(C))\]

This adjunction automatically implies $d_{\mathrm{bar}}^2 = 0$ by the universal property.
\end{theorem}

\begin{proof}
The proof assembles three previously established results into the claimed adjunction.

\emph{Functoriality.}
By Theorem~\ref{thm:bar-functorial}, the bar construction $B$ is functorial: a morphism $f: \mathcal{A} \to \mathcal{B}$ of chiral algebras induces $B(f): B(\mathcal{A}) \to B(\mathcal{B})$ compatible with the coalgebra structures.  The category $\mathrm{ChAlg}^{\mathrm{central}}$ consists of chiral algebras whose curvature $\mu_0$ lies in the center $Z(\mathcal{A})$.

\emph{Adjunction.}
The bar-cobar adjunction $B \dashv \Omega$ is established in Theorem~\ref{thm:bar-cobar-adjunction}.  For a coalgebra~$C$ and a chiral algebra~$\mathcal{A}$, a coalgebra map $\phi: B(\mathcal{A}) \to C$ determines a chiral algebra map $\tilde{\phi}: \mathcal{A} \to \Omega(C)$ by the universal property of the tensor coalgebra: $\phi$ restricts to a map on cogenerators $\bar{\phi}: s\mathcal{A} \to \bar{C}$, and $\tilde{\phi} = s^{-1}\bar{\phi}$ is the corresponding algebra map.  Conversely, $\tilde{\phi}: \mathcal{A} \to \Omega(C)$ extends uniquely to $\phi: B(\mathcal{A}) \to C$ by the cofree property of the tensor coalgebra.

\emph{$d^2 = 0$ from the adjunction.}
The bar differential $d_B$ on $B(\mathcal{A})$ satisfies $d_B^2 = 0$ when the curvature $\mu_0$ is central: by Theorem~\ref{thm:central-implies-strict}, the centrality $\mu_0 \in Z(\mathcal{A})$ implies that the curvature-induced term $[\mu_0, -]$ in $d_B^2$ vanishes identically.  From the Grothendieck perspective, this is automatic: $d_B$ is the unique coderivation on $B(\mathcal{A})$ whose projection to cogenerators recovers the chiral algebra structure maps.  The condition $d_B^2 = 0$ is equivalent to the $A_\infty$ relations on $\mathcal{A}$, which reduce to the associativity of the chiral product when $\mu_0$ is central (the curved $A_\infty$ relations $\sum m_{r+1+t}(\mathrm{id}^{\otimes r} \otimes m_s \otimes \mathrm{id}^{\otimes t}) = 0$ collapse because the terms involving $m_0 = \mu_0$ commute with everything).
\end{proof}

\begin{remark}[Dependencies]
This proof uses three ingredients:
\begin{enumerate}[label=(D\arabic*)]
\item Concrete bar-functor functoriality (Theorem~\ref{thm:bar-functorial}).
\item The bar-cobar adjunction framework (Theorem~\ref{thm:bar-cobar-adjunction}).
\item Central curvature implies strict nilpotence (Theorem~\ref{thm:central-implies-strict}).
\end{enumerate}
\end{remark}

Once centrality of $\mu_0$ ensures the adjunction exists, the nilpotence $d^2 = 0$ follows formally.

\subsection{Strict versus homotopy nilpotence}

\begin{remark}[Strict vs homotopy nilpotence]
For chiral algebras $\mathcal{A}$ arising from vertex operator algebras
(satisfying the vacuum axiom $Y(\mathbf{1},z) = \mathrm{id}$) or from
conformal field theories whose anomalies are central, we have $d_{\text{bar}}^2 = 0$ strictly at all genera. The ingredients are: (1) the curvature $\mu_0$ is proportional to the vacuum $\mathbf{1}$, hence central by the vacuum axiom (Remark~\ref{rem:voa-central-curvature}); (2) Arnold relations ensure $d_{\mathrm{residue}}^2 = 0$; (3) the Leibniz rule gives the anticommutation of $d_{\mathrm{internal}}$ and $d_{\mathrm{residue}}$; (4) the modular operad boundary relations on $\overline{\mathcal{M}}_{g,n}$ extend this to all genera (Theorem~\ref{thm:genus-induction-strict}).

This covers all standard examples (Heisenberg, Kac--Moody, Virasoro, $W$-algebras, free field systems). For non-unital chiral algebras or theories with non-central anomalies, $d^2 = 0$ may hold only up to homotopy.
\end{remark}


\section{Non-quadratic chiral algebras}
\label{sec:filtered-vs-curved-comprehensive}

\label{thm:bar-convergence}% alias: canonical theorem inside commented-out block below
The four-regime hierarchy (Convention~\ref{conv:regime-tags})
classifies chiral algebras by their bar-complex behavior:
quadratic (strict $d^2=0$), curved-central ($\dfib^2 = \kappa\cdot\omega_g$
with central curvature), filtered-complete (requiring I-adic completion),
and programmatic (infinite generators, MC4 frontier).
The examples of each regime are given in
\S\ref{subsec:three-regimes} above.

\iffalse
%%% BEGIN COMMENTED-OUT: Four-class exposition (~600 lines)
%%% This section duplicated the three-regime classification at
%%% \S\ref{subsec:three-regimes} with a fourth "general" class.
%%% The canonical classification is Convention~\ref{conv:regime-tags}.
Four increasingly general classes form a hierarchy
\begin{equation}
\text{Quadratic} \subset \text{Curved} \subset \text{Filtered} \subset \text{General},
\end{equation}
and the techniques needed at each level differ sharply: quadratic
algebras admit direct bar-cobar duality without completion; curved
algebras require completed bar-cobar but preserve the adjunction;
filtered algebras always need nilpotent completion
(Appendix~\ref{app:nilpotent-completion}); and general algebras may
lie beyond the present theorematic reach (e.g., $\mathcal{W}_\infty$,
whose standard tower is theorematic while the remaining frontier is the
exact MC4 H-level / coefficient-identification package).
The passage from quadratic to general is also the passage from $E_\infty$
to $E_1$: the Yangian $Y(\mathfrak{g})$, which is $E_1$-chiral but
not $E_\infty$, falls in the filtered class and requires the
completed framework of Theorem~\ref{thm:en-koszul-duality}.

\subsection{Definitions: four classes of chiral algebras}

\subsubsection{Class I: quadratic chiral algebras}

\begin{definition}[Quadratic chiral algebra]\label{def:quadratic-chiral}
The Heisenberg algebra $\mathcal{H}_k$ is the prototype: a single generator $\alpha$ with a purely quadratic OPE $\alpha(z)\alpha(w) = k/(z-w)^2 + \mathrm{reg.}$ (\S\ref{sec:frame-bar-deg1}).  The following definition captures this structure in general.

A chiral algebra $\mathcal{A}$ is \emph{quadratic} if it admits a presentation:
\begin{equation}
\mathcal{A} = \text{Free}_{\text{ch}}(V) / (R)
\end{equation}
where:
\begin{itemize}
\item $V$ is a graded vector space of generators
\item $R \subset V \otimes V$ consists of \emph{quadratic} relations only
\item No higher-order relations (no terms in $V^{\otimes n}$ for $n \geq 3$)
\end{itemize}
\end{definition}

\begin{example}[Heisenberg: curved quadratic]\label{ex:heisenberg-quadratic}
The Heisenberg algebra $\mathcal{H}_k$ has a single generator $J$ (the current) with OPE $J(z)J(w) \sim k/(z-w)^2$ and no simple pole (the Lie bracket vanishes).

The double pole produces a \emph{curved} quadratic structure: the relation $J \otimes J \sim k \cdot \mathbf{1}$ is inhomogeneous (the right side is degree~0, not degree~2), so the bar complex has curvature $m_0 = -k \cdot c \neq 0$ for $k \neq 0$.  In the language of Definition~\ref{def:quadratic-chiral}, the Heisenberg is ``quadratic'' in the sense that the OPE involves at most two generators on the left side, but it is \emph{not} strictly quadratic (i.e., $m_0 \neq 0$).

\emph{Koszul dual.}
\[\mathcal{H}_k^! = \mathrm{Sym}^{\mathrm{ch}}(V^*) \quad \text{(commutative chiral algebra on the dual space, with curvature)}\]

Since $\mathfrak{h}$ is abelian, the chiral Koszul dual is $\mathrm{Sym}^{\mathrm{ch}}(V^*)$ with $d = 0$ and curvature $m_0 = -k \cdot c$ (level negates; $h^\vee = 0$).  This is a \emph{commutative} chiral algebra with curvature, not the Heisenberg at negated level.  The bar differential vanishes (Example~\ref{ex:ope-to-residue}), so $H^*(\bar{B}(\mathcal{H}_k)) \simeq \bar{B}(\mathcal{H}_k)$ with CE cooperad structure.

\emph{Contrast with Free Fermions}:
\begin{center}
\begin{tabular}{|l|l|l|l|}
\hline
\emph{Algebra} & \textbf{OPE Pole} & \textbf{Bar Differential} & \textbf{Effect} \\
\hline
Free Fermion $\psi$ & Simple: $\frac{1}{z-w}$ & Non-zero & contracts fields \\
Heisenberg $J$ & Double: $\frac{k}{(z-w)^2}$ & Zero (triple pole) & residue vanishes \\
\hline
\end{tabular}
\end{center}

No completion is needed. The bar complex is conilpotent (finite-dimensional in each bar degree) and the differential is zero except for the curvature term $m_0$. Convergence of the bar spectral sequence follows: the filtration by bar degree is bounded below ($n \geq 0$) and exhaustive, with finite-dimensional graded pieces, so the spectral sequence converges to the bar homology by the standard bounded-below complete convergence theorem (\cite{LV12}, Theorem~2.2.3).

\emph{Reference}: See \cite{GLZ22} Section 6, Proposition 6.2 for the identification
of the twisted chiral enveloping algebra with CE algebra structure.
\end{example}

\begin{example}[Affine Kac--Moody: quadratic]\label{ex:km-quadratic}
For $\widehat{\mathfrak{g}}_k$ (affine Lie algebra at level $k$), the generators are $V = \mathfrak{g}$ with relations $R = \{J^a \otimes J^b - f^{abc}J^c - k \delta^{ab} \mathbf{1}\}$.

The OPE is:
\[J^a(z)J^b(w) = \frac{k \delta^{ab}}{(z-w)^2} + \frac{f^{abc}J^c(w)}{z-w} + \text{regular}\]

This is quadratic: only products of two currents appear, and the structure constants $f^{abc}$ are linear in $J^c$.

\emph{Koszul dual.}
\[\bar{B}^{\mathrm{ch}}(\widehat{\mathfrak{g}}_k) \simeq \mathrm{CE}^{\mathrm{ch},c}(\mathfrak{g}_{-k-2h^\vee})
\quad \text{(curved CE coalgebra at shifted level)}\]
and applying the cobar functor recovers the vertex algebra $\widehat{\mathfrak{g}}_{-k-2h^\vee}$.

No completion is needed.
\end{example}

\subsubsection{Class II: curved (non-quadratic) chiral algebras}

\begin{definition}[Curved chiral algebra]\label{def:curved-chiral-detailed}
A chiral algebra $\mathcal{A}$ is \emph{curved} (but not necessarily quadratic) if:
\begin{enumerate}
\item It has a presentation $\mathcal{A} = \text{Free}_{\text{ch}}(V) / (R)$
\item The relations $R$ may involve terms in $V^{\otimes n}$ for $n \geq 3$
\item There exists a \emph{central curvature element} $\mu_0 \in Z(\mathcal{A})^2$
\item The curvature satisfies the MC equation: 
\[\sum_{k=0}^\infty \frac{1}{k!} m_k(\mu_0^{\otimes k}) = 0\]
\end{enumerate}
\end{definition}

\begin{example}[Virasoro: curved, non-quadratic]\label{ex:virasoro-curved}
The Virasoro algebra $\text{Vir}_c$ has:
\begin{itemize}
\item Generators: $V = \mathbb{C} \cdot T$ (stress tensor)
\item Quadratic part: $T \otimes T \sim \frac{c}{(z-w)^4} + \frac{2T}{(z-w)^2}$
\item \emph{Cubic term}: $m_3(T \otimes T \otimes T) \neq 0$ (Schwarzian derivative)
\end{itemize}

The OPE is:
\[T(z)T(w) = \frac{c/2}{(z-w)^4} + \frac{2T(w)}{(z-w)^2} + \frac{\partial T(w)}{z-w} + \text{regular}\]

The algebra is curved because the central charge $c$ is a curvature
($\mu_0 = c \cdot \mathbf{1}$) which is central ($[c, T] = 0$) and
satisfies $m_1(c) = 0$. It is non-quadratic because the stress tensor
has a non-zero triple product $T(z)T(w)T(u) \sim \cdots$ encoded by the
Schwarzian derivative.

In the present completed formalism, one tracks the same-family partner
determined by the central-charge involution
$c \leftrightarrow 26-c$
\textup{(}Example~\textup{\ref{ex:virasoro-koszul-dual}}\textup{)}.
This supplies the model/shadow object used in the genus calculations,
while the stronger infinite-generator realization is delegated to the
exact $W_\infty$ MC4 package of filtered H-level targets, residue
identities, and finite detection on $\mathcal{I}_N$, whose first live
higher-spin stages are already reduced on the theorem surface to the
stage-$3$ fifteen-coefficient packet and the exact stage-$4$ six-entry
identity packet on~$\mathcal{I}_4$, organized as three local OPE blocks
and splitting into four higher-spin channels together with two
Virasoro-target channels whose principal Drinfeld--Sokolov values are
fixed and whose residue-side contraction is the extra input isolated by
Corollary~\ref{cor:winfty-stage4-residue-four-channel}.
Completion is needed: the non-quadratic relations require
\[\widehat{\bar{B}}(\text{Vir}_c) = \varprojlim_n \bar{B}(\text{Vir}_c)/I^n\]
where $I$ is the augmentation ideal.
\end{example}

\subsubsection{Class III: filtered chiral algebras}

\begin{definition}[Filtered chiral algebra]\label{def:filtered-chiral}
A chiral algebra $\mathcal{A}$ is \emph{filtered} if it carries a filtration:
\begin{equation}
F_0\mathcal{A} \subset F_1\mathcal{A} \subset F_2\mathcal{A} \subset \cdots \subset \mathcal{A}
\end{equation}
satisfying:
\begin{enumerate}
\item \emph{Multiplicativity}: $m_2(F_i\mathcal{A} \otimes F_j\mathcal{A}) \subset F_{i+j}\mathcal{A}$
\item \emph{Exhaustive}: $\mathcal{A} = \bigcup_{i=0}^\infty F_i\mathcal{A}$
\item \emph{Separated}: $\bigcap_{i=0}^\infty F_i\mathcal{A} = 0$
\item \emph{Complete}: $\mathcal{A} \cong \varprojlim_n \mathcal{A}/F_n\mathcal{A}$
\end{enumerate}

The \emph{associated graded} is:
\begin{equation}
\text{gr}(\mathcal{A}) = \bigoplus_{i=0}^\infty F_i\mathcal{A}/F_{i-1}\mathcal{A}
\end{equation}
\end{definition}

\begin{example}[\texorpdfstring{$W_3$}{W3}: filtered, non-curved]\label{ex:w3-filtered}
The $W_3$ algebra has generators $L$ (dimension~2, generating the Virasoro subalgebra) and $W$ (dimension~3, a primary field).

\emph{Filtration by operator dimension.}
\begin{align}
F_0W_3 &= \mathbb{C} \cdot \mathbf{1} \\
F_1W_3 &= \mathbb{C} \cdot \mathbf{1} \oplus \mathbb{C} \cdot \partial L \oplus \cdots \\
F_2W_3 &= F_1W_3 \oplus \mathbb{C} \cdot L \oplus \mathbb{C} \cdot \partial^2 L \oplus \cdots \\
F_3W_3 &= F_2W_3 \oplus \mathbb{C} \cdot W \oplus \mathbb{C} \cdot (L \cdot L) \oplus \cdots
\end{align}

The OPE is non-linear:
\[W(z)W(w) = \frac{\cdots}{(z-w)^6} + \cdots + \frac{\Lambda(L \cdot L)(w)}{(z-w)^2} + \cdots\]
where $\Lambda(L \cdot L)$ is a composite operator, not a single generator. The algebra is filtered rather than curved because it is not generated by a finite-dimensional space $V$: composite operators like $(L \cdot L)$ appear at all levels, and the filtration is infinite-dimensional at each level. The Koszul dual is
\[W_3^! = \widehat{\mathrm{CoW}_3}
\quad \text{(completed cooperad structure)}.\]
The bar construction must be completed:
\[\widehat{\bar{B}}(W_3) = \varprojlim_n \bar{B}(W_3)/F_n\]
where $F_n$ is the filtration by operator dimension.
\end{example}

\subsubsection{Class IV: general (exact MC4 frontier)}

\begin{example}[\texorpdfstring{$W_\infty$}{W_infinity}: exact MC4 frontier]\label{ex:winfty-completion-frontier-comprehensive}
The $W_\infty$ algebra has:
\begin{itemize}
\item Generators: $W^{(n)}$ for all $n \geq 2$ (infinitely many generators)
\item Relations: Infinitely many non-linear relations
\item No finite presentation
\end{itemize}

No monograph-wide universal H-level duality theorem is currently
proved: the generating space $V$ is infinite-dimensional, and the
manuscript does not yet construct the filtered factorization object
whose finite quotients recover the principal stages and satisfy the
exact residue-identification package.

The frontier is now narrower.  For the standard principal-stage tower,
Corollary~\ref{cor:winfty-standard-mc4-package} already gives the
standard-tower completed M-level bar-cobar package.  The remaining
problem is not to construct a first completed bar object, but to
compare that standard completion with stronger H-level
factorization-theoretic or physical realizations of $W_\infty$ by
proving the explicit identities
$\mathsf{C}^{\mathrm{res}}_{s,t;u;m,n}(N)
=\mathsf{C}^{\mathrm{DS}}_{s,t;u;m,n}(N)$
and reducing them to the finite-detection packet $\mathcal{I}_N$, whose
first live higher-spin stages are already reduced on the theorem
surface to the stage-$3$ fifteen-coefficient packet and the exact
stage-$4$ six-entry identity packet on~$\mathcal{I}_4$, organized as
three local OPE blocks and splitting into four higher-spin channels
together with two Virasoro-target channels whose principal values are
fixed and whose residue-side contraction is the extra input isolated by
Corollary~\ref{cor:winfty-stage4-residue-four-channel}.

\emph{Physical interpretation.}
$W_\infty$ describes non-local interactions in 2D gravity.  The open problem is
to construct the correct filtered H-level dual description and verify
its exact finite-stage coefficient identities, not to prove that no such
description can exist.
\end{example}

\subsection{Comparison of the four classes}

\begin{center}
\renewcommand{\arraystretch}{1.5}
\begin{tabular}{|>{\raggedright\arraybackslash}p{2.5cm}|>{\raggedright\arraybackslash}p{2.5cm}|>{\raggedright\arraybackslash}p{2.5cm}|>{\raggedright\arraybackslash}p{2.5cm}|>{\raggedright\arraybackslash}p{3cm}|}
\hline
\textbf{Class} & \textbf{Generators} & \textbf{Relations} & \textbf{Completion?} & \emph{Examples} \\
\hline
\hline
\emph{Quadratic} & 
Finite-dim $V$ & 
$R \subset V^{\otimes 2}$ only & 
NO &
$\mathcal{H}_k$, $\widehat{\mathfrak{g}}_k$ \\
\hline
\emph{Curved} & 
Finite-dim $V$ & 
$R \subset \bigoplus_{n \geq 2} V^{\otimes n}$, $\mu_0 \in Z(\mathcal{A})$ & 
SOMETIMES &
$\text{Vir}_c$ \\
\hline
\emph{Filtered} & 
Infinite-dim, graded & 
All orders, composite ops & 
YES &
$W_3$, $W_N$ \\
\hline
\textbf{General / frontier} & 
Infinite-dim, ungraded & 
No uniform finite-type control & 
OPEN / frontier &
$W_\infty$ \\
\hline
\end{tabular}
\end{center}

\subsection{Theoretical framework: filtered cooperads}

Following Gui--Li--Zeng \cite{GLZ22}, we develop the theory of filtered cooperads.

\begin{definition}[Filtered cooperad]\label{def:filtered-cooperad}
A \emph{filtered cooperad} $\mathcal{C}$ is a cooperad equipped with a filtration:
\begin{equation}
F^0\mathcal{C} \supset F^1\mathcal{C} \supset F^2\mathcal{C} \supset \cdots
\end{equation}
(decreasing) satisfying:
\begin{enumerate}
\item \emph{Coalgebra compatibility}: 
\[\Delta(F^k\mathcal{C}) \subset \sum_{i+j=k} F^i\mathcal{C} \otimes F^j\mathcal{C}\]
\item \emph{Separated}: $\bigcap_{k=0}^\infty F^k\mathcal{C} = 0$
\item \emph{Complete}: $\mathcal{C} \cong \varprojlim_k \mathcal{C}/F^k\mathcal{C}$
\end{enumerate}
\end{definition}

\begin{theorem}[Filtered Koszul duality (GLZ) {\cite{GLZ22}}; \ClaimStatusProvedElsewhere]\label{thm:filtered-koszul-glz}
Let $\mathcal{A}$ be a filtered chiral algebra with:
\begin{itemize}
\item Associated graded $\text{gr}(\mathcal{A})$ is quadratic
\item Filtration is compatible with chiral product
\item Completion $\widehat{\mathcal{A}} = \varprojlim_n \mathcal{A}/F_n\mathcal{A}$ exists
\end{itemize}

Then the completed bar construction:
\begin{equation}
\widehat{\bar{B}}(\mathcal{A}) := \varprojlim_n \bar{B}(\mathcal{A})/F_n
\end{equation}
computes a \emph{filtered Koszul dual} $\mathcal{A}^!_{\text{filt}}$ with:
\begin{equation}
\Omega(\widehat{\bar{B}}(\mathcal{A})) \simeq \widehat{\mathcal{A}}
\end{equation}
as filtered chiral algebras.
\end{theorem}

\begin{proof}[Proof outline (following GLZ)]
The argument proceeds through four steps.

\emph{Step 1: Associated graded.}
Since $\mathrm{gr}(\mathcal{A})$ is quadratic, Theorem~\ref{thm:quadratic-koszul} gives
$\Omega(B(\mathrm{gr}(\mathcal{A}))) \simeq \mathrm{gr}(\mathcal{A})$.

\emph{Step 2: Spectral sequence.}
The filtration induces
$E_1^{p,q} = H^q(B(F_p\mathcal{A}/F_{p-1}\mathcal{A})) \Rightarrow H^{p+q}(\widehat{\bar{B}}(\mathcal{A}))$.
The $E_1$ page reduces to the associated graded computation of Step~1.

\emph{Step 3: Convergence.}
The inverse system $\{H^*(\bar{B}(\mathcal{A})/F^n)\}_{n \geq 0}$ satisfies the Mittag-Leffler condition because the filtration is exhaustive and the successive quotients $F^n/F^{n+1}$ are finite-dimensional (by the finite-type hypothesis on $\mathcal{A}$), so the image sequences $\mathrm{im}(H^*(F^{n+k}/F^n) \to H^*(F^n/F^n))$ stabilize.  Therefore $\varprojlim^1 = 0$ and $\varprojlim_n H^*(\bar{B}/F^n) = H^*(\widehat{\bar{B}})$.

\emph{Step 4: Cobar recovery.}
The cobar-bar adjunction $\Omega \dashv B$ restricts to an equivalence on pro-nilpotent objects by the filtered analogue of Theorem~\ref{thm:bar-cobar-isomorphism-main}; the completion $\widehat{\bar{B}}(\mathcal{A})$ is pro-nilpotent by construction, so $\Omega(\widehat{\bar{B}}(\mathcal{A})) \simeq \widehat{\mathcal{A}}$.  See Positselski~\cite{Positselski11} for the general framework of curved Koszul duality with completions.
\end{proof}

\subsection{When does filtering degenerate to curved?}

\begin{proposition}[Filtered implies curved; \ClaimStatusProvedHere]\label{prop:filtered-to-curved}
A filtered chiral algebra $\mathcal{A}$ has an associated \emph{curved structure} if:
\begin{enumerate}
\item The filtration is \emph{finite-dimensional at each level}:
$\dim(F_k\mathcal{A}/F_{k-1}\mathcal{A}) < \infty$ for all $k$
\item The associated graded $\mathrm{gr}(\mathcal{A})$ is generated by $\mathrm{gr}^1(\mathcal{A})$
\item All higher relations are \emph{consequences} of lower ones plus curvature
\end{enumerate}

In this case, the filtered structure \emph{degenerates} to a curved structure with:
\[\mu_0 \in F_2\mathcal{A}\]
encoding the deviation from quadratic.
\end{proposition}

\begin{proof}
By condition~(2), $\mathrm{gr}(\mathcal{A})$ is generated in degree~1.  The free chiral algebra $\mathrm{Free}(\mathrm{gr}^1(\mathcal{A}))$ surjects onto $\mathrm{gr}(\mathcal{A})$, with kernel generated by the quadratic relations $R \subset \mathrm{gr}^1 \boxtimes \mathrm{gr}^1$.

The filtered algebra $\mathcal{A}$ deforms $\mathrm{gr}(\mathcal{A})$.  By condition~(3), the deformation is controlled by a single curvature element: choose generators $a_1, \ldots, a_r$ lifting $\mathrm{gr}^1(\mathcal{A})$ to $F_1\mathcal{A}$.  The product $\mu(a_i, a_j)$ in $\mathcal{A}$ differs from the product in $\mathrm{gr}(\mathcal{A})$ by an element in $F_2\mathcal{A}$:
\[
\mu_\mathcal{A}(a_i, a_j) = \mu_{\mathrm{gr}}(a_i, a_j) + \mu_0^{ij}
\]
where $\mu_0^{ij} \in F_2\mathcal{A}$.  By condition~(3), the higher-order corrections are determined by these quadratic corrections.  Collecting $\mu_0 = \sum_{ij} \mu_0^{ij} \in F_2\mathcal{A}$, we obtain the curved structure: the bar differential satisfies $d^2 = [\mu_0, -]$, where $\mu_0$ encodes the total deviation from quadratic.

Condition~(1) ensures that this process terminates at each filtration level, giving a well-defined curved $A_\infty$ structure.
\end{proof}

\begin{example}[Virasoro: filtered degenerates to curved]\label{ex:vir-filtered-to-curved}
The Virasoro algebra can be viewed as:

\emph{Option 1 - Filtered.}
\begin{align}
F_0\text{Vir} &= \mathbb{C} \cdot \mathbf{1} \\
F_1\text{Vir} &= F_0 \oplus \mathbb{C} \cdot \partial T \\
F_2\text{Vir} &= F_1 \oplus \mathbb{C} \cdot T \\
F_3\text{Vir} &= F_2 \oplus \mathbb{C} \cdot \partial^2 T \\
&\vdots
\end{align}

\emph{Option 2 - Curved.} Generator $V = \mathbb{C} \cdot T$, curvature $\mu_0 = c \cdot \mathbf{1}$, and higher operations $m_3(T \otimes T \otimes T)$ from the Schwarzian.

The filtration $F_k$ is generated by $T$ and its derivatives up to order $k-2$. All composite
operators like $\partial^n T$ are derivatives of the single generator $T$, so the algebra is
effectively curved rather than truly filtered.

The curvature $\mu_0 = c$ captures the failure of $T$ to be a quadratic generator.
\end{example}

\begin{example}[\texorpdfstring{$W_3$}{W3}: truly filtered, not curved]\label{ex:w3-not-curved}
The $W_3$ algebra is \emph{genuinely filtered}: its generators $L$ (dimension~2) and $W$ (dimension~3) produce composite operators $(L \cdot L)$, $(L \cdot W)$, etc.\ in the OPE that are \emph{not} derivatives of $L$ or $W$.

Therefore $W_3$ cannot be reduced to a curved algebra with finite-dimensional generators. 
It requires the full filtered framework.

The distinction is that for Virasoro, $T$ and all $\partial^n T$ are ``the same'' generator (derivatives),
whereas for $W_3$, the fields $L$, $W$, and $(L \cdot L)$ are \emph{independent} generators.
This is why $W_3$ requires completion while Heisenberg and Kac--Moody do not.
\end{example}

\subsection{Explicit calculations: three examples}

\subsubsection{Heisenberg (quadratic): no completion}

\begin{example}[Heisenberg: explicit bar complex]\label{ex:heisenberg-bar-explicit}
For $\mathcal{H}_k$ with generator $J$:

\emph{Bar complex.}
\begin{align}
\bar{B}^0(\mathcal{H}_k) &= \mathbb{C} \cdot \mathbf{1} \\
\bar{B}^1(\mathcal{H}_k) &= \mathbb{C} \cdot J \\
\bar{B}^2(\mathcal{H}_k) &= \mathbb{C} \cdot (J \otimes J) \\
\bar{B}^3(\mathcal{H}_k) &= \mathbb{C} \cdot (J \otimes J \otimes J) \\
&\vdots
\end{align}

\emph{Bar differential.}

Step 1: Write the differential explicitly:
\[d: \bar{B}_1 \to \bar{B}_0\]
\[d(J(z_1) \otimes J(z_2) \otimes \eta_{12}) = \text{Res}_{z_1=z_2}[J(z_1)J(z_2) \cdot d\log(z_1-z_2)]\]

Step 2: Insert the OPE:
\begin{align}
J(z_1)J(z_2) \cdot d\log(z_1-z_2) &= \frac{k}{(z_1-z_2)^2} \cdot \frac{dz_1}{z_1-z_2} \\
&= \frac{k \, dz_1}{(z_1-z_2)^3}
\end{align}

Step 3: Compute the residue:
\[\text{Res}_{z_1=z_2}\left[\frac{k \, dz_1}{(z_1-z_2)^3}\right] = 0\]

The triple pole has zero residue, so $d = 0$ at this level.

\begin{align}
d(J) &= 0 \\
d(J \otimes J) &= 0 \quad \text{(not } k \cdot \mathbf{1}\text{; the triple pole residue vanishes)} \\
d(J \otimes J \otimes J) &= 0 \quad \text{(all pairwise OPEs have poles of order $\geq 2$, hence zero residue)}
\end{align}

\emph{Cohomology.}
\begin{align}
H^0(\bar{B}(\mathcal{H}_k)) &= \mathbb{C} \cdot \mathbf{1} \quad \text{(vacuum)} \\
H^1(\bar{B}(\mathcal{H}_k)) &\neq 0 \quad \text{(survives: the double pole means } d = 0\text{)} \\
H^n(\bar{B}(\mathcal{H}_k)) &\text{ has CE cooperad structure}
\end{align}

\emph{Koszul dual.}
\[\mathcal{H}_k^! = \mathrm{Sym}^{\mathrm{ch}}(V^*)\]
the commutative chiral algebra on the dual space $V^* = \mathrm{Hom}(V, \mathbb{C})$.

\emph{Cohomology.}
\begin{itemize}
\item $H^0 = \mathbb{C} \cdot \mathbf{1}$ (vacuum)
\item $H^1 = \bar{B}_1$ itself (since $d = 0$)
\item $H^n$ for $n \geq 2$ follow similar pattern
\end{itemize}

The full cohomology $H^*(\bar{B}(\mathcal{H}_k))$ has the structure of the CE cooperad:
\[H^*(\bar{B}(\mathcal{H}_k)) \simeq \text{CE}^!(\mathfrak{h}_k)\]

Taking the cobar dual:
\[\Omega(H^*(\bar{B}(\mathcal{H}_k))) \simeq \Omega(\text{CE}^!(\mathfrak{h}_k)) \simeq \mathcal{H}_k\]

\emph{Physical interpretation.}
\begin{itemize}
\item The Heisenberg algebra $\mathcal{H}_k$ describes a free boson
\item Its Koszul dual $\mathrm{Sym}^{\mathrm{ch}}(V^*)$ is the commutative chiral algebra on $V^*$, with curvature $m_0 = -k \cdot c$
\item Physically, this commutative dual describes the \emph{ghost system} for gauging the $U(1)$ symmetry
\item This matches the BV-BRST quantization of the gauged theory
\end{itemize}

No completion is needed: the bar complex is finite-dimensional at each degree and
converges immediately.
\end{example}

\subsubsection{Virasoro (curved): sometimes completion}

\begin{example}[Virasoro: bar complex requires completion]\label{ex:vir-bar-completion}
For $\text{Vir}_c$ with generator $T$:

\emph{Bar complex (before completion).}
\begin{align}
\bar{B}^0(\text{Vir}) &= \mathbb{C} \cdot \mathbf{1} \\
\bar{B}^1(\text{Vir}) &= \mathbb{C} \cdot T \oplus \mathbb{C} \cdot \partial T \oplus \cdots \\
\bar{B}^2(\text{Vir}) &= (\mathbb{C} \cdot T \oplus \cdots)^{\otimes 2} \\
&\vdots
\end{align}

\emph{Issue.} The space $\bar{B}^1$ is \emph{infinite-dimensional} because it includes 
all derivatives $\partial^n T$ for $n \geq 0$.

\emph{Completion.} Define the augmentation ideal:
\[I = \langle T, \partial T, \partial^2 T, \ldots \rangle\]

Complete with respect to $I$:
\[\widehat{\bar{B}}(\text{Vir}) = \varprojlim_n \bar{B}(\text{Vir})/I^n\]

\emph{Completed differential.}
\[\widehat{d}(T \otimes T) = \text{Res}(T(z)T(w)) + c \cdot \mathbf{1}\]

The curvature $c$ ensures $\widehat{d}^2 = 0$ on the completed complex.

\emph{Cohomology.}
\[H^*(\widehat{\bar{B}}(\text{Vir})) = \widehat{U(\text{Vir})}^*_{-c}\]
is the completed dual universal enveloping algebra.

Completion is essential. Without completion, the bar complex does not converge and the
Koszul dual is not well-defined.
\end{example}

\subsubsection{\texorpdfstring{$W_3$ (filtered): always completion}{W-3 (filtered): always completion}}

\begin{example}[\texorpdfstring{$W_3$}{W3}: bar complex must be completed]\label{ex:w3-bar-completion}
For $W_3$ with generators $L$ (dimension 2) and $W$ (dimension 3):

\emph{Bar complex (before completion).}
\begin{align}
\bar{B}^0(W_3) &= \mathbb{C} \cdot \mathbf{1} \\
\bar{B}^1(W_3) &= \mathbb{C} \cdot L \oplus \mathbb{C} \cdot W \oplus \text{(derivatives and composites)} \\
\bar{B}^2(W_3) &= \text{(all pairs)} \\
&\vdots
\end{align}

\emph{Problem.} Already at degree 1, we have:
\begin{itemize}
\item Generators: $L$, $W$
\item First derivatives: $\partial L$, $\partial W$
\item Second derivatives: $\partial^2 L$, $\partial^2 W$
\item Composites: $(L \cdot L)$, $(L \cdot W)$, $(W \cdot W)$
\item Higher composites: $(\partial L \cdot L)$, etc.
\end{itemize}

This is infinite-dimensional even before taking products.

\emph{Filtration.} Filter by total operator dimension:
\begin{align}
F_0 &= \mathbb{C} \cdot \mathbf{1} \\
F_2 &= F_0 \oplus \mathbb{C} \cdot L \\
F_3 &= F_2 \oplus \mathbb{C} \cdot W \oplus \mathbb{C} \cdot \partial L \\
F_4 &= F_3 \oplus \mathbb{C} \cdot \partial W \oplus \mathbb{C} \cdot \partial^2 L 
\oplus \mathbb{C} \cdot (L \cdot L) \\
&\vdots
\end{align}

\emph{Completed bar complex.}
\[\widehat{\bar{B}}(W_3) = \varprojlim_n \bar{B}(W_3)/F_n\]

\emph{Completed differential.}
\[\widehat{d}(W \otimes W) = \text{Res}(W(z)W(w)) + \text{(composite terms)} + c \cdot \mathbf{1}\]

The composite terms involve $(L \cdot L)$ and higher, which are not in the span of $\{L, W\}$.

\emph{Cohomology.}
\[H^*(\widehat{\bar{B}}(W_3)) = \widehat{\text{CoW}_3}\]
is the completed cooperad structure dual to $W_3$.

Completion is essential. Without it, the bar construction does not converge.
\end{example}

\subsection{Convergence criteria}

\begin{theorem}[Convergence of bar construction; \ClaimStatusProvedHere]\label{thm:bar-convergence}
Let $\mathcal{A}$ be a chiral algebra and write
$\bar{B}(\mathcal{A})=\bigoplus_{n\ge0}\bar{B}^n(\mathcal{A})$ with its bar-degree
filtration $F_p:=\bigoplus_{n\le p}\bar{B}^n(\mathcal{A})$.
Assume $\dim\bar{B}^n(\mathcal{A})<\infty$ for all $n$. Then the following are equivalent:
\begin{enumerate}
\item $\bar{B}(\mathcal{A})$ converges without completion, i.e. the natural map
$\bigoplus_n\bar{B}^n(\mathcal{A})\to\prod_n\bar{B}^n(\mathcal{A})$ is a
quasi-isomorphism for the bar differential.
\item \emph{Growth bound:}
$\limsup_{n\to\infty}\dim(\bar{B}^n(\mathcal{A}))^{1/n}<\infty$.
\item \emph{Filtration compatibility:}
$d(F_p)\subseteq F_p$ and $d$ lowers bar degree by one on $\gr_F\bar{B}(\mathcal{A})$.
\end{enumerate}
If (2) or (3) fails, one must use the completed bar complex
$\widehat{\bar{B}}(\mathcal{A})$.
\end{theorem}

\begin{proof}
\emph{(2)+(3) $\Rightarrow$ (1).}
Condition (3) gives a filtered complex structure on the direct sum complex and
the associated spectral sequence has
\[
E_0^{p,\bullet}=\bar{B}^p(\mathcal{A}), \qquad d_0:\bar{B}^p\to\bar{B}^{p-1}.
\]
Because each $\bar{B}^p$ is finite-dimensional and condition (2) gives uniform
at-most-exponential growth, each total degree receives only finite-dimensional
contributions, so the filtered spectral sequence is well-defined and converges
to $H^*(\bar{B}(\mathcal{A}))$. Thus no topological completion is required.
Together with Theorem~\ref{thm:conilpotency-convergence}, this means every
iterated coproduct/differential expression is finite on each element.

\emph{(1) $\Rightarrow$ (3).}
If convergence holds in the uncompleted direct sum, $d$ cannot create infinitely
many bar-degree components from a single homogeneous input; equivalently,
$d(F_p)\subseteq F_p$ and the induced differential on $\gr_F$ is degree $-1$.
Otherwise $d$ is only defined in the product completion.

\emph{(1) $\Rightarrow$ (2).}
By the Eilenberg--Moore comparison theorem \cite[Theorem~7.14]{Weibel94}, the canonical map from the direct-sum totalization to the product totalization is a quasi-isomorphism if and only if the filtration is exhaustive and complete; the growth condition~(2) ensures that the completed bar filtration is exhaustive.  Conversely, if the growth bound fails, the filtration is not exhaustive: there exist infinite-bar-degree elements in the completion that project to zero in every finite truncation, producing cohomology classes in the completion not represented by any finite bar element.  These classes obstruct the comparison quasi-isomorphism, contradicting~(1).  Hence (2) is necessary.

\emph{Quadratic finite-type case.}
For $\mathcal{A}=T(V)/(R)$ with $V$ finite-dimensional,
$\bar{B}^n(\mathcal{A})\cong (sV)^{\otimes n}$, so (2) holds with growth constant
$\dim V$, and (3) holds by definition of the bar filtration. Therefore quadratic
finite-type algebras converge without completion.
\end{proof}

\begin{remark}[Dependencies]
This proof uses:
\begin{enumerate}[label=(D\arabic*)]
\item conilpotent convergence control (Theorem~\ref{thm:conilpotency-convergence});
\item filtered-to-curved reduction inputs for non-quadratic regimes
(Theorem~\ref{thm:filtered-to-curved});
\item filtered/complete Koszul duality context (Theorem~\ref{thm:filtered-koszul-glz}).
\end{enumerate}
\end{remark}

\subsection{Physical interpretation}

Quadratic algebras correspond to free field theories (Heisenberg $\leftrightarrow$ free boson, Kac--Moody $\leftrightarrow$ WZW).  Curved algebras correspond to interacting theories with anomalies (Virasoro, with central charge $c$ measuring quantum breaking of scale invariance).  Filtered algebras correspond to theories with composite operators ($W_3$ and Toda field theory).  General algebras correspond to non-local theories ($W_\infty$), where the exact MC4 package of filtered H-level realization plus finite-stage coefficient identification remains open.

From Kontsevich's geometric viewpoint, the filtration level corresponds to the codimension of collision loci: quadratic algebras see only pairwise collisions, curved algebras see central terms from $n$-point collisions, and filtered algebras require higher-codimension strata.  The completion $\widehat{\bar{B}}(\mathcal{A})$ is the formal neighborhood of the diagonal in configuration space.

\subsection{Summary and decision tree}

\begin{remark}[Summary for practitioners]\label{rem:practitioner-takeaway}
\emph{Before computing the Koszul dual, one should determine:}
\begin{enumerate}
\item Is the algebra quadratic? $\Rightarrow$ Proceed directly
\item Is it curved with central curvature? $\Rightarrow$ Check if $\dim(\bar{B}^1) < \infty$
  \begin{itemize}
  \item If yes: No completion
  \item If no: Completion is required
  \end{itemize}
\item Does it have composite operators? $\Rightarrow$ Must complete
\item Is the generating space infinite-dimensional? $\Rightarrow$ treat as an exact MC4 frontier: first construct the filtered H-level target, then identify its finite quotients by explicit coefficient identities and finite detection
\end{enumerate}

Most vertex algebras from CFT are either quadratic or curved with finite-dimensional
$\bar{B}^1$, so Koszul duality applies.
\end{remark}

The bar complex $\bar{B}^{\mathrm{ch}}(\cA)$ is a chiral coalgebra; when the curvature is nonzero, it is a \emph{curved} coalgebra requiring Positselski's coderived categories for its homological algebra.
\fi
%%% END COMMENTED-OUT: Four-class exposition

\section{Chiral coalgebra homological algebra}
\label{sec:chiral-coalgebra-homalg}
\index{chiral coalgebra!homological algebra|textbf}
\index{Positselski!chiral lift|textbf}

The bar complex $\bar{B}^{\mathrm{ch}}(\cA)$ is a chiral coalgebra
(Theorem~\ref{thm:bar-chiral}).  At genus~$0$, it is a
DG chiral coalgebra ($\dzero^2 = 0$); at genus~$g \geq 1$ with
$\kappa(\cA) \neq 0$, it is a \emph{curved} DG chiral coalgebra
($\dfib^{\,2} = \mcurv{g} \neq 0$; Convention~\ref{conv:higher-genus-differentials}).  In the curved case, the ordinary
derived category is trivial
(Proposition~\ref{prop:curved-bar-acyclicity}), and a richer
homological algebra is needed.

The chiral analogues of Positselski's coalgebra homological
algebra~\cite{Positselski11} replace the ground field~$k$ by
$\mathcal{D}_X$ and the tensor product by $\chirtensor$; all
exactness properties are inherited from configuration space geometry.

\subsection{Chiral CDG-coalgebras}
\label{subsec:chiral-CDG-coalgebra}

\begin{definition}[Chiral CDG-coalgebra]\label{def:chiral-CDG-coalgebra}
\index{CDG-coalgebra!chiral}
A \emph{chiral CDG-coalgebra} over a smooth curve $X$ is a triple
$(C, d, h)$ where:
\begin{enumerate}
\item $C$ is a graded chiral coalgebra on $X$: a graded
  $\mathcal{D}_X$-module equipped with a comultiplication
  $\Delta\colon C \to C \chirtensor C$
  (defined via configuration space splitting maps,
  cf.\ Theorem~\ref{thm:bar-chiral}),
  a counit $\varepsilon\colon C \to \omega_X$, coassociativity,
  and the counit axioms.
\item $d\colon C \to C$ is a homogeneous endomorphism of degree~$1$
  that is an odd coderivation: the composition
  $\Delta \circ d = (d \otimes \mathrm{id} +
  \mathrm{id} \otimes d) \circ \Delta$
  holds, and $\varepsilon \circ d = 0$ on $C^{>-2}$
  (the counit annihilates $d$ except possibly in the lowest degree).
\item $h\colon C \to \omega_X$ is a homogeneous linear functional
  of degree~$2$ (the \emph{curvature}), satisfying:
  \begin{enumerate}
  \item $d^2(c) = h \ast c - c \ast h$ for all $c \in C$,
    where $h \ast c$ denotes the coaction of the curvature
    (defined via the comodule structure of $C$ over itself);
  \item $h \circ d(c) = 0$ for all $c \in C$
    (the curvature is a cocycle for the coderivation~$d$).
  \end{enumerate}
\end{enumerate}
When $h = 0$, we recover an ordinary DG chiral coalgebra.
\end{definition}

\begin{example}[The bar complex as chiral CDG-coalgebra]
\label{ex:bar-as-CDG}
\index{bar complex!as CDG-coalgebra}
For a Koszul chiral algebra $\cA$ on $X$:
\begin{itemize}
\item At genus~$0$: $\bar{B}^{(0)}(\cA) = (\bar{B}^{\mathrm{ch}}(\cA), d_B, 0)$
  is a DG chiral coalgebra ($h = 0$).
\item At genus~$g \geq 1$: $\bar{B}^{(g)}(\cA) = (\bar{B}^{\mathrm{ch}}(\cA), d_B^{(g)}, h_g)$
  is a CDG chiral coalgebra with curvature
  $h_g = m_0^{(g)} = \kappa(\cA) \cdot \lambda_g$
  (Theorem~\ref{thm:genus-universality}), where $\kappa(\cA)$ is
  the obstruction coefficient and $\lambda_g$ is the Faber--Pandharipande
  class on $\overline{\mathcal{M}}_g$.
\end{itemize}
The CDG axioms are verified as follows.  Axiom (2a): the identity
$d^2 = [m_0, -]$ is the curved $A_\infty$ relation
(Corollary~\ref{cor:critical-level-universality}).
Axiom (2b): $h \circ d = 0$ follows from $m_0$ being a scalar
(central element), so $d(m_0) = 0$.
\end{example}

\begin{remark}[Chiral CDG-coalgebra morphisms]
\label{rem:CDG-morphisms}
A morphism of chiral CDG-coalgebras $(C, d_C, h_C)$
$\to$ $(D, d_D, h_D)$ is a pair $(f, a)$ where
$f\colon C \to D$ is a morphism of graded chiral
coalgebras and $a\colon C \to \omega_X$
is a homogeneous linear function of degree~$1$ such
that
\begin{align*}
d_D(f(c)) &= f(d_C(c)) + f(a \ast c) - (-1)^{|c|}f(c \ast a), \\
h_D(f(c)) &= h_C(c) + a(d_C(c)) + a^2(c)
\end{align*}
for all $c \in C$.
Composition of morphisms is defined by
$(g, b) \circ (f, a) = (g \circ f,\allowbreak b \circ f + a)$.
This follows Positselski~\cite[\S4.1]{Positselski11}.
\end{remark}

\subsection{Chiral CDG-comodules and CDG-contramodules}
\label{subsec:chiral-CDG-comod-contra}

\begin{definition}[Chiral CDG-comodule]\label{def:chiral-CDG-comodule}
\index{CDG-comodule!chiral}
Let $(C, d, h)$ be a chiral CDG-coalgebra.  A \emph{left chiral
CDG-comodule} over $C$ is a graded left $C$-comodule $M$
(i.e., a graded $\mathcal{D}_X$-module with coaction
$\lambda\colon M \to C \chirtensor M$ satisfying coassociativity
and counit axioms) endowed with a homogeneous endomorphism
$d_M\colon M \to M$ of degree~$1$ such that:
\begin{enumerate}
\item The coaction commutes with differentials:
  $\lambda \circ d_M = (d \otimes \mathrm{id} + \mathrm{id} \otimes d_M) \circ \lambda$.
\item The curvature relation:
  $d_M^2(x) = h \ast x$ for all $x \in M$,
  where $h \ast x$ denotes the action of the curvature
  via the coaction map.
\end{enumerate}

A \emph{right chiral CDG-comodule} is defined analogously with
coaction $\rho\colon N \to N \chirtensor C$ and
$d_N^2(y) = -y \ast h$.

Morphisms of chiral CDG-comodules are homogeneous $\mathcal{D}_X$-linear
maps of degree~$0$ that are compatible with the coaction
and commute with the differential (in the usual graded sense).
\end{definition}

\begin{definition}[Chiral CDG-contramodule]\label{def:chiral-CDG-contramodule}
\index{CDG-contramodule!chiral}
Let $(C, d, h)$ be a chiral CDG-coalgebra.  A \emph{left chiral
CDG-contramodule} over $C$ is a graded left $C$-contramodule $P$
(i.e., a graded $\mathcal{D}_X$-module with contraaction
$\pi_P\colon \mathrm{Hom}^{\mathrm{ch}}(C, P) \to P$
satisfying the contraassociativity and contraunit axioms)
endowed with a
homogeneous endomorphism $d_P\colon P \to P$ of degree~$1$ such that:
\begin{enumerate}
\item The contraaction commutes with differentials on
  $\mathrm{Hom}^{\mathrm{ch}}(C, P)$ and $P$.
\item The curvature relation:
  $d_P^2(p) = h \ast p$ for all $p \in P$,
  where $h \ast p = \pi_P(h \otimes p)$ denotes the contraaction
  of the curvature.
\end{enumerate}
Morphisms of chiral CDG-contramodules are defined analogously.
\end{definition}

\begin{remark}[Contramodules in the chiral setting]
\label{rem:chiral-contramodules}
\index{contramodule!chiral}
The chiral contraaction $\pi_P\colon \mathrm{Hom}^{\mathrm{ch}}(C, P) \to P$
replaces Positselski's $\mathrm{Hom}_k(C, P) \to P$, where
$\mathrm{Hom}^{\mathrm{ch}}(C, P) := \mathrm{Hom}_{\mathcal{D}_X}(C, P)$
is the internal Hom in the chiral category.  For a conilpotent chiral
coalgebra $C$ with finite-dimensional graded pieces (as holds for
$\bar{B}^{\mathrm{ch}}(\cA)$ when $\cA$ is finitely generated over $\mathcal{D}_X$),
$\mathrm{Hom}^{\mathrm{ch}}(C, P) \cong C^* \chirtensor P$
where $C^* = \mathrm{Hom}_{\mathcal{D}_X}(C, \omega_X)$ is the
graded dual.  In this finite-type case, contramodules over $C$ are
equivalent to complete modules over the dual algebra $C^*$, recovering
the correspondence between conilpotent comodules and complete modules
already used in Theorem~\ref{thm:e1-module-koszul-duality}.
\end{remark}

\begin{remark}[No cohomology for CDG-comodules]
\label{rem:no-CDG-cohomology}
Since $d^2 \neq 0$ in a CDG-comodule (when $h \neq 0$), there is no
notion of cohomology.  This is the fundamental reason why the ordinary
derived category is inadequate at higher genus: one cannot define
acyclic objects via vanishing cohomology.  The correct
substitute is the coderived category (Definition~\ref{def:chiral-coderived}).
\end{remark}

\subsection{Chiral cotensor and cohomomorphism complexes}
\label{subsec:chiral-cotensor}

\begin{definition}[Chiral cotensor product]\label{def:chiral-cotensor}
\index{cotensor product!chiral}
Let $(C, d, h)$ be a chiral CDG-coalgebra, $N$ a right chiral
CDG-comodule, and $M$ a left chiral CDG-comodule.  The
\emph{chiral cotensor product} $N \square_C^{\mathrm{ch}} M$ is
the equalizer:
\[
N \square_C^{\mathrm{ch}} M =
\ker\bigl(\rho \chirtensor \mathrm{id} - \mathrm{id} \chirtensor \lambda\colon
N \chirtensor M \rightrightarrows N \chirtensor C \chirtensor M\bigr)
\]
in the category of graded $\mathcal{D}_X$-modules.
The differential on $N \chirtensor M$ restricts to
$N \square_C^{\mathrm{ch}} M$ (this uses the compatibility of coactions
with differentials), and the restricted differential has
$d^2 = h_N \ast (-) - (-) \ast h_M$ on the cotensor product, making
$N \square_C^{\mathrm{ch}} M$ a complex (i.e., $d^2 = 0$) precisely
when $h_N = h_M$ (the curvatures match), which holds when $N$ and $M$
are CDG-comodules over the same CDG-coalgebra~$C$.
\end{definition}

\begin{definition}[Chiral cohomomorphisms]\label{def:chiral-cohom}
\index{cohomomorphism!chiral}
For left chiral CDG-comodules $L$ and $M$ over $C$, the
\emph{chiral cohomomorphism complex} is:
\[
\mathrm{Cohom}_C^{\mathrm{ch}}(L, M) =
\{f \in \mathrm{Hom}_{\mathcal{D}_X}(L, M) :
\lambda_M \circ f = (\mathrm{id}_C \chirtensor f) \circ \lambda_L\}
\]
with differential $(d f)(x) = d_M(f(x)) - (-1)^{|f|} f(d_L(x))$.
This differential has a zero square (using the CDG-comodule axioms),
making $\mathrm{Cohom}_C^{\mathrm{ch}}(L, M)$ an honest complex.
\end{definition}

\begin{definition}[Chiral contratensor product]\label{def:chiral-contratensor}
\index{contratensor product!chiral}
For a right chiral CDG-comodule $N$ and a left chiral
CDG-contramodule $P$ over $C$, the \emph{chiral contratensor
product} $N \odot_C^{\mathrm{ch}} P$ is the coequalizer:
\[
N \odot_C^{\mathrm{ch}} P =
\mathrm{coker}\bigl(
N \chirtensor \mathrm{Hom}^{\mathrm{ch}}(C, P) \rightrightarrows N \chirtensor P
\bigr)
\]
where the two maps are
$\rho \chirtensor \pi_P$ and $\mathrm{id}_N \chirtensor \pi_P$.
The induced differential on $N \odot_C^{\mathrm{ch}} P$ has zero
square by the same curvature cancellation as for the cotensor product.
\end{definition}

\begin{lemma}[Adjunction; \ClaimStatusProvedHere]\label{lem:chiral-co-contra-adjunction}
\index{adjunction!chiral cotensor--contratensor}
There is a canonical isomorphism of complexes:
\[
\mathrm{Hom}_{\mathcal{D}_X}(N \square_C^{\mathrm{ch}} M,\, V)
\;\cong\;
\mathrm{Cohom}_C^{\mathrm{ch}}\bigl(M,\,
\mathrm{Hom}_{\mathcal{D}_X}(N, V)\bigr)
\]
for any graded $\mathcal{D}_X$-module $V$, and dually:
\[
\mathrm{Hom}_{\mathcal{D}_X}(C \chirtensor V,\, P)
\;\cong\;
\mathrm{Hom}_{\mathcal{D}_X}(V,\,
\mathrm{Hom}^{\mathrm{ch}}(C, P))
\]
establishing $\Phi_C \dashv \Psi_C$
(see Definition~\textup{\ref{def:chiral-Phi-Psi}} below).
\end{lemma}

\begin{proof}
Both isomorphisms are instances of the tensor-hom adjunction in the
chiral category.  For the first: an element of the left side is a
$\mathcal{D}_X$-linear map $f\colon N \square_C^{\mathrm{ch}} M \to V$.
By the universal property of the equalizer, this is equivalent to a
$\mathcal{D}_X$-linear map $\tilde{f}\colon N \chirtensor M \to V$
that equalizes the two coaction maps.  By the chiral tensor-hom
adjunction (which holds for holonomic $\mathcal{D}_X$-modules
by~\cite[Proposition~3.4.6]{BD04}), $\tilde{f}$ corresponds to a map
$M \to \mathrm{Hom}_{\mathcal{D}_X}(N, V)$ that is
$C$-colinear, i.e., an element of
$\mathrm{Cohom}_C^{\mathrm{ch}}(M, \mathrm{Hom}_{\mathcal{D}_X}(N, V))$.
The second isomorphism is analogous.
\end{proof}

\subsection{Coderived and contraderived categories}
\label{subsec:chiral-coderived-contraderived}

\begin{definition}[Coderived and contraderived categories]
\label{def:chiral-coderived}
\index{coderived category!chiral|textbf}
\index{contraderived category!chiral|textbf}
Let $(C, d, h)$ be a chiral CDG-coalgebra.

\emph{Exact triples.}
An exact triple of left chiral CDG-comodules is a short sequence
$0 \to L \to M \to N \to 0$ that is exact in the abelian category
of graded $C$-comodules (with closed, degree-zero morphisms).
The analogous notion applies to CDG-contramodules.

\emph{Coacyclic objects.}
A left chiral CDG-comodule $M$ is \emph{coacyclic} if it belongs to
the minimal thick subcategory of $\mathrm{Hot}(C\text{-}\mathrm{comod}^{\mathrm{ch}})$
that contains the totalizations of all exact triples and is closed
under infinite direct sums.  Denote by
$\mathrm{Acycl}^{\mathrm{co}}(C\text{-}\mathrm{comod}^{\mathrm{ch}})$
the thick subcategory of coacyclic CDG-comodules.

\emph{Contraacyclic objects.}
A left chiral CDG-contramodule $P$ is \emph{contraacyclic} if it
belongs to the minimal thick subcategory of
$\mathrm{Hot}(C\text{-}\mathrm{contra}^{\mathrm{ch}})$
that contains the totalizations of all exact triples and is closed
under infinite direct products.  Denote by
$\mathrm{Acycl}^{\mathrm{ctr}}(C\text{-}\mathrm{contra}^{\mathrm{ch}})$
the thick subcategory of contraacyclic CDG-contramodules.

\emph{The exotic derived categories.}
\begin{align}
D^{\mathrm{co}}(C\text{-}\mathrm{comod}^{\mathrm{ch}})
&:= \mathrm{Hot}(C\text{-}\mathrm{comod}^{\mathrm{ch}})
\,/\, \mathrm{Acycl}^{\mathrm{co}}(C\text{-}\mathrm{comod}^{\mathrm{ch}})
\label{eq:chiral-coderived}\\
D^{\mathrm{ctr}}(C\text{-}\mathrm{contra}^{\mathrm{ch}})
&:= \mathrm{Hot}(C\text{-}\mathrm{contra}^{\mathrm{ch}})
\,/\, \mathrm{Acycl}^{\mathrm{ctr}}(C\text{-}\mathrm{contra}^{\mathrm{ch}})
\label{eq:chiral-contraderived}
\end{align}
\end{definition}

\begin{remark}[Comparison with ordinary derived category]
\label{rem:exotic-vs-ordinary}
When $h = 0$ (genus~$0$), CDG-comodules are ordinary DG-comodules
and one can also form the standard derived category
$D(C\text{-}\mathrm{comod}^{\mathrm{ch}})$
by localizing at acyclic objects (those with $H^*(M) = 0$).  In general,
$\mathrm{Acycl}^{\mathrm{co}} \supset \mathrm{Acycl}$,
so $D^{\mathrm{co}}$ is a quotient of $D$.  When $C$ has
finite homological dimension (e.g., when $\cA$ is Koszul with finite
global dimension), the two categories coincide
(Positselski~\cite[\S4.5]{Positselski11}).  When $h \neq 0$, the standard
derived category is trivial (no cohomology exists), while
$D^{\mathrm{co}}$ retains rich structure.
\end{remark}

\subsection{Injective comodules and projective contramodules}
\label{subsec:chiral-inj-proj}

\begin{definition}[Cofree chiral CDG-comodule]
\label{def:cofree-chiral-comod}
\index{cofree comodule!chiral}
For a graded $\mathcal{D}_X$-module $V$, the \emph{cofree graded
chiral $C$-comodule} on $V$ is:
\[
J(V) := C \chirtensor V
\]
with coaction $\Delta \chirtensor \mathrm{id}_V\colon
C \chirtensor V \to C \chirtensor C \chirtensor V$.
The CDG-comodule differential is
$d_{J(V)}(c \otimes v) = d(c) \otimes v$,
and the curvature relation $d_{J(V)}^2 = h \ast (-)$ is inherited
from~$C$.
\end{definition}

\begin{proposition}[Injective and projective resolutions;
\ClaimStatusProvedHere]\label{prop:chiral-inj-proj-resolutions}
\index{injective resolution!chiral CDG-comodule}
\index{projective resolution!chiral CDG-contramodule}
Let $(C, d, h)$ be a conilpotent chiral CDG-coalgebra with
finite-dimensional graded pieces.
\begin{enumerate}[label=\textup{(\alph*)}]
\item For any left chiral CDG-comodule $M$, the cofree chiral
  comodule $J(M) = C \chirtensor M$ is \emph{injective} in the
  abelian category of graded $C^{\#}$-comodules (where $C^{\#}$
  denotes $C$ with forgotten differential and curvature).
\item For any left chiral CDG-contramodule $P$, the free chiral
  contramodule $F(P) = \mathrm{Hom}^{\mathrm{ch}}(C, P)$ is
  \emph{projective} in the abelian category of graded
  $C^{\#}$-contramodules.
\item The homotopy categories admit semiorthogonal decompositions:
  the coderived category $D^{\mathrm{co}}(C\text{-}\mathrm{comod}^{\mathrm{ch}})$
  is equivalent to the homotopy category of CDG-comodules
  with injective underlying graded comodules, and
  $D^{\mathrm{ctr}}(C\text{-}\mathrm{contra}^{\mathrm{ch}})$
  is equivalent to the homotopy category of CDG-contramodules
  with projective underlying graded contramodules.
\end{enumerate}
\end{proposition}

\begin{proof}
\emph{Part (a).}
Let $0 \to L \to M \to N \to 0$ be an exact sequence of graded
$C^{\#}$-comodules.  We must show that
$0 \to \mathrm{Hom}_{C^{\#}}(N, J(V)) \to \mathrm{Hom}_{C^{\#}}(M, J(V))
\to \mathrm{Hom}_{C^{\#}}(L, J(V)) \to 0$
is exact.  By the tensor-hom adjunction in the chiral category:
$\mathrm{Hom}_{C^{\#}}(M, C \chirtensor V)
\cong \mathrm{Hom}_{\mathcal{D}_X}(M, V)$,
and the functor $\mathrm{Hom}_{\mathcal{D}_X}(-, V)$ is exact
on holonomic $\mathcal{D}_X$-modules (since holonomic modules have
finite length).

\emph{Part (b).}
Dual to (a): for an exact sequence
$0 \to P \to Q \to R \to 0$ of graded $C^{\#}$-contramodules,
$\mathrm{Hom}_{C^{\#}}(\mathrm{Hom}^{\mathrm{ch}}(C, V), -)$
is exact because
$\mathrm{Hom}_{C^{\#}}(\mathrm{Hom}^{\mathrm{ch}}(C, V), P)
\cong \mathrm{Hom}_{\mathcal{D}_X}(V, P)$.

\emph{Part (c).}
For any CDG-comodule $M$, the cobar resolution provides a
functorial injective resolution: form
$J = \bigoplus_{n=0}^{\infty} C^{\chirtensor (n+1)} \chirtensor M$
with the standard bar-type differential (the chiral analogue of
Positselski~\cite[\S4.4, proof of Theorem]{Positselski11}).
The closed morphism $M \to J$ has the property that
$\mathrm{cone}(M \to J)$ is coacyclic (each quotient
$\tau_{\leq n}J / \tau_{\leq n-1}J$ is a cofree CDG-comodule, and
the cone is built from these by exact triangles and countable direct
sums).  The projectivity statement for contramodules is dual:
take the bar resolution
$\cdots \to \mathrm{Hom}^{\mathrm{ch}}(C^{\chirtensor 2}, P)
\to \mathrm{Hom}^{\mathrm{ch}}(C, P) \to P$
and observe that the cone is contraacyclic (built from free
contramodules by exact triangles and countable direct products).

The semiorthogonal decomposition follows formally: the localization
functor from the full homotopy category to $D^{\mathrm{co}}$
restricts to an equivalence on the full subcategory of objects with
injective underlying graded comodule (since coacyclic objects with
injective underlying comodule are contractible, by the same argument as
Positselski~\cite[\S4.4]{Positselski11}, using the injective
resolution of the previous paragraph).  The dual statement holds for
$D^{\mathrm{ctr}}$.
\end{proof}

\begin{proposition}[Explicit CDG Hom-complex; \ClaimStatusProvedHere]
\label{prop:cdg-hom-complex}
\index{CDG-comodule!Hom-complex|textbf}
Let $(C, d_C, h)$ be a chiral CDG-coalgebra and let $M$, $N$ be
left chiral CDG-comodules over~$C$.  The \emph{CDG Hom-complex}
$\underline{\mathrm{Hom}}_C^{\mathrm{ch}}(M, N)$ is the graded
$\mathcal{D}_X$-module of graded $C^{\#}$-comodule morphisms,
with differential
\begin{equation}\label{eq:cdg-hom-differential}
d_{\mathrm{Hom}}(f) \;=\; d_N \circ f
\;-\; (-1)^{|f|}\, f \circ d_M.
\end{equation}
This differential satisfies $d_{\mathrm{Hom}}^2(f) = [h_N, f] - [h_M, f]$
where $h_M$, $h_N$ denote the curvature endomorphisms
$h \ast (-)$ on $M$ and $N$ respectively.  In particular:
\begin{enumerate}[label=\textup{(\alph*)}]
\item When $h = 0$ \textup{(}genus~$0$\textup{)},
  $d_{\mathrm{Hom}}^2 = 0$ and the Hom-complex is a genuine
  DG-module.
\item When $h \neq 0$, the Hom-complex is a \emph{CDG-module}
  with curvature $h_{\mathrm{Hom}}(f) = h_N \circ f - f \circ h_M$.
  For endomorphisms ($M = N$),
  $d_{\mathrm{Hom}}^2(f) = [h_M, f]$, so $d_{\mathrm{Hom}}$ is
  a genuine differential on the \emph{center}
  $Z(\mathrm{End}_C(M)) = \{f : h_M \circ f = f \circ h_M\}$.
\end{enumerate}
\end{proposition}

\begin{proof}
Direct computation.  For $f \in \underline{\mathrm{Hom}}_C^{\mathrm{ch}}(M, N)^k$:
\begin{align*}
d_{\mathrm{Hom}}^2(f)
&= d_N \circ \bigl(d_N \circ f - (-1)^k f \circ d_M\bigr)
  - (-1)^{k+1}\bigl(d_N \circ f - (-1)^k f \circ d_M\bigr) \circ d_M \\
&= d_N^2 \circ f + (-1)^{k+1}\, d_N \circ f \circ d_M
  + (-1)^{k+1}\, d_N \circ f \circ d_M + f \circ d_M^2 \\
&= d_N^2 \circ f - f \circ d_M^2 \\
&= (h_N \ast -) \circ f - f \circ (h_M \ast -),
\end{align*}
where the last step uses the CDG relation $d^2 = h \ast (-)$ on
$M$ and~$N$.  When $M = N$, this is $[h_M, f]$, which vanishes
on central endomorphisms.
\end{proof}

\begin{corollary}[Contractibility of coacyclic injectives; \ClaimStatusProvedHere]
\label{cor:coacyclic-injective-contractible}
\index{coacyclic!contractibility criterion|textbf}
Let $(C, d, h)$ be a conilpotent chiral CDG-coalgebra with
finite-dimensional graded pieces, and let $M$ be a chiral
CDG-comodule with injective underlying graded $C^{\#}$-comodule.
Then $M$ is coacyclic if and only if it admits a
\emph{contracting homotopy}: a degree~$-1$ graded
$C^{\#}$-comodule endomorphism $s \colon M \to M$ satisfying
\begin{equation}\label{eq:contracting-homotopy-cdg}
d_M \circ s + s \circ d_M \;=\; \mathrm{id}_M.
\end{equation}
\end{corollary}

\begin{proof}
If such an $s$ exists, then $\mathrm{id}_M = d_{\mathrm{Hom}}(s)$
in the CDG Hom-complex of
Proposition~\ref{prop:cdg-hom-complex} (with $M = N$), so $M$ is
null-homotopic and \emph{a fortiori} coacyclic.

Conversely, suppose $M$ is coacyclic with injective underlying
graded comodule.  By the semiorthogonal decomposition in
Proposition~\ref{prop:chiral-inj-proj-resolutions}(c),
$D^{\mathrm{co}}(C\text{-}\mathrm{comod}^{\mathrm{ch}})$
is equivalent to
$\mathrm{Hot}(C\text{-}\mathrm{comod}^{\mathrm{ch}}_{\mathrm{inj}})$.
A coacyclic object that already has injective underlying comodule
represents the zero object in this homotopy category, hence is
null-homotopic: there exists a degree~$-1$ closed morphism
$s$ with $d_{\mathrm{Hom}}(s) = \mathrm{id}_M$.
Expanding via~\eqref{eq:cdg-hom-differential} gives~\eqref{eq:contracting-homotopy-cdg}.
\end{proof}

\subsection{The correspondence functors}
\label{subsec:chiral-correspondence-functors}

\begin{definition}[Functors \texorpdfstring{$\Phi_C^{\mathrm{ch}}$}{Phi_C^ch} and \texorpdfstring{$\Psi_C^{\mathrm{ch}}$}{Psi_C^ch}]
\label{def:chiral-Phi-Psi}
\index{Positselski!functors PhiPsi@$\Phi_C$, $\Psi_C$}
Let $(C, d, h)$ be a chiral CDG-coalgebra.
\begin{itemize}
\item For a left chiral CDG-contramodule $P$, define
  $\Phi_C^{\mathrm{ch}}(P)$ to be the left chiral CDG-comodule
  whose underlying graded $\mathcal{D}_X$-module is the
  chiral contratensor product $C \odot_C^{\mathrm{ch}} P$.
  Concretely, $\Phi_C^{\mathrm{ch}}(P)$ is a quotient of
  $C \chirtensor P$; the comodule structure comes from the
  comultiplication on the left factor~$C$, and the differential is
  induced by the differentials on $C$ and $P$ via the standard
  formula for the tensor product of graded $\mathcal{D}_X$-modules
  with differentials.
\item For a left chiral CDG-comodule $M$, define
  $\Psi_C^{\mathrm{ch}}(M)$ to be the left chiral CDG-contramodule
  whose underlying graded $\mathcal{D}_X$-module is
  $\mathrm{Hom}_C^{\mathrm{ch}}(C, M) = \mathrm{Cohom}_C^{\mathrm{ch}}(C, M)$.
  The contramodule structure comes from the comultiplication
  on $C$, and the differential is induced from $d_C$ and $d_M$.
\end{itemize}
These define DG-functors:
\begin{align*}
\Phi_C^{\mathrm{ch}}\colon \mathrm{DG}(C\text{-}\mathrm{contra}^{\mathrm{ch}})
&\longrightarrow \mathrm{DG}(C\text{-}\mathrm{comod}^{\mathrm{ch}}), \\
\Psi_C^{\mathrm{ch}}\colon \mathrm{DG}(C\text{-}\mathrm{comod}^{\mathrm{ch}})
&\longrightarrow \mathrm{DG}(C\text{-}\mathrm{contra}^{\mathrm{ch}}).
\end{align*}
\end{definition}

\begin{lemma}[Key properties of \texorpdfstring{$\Phi_C^{\mathrm{ch}}$}{Phi_C^ch} and
\texorpdfstring{$\Psi_C^{\mathrm{ch}}$}{Psi_C^ch}; \ClaimStatusProvedHere]
\label{lem:Phi-Psi-properties}
\begin{enumerate}[label=\textup{(\alph*)}]
\item $\Phi_C^{\mathrm{ch}}$ is left adjoint to $\Psi_C^{\mathrm{ch}}$
  on the homotopy categories.
\item $\Phi_C^{\mathrm{ch}}$ maps the free graded contramodule
  $\mathrm{Hom}^{\mathrm{ch}}(C, V)$ to the cofree graded comodule
  $C \chirtensor V$, for any graded $\mathcal{D}_X$-module $V$.
\item $\Psi_C^{\mathrm{ch}}$ maps the cofree graded comodule
  $C \chirtensor V$ to the free graded contramodule
  $\mathrm{Hom}^{\mathrm{ch}}(C, V)$.
\item $\Phi_C^{\mathrm{ch}}$ maps projective graded contramodules
  to injective graded comodules, and $\Psi_C^{\mathrm{ch}}$ maps
  injective graded comodules to projective graded contramodules.
\end{enumerate}
\end{lemma}

\begin{proof}
\emph{Part (a).}
The adjunction $\Phi_C^{\mathrm{ch}} \dashv \Psi_C^{\mathrm{ch}}$
follows from the standard isomorphism
$\mathrm{Hom}_{\mathcal{D}_X}(C \chirtensor V, P)
\cong \mathrm{Hom}_{\mathcal{D}_X}(V, \mathrm{Hom}^{\mathrm{ch}}(C, P))$
(the chiral tensor-hom adjunction), which makes
$\Phi_C^{\mathrm{ch}}$ left adjoint to $\Psi_C^{\mathrm{ch}}$ as
DG-functors (Lemma~\ref{lem:chiral-co-contra-adjunction}).

\emph{Parts (b) and (c).}
The free graded contramodule on $V$ is
$\mathrm{Hom}^{\mathrm{ch}}(C, V) = \mathrm{Hom}_{\mathcal{D}_X}(C, V)$.
Then:
\[
\Phi_C^{\mathrm{ch}}(\mathrm{Hom}^{\mathrm{ch}}(C, V))
= C \odot_C^{\mathrm{ch}} \mathrm{Hom}^{\mathrm{ch}}(C, V).
\]
By the contramodule-comodule correspondence at the graded level
(cf.\ the example in
Positselski~\cite[\S5.1]{Positselski11}), the
contratensor product of $C$ with the free contramodule on $V$ is
canonically isomorphic to the cofree comodule $C \chirtensor V$.
In the chiral setting, this isomorphism follows from the
tensor-hom adjunction for $\mathcal{D}_X$-modules:
\begin{align*}
C \odot_C^{\mathrm{ch}} \mathrm{Hom}^{\mathrm{ch}}(C, V)
&\cong \mathrm{coker}\bigl(
C \chirtensor \mathrm{Hom}^{\mathrm{ch}}(C, \mathrm{Hom}^{\mathrm{ch}}(C, V))
\rightrightarrows C \chirtensor \mathrm{Hom}^{\mathrm{ch}}(C, V)
\bigr) \\
&\cong C \chirtensor V,
\end{align*}
where the last step uses the counit of the adjunction
$\mathrm{ev}\colon \mathrm{Hom}^{\mathrm{ch}}(C, V) \chirtensor C
\to V$ and the coassociativity of $C$.  Part~(c) is the dual computation.

\emph{Part (d).}
Combine (b)--(c) with
Proposition~\ref{prop:chiral-inj-proj-resolutions}(a)--(b):
free graded contramodules are projective and cofree graded comodules
are injective, so $\Phi_C^{\mathrm{ch}}$ sends projective objects
to injective objects and vice versa.
\end{proof}

\subsection{The chiral comodule-contramodule correspondence}
\label{subsec:chiral-co-contra-correspondence}

\begin{theorem}[Chiral comodule-contramodule correspondence;
\ClaimStatusProvedHere]
\label{thm:chiral-co-contra-correspondence}
\index{comodule-contramodule correspondence!chiral|textbf}
\index{Positselski!chiral correspondence|textbf}
Let $(C, d, h)$ be a conilpotent chiral CDG-coalgebra on a smooth
curve $X$ with finite-dimensional graded pieces.  The derived functors
\begin{align*}
\mathbb{L}\Phi_C^{\mathrm{ch}}\colon
D^{\mathrm{ctr}}(C\text{-}\mathrm{contra}^{\mathrm{ch}})
&\longrightarrow
D^{\mathrm{co}}(C\text{-}\mathrm{comod}^{\mathrm{ch}}), \\
\mathbb{R}\Psi_C^{\mathrm{ch}}\colon
D^{\mathrm{co}}(C\text{-}\mathrm{comod}^{\mathrm{ch}})
&\longrightarrow
D^{\mathrm{ctr}}(C\text{-}\mathrm{contra}^{\mathrm{ch}})
\end{align*}
are mutually inverse equivalences of triangulated categories.
\end{theorem}

\begin{proof}
The proof follows Positselski~\cite[\S5.2]{Positselski11}, with the
ordinary tensor product replaced by the chiral tensor product
$\chirtensor$ and exactness properties established via the geometry
of $\mathcal{D}_X$-modules on configuration spaces.  We verify
each step.

\emph{Step~1: Derived functors.}
By Proposition~\ref{prop:chiral-inj-proj-resolutions}(c),
the contraderived and coderived categories
\begin{align*}
D^{\mathrm{ctr}}(C\text{-}\mathrm{contra}^{\mathrm{ch}})
&\simeq K_{\mathrm{proj}}(C\text{-}\mathrm{contra}^{\mathrm{ch}}),\\
D^{\mathrm{co}}(C\text{-}\mathrm{comod}^{\mathrm{ch}})
&\simeq K_{\mathrm{inj}}(C\text{-}\mathrm{comod}^{\mathrm{ch}})
\end{align*}
are equivalent to the homotopy categories of
CDG-contra\-modules (resp.\ CDG-comodules) with
projective (resp.\ injective) underlying graded module.
The left derived functor~$\mathbb{L}\Phi_C^{\mathrm{ch}}$
restricts to projective contramodules;
$\mathbb{R}\Psi_C^{\mathrm{ch}}$ restricts to injective comodules.

\emph{Step~2: Mutual inverse on generators.}
By Lemma~\ref{lem:Phi-Psi-properties}(b)--(c), $\Phi_C^{\mathrm{ch}}$
sends the free graded contramodule $\mathrm{Hom}^{\mathrm{ch}}(C, V)$
to the cofree graded comodule $C \chirtensor V$, and
$\Psi_C^{\mathrm{ch}}$ sends $C \chirtensor V$ back to
$\mathrm{Hom}^{\mathrm{ch}}(C, V)$.  These are the generating objects
for the respective categories (any CDG-comodule has an injective
resolution by cofree comodules, and any CDG-contramodule has a
projective resolution by free contramodules, by
Proposition~\ref{prop:chiral-inj-proj-resolutions}).

\emph{Step~3: Equivalence on full triangulated subcategories.}
We claim that the restrictions of $\Phi_C^{\mathrm{ch}}$ and
$\Psi_C^{\mathrm{ch}}$ to the full triangulated subcategories
\begin{align*}
\mathrm{Hot}(C\text{-}\mathrm{contra}^{\mathrm{ch}}_{\mathrm{proj}})
&\subset \mathrm{Hot}(C\text{-}\mathrm{contra}^{\mathrm{ch}}), \\
\mathrm{Hot}(C\text{-}\mathrm{comod}^{\mathrm{ch}}_{\mathrm{inj}})
&\subset \mathrm{Hot}(C\text{-}\mathrm{comod}^{\mathrm{ch}})
\end{align*}
are mutually inverse equivalences.

To see this, observe that the unit and counit natural transformations
\[
\eta\colon \mathrm{id} \to \Psi_C^{\mathrm{ch}} \circ \Phi_C^{\mathrm{ch}},
\qquad
\varepsilon\colon \Phi_C^{\mathrm{ch}} \circ \Psi_C^{\mathrm{ch}} \to \mathrm{id}
\]
are isomorphisms on free contramodules and cofree comodules
(Step~2).  Since these objects generate the respective
homotopy categories (every object is built from generators by
shifts, cones, and countable (co)products),
and both $\Phi_C^{\mathrm{ch}}$ and $\Psi_C^{\mathrm{ch}}$ preserve
exact triangles and countable (co)products, the natural
transformations $\eta$ and $\varepsilon$ are isomorphisms
on all objects with projective (resp.\ injective) underlying graded
structure.

\emph{Step~4: Descent to exotic derived categories.}
To pass from the full homotopy categories to $D^{\mathrm{co}}$ and
$D^{\mathrm{ctr}}$, we verify that $\Phi_C^{\mathrm{ch}}$ maps
contraacyclic CDG-contramodules to coacyclic CDG-comodules, and
$\Psi_C^{\mathrm{ch}}$ maps coacyclic CDG-comodules to
contraacyclic CDG-contramodules.

For $\Phi_C^{\mathrm{ch}}$: since $\Phi_C^{\mathrm{ch}}$ is a
left adjoint, it preserves (all) colimits, in particular coproducts
and cokernels.  A contraacyclic CDG-contramodule is built from
totalizations of exact triples and direct products.  Now,
$\Phi_C^{\mathrm{ch}}$ preserves exact triples (being a left adjoint
of an exact functor, hence right exact; and by
Lemma~\ref{lem:Phi-Psi-properties}(d) it sends projective
contramodules to injective comodules, so the left derived functor
is exact).  The image of a totalization under $\Phi_C^{\mathrm{ch}}$
is a totalization, and the image of a direct product is computed
using the cofree resolution from Step~2.  Thus
$\Phi_C^{\mathrm{ch}}$ preserves contraacyclic $\to$ coacyclic.
The argument for $\Psi_C^{\mathrm{ch}}$ is dual.

\emph{Step~5: Conclusion.}
Combining Steps~3 and~4:  $\mathbb{L}\Phi_C^{\mathrm{ch}}$ and
$\mathbb{R}\Psi_C^{\mathrm{ch}}$ are well-defined triangulated
functors between $D^{\mathrm{ctr}}$ and $D^{\mathrm{co}}$, they
are mutually inverse on the generating subcategories (projective
contramodules $\leftrightarrow$ injective comodules), and both
categories are compactly generated by finite-dimensional objects
(using conilpotency and finite-type grading; the argument parallels
Positselski~\cite[\S5.5]{Positselski11}, with the key input that any
graded $C$-comodule is the union of its finite-dimensional
subcomodules, which holds for conilpotent chiral coalgebras because
the conformal weight grading bounds the coaction filtration).
By the uniqueness of extensions of equivalences from compact generators
to the full compactly generated category
(cf.~\cite[\S5.5]{Positselski11}), the equivalence extends to all objects.
\end{proof}

\begin{remark}[Hypotheses]\label{rem:co-contra-hypotheses}
The two key hypotheses---conilpotency and finite-type grading---are
satisfied by $C = \bar{B}^{\mathrm{ch}}(\cA)$ whenever $\cA$ is a
finitely generated chiral algebra with finite-dimensional conformal
weight spaces (which includes all Koszul chiral algebras arising in
this monograph: Kac--Moody, Virasoro, $W$-algebras, Heisenberg, free
fields, and Yangians).  Conilpotency was established in
Theorem~\ref{thm:coalgebra-via-NAP}(4); finite-type grading follows
from the finite-generation hypothesis on~$\cA$.
\end{remark}

\subsection{Application: the Positselski equivalence for chiral algebras}
\label{subsec:positselski-chiral-equivalence}

The chiral comodule-contramodule correspondence yields the derived module category equivalence at all genera.

\begin{theorem}[Positselski equivalence for chiral algebras;
\ClaimStatusProvedHere]
\label{thm:positselski-chiral-proved}
\index{Positselski!chiral equivalence|textbf}
\index{comodule-contramodule correspondence!chiral Koszul|textbf}
For a Koszul chiral algebra $\cA$ on a smooth curve $X$ with
curved bar complex $C = \bar{B}^{\mathrm{ch}}(\cA)$
(a conilpotent chiral CDG-coalgebra by
Example~\textup{\ref{ex:bar-as-CDG}}), there is an equivalence of
triangulated categories:
\begin{equation}\label{eq:positselski-chiral-proved}
D^{\mathrm{co}}(\cA\text{-}\mathrm{CoMod}^{\mathrm{conil,\,ch}})
\;\simeq\;
D^{\mathrm{ctr}}(\cA^!\text{-}\mathrm{ContraMod}^{\mathrm{ch}})
\end{equation}
between the coderived category of conilpotent chiral
$\cA$-comodules and the contraderived category of chiral
$\cA^!$-contramodules.

At genus~$0$ (where $h = 0$), this reduces to the module Koszul
duality of
Theorem~\textup{\ref{thm:e1-module-koszul-duality}}.
At genus~$g \geq 1$ (where $h = m_0^{(g)} \neq 0$), the ordinary
derived categories are trivial, and the coderived/contraderived
categories provide the correct framework.
\end{theorem}

\begin{proof}
Apply Theorem~\ref{thm:chiral-co-contra-correspondence} with
$C = \bar{B}^{\mathrm{ch}}(\cA)$.

\emph{Identification of comodule category.}
A conilpotent chiral CDG-comodule over
$C = \bar{B}^{\mathrm{ch}}(\cA)$ is a graded $\mathcal{D}_X$-module~$M$
with coaction $M \to \bar{B}^{\mathrm{ch}}(\cA) \chirtensor M$ and
differential satisfying $\dfib^{\,2} = \mcurv{g} \ast (-)$ (Convention~\ref{conv:higher-genus-differentials}).  This is
precisely a (curved) $\cA$-module in the sense that the coaction
encodes the $\cA$-module structure via the bar construction:
$\cA\text{-}\mathrm{CoMod}^{\mathrm{conil,\,ch}}
\cong \bar{B}^{\mathrm{ch}}(\cA)\text{-}\mathrm{comod}^{\mathrm{ch}}$
(this is the module-level bar construction,
cf.\ the proof of Theorem~\ref{thm:e1-module-koszul-duality},
Step~1).

\emph{Identification of contramodule category.}
The dual algebra $C^* = \bar{B}^{\mathrm{ch}}(\cA)^*$ is canonically
identified with $\cA^!$ (the Koszul dual chiral algebra, via
Theorem~\ref{thm:bar-computes-dual}).  A chiral CDG-contramodule
over $C$ is a graded $\mathcal{D}_X$-module $P$ with contraaction
$\mathrm{Hom}^{\mathrm{ch}}(\bar{B}^{\mathrm{ch}}(\cA), P) \to P$,
which (by the finite-type duality of Remark~\ref{rem:chiral-contramodules})
is equivalent to a complete (pro-nilpotent) module over~$\cA^!$:
$\cA^!\text{-}\mathrm{ContraMod}^{\mathrm{ch}}
\cong \bar{B}^{\mathrm{ch}}(\cA)\text{-}\mathrm{contra}^{\mathrm{ch}}$.

\emph{Conclusion.}
The hypotheses of Theorem~\ref{thm:chiral-co-contra-correspondence}
are satisfied: $\bar{B}^{\mathrm{ch}}(\cA)$ is conilpotent
(Theorem~\ref{thm:coalgebra-via-NAP}(4)) with finite-dimensional
graded pieces (since $\cA$ is finitely generated with
finite-dimensional conformal weight spaces).  The equivalence
$D^{\mathrm{co}}(C\text{-}\mathrm{comod}^{\mathrm{ch}})
\simeq D^{\mathrm{ctr}}(C\text{-}\mathrm{contra}^{\mathrm{ch}})$
of Theorem~\ref{thm:chiral-co-contra-correspondence} therefore yields
\eqref{eq:positselski-chiral-proved}.
\end{proof}

\begin{theorem}[Full derived module equivalence; \ClaimStatusProvedHere]
\label{thm:full-derived-module-equiv-proved}
\index{derived category!full module equivalence|textbf}
\index{Positselski!full derived equivalence|textbf}
For a Koszul chiral algebra $\cA$ with curved bar complex
$\bar{B}(\cA)$, there is a triangulated equivalence:
\begin{equation}\label{eq:full-derived-equiv-proved}
D^b(\mathrm{Mod}^{\mathrm{compl}}(\cA))
\;\simeq\;
D^{\mathrm{co}}(\mathrm{CoMod}^{\mathrm{conil}}(\cA^!))
\end{equation}
where $\mathrm{Mod}^{\mathrm{compl}}(\cA)$ denotes the category
of complete \textup{(}pro-nilpotent\textup{)} $\cA$-modules
and $\mathrm{CoMod}^{\mathrm{conil}}(\cA^!)$ denotes the
category of conilpotent $\cA^!$-comodules.

This equivalence induces the Ext duality:
\begin{equation}\label{eq:ext-duality-proved}
\mathrm{Ext}^n_{\cA}(M_1, M_2)
\;\cong\;
\mathrm{Ext}^{d-n}_{\cA^!}(\Phi(M_1), \Phi(M_2))^{\vee}
\end{equation}
where $d$ is the global dimension of $\cA$ and
$\Phi = \bar{B}_{\cA}$ is the bar functor.
\end{theorem}

\begin{proof}
\emph{Step~1.}
By Theorem~\ref{thm:positselski-chiral-proved}:
\[
D^{\mathrm{co}}(\cA\text{-}\mathrm{CoMod}^{\mathrm{conil,\,ch}})
\simeq
D^{\mathrm{ctr}}(\cA^!\text{-}\mathrm{ContraMod}^{\mathrm{ch}}).
\]

\emph{Step~2.}
By Remark~\ref{rem:chiral-contramodules}, complete $\cA^!$-modules
are equivalent to chiral CDG-contramodules over $\bar{B}^{\mathrm{ch}}(\cA)$.
Moreover, the contraderived category of such contramodules is
identified with the bounded derived category of complete
$\cA$-modules:
$D^{\mathrm{ctr}}(\cA^!\text{-}\mathrm{ContraMod}^{\mathrm{ch}})
\simeq D^b(\mathrm{Mod}^{\mathrm{compl}}(\cA))$.
This identification uses the finite homological dimension:
since $\cA$ is Koszul, the bar complex
$\bar{B}^{\mathrm{ch}}(\cA)$ has finite homological dimension
(bounded by the global dimension of the underlying operadic
structure), so the contraderived and ordinary derived categories
coincide on the algebra side
(Positselski~\cite[\S4.5]{Positselski11}).

\emph{Step~3.}
Combining: $D^b(\mathrm{Mod}^{\mathrm{compl}}(\cA)) \simeq
D^{\mathrm{co}}(\mathrm{CoMod}^{\mathrm{conil}}(\cA^!))$.

\emph{Ext duality.}
The equivalence $\Phi = \bar{B}_{\cA}$ sends
$\mathrm{Hom}^*_{\cA}(M_1, M_2)$ to
$\mathrm{Cohom}^*_{\cA^!}(\Phi(M_1), \Phi(M_2))$.
By the comodule-contramodule correspondence
(Theorem~\ref{thm:chiral-co-contra-correspondence}, \S5.3 of
\cite{Positselski11}), the cohomomorphism complex is identified
with the Ext complex on the dual side, with the dimension shift
coming from Verdier duality on configuration spaces:
$\mathrm{Ext}^n_{\cA}(M_1, M_2) \cong
\mathrm{Ext}^{d-n}_{\cA^!}(\Phi(M_1), \Phi(M_2))^{\vee}$.
\end{proof}

\begin{remark}[Relation to genus-0 module Koszul duality]
\label{rem:genus-0-recovery}
At genus~$0$, the curvature $h = 0$, so CDG-comodules reduce to
DG-comodules, the coderived category reduces to the ordinary derived
category, and Theorem~\ref{thm:full-derived-module-equiv-proved}
recovers the module Koszul duality of
Theorem~\ref{thm:e1-module-koszul-duality}:
\[
D^b(\mathrm{Mod}^{\mathrm{compl}}(\cA))
\simeq D(\mathrm{CoMod}^{\mathrm{conil}}(\cA^!)).
\]
The virtue of the coderived/contraderived framework is that it
provides the correct generalization to genus~$g \geq 1$, where the
curvature $m_0^{(g)} \neq 0$ makes the ordinary derived category
trivial.
\end{remark}

\begin{remark}[The role of the configuration space geometry]
\label{rem:config-space-role}
The proof of the chiral comodule-contramodule correspondence
(Theorem~\ref{thm:chiral-co-contra-correspondence}) uses the
geometry of configuration spaces in three places:
\begin{enumerate}
\item \emph{Exactness of the chiral tensor product.}
  The tensor-hom adjunction for holonomic $\mathcal{D}_X$-modules
  on $\overline{C}_n(X)$ ensures that cofree comodules are injective
  and free contramodules are projective
  (Proposition~\ref{prop:chiral-inj-proj-resolutions}).
\item \emph{Conilpotency from conformal weight.}
  The conformal weight grading on $\bar{B}^{\mathrm{ch}}(\cA)$
  (inherited from $\cA$) ensures conilpotency
  (Theorem~\ref{thm:coalgebra-via-NAP}(4)), which is needed for
  the compact generation argument in Step~5 of the proof.
\item \emph{Finite-dimensionality from holonomicity.}
  The holonomicity of the bar complex on each configuration space
  stratum (Lemma~\ref{lem:bar-holonomicity}) ensures
  finite-dimensional graded pieces, which is
  needed for the duality between contramodules and complete modules
  (Remark~\ref{rem:chiral-contramodules}).
\end{enumerate}
These geometric inputs are specific to the chiral setting and have
no analogue in Positselski's ground-field framework.
\end{remark}

%================================================================
% SECTION: BAR-COBAR INVERSION - COMPLETE QUASI-ISOMORPHISM
%================================================================

\section{Bar-cobar inversion}
\label{sec:bar-cobar-inversion-quasi-iso}

\subsection{Statement of the main result}

\begin{theorem}[Bar-cobar inversion is quasi-isomorphism; \ClaimStatusProvedHere]\label{thm:bar-cobar-inversion-qi}
\textup{[Regime: quadratic at genus~$0$; curved-central at genus~$g \geq 1$
on the Koszul locus; filtered-complete off it
\textup{(}cf.\ Convention~\textup{\ref{conv:regime-tags})}.]}

The Heisenberg inversion $\Omega^{\mathrm{ch}}(\bar{B}^{\mathrm{ch}}(\mathcal{H}_k)) \xrightarrow{\sim} \mathcal{H}_k$ (\S\ref{sec:frame-inversion}) generalizes to all Koszul chiral algebras.

Let $\mathcal{A}$ be a \emph{Koszul} chiral algebra on a Riemann surface $X$
(Definition~\ref{def:chiral-koszul-pair}), with
$\bar{B}^{\mathrm{ch}}(\mathcal{A})$ conilpotent or $\mathcal{A}$ complete
with respect to its augmentation ideal
(\S\ref{sec:i-adic-completion}).  Then the natural map:
\[\psi: \Omega(\bar{B}(\mathcal{A})) \longrightarrow \mathcal{A}\]
induced by the bar-cobar adjunction is a \emph{quasi-isomorphism}, not merely
an isomorphism in cohomology.

More precisely:
\begin{enumerate}
\item The map $\psi$ is a morphism of chiral algebras (respects all structure)

\item At each genus $g$, the genus-$g$ component:
      \[\psi_g: \Omega_g(\bar{B}_g(\mathcal{A})) \longrightarrow \mathcal{A}\]
      is a quasi-isomorphism

\item The full genus-graded map:
      \[\psi = \bigoplus_{g=0}^\infty \psi_g: \Omega(\bar{B}(\mathcal{A})) \longrightarrow \mathcal{A}\]
      converges and is a quasi-isomorphism in the $\hbar$-adic completion

\item There exists a spectral sequence converging to $H^\bullet(\mathcal{A})$
      with $E_1$-page given by the bar-cobar complex; it collapses at $E_2$
      by the Koszul property (Theorem~\ref{thm:spectral-sequence-collapse})
\end{enumerate}

\smallskip\noindent\emph{Scope.}
The Koszul hypothesis is essential: for admissible-level affine algebras
and minimal-model Virasoro/W-algebras the bar spectral sequence has
nontrivial higher differentials and the counit $\psi$ is \emph{not} a
quasi-isomorphism (see Remark~\ref{rem:sl2-admissible}
and~\ref{rem:virasoro-module-koszul-minimal}).  In those regimes one must work
with the completed or coderived bar-cobar correspondence
(Corollary~\ref{cor:bar-cobar-inverse},
\S\ref{sec:i-adic-completion}).
\end{theorem}

\begin{proof}[Dependency-closed proof]
We verify each numbered clause by previously established results in this chapter.
\begin{enumerate}[label=(D\arabic*)]
\item \emph{Morphism of chiral algebras.}
The map $\psi$ is the counit of the bar-cobar adjunction from
Theorem~\ref{thm:bar-cobar-adjunction}; therefore it is canonical and respects
the algebraic structure by functoriality of bar and cobar
(Theorem~\ref{thm:bar-functorial}, Theorem~\ref{thm:cobar-functorial}).

\item \emph{Genuswise quasi-isomorphism.}
The genus-$0$ base case is proved independently by
Theorem~\ref{thm:bar-nilpotency-complete} and
Theorem~\ref{thm:chiral-koszul-duality};
equivalently, it is the implication
\ref{ftm:koszul}$\Rightarrow$\ref{ftm:counit} of
the fundamental theorem of chiral twisting morphisms
(Theorem~\ref{thm:fundamental-twisting-morphisms}).
The genus-$g$ components $\psi_g$ for $g \geq 1$ are quasi-isomorphisms by
induction on genus via
Theorem~\ref{thm:genus-graded-convergence}, item~(2), together with the
higher-genus inversion package (Theorem~\ref{thm:higher-genus-inversion}).

\item \emph{Full genus-graded convergence and quasi-isomorphism.}
This is Theorem~\ref{thm:genus-graded-convergence}, item (3), which identifies
the completed genus series $\psi=\sum_{g\ge 0}\hbar^{2g-2}\psi_g$ and proves convergence.

\item \emph{Spectral sequence statement.}
Existence and explicit page description are exactly
Theorem~\ref{thm:bar-cobar-spectral-sequence}; collapse for Koszul input is
Theorem~\ref{thm:spectral-sequence-collapse}.
\end{enumerate}
Combining (D1)--(D4) gives all four clauses of the theorem.
\end{proof}

\begin{remark}[Relation to the fundamental theorem \texorpdfstring{$\mathrm{A}_0$}{A_0}]
\label{rem:inversion-vs-fundamental}
At genus~$0$, Theorem~\ref{thm:bar-cobar-inversion-qi} is a
consequence of the fundamental theorem of chiral twisting morphisms
(Theorem~\ref{thm:fundamental-twisting-morphisms}):
the counit $\Omega_X(\bar{B}_X(\cA)) \to \cA$ being a
quasi-isomorphism is one of the four equivalent characterizations
of Koszulity established there.  The present theorem extends this
to all genera via the genus induction of clause~(D2).
\end{remark}

\begin{remark}[Quasi-isomorphism of chain complexes vs chiral algebras]\label{rem:qi-vs-homology-iso}
The distinction is between two levels of structure:

\emph{Chain complex quasi-isomorphism.} A chain map $\psi: C^\bullet \to D^\bullet$ inducing an isomorphism $H^\bullet(\psi): H^\bullet(C) \xrightarrow{\cong} H^\bullet(D)$ on cohomology. This is the standard notion; a ``homology isomorphism'' and a ``quasi-isomorphism'' are the \emph{same thing} at the chain complex level.

\emph{Chiral algebra quasi-isomorphism.} An $A_\infty$ (or chiral algebra) morphism $\psi\colon \mathcal{A} \to \mathcal{B}$ whose linear component $\psi_1\colon \mathcal{A} \to \mathcal{B}$ is a chain complex quasi-isomorphism. Such a map preserves the full $A_\infty$ structure up to coherent homotopy. The distinction matters:
A chain complex quasi-isomorphism tells us only about $H^\bullet$ and need not respect multiplicative or $A_\infty$ structure, whereas a chiral algebra quasi-isomorphism gives a full equivalence in the derived category of chiral algebras, preserving all homotopy-theoretic information.  For Koszul duality, the map $\Omega(B(\mathcal{A})) \to \mathcal{A}$ must be a quasi-isomorphism of \emph{chiral algebras}, not merely of chain complexes, to establish derived equivalences.

\emph{Example.}
Consider a dga $(C^\bullet, d, \mu)$ with:
\[\cdots \to 0 \to \mathbb{C} \xrightarrow{0} \mathbb{C} \to 0 \to \cdots\]

This has $H^0 = \mathbb{C}$, $H^i = 0$ for $i \neq 0$.

Compare with the algebra $(D^\bullet, 0, \nu) = (\mathbb{C}, 0, 0)$ in degree 0.
The projection $\pi: C^\bullet \to D^\bullet$ is a quasi-isomorphism of chain
complexes (both have $H^0 = \mathbb{C}$), but $\pi$ does \emph{not} respect the
multiplication ($\pi$ need not intertwine $\mu$ and $\nu$ at the chain level).
To obtain a quasi-isomorphism of \emph{algebras}, one must work with $A_\infty$
morphisms, which encode the higher homotopies relating the two multiplicative
structures.
\end{remark}

\subsection{Proof strategy and filtration}

\begin{definition}[Bar-cobar filtration]\label{def:bar-cobar-filtration}
Define a decreasing filtration on $\Omega(\bar{B}(\mathcal{A}))$ by:
\[F^p\Omega(\bar{B}(\mathcal{A})) = \bigoplus_{n \geq p} \Omega^n(\bar{B}^n(\mathcal{A}))\]

This is the filtration by \emph{bar degree} (= cobar arity).

\emph{Geometric meaning.} $F^p$ consists of elements involving at least $p$ points
in configuration space. As $p \to \infty$, we are considering increasingly complicated
configurations.

\emph{Properties.}
\begin{enumerate}
\item $F^0 \supseteq F^1 \supseteq F^2 \supseteq \cdots$
\item $\bigcap_{p=0}^\infty F^p = 0$ (completeness)
\item The differential respects filtration: $d(F^p) \subseteq F^p$
\item The natural map factors through the filtration
\end{enumerate}
\end{definition}

\begin{lemma}[Associated graded; \ClaimStatusProvedHere]\label{lem:bar-cobar-associated-graded}
The associated graded of the bar-cobar filtration is:
\[\text{Gr}^p\Omega(\bar{B}(\mathcal{A})) = \Omega^p(\bar{B}^p(\mathcal{A}))\]

The induced differential on $\text{Gr}^\bullet$ is the \emph{internal differential} $d_{\text{internal}}$ alone (i.e., the differential inherited from $\mathcal{A}$).
\end{lemma}

\begin{proof}
By definition of associated graded:
\[\text{Gr}^p = F^p / F^{p+1} = \Omega^p(\bar{B}^p(\mathcal{A}))\]

The full differential on $\Omega(\bar{B}(\mathcal{A}))$ decomposes as $d = d_{\text{internal}} + d_{\text{bar}} + d_{\text{cobar}}$, where $d_{\text{internal}}$ preserves the filtration degree (it acts within each tensor factor), $d_{\text{bar}}$ preserves or lowers filtration degree (collisions reduce arity), and $d_{\text{cobar}}$ raises filtration degree by~$1$ (comultiplication increases cobar arity).  On the associated graded $\text{Gr}^p = F^p/F^{p+1}$, only the \emph{filtration-preserving} component survives: $d_{\text{cobar}}$ maps $F^p$ into $F^{p+1}$ and so vanishes on $\text{Gr}^p$, while $d_{\text{bar}}$ maps $F^p$ into $F^{p-1} \subset F^p$ and so also acts trivially on $\text{Gr}^p$.  Thus $d_{\text{gr}} = d_{\text{internal}}$, which is the $d_0$ differential of the spectral sequence (cf.\ Theorem~\ref{thm:bar-cobar-spectral-sequence}).
\end{proof}

\subsection{Spectral sequence construction}

\begin{theorem}[Bar-cobar spectral sequence; \ClaimStatusProvedHere]\label{thm:bar-cobar-spectral-sequence}
\label{thm:koszul-spectral-sequence}
The filtration from Definition~\ref{def:bar-cobar-filtration} induces a spectral 
sequence:
\[E_0^{p,q} = \left(F^p / F^{p+1}\right)^{p+q} \implies H^{p+q}(\Omega(\bar{B}(\mathcal{A})))\]
converging to the cohomology of $\Omega(\bar{B}(\mathcal{A}))$.  Here $p$ is the filtration degree (cobar arity) and $q$ is the complementary degree (from the internal and bar gradings), so the total degree is $p + q$.

\emph{Explicit description of pages.}
\begin{align*}
E_0^{p,q} &= \left(F^p\Omega^{p+q}(\bar{B}(\mathcal{A})) \big/ F^{p+1}\Omega^{p+q}(\bar{B}(\mathcal{A}))\right) \quad \text{(associated graded)} \\
E_1^{p,q} &= H^{p+q}\!\left(E_0^{p,\bullet},\, d_{\text{internal}}\right)
            \quad \text{(internal cohomology at fixed cobar degree $p$)} \\
E_2^{p,q} &= H^p\!\left(E_1^{\bullet,q},\, d_{\text{bar}}\right)
            \quad \text{(bar cohomology)} \\
E_\infty^{p,q} &= \mathrm{Gr}^p H^{p+q}(\Omega(\bar{B}(\mathcal{A}))) \quad \text{(limiting page)}
\end{align*}
\end{theorem}

\begin{proof}
This is a standard spectral sequence associated to a filtered complex (Weibel~\cite{Weibel94}, Chapter~5).  We verify the key properties:

\emph{Step 1: $E_0$ page.} This is just the raw complex with its bigrading:
\[E_0^{p,q} = F^p\Omega^{p+q}(\bar{B}(\mathcal{A})) / F^{p+1}\Omega^{p+q}(\bar{B}(\mathcal{A}))\]

By definition of the filtration, this isolates the piece in cobar degree~$p$ and total degree $p+q$.

\emph{Step 2: $d_0$ differential.} On the $E_0$ page:
\[d_0: E_0^{p,q} \to E_0^{p,q+1}\]
is the \emph{internal differential} $d_{\text{internal}}$ (from the differential 
on $\mathcal{A}$ itself).

Taking cohomology gives the $E_1$ page.

\emph{Step 3: $d_1$ differential.} On the $E_1$ page:
\[d_1: E_1^{p,q} \to E_1^{p+1,q}\]
is induced by the \emph{bar differential} $d_{\text{bar}}$ (collisions in 
configuration space).

Taking cohomology gives the $E_2$ page.

\emph{Step 4: Higher differentials.} For $r \geq 2$:
\[d_r: E_r^{p,q} \to E_r^{p+r,q-r+1}\]

These differentials encode higher-order interactions between bar and cobar operations.

\emph{Step 5: Convergence.} The spectral sequence converges because:
\begin{enumerate}
\item The filtration is complete: $\bigcap_p F^p = 0$
\item The filtration is exhaustive: $\bigcup_p F^p = \Omega(\bar{B}(\mathcal{A}))$
\item The complex is bounded in each column (fixed $p$)
\end{enumerate}

By standard spectral sequence theory (Weibel \cite{Weibel94}, Chapter 5), this ensures:
\[E_\infty^{p,q} = \text{Gr}^p H^{p+q}(\Omega(\bar{B}(\mathcal{A})))\]
\end{proof}

\begin{theorem}[Collapse at \texorpdfstring{$E_2$}{E2}; \ClaimStatusProvedHere]\label{thm:spectral-sequence-collapse}
\textup{[Regime: quadratic, Koszul locus
\textup{(}Convention~\textup{\ref{conv:regime-tags})}.]}

For the Heisenberg algebra, the spectral sequence collapsed trivially because $d_{\mathrm{bracket}} = 0$ (no simple poles): the bar cohomology was concentrated in degrees $0$ and $2$ without further work (\S\ref{sec:frame-bar-all}).  For non-abelian algebras, the collapse is a consequence of the Koszul property.

For a Koszul chiral algebra $\mathcal{A}$ (Remark~\ref{rem:koszul-chiral}), the spectral sequence from
Theorem~\ref{thm:bar-cobar-spectral-sequence} collapses at the $E_2$ page:
\[E_2^{p,q} = E_\infty^{p,q}.\]
\end{theorem}

\begin{proof}
The Koszul property of $\mathcal{A}$ means that the natural inclusion $\iota: \mathcal{A}^! \hookrightarrow \bar{B}(\mathcal{A})$ of the quadratic dual into the bar complex is a quasi-isomorphism (see \cite[Theorem~3.4.1]{LV12} for the classical case and Chapter~\ref{chap:koszul-pairs} for the chiral adaptation). Since $\mathcal{A}^!$ is generated in weight~1 with quadratic relations, its image under $\iota$ lies on the diagonal $p + q = \text{const}$ in the bigrading $(p,q)$ by bar degree and internal degree.

At the $E_1$ page, we compute the cohomology of each column with respect to $d_0 = d_{\text{int}}$. The $E_1$ page of the bar-cobar spectral sequence is computed by the cohomology of the bar complex: each column $E_0^{p,\bullet}$ with its internal differential $d_0$ is precisely the bar complex $\bar{B}^p(\mathcal{A})$ in cobar degree~$p$, so $E_1^{p,q} = H^q(\bar{B}^p(\mathcal{A}), d_{\mathrm{int}})$. By the Koszul quasi-isomorphism, the bar cohomology is concentrated on the diagonal, giving $E_1^{p,q} = 0$ whenever $q \neq 0$ (after reindexing by the bar filtration degree).

At the $E_2$ page, we take cohomology with respect to $d_1$ (induced by $d_{\text{bar}}$). The $E_2$ groups inherit the diagonal concentration.

For $r \geq 2$, the differential $d_r: E_r^{p,q} \to E_r^{p+r, q-r+1}$ shifts the bidegree by $(r, 1-r)$. Since $E_2^{p,q} = 0$ for $q \neq 0$, either the source or the target of $d_r$ vanishes for every $r \geq 2$. Therefore $d_r = 0$ for all $r \geq 2$, and the spectral sequence collapses at $E_2$.
\end{proof}

\subsection{Convergence at all genera}

\begin{theorem}[Genus-graded convergence; \ClaimStatusProvedHere]\label{thm:genus-graded-convergence}
\textup{[Regime: curved-central on the Koszul locus
\textup{(}Convention~\textup{\ref{conv:regime-tags})}.]}

Let $\mathcal{A}$ be a Koszul chiral algebra satisfying the hypotheses of
Theorem~\ref{thm:bar-cobar-inversion-qi}.  The bar-cobar inversion
$\psi: \Omega(\bar{B}(\mathcal{A})) \to \mathcal{A}$
converges at each genus $g$, and the full genus-graded sum converges in the
$\hbar$-adic completion.

More precisely:
\begin{enumerate}
\item \emph{Genus zero:} 
      \[\psi_0: \Omega_0(\bar{B}_0(\mathcal{A})) \xrightarrow{\sim} \mathcal{A}\]
      is a quasi-isomorphism (classical result, BD §3.7)
      
\item \emph{Fixed genus $g$:}
      \[\psi_g: \Omega_g(\bar{B}_g(\mathcal{A})) \to \mathcal{A}_g\]
      is a quasi-isomorphism after appropriate quantum corrections (here $\mathcal{A}_g$ denotes the genus-$g$ component of the genus-graded chiral algebra)
      
\item \emph{Genus series:}
      \[\psi = \sum_{g=0}^\infty \hbar^{2g-2} \psi_g\]
      converges in the $\hbar$-adic completion for $|\hbar| < R$ (radius determined 
      by growth of moduli spaces)
\end{enumerate}
\end{theorem}

\begin{proof}
We prove each case separately.

\emph{Case 1: Genus zero (classical).}

At genus zero, we work with rational curves $\mathbb{P}^1$. The bar complex is:
\[\bar{B}_0^n(\mathcal{A}) = \Gamma\left(\overline{C}_n(\mathbb{P}^1), 
\mathcal{A}^{\boxtimes n} \otimes \Omega^\bullet\right)\]

Beilinson--Drinfeld proved \cite{BD04} Theorem 3.7.11:
\[\Omega_0(\bar{B}_0(\mathcal{A})) \xrightarrow{\sim} \mathcal{A}\]

Their proof uses:
\begin{itemize}
\item Chevalley--Cousin resolution
\item Ran space formalism
\item Descent from configuration spaces
\end{itemize}

We have verified (Theorem~\ref{thm:BD-extension-higher-genus}) that all technical 
conditions hold at genus zero.

\emph{Case 2: Fixed genus $g \geq 1$.}

At higher genus, configuration spaces fiber over moduli space:
\[\pi: \overline{C}_n(X) \to \overline{\mathcal{M}}_g\]

The bar complex becomes:
\[\bar{B}_g^n(\mathcal{A}) = R\pi_*\left(\mathcal{A}^{\boxtimes n} \otimes \Omega^\bullet_{\log}\right)\]
(equivalently, $\int_{\overline{\mathcal{M}}_{g,n}} \mathcal{A}^{\boxtimes n} \otimes \Omega^\bullet_{\log}$ in the factorization homology notation of Lemma~\ref{lem:bar-as-fact-hom-AF}).

The pushforward $\pi_*$ preserves quasi-isomorphisms:

\begin{lemma}[Derived pushforward preserves QI; \ClaimStatusProvedHere]\label{lem:pushforward-preserves-qi}
For a proper morphism $\pi: Y \to Z$ and a quasi-isomorphism $f: \mathcal{F} \xrightarrow{\sim} \mathcal{G}$
of bounded-below complexes of sheaves on $Y$, the derived pushforward:
\[R\pi_*(f): R\pi_*\mathcal{F} \xrightarrow{\sim} R\pi_*\mathcal{G}\]
is a quasi-isomorphism in $D^+(Z)$.
\end{lemma}

\begin{proof}[Proof of Lemma]
Since $R\pi_*$ is a derived functor, it preserves quasi-isomorphisms by definition (a quasi-isomorphism between complexes of sheaves becomes an isomorphism in the derived category, and $R\pi_*$ is a functor on derived categories).  The underived $\pi_*$ does \emph{not} preserve quasi-isomorphisms in general.

In our application, the complexes involved are complexes of locally free sheaves on smooth proper varieties, so the Leray spectral sequence $H^p(Z, R^q\pi_*\mathcal{F}) \Rightarrow H^{p+q}(Y, \mathcal{F})$ applies, and the statement follows from the proper base change theorem.
\end{proof}

Applying this lemma: Since $\psi$ is a quasi-isomorphism fiberwise (over each
point of $\overline{\mathcal{M}}_g$), the derived pushforward $R\pi_*(\psi)$ is also a quasi-isomorphism (the FM compactification $\overline{C}_n(X)$ is proper over $\overline{\mathcal{M}}_g$, and the complexes are bounded, so proper base change applies).

\emph{Quantum corrections.} At genus $g \geq 1$, we must account for:
\begin{itemize}
\item Central charge contributions: $\sim \int_{\overline{\mathcal{M}}_g} \lambda_1$ 
      (Hodge class)
\item Modular form corrections: Period integrals over $H^1(\Sigma_g)$
\item Anomaly cancellation: Ensures $\dfib^{\,2} = 0$ (fiberwise nilpotence)
\end{itemize}

With these corrections included (see Theorem~\ref{thm:higher-genus-inversion}), 
$\psi_g$ is a quasi-isomorphism.

\emph{Case 3: Genus series convergence.}

Consider the formal series:
\[\psi(\hbar) = \sum_{g=0}^\infty \hbar^{2g-2} \psi_g\]

\emph{Growth estimate.} By Mirzakhani's results \cite{Mirzakhani}, Weil--Petersson volumes satisfy
\[\mathrm{Vol}(\overline{\mathcal{M}}_g) \sim (2g)!\]
with factorial growth.  Consequently the formal genus series has zero radius of analytic convergence, consistent with the non-Borel-summability of the perturbative string expansion.

\emph{$\hbar$-adic completion.} For rigorous formal computations, work in:
\[\widehat{\Omega}(\bar{B}(\mathcal{A}))_\hbar = \varprojlim_n 
\Omega(\bar{B}(\mathcal{A})) / \hbar^n\]

In this completion, the series converges unconditionally.
\end{proof}

\subsection{The counit of the adjunction}

\begin{proposition}[Counit is quasi-isomorphism; \ClaimStatusProvedHere]\label{prop:counit-qi}
The counit of the bar-cobar adjunction:
\[\epsilon: \bar{B}(\Omega(\mathcal{C})) \longrightarrow \mathcal{C}\]
for a chiral coalgebra $\mathcal{C}$ is also a quasi-isomorphism (dual statement).
\end{proposition}

\begin{proof}
By Verdier duality (Theorem~\ref{thm:verdier-bar-cobar}), the counit $\epsilon_{\mathcal{C}} \colon \bar{B}^{\mathrm{ch}}(\Omega^{\mathrm{ch}}(\mathcal{C})) \to \mathcal{C}$ is the Verdier dual of the unit $\eta_{\mathcal{A}^!}$, which is a quasi-isomorphism by Theorem~\ref{thm:geom-unit}.  Since Verdier duality on the Fulton--MacPherson compactifications preserves quasi-isomorphisms (Theorem~\ref{thm:verdier-bar-cobar}), the result follows.
\end{proof}

\subsection{Functoriality of the quasi-isomorphism}

\begin{theorem}[Functoriality; \ClaimStatusProvedHere]\label{thm:bar-cobar-inversion-functorial}
\label{thm:cobar-functorial}
The quasi-isomorphism $\psi: \Omega(\bar{B}(\mathcal{A})) \xrightarrow{\sim} \mathcal{A}$ 
is \emph{functorial}: for any morphism $f: \mathcal{A} \to \mathcal{A}'$ of 
chiral algebras, the diagram commutes:

\begin{center}
\begin{tikzcd}
\Omega(\bar{B}(\mathcal{A})) \ar[r, "\psi"] \ar[d, "\Omega(\bar{B}(f))"] 
& \mathcal{A} \ar[d, "f"] \\
\Omega(\bar{B}(\mathcal{A}')) \ar[r, "\psi'"] 
& \mathcal{A}'
\end{tikzcd}
\end{center}
\end{theorem}

\begin{proof}
This follows from the functoriality of the bar construction (Theorem~\ref{thm:bar-functorial}) and the universal property of the bar-cobar adjunction (Theorem~\ref{thm:bar-cobar-adjunction}).

\emph{Step 1:} The bar construction is functorial:
\[\bar{B}(f): \bar{B}(\mathcal{A}) \to \bar{B}(\mathcal{A}')\]

\emph{Step 2:} The cobar construction is functorial:
\[\Omega(g): \Omega(\mathcal{C}) \to \Omega(\mathcal{C}')\]
for any coalgebra morphism $g$.

\emph{Step 3:} The natural transformation $\psi$ is defined universally via the 
adjunction, hence commutes with all morphisms.

\emph{Step 4:} At each genus, $\psi_g$ is natural in $\mathcal{A}$, so the 
genus-graded sum is also natural.
\end{proof}

\subsection{Applications to derived equivalences}

\begin{corollary}[Derived equivalence; \ClaimStatusProvedHere]\label{cor:derived-equivalence-bar-cobar}
\label{cor:derived-equivalence-koszul}
Let $\mathcal{A}$ be a Koszul chiral algebra on a smooth curve $X$ with Koszul dual
coalgebra $\mathcal{A}^!$.  Make the following assumptions:
\begin{enumerate}[label=(\roman*)]
\item $\mathcal{A}$ is augmented over $\mathcal{O}_X$ and locally finitely generated
as a $\mathcal{D}_X$-module.
\item ``Module'' means \emph{chiral module} in the sense of Beilinson--Drinfeld
\cite{BD04}: a $\mathcal{D}_X$-module $M$ with a chiral action map
$\mu_M: j_*j^*(\mathcal{A} \boxtimes M) \to \Delta_! M$.
\item ``Comodule'' means a conilpotent $\mathcal{A}^!$-comodule
$N$ with coaction $\Delta_N: N \to \mathcal{A}^! \chirtensor N$ satisfying
the usual coassociativity and counit axioms.
\item The derived categories $\mathcal{D}^b$ are the bounded derived categories
of complexes with finitely generated cohomology.
\end{enumerate}
Under these hypotheses, the bar and cobar constructions induce an equivalence
of triangulated categories:
\[\mathcal{D}^b(\text{Mod}^{\mathrm{fg}}(\mathcal{A})) \simeq
\mathcal{D}^b(\text{Comod}^{\mathrm{conil}}(\mathcal{A}^!))\]
via the functors $\bar{B}(-) = \mathcal{A}^! \otimes^{\tau}_{\mathcal{A}} (-)$
(bar extension of scalars) and $\Omega(-) = \mathcal{A} \otimes^{\tau}_{\mathcal{A}^!} (-)$
(cobar extension of scalars), where $\tau: \mathcal{A}^! \to \mathcal{A}$ is the
universal twisting cochain.
\end{corollary}

\begin{proof}
The quasi-isomorphisms $\psi: \Omega(\bar{B}(\mathcal{A})) \xrightarrow{\sim} \mathcal{A}$
and $\epsilon: \bar{B}(\Omega(\mathcal{A}^!)) \xrightarrow{\sim} \mathcal{A}^!$
establish:

\begin{center}
\begin{tikzcd}
\text{Mod}^{\mathrm{fg}}(\mathcal{A}) \ar[r, "\bar{B}", shift left=1ex]
& \text{Comod}^{\mathrm{conil}}(\mathcal{A}^!) \ar[l, "\Omega", shift left=1ex]
\end{tikzcd}
\end{center}

with $\Omega \circ \bar{B} \simeq \text{id}$ and $\bar{B} \circ \Omega \simeq \text{id}$
up to quasi-isomorphism.  The finiteness hypothesis~(i) ensures the bar construction
preserves bounded complexes; the conilpotence hypothesis~(iii) ensures the cobar
construction converges.  Together these give the stated equivalence on derived categories.
\end{proof}

\begin{remark}[Scope and limitations]\label{rem:derived-equiv-scope}
The hypotheses above are essential:
\begin{itemize}
\item Without finite generation, the bar construction may produce unbounded complexes.
\item Without conilpotence, the cobar series $\Omega(\mathcal{A}^!) = \bigoplus_{n \geq 1}
\mathcal{A}^{!\otimes n}$ may not converge, and the functor $\Omega$ is not well-defined.
\item For non-quadratic algebras, one must replace $\mathcal{A}^!$ by its nilpotent
completion (Appendix~\ref{app:nilpotent-completion}) and work with pro-nilpotent
comodules.
\end{itemize}
Even in ordinary algebra, Koszul duality yields equivalences between
\emph{specific subcategories}, not the entire derived category without
finiteness conditions; cf.~Beilinson--Ginzburg--Soergel \cite{BGS96}.
\end{remark}

\begin{remark}[Physical significance]\label{rem:qi-matters-physics}
From the physics perspective, the distinction between homology isomorphism and 
quasi-isomorphism corresponds to:

A homology isomorphism gives only on-shell equivalence: physical states match, but one cannot compute scattering amplitudes or access quantum corrections.  A full quasi-isomorphism gives off-shell equivalence: the entire QFT matches, correlation functions and amplitudes are computable, quantum corrections are encoded in the higher homotopies, and the path integral measure is determined.  This is why establishing the quasi-isomorphism---not just the homology isomorphism---is essential for physical applications.
\end{remark}

%================================================================
% RECOGNIZING KOSZUL DUALS IN PRACTICE
%================================================================

\section{Recognizing Koszul duals in practice}
\label{sec:recognizing-koszul-duals}

\begin{remark}[Identifying \texorpdfstring{$\mathcal{A}^!$}{A!} in practice]\label{rem:identify-koszul-wild}
When encountering a coalgebra $\widehat{\mathcal{C}}$ in geometry or physics, use 
the following checklist to determine if it is a Koszul dual:

\emph{Step 1: Check necessary conditions} (Theorem~\ref{thm:essential-image-koszul}):
\begin{itemize}
\item[$\square$] Conilpotent? ($\bigcap_n \text{coker}(\Delta^n) = 0$)
\item[$\square$] Connected? ($\epsilon: \widehat{\mathcal{C}} \twoheadrightarrow \mathbb{C}$)
\item[$\square$] Geometrically representable? (arises from configuration spaces)
\item[$\square$] Curvature central? (if curved)
\item[$\square$] Formally complete? (with respect to coaugmentation)
\end{itemize}

\emph{Step 2: Compute candidate algebra:}
\[\mathcal{A}_{\text{candidate}} = \Omega(\widehat{\mathcal{C}})\]

\emph{Step 3: Verify bar-cobar inversion:} Compute $\bar{B}(\mathcal{A}_{\text{candidate}})$ and check if $\bar{B}(\mathcal{A}_{\text{candidate}}) \simeq \widehat{\mathcal{C}}$.

\emph{Step 4: If yes:}
\[\widehat{\mathcal{C}} = \mathcal{A}_{\text{candidate}}^!\]

\emph{Examples where this works.}
\begin{itemize}
\item Heisenberg coalgebra $\to$ Heisenberg algebra
\item Exterior coalgebra $\to$ Free fermion $\beta\gamma$
\item Langlands dual Kac--Moody $\to$ Original Kac--Moody
\item Certain W-algebra coalgebras $\to$ W-algebras at special central charges
\end{itemize}

\emph{Examples where this fails:} non-conilpotent coalgebras (cannot be Koszul duals) and geometrically non-representable coalgebras (not from configuration spaces).
\end{remark}

\section{The Ran-space perspective}
\label{sec:ran-space-perspective}
\index{Ran space!bar-cobar perspective}

The $\infty$-categorical framework of factorization
algebras on $\operatorname{Ran}(X)$ provides the conceptual home for
the bar-cobar adjunction.

\begin{remark}[From configuration spaces to Ran space]\label{rem:config-to-ran}
The configuration spaces $C_n(X)$ and their FM compactifications
$\overline{C}_n(X)$ are finite-dimensional approximations to the
\emph{Ran space} $\operatorname{Ran}(X) = \operatorname{colim}_n\,
X^n / \Sigma_n$, the space of all finite non-empty subsets of~$X$.
The bar complex $\bar{B}(\cA) = \bigoplus_n \bar{B}^n(\cA)$, with
$\bar{B}^n(\cA)$ living on $\overline{C}_n(X)$, assembles into a
single object on $\operatorname{Ran}(X)$.  The bar differential,
given by residues at collision divisors
(Definition~\ref{def:bar-differential-complete}), becomes the
factorization structure map: when two points collide, the OPE
produces the ``fusion'' of the corresponding tensor factors.

In this language, the Arnold relations that ensure $d^2 = 0$
(Theorem~\ref{thm:bar-nilpotency-complete}) are the coherence conditions
for the factorization algebra axiom: the result of a three-point
collision does not depend on the order in which pairs collide,
because the relevant boundary strata of $\overline{C}_3(X)$ are
related by the Arnold relation
$\eta_{12} \wedge \eta_{23} + \eta_{23} \wedge \eta_{31}
+ \eta_{31} \wedge \eta_{12} = 0$.
\end{remark}

\begin{remark}[Verdier duality as the engine of Koszul duality]\label{rem:verdier-engine}
\index{Verdier duality!Koszul duality engine}
The bar-cobar adjunction
\[
\bar{B}_X \colon
\operatorname{Alg}_{\mathrm{aug}}(\operatorname{Fact}(X))
\;\rightleftarrows\;
\operatorname{CoAlg}_{\mathrm{conil}}(\operatorname{Fact}(X))
\,:\!\Omega_X
\]
is controlled by Verdier duality on $\operatorname{Ran}(X)$.
The key identification is:
\[
\mathbb{D}_{\operatorname{Ran}} \bar{B}_X(\cA) \;\simeq\; \bar{B}_X(\cA^!).
\]
At the chain level, this identity is the content of
Theorem~\ref{thm:bar-cobar-isomorphism-main}: the Verdier dual of the
bar construction of $\cA$ is the bar construction of the Koszul dual
$\cA^!$.  The FM propagator $\eta_{ij}$ and its Verdier dual
$\eta_{ij}^\vee$ implement this duality on configuration spaces:
the propagator mediates the passage from algebra to coalgebra, and
Verdier duality exchanges the two.

This perspective makes the Poincar\'e--Koszul duality of
Ayala--Francis~\cite{AF15} manifest: their equivalence
$\int_M A \simeq (\int_{-M} A^!)^\vee$ is the shadow, at the level
of factorization homology, of the chain-level Verdier identification
above.
\end{remark}

\begin{remark}[Curvature and coderived categories]\label{rem:curvature-coderived}
\index{coderived category!bar-cobar}
\index{curvature!coderived category}
When curvature is present---i.e., when the genus-$1$ obstruction
$m_0^{(1)} = \kappa(\cA) \neq 0$ (as it does for all non-trivial
Koszul chiral algebras; see
Theorem~\ref{thm:genus-universality})---the bar complex satisfies
$\dfib^{\,2} = \kappa \cdot \omega_1 \cdot \mathrm{id}$ rather than $\dfib^{\,2} = 0$
(see Convention~\ref{conv:higher-genus-differentials} for the notation).
The ordinary derived category of $\cA$-modules is then the wrong
quotient: it erases the distinction between ``acyclic because of
$\dfib^{\,2} \neq 0$'' and ``acyclic because of genuine exactness.''

The correct ambient category is the \emph{coderived category}
$\mathcal{D}^{\mathrm{co}}(\cA)$ in the sense of
Positselski~\cite{Positselski11}.  In this category, a curved complex
$(M, d_M)$ with $d_M^2 = \kappa \cdot \mathrm{id}_M$ is
\emph{not} set to zero; it represents a genuine object encoding the
curvature.  The bar-cobar adjunction lifts to an equivalence
\[
\bar{B}_X \colon
\mathcal{D}^{\mathrm{co}}(\operatorname{Mod}(\cA))
\;\simeq\;
\mathcal{D}^{\mathrm{ctr}}(\operatorname{Comod}(\cA^!))
\]
between the coderived category of modules and the contraderived category
of comodules.  On the Koszul locus---where the curvature vanishes or
is absorbed into the differential---the coderived and derived categories
agree, and the equivalence reduces to
Corollary~\ref{cor:derived-equivalence-koszul}.  Off this locus, the
coderived formalism captures the ``acyclic-but-not-contractible''
phenomena that higher-genus curvature creates.

This perspective clarifies the physical meaning of curvature: the
cosmological constant $\Lambda \sim \kappa(\cA)$ in the AdS$_3$/CFT$_2$
correspondence (Conjecture~\ref{conj:ads-cft-bar}) is the obstruction
to reducing the coderived category to the derived category.  Anomaly
cancellation at $c = 26$ ($\kappa = 0$) is precisely the condition
under which the coderived and derived categories agree, i.e., under
which the string theory is ``on-shell'' in the derived-categorical
sense.
\end{remark}

The bar-cobar adjunction for chiral algebras is now fully constructed:
the bar differential is a residue on Fulton--MacPherson
compactifications, $d^2 = 0$ follows from the Arnold relations, and
Verdier duality on $\operatorname{Ran}(X)$ intertwines bar and cobar
(Theorem~\ref{thm:verdier-bar-cobar}).  For Koszul pairs, the
adjunction is an equivalence (Theorem~\ref{thm:bar-cobar-isomorphism-main}).
This is the genus-$0$ theory.  Everything built so far lives on
configuration spaces of the fixed curve~$X$, with no reference to
moduli of curves.

\subsection{Bar-cobar as factorization homology}
\label{subsec:bar-cobar-fh}

The Ran-space remarks above assemble the bar complex into an
object on $\operatorname{Ran}(X)$.  We now show that this
assembly is not a reformulation but an \emph{identification}:
the bar construction \emph{is} factorization homology, and the
bar-cobar adjunction \emph{is} the adjunction of factorization
(co)homology functors.

\begin{proposition}[Bar construction as factorization homology;
\ClaimStatusProvedHere]
\label{prop:bar-fh}
\index{factorization homology!bar construction}
\index{bar construction!as factorization homology}
Let $\cA$ be an augmented $\Eone$-chiral algebra on a smooth
curve~$X$, and let $\overline{\cA} = \ker(\epsilon \colon \cA \to k)$
be the augmentation ideal.  Then there is a natural equivalence
\[
\Bbar^{\mathrm{geom}}_X(\cA)
\;\simeq\;
\int_X \overline{\cA}^{\mathrm{Lie}}
\]
where $\overline{\cA}^{\mathrm{Lie}}$ denotes $\overline{\cA}$ viewed as
an $\Eone$-Lie algebra via the commutator bracket, and the right-hand
side is factorization homology in the sense of
\textup{\cite{AF15, BD04}}.  More precisely:
\begin{enumerate}[label=(\roman*)]
\item \textbf{Finite approximations}: For each $n \geq 1$,
the degree-$n$ piece $\Bbar^n_X(\cA) =
(s^{-1}\overline{\cA})^{\otimes n}_{\overline{C}_n(X)}$
(the bar complex on~$n$ points, supported on the
Fulton--MacPherson compactification $\overline{C}_n(X)$) is the
restriction of the factorization homology $\int_X\overline{\cA}$
to $n$-point configurations.
\item \textbf{Assembly}: The bar differential
$d_{\mathrm{bar}} = \sum_{i < j} d_{ij}$ (residue along collision
divisors, Definition~\ref{def:bar-differential-complete}) is the
structure map of the factorization algebra:
the collision $z_i \to z_j$ implements the factorization product
$\int_{U_i}\overline{\cA} \otimes
\int_{U_j}\overline{\cA} \to
\int_{U_i \cup U_j}\overline{\cA}$
for shrinking discs $U_i \ni z_i$, $U_j \ni z_j$.
\item \textbf{Colimit}: The total bar complex
$\Bbar_X(\cA) = \bigoplus_{n \geq 0}\Bbar^n_X(\cA)$
is the colimit over $n$ of the finite approximations,
i.e., the factorization homology
$\int_X\overline{\cA} = \operatorname{colim}_{S \in
\mathrm{Fin}_{/X}} \bigotimes_{s \in S}\overline{\cA}_s$.
\end{enumerate}
\end{proposition}

\begin{proof}
The factorization homology of an $E_n$-algebra $A$ on a manifold
$M$ is defined as
\[
\textstyle\int_M A \;=\; \operatorname{colim}_{U_1 \sqcup \cdots
\sqcup U_k \hookrightarrow M} A(U_1) \otimes \cdots
\otimes A(U_k)\,,
\]
where the colimit is over rectilinear embeddings of disjoint discs
\cite[Definition~3.1]{AF15}.  For $n = 1$ on a curve~$X$, this is
the colimit over configurations of disjoint discs
$D_1 \sqcup \cdots \sqcup D_k \hookrightarrow X$, with
$A(D_i) = \overline{\cA}_{z_i}$ (the fiber of $\overline{\cA}$
at the center $z_i$ of $D_i$).

Part (i): The term $\Bbar^n_X(\cA)$ places one copy of
$s^{-1}\overline{\cA}$ at each of $n$ distinct points
$z_1, \ldots, z_n \in X$.  The locus of distinct points is the
open configuration space $C_n(X) = X^n \smallsetminus \Delta$,
and the extension to $\overline{C}_n(X)$ encodes the factorization
product at coincident points.

Part (ii): When $z_i \to z_j$, the factorization structure map
composes the two copies of $\overline{\cA}$ using the chiral
product.  At the chain level, this composition is the OPE residue
$\mathrm{Res}_{z_i = z_j}(\mu^{\mathrm{ch}}(\bar{a}_i, \bar{a}_j)
\cdot \eta_{ij})$, which is exactly the bar differential $d_{ij}$.

Part (iii): The colimit over $n$ in the bar complex is the colimit
over the category $\mathrm{Fin}_{/X}$ of finite subsets of $X$,
which is the Ran space $\operatorname{Ran}(X) =
\operatorname{colim}_n X^n/\Sigma_n$.  The factorization homology
$\int_X\overline{\cA}$ is this colimit by definition.
\end{proof}

\begin{proposition}[Cobar as factorization cohomology;
\ClaimStatusProvedHere]
\label{prop:cobar-fh}
\index{factorization homology!cobar construction}
\index{cobar construction!as factorization cohomology}
The cobar construction $\Omega_X(\cC)$ of a conilpotent
$\Eone$-chiral coalgebra $\cC$ is naturally equivalent to the
\emph{factorization cohomology}:
\[
\Omega_X(\cC) \;\simeq\;
\mathbb{R}\mathrm{Hom}_{\operatorname{Fact}(X)}
\bigl(k,\, \cC\bigr)
\;=:\;
\int^X \cC\,,
\]
the derived internal hom from the unit factorization algebra $k$
to~$\cC$ in the $\infty$-category of factorization coalgebras on~$X$.
Under Verdier duality $\mathbb{D}_{\mathrm{Ran}}$, this becomes:
\[
\Omega_X(\cC)
\;\simeq\;
\Bigl(\int_X \mathbb{D}_{\mathrm{Ran}}(\cC)\Bigr)^{\!\vee}\,,
\]
the linear dual of the factorization homology of the Verdier dual.
For $\cC = \Bbar_X(\cA)$, this gives
\[
\Omega_X\Bbar_X(\cA) \;\simeq\;
\Bigl(\int_X \Bbar_X(\cA^!)\Bigr)^{\!\vee}\,,
\]
which identifies the bar-cobar counit $\psi \colon
\Omega_X\Bbar_X(\cA) \to \cA$ with the Poincar\'e--Koszul duality
map of Ayala--Francis~\textup{\cite{AF15}}.
\end{proposition}

\begin{proof}
The cobar construction $\Omega_X(\cC) = T(s^{-1}\,\overline{\cC})$ with
differential $d_{\mathrm{cobar}} = d_{\mathrm{coop}} +
d_{\mathrm{tw}}$ is the free algebra on $s^{-1}\,\overline{\cC}$, with
the differential encoding the coalgebra structure of~$\cC$.
The free algebra on a module $V$ is
$\mathbb{R}\mathrm{Hom}(k, T^c(V))$, and for
$V = s^{-1}\,\overline{\cC}$ supported on $\operatorname{Ran}(X)$, this
is the factorization cohomology $\int^X\cC$ by
\cite[Proposition~5.5.2.5]{HA}.

The Verdier duality identification follows from
$\mathbb{D}_{\mathrm{Ran}}(\Bbar_X(\cA)) \simeq \Bbar_X(\cA^!)$
(Theorem~\ref{thm:verdier-bar-cobar}) and the general adjunction
$\mathrm{Hom}(k, \cC) \simeq (\int_X \mathbb{D}\cC)^\vee$.
The identification with Poincar\'e--Koszul duality follows from
\cite[Theorem~4.13]{AF15}: their equivalence
$\int_M A \simeq (\int_{-M} A^\vee)^\vee$ is precisely the
composition of Verdier duality on $\operatorname{Ran}(X)$ with
the bar-cobar adjunction.
\end{proof}

\begin{remark}[The bar-cobar adjunction as a duality of integration
theories]
\label{rem:bar-cobar-integration}
\index{factorization homology!duality of integration theories}
Propositions~\ref{prop:bar-fh} and~\ref{prop:cobar-fh} together
identify the bar-cobar adjunction
\[
\Bbar_X \colon
\operatorname{Alg}_{\mathrm{aug}}(\operatorname{Fact}(X))
\;\rightleftarrows\;
\operatorname{CoAlg}_{\mathrm{conil}}(\operatorname{Fact}(X))
\,:\!\Omega_X
\]
as the adjunction between two integration theories:
\begin{itemize}
\item The bar functor $\Bbar_X$ is \emph{factorization integration}:
it computes $\int_X(\cdot)$ by summing over configurations of
points on~$X$, with the OPE providing the integrand.
\item The cobar functor $\Omega_X$ is \emph{factorization
co-integration}: it computes $\int^X(\cdot)$ by taking the
internal hom from the unit, with the coproduct providing the
``inverse transform.''
\item Koszul duality $\cA \mapsto \cA^!$ is the
\emph{Verdier transform} mediating between the two:
$\Bbar(\cA)$ and $\Bbar(\cA^!)$ are Verdier dual, so the
bar of one algebra is the co-integration data for the dual.
\end{itemize}

This perspective reveals the bar-cobar machine as a
\emph{nonlinear Fourier duality} (cf.\
Remark~\ref{rem:koszul-fourier} and
Chapter~\ref{ch:en-koszul-duality}): the OPE kernel
$\eta_{ij} = d\log(z_i - z_j)$ plays the role of the Fourier
kernel $e^{2\pi i x\xi}$, the bar construction is the ``forward
transform'' (integration against the kernel), and the cobar is the
``inverse transform.''  Koszulness is the condition that the
transform is invertible (Fourier inversion), and the bar-cobar
counit $\Omega\Bbar(\cA) \to \cA$ is the Fourier inversion formula.
\end{remark}

\begin{remark}[The master diagram at genus~\texorpdfstring{$0$}{0}]
\label{rem:master-diagram-genus0}
\index{master diagram!genus zero}
The results of this chapter instantiate the master
diagram~\eqref{eq:master-square} at genus~$0$:
\[
\begin{tikzcd}[row sep=large, column sep=large]
\cA \arrow[r, "\barB_X"] \arrow[d, "(\cdot)^!"']
& \barB_X(\cA) \arrow[d, "\mathbb{D}_{\mathrm{Ran}}"] \\
\cA^! \arrow[r, "\barB_X"']
& \barB_X(\cA^!)
\end{tikzcd}
\]
The top row is the bar construction of this chapter; the right
column is Verdier duality
(Theorem~\ref{thm:verdier-bar-cobar}); the commutativity
$\mathbb{D}_{\mathrm{Ran}}\barB_X(\cA) \simeq \barB_X(\cA^!)$
is Theorem~A.  The cobar functor $\Omega_X$ is left adjoint
to~$\barB_X$, and the counit
$\Omega_X\barB_X(\cA) \xrightarrow{\sim} \cA$ is an
equivalence for Koszul algebras.

In the nilpotence-periodicity framework
(Remark~\ref{rem:nilpotence-periodicity}), this is the
\emph{single-valued regime} of the categorical logarithm:
$d^2 = 0$ exactly (the Arnold relations give single-valuedness),
so the logarithm has no monodromy and the diagram commutes
strictly.  At genus~$g \geq 1$
(Chapter~\ref{chap:higher-genus}), the logarithm acquires
monodromy---the square acquires
curvature~$\kappa(\cA) \cdot \omega_g$---and strict commutativity
is replaced by commutativity in the coderived category:
the passage from genus~$0$ to genus~$g$ is the passage from
single-valued to multi-valued logarithm, and the coderived
category is the analytic continuation that makes the multi-valued
logarithm well-defined.

The factorization homology perspective
(Propositions~\ref{prop:bar-fh},~\ref{prop:cobar-fh}) upgrades the
diagram from a commutative square of chain complexes to a
commutative square of \emph{factorization algebra objects} on
$\operatorname{Ran}(X)$.  The Verdier duality
$\mathbb{D}_{\mathrm{Ran}}$ becomes an $\infty$-categorical
involution, and the commutativity is an equivalence of factorization
algebras, not merely a quasi-isomorphism.  This
$\infty$-categorical enhancement is the H-level datum that
organizes the M-level constructions.
\end{remark}

\section{The categorical logarithm}
\label{sec:categorical-logarithm}
\index{categorical logarithm|textbf}
\index{nilpotence-periodicity correspondence!categorical logarithm}

The constructions of this chapter---bar, cobar, Verdier intertwining,
bar-cobar inversion, factorization homology---are not independent
results.  They are aspects of a single structure: the
\emph{categorical logarithm} for factorization algebras on algebraic
curves.  In this section we develop this identification
systematically, extracting from it both the conceptual unity of the
four main theorems and concrete predictions for higher-genus theory.

\subsection{The exponential-logarithm correspondence}
\label{subsec:exp-log-correspondence}

\begin{definition}[Categorical logarithm and exponential]
\label{def:categorical-log-exp}
\index{categorical logarithm!definition}
\index{categorical exponential!definition}
Let $\cA$ be an augmented $\Eone$-chiral algebra on a smooth
curve~$X$.  The \emph{categorical logarithm} of~$\cA$ is the bar
construction
\[
\log^{\mathrm{cat}}(\cA) \;:=\; \barB_X(\cA)\,,
\]
viewed as a factorization coalgebra on $\operatorname{Ran}(X)$
(Proposition~\ref{prop:bar-fh}).  The \emph{categorical exponential}
is the cobar construction
\[
\exp^{\mathrm{cat}}(\cC) \;:=\; \Omega_X(\cC)\,,
\]
viewed as a factorization algebra on $\operatorname{Ran}(X)$
(Proposition~\ref{prop:cobar-fh}).
\end{definition}

This terminology is justified by the following dictionary, each entry
of which is a theorem proved in this chapter or in
Chapter~\ref{chap:higher-genus}:

\begin{center}
\small
\renewcommand{\arraystretch}{1.3}
\begin{tabular}{p{5.5cm}p{5.5cm}}
\textbf{Classical logarithm} & \textbf{Categorical logarithm} \\
\hline
$\log \colon (\mathbb{C}^\times, \cdot) \to (\CC, +)$ &
$\barB \colon \mathrm{Alg}_{\mathrm{aug}} \to \mathrm{CoAlg}_{\mathrm{conil}}$ \\[2pt]
Maps multiplicative to additive &
Maps OPE (multiplicative) to bar differential (nilpotent/additive) \\[2pt]
$\exp(\log(z)) = z$ on convergence disk &
$\Omega(\barB(\cA)) \xrightarrow{\sim} \cA$ on Koszul locus
(Theorem~\ref{thm:bar-cobar-inversion-qi}) \\[2pt]
Power series $\log(1+z) = z - z^2/2 + \cdots$ &
PBW spectral sequence:
$E_1^{p,q} = H^q(\mathrm{gr}^p \barB(\cA))$ \\[2pt]
Convergence radius &
Koszul locus: $E_2$~degeneration
(Theorem~\ref{thm:spectral-sequence-collapse}) \\[2pt]
Analytic continuation past radius &
Coderived category completion
(\S\ref{subsec:chiral-coderived-contraderived}) \\[2pt]
$\log(ab) = \log(a) + \log(b)$ &
$Q_g(\cA) + Q_g(\cA^!) = H^*(\overline{\mathcal{M}}_g, Z_\cA)$
(Theorem~\ref{thm:quantum-complementarity-main}) \\[2pt]
Single-valued on simply connected domain &
$\dzero^2 = 0$ at genus~$0$
(Theorem~\ref{thm:bar-nilpotency-complete}) \\[2pt]
Multi-valued on $\mathbb{C}^\times$; monodromy $2\pi i$ &
$\dfib^{\,2} = \kappa(\cA) \cdot \omega_g$ at genus~$g \geq 1$ \\[2pt]
Residue at branch point &
Curvature $\kappa(\cA)$ \\[2pt]
Leading coefficient of $\log(1+z) = z + \cdots$ &
Modular characteristic $\kappa(\cA)$
(Theorem~\ref{thm:modular-characteristic})
\end{tabular}
\end{center}

\begin{remark}[The logarithmic kernel]
\label{rem:logarithmic-kernel}
\index{logarithmic kernel}
The fact that the bar construction uses a \emph{logarithmic} kernel
$\eta_{ij} = d\log(z_i - z_j)$ is the surface manifestation of the
categorical identification.  The classical logarithm of a product is a
sum: $\log(z_1 z_2) = \log z_1 + \log z_2$.  The bar differential
implements precisely this: the OPE (a product structure at colliding
points) is converted, via the residue $\mathrm{Res}_{z_i = z_j}
(\mu^{\mathrm{ch}} \cdot \eta_{ij})$, into a sum (the bar
differential is a sum over pairwise collisions).  The kernel $\eta$ is
the infinitesimal version of this conversion.
\end{remark}

\subsection{The Taylor expansion of the logarithm}
\label{subsec:taylor-expansion-log}
\index{PBW spectral sequence!as Taylor expansion}

The PBW spectral sequence of a chiral algebra~$\cA$
(Definition~\ref{def:chiral-koszul-pair}) is the Taylor expansion of
the categorical logarithm.

More precisely: the weight filtration on the bar complex
$\barB(\cA) = \bigoplus_{n \geq 0} (s^{-1}\overline{\cA})^{\otimes n}$
induces a spectral sequence
\[
E_1^{p,q} \;=\; H^q\bigl(\mathrm{gr}^p\, \barB(\cA)\bigr)
\;\Longrightarrow\;
H^{p+q}(\barB(\cA))\,.
\]
In the exp/log framework, the filtration degree~$p$ is the
\emph{order of the Taylor expansion}: the $p$-th filtered piece
$F^p\barB(\cA)$ retains contributions from at most $p$~points in
configuration space, i.e., the first $p$~terms of the logarithmic
series.  The differential $d_1$ on the $E_1$-page extracts the
\emph{leading term} of the logarithm (the binary OPE residue); higher
differentials $d_r$ capture \emph{higher-order corrections} from
$r$-fold collisions.

\begin{proposition}[Convergence criterion; \ClaimStatusProvedElsewhere{}
Theorems~\ref{thm:spectral-sequence-collapse}
and~\ref{thm:bar-cobar-inversion-qi}]
\label{prop:log-convergence}
\index{convergence!of categorical logarithm}
The categorical logarithm \emph{converges}---meaning $\exp \circ \log
= \mathrm{id}$---if and only if the PBW spectral sequence degenerates
at~$E_2$.  The $E_2$~degeneration is equivalent to:
\begin{enumerate}[label=(\roman*)]
\item The chiral algebra~$\cA$ lies on the \emph{Koszul locus}
(Definition~\ref{def:chiral-koszul-pair});
\item The Taylor expansion has finitely many nonzero corrections at
each filtration degree: the logarithm is ``polynomial'' in the
filtration variable;
\item The counit $\psi \colon \Omega\barB(\cA) \to \cA$ is a
quasi-isomorphism
(Theorem~\ref{thm:bar-cobar-inversion-qi}).
\end{enumerate}
Off the Koszul locus, the Taylor series diverges in the ordinary
derived category.  The coderived category
(\S\ref{subsec:chiral-coderived-contraderived}) provides the analytic
continuation: it is the completion of the derived category that makes
the divergent logarithm well-defined, exactly as the maximal analytic
continuation of $\log z$ requires passage from $\mathbb{C}^\times$ to its
universal cover.
\end{proposition}

\begin{remark}[Convergence radius from OPE data]
\label{rem:convergence-radius-ope}
\index{convergence!radius from OPE poles}
\index{Koszul locus!characterization via OPE}
The analogy with power series provides a heuristic for the following
conjectural expectation: the convergence radius of the categorical
logarithm should be computable from the input data.
For the classical logarithm $\log(1+z) = \sum (-1)^{n+1} z^n/n$, the
radius is determined by the growth rate of the coefficients.  For the
categorical logarithm, the ``coefficients'' are the OPE residues, and
the ``growth rate'' is the complexity of multi-point collisions.

In concrete terms: Kac--Moody algebras have first-order OPE poles
(linear growth), and the categorical logarithm converges
(Theorem~\ref{thm:master-pbw}).  The Virasoro algebra has a
second-order pole from the central term, but the Sugawara constraint
tames the growth, and convergence holds
(Theorem~\ref{thm:virasoro-chiral-koszul}).  For non-Koszul algebras
(e.g., admissible-level affine algebras), the OPE singularity structure
generates higher differentials in the spectral sequence---the Taylor
coefficients grow too fast for convergence.  A precise bound relating
the OPE pole orders to the convergence radius of the categorical
logarithm is a natural target for future work.
\end{remark}

\subsection{Additivity and the branch structure}
\label{subsec:additivity-branch}
\index{categorical logarithm!additivity}
\index{complementarity!as additivity of logarithm}

The fundamental property of any logarithm is that it converts products
to sums.  For the categorical logarithm, this is the content of the
complementarity theorem.

\begin{proposition}[Additivity of the categorical logarithm;
\ClaimStatusProvedElsewhere{} Theorem~\ref{thm:quantum-complementarity-main}]
\label{prop:log-additivity}
Let $(\cA, \cA^!)$ be a Koszul pair.  The quantum corrections to the
bar complex satisfy
\[
Q_g(\cA) \;+\; Q_g(\cA^!)
\;=\; H^*(\overline{\mathcal{M}}_g,\, Z_\cA)\,.
\]
In the exp/log framework, this is the statement
$\log^{\mathrm{cat}}(\cA \otimes \cA^!) =
\log^{\mathrm{cat}}(\cA) + \log^{\mathrm{cat}}(\cA^!)$: the
logarithm of the Koszul product is the sum of the logarithms, and
the right-hand side is the total logarithmic data
$\log^{\mathrm{cat}}(\cA \otimes \cA^!)$ computed from the center~$Z_\cA$
via the Hodge cohomology of~$\overline{\mathcal{M}}_g$.

The Lagrangian decomposition
$C_g(\cA) = Q_g(\cA) \oplus Q_g(\cA^!)$ is the decomposition of the
total logarithm into its two \emph{branches}: the sheet associated
to~$\cA$ and the sheet associated to its Koszul dual.  The pairing
between them is the monodromy acquired when analytically continuing
from one branch to the other.
\end{proposition}

\begin{remark}[The categorical Baker--Campbell--Hausdorff formula]
\label{rem:categorical-bch}
\index{Baker--Campbell--Hausdorff formula!categorical analogue}
For the classical logarithm, $\log(e^X e^Y) = X + Y + \tfrac{1}{2}[X,Y] + \cdots$
(the BCH formula).  The categorical analogue is the twisted tensor
product decomposition: for an extension
$0 \to \cA' \to \cA \to \cA'' \to 0$ of chiral algebras, the bar
complex decomposes as
$\barB(\cA) \simeq \barB(\cA'') \otimes_\tau \barB(\cA')$,
where the twisting cochain~$\tau$ encodes the extension data through
iterated brackets, in direct analogy with the commutator corrections
of BCH\@.  This follows from the standard filtered bar construction
applied to the extension filtration, and the twisting cochain~$\tau$
can be computed order by order from the connecting maps of the
extension.  A systematic development of the categorical BCH, including
explicit formulas for the first correction terms and their geometric
content, is deferred to \S\ref{sec:maurer-cartan-curved}.
\end{remark}

\subsection{From single-valued to multi-valued: the genus transition}
\label{subsec:genus-transition-log}
\index{categorical logarithm!genus transition}
\index{monodromy!of categorical logarithm}

The deepest structural consequence of the exp/log identification is the
genus transition.

\begin{center}
\small
\renewcommand{\arraystretch}{1.3}
\begin{tabular}{p{5.5cm}p{5.5cm}}
\textbf{Classical complex analysis} & \textbf{Categorical logarithm} \\
\hline
Simply connected domain &
Genus~$0$: $\pi_1(C_2(\mathbb{P}^1)) = 0$ \\[2pt]
$\log$ is single-valued, holomorphic &
$\dzero^2 = 0$ (Arnold relations), no curvature \\[2pt]
Multiply connected domain ($\mathbb{C}^\times$) &
Genus~$g \geq 1$: $\pi_1(\Sigma_g) \neq 0$ \\[2pt]
$\log$ acquires monodromy $2\pi i n$ &
$\dfib^{\,2} = \kappa(\cA) \cdot \omega_g$ \\[2pt]
Residue at branch point: $\mathrm{Res}_{z=0}(dz/z) = 1$ &
Curvature: $\kappa(\cA) \in H^1(X, \cA) \otimes H^0(\overline{\mathcal{M}}_g, \omega)$ \\[2pt]
Total differential $d\log z$ is exact on universal cover &
Total differential $\Dg{g}^2 = 0$ after Gauss--Manin correction \\[2pt]
Period matrix $\tau_{ij} = \oint_{B_j} \omega_i$ &
Modular package $(\kappa, \Delta, \Theta, \Pi)$ \\[2pt]
Torelli: $\tau$ determines the curve &
Categorical Torelli: monodromy of $\barB_g$ determines~$\cA$
up to quasi-isomorphism (via Theorem~\ref{thm:bar-cobar-inversion-qi})
\end{tabular}
\end{center}

\begin{remark}[The chiral Riemann--Hilbert correspondence, revisited]
\label{rem:chiral-rh-revisited}
\index{Riemann--Hilbert correspondence!chiral}
The Riemann--Hilbert correspondence identifies flat connections
(nilpotent data: $\nabla^2 = 0$) with monodromy representations
(periodic data: $\rho(\gamma^N) = 1$ for finite-order monodromy).
The categorical logarithm instantiates this at the level of
factorization algebras:
\begin{enumerate}[label=(\roman*)]
\item The \emph{flat connection} is the total corrected differential
$\Dg{g}$ with $\Dg{g}^2 = 0$, a flat connection on the bar family
$\barB_g(\cA) \to \overline{\mathcal{M}}_g$.
\item The \emph{curvature} of the fiberwise differential
$\dfib^{\,2} = \kappa \cdot \omega_g$ is the infinitesimal generator
of the monodromy.
\item The \emph{monodromy representation} is the action of the mapping
class group $\Gamma_g$ on $H^*(\barB_g(\cA))$, whose leading
eigenvalues are $\exp(2\pi i \kappa(\cA))$.
\end{enumerate}
The three differentials $\dzero$, $\dfib$, $\Dg{g}$
(Convention~\ref{conv:higher-genus-differentials}) are three levels of
the logarithm: the germ ($\dzero$, single-valued), the
multi-valued lift ($\dfib$, curvature $= $ residue), and the
trivialization on the universal cover ($\Dg{g}$, flat).
\end{remark}

\subsection{The characteristic hierarchy as logarithmic data}
\label{subsec:characteristic-log}
\index{modular characteristic!as logarithmic data}

The modular package $(\kappa, \Delta, \Theta, \Pi)$ of a chiral
algebra is the data of the categorical logarithm, organized by order:

\begin{enumerate}[label=\textbf{Level \arabic*}:, leftmargin=3.5em]
\item \textbf{Scalar logarithmic data.}  The modular
characteristic $\kappa(\cA)$ is the leading coefficient of the
categorical logarithm---the residue of the curvature.  Its
exponential is the scalar period~$N$: the smallest integer for which
$\exp(2\pi i N c/24) = 1$.  This is the content of Theorem~D
(Theorem~\ref{thm:modular-characteristic}).

\item \textbf{Spectral logarithmic data.}  The spectral
discriminant~$\Delta(\cA)$ encodes the eigenvalues of the monodromy of
the categorical logarithm.  Its exponential is the periodicity
profile~$\Pi(\cA)$: the collection of periods of individual T-matrix
eigenvalues $\lambda_i = \exp(2\pi i(h_i - c/24))$.

\item \textbf{Full logarithmic data.}  The universal Maurer--Cartan
class $\Theta_\cA \in H^2(\overline{\mathcal{M}}_g, Z_\cA)$
(Theorem~\ref{thm:mc2-full-resolution}) is the complete categorical
logarithm: it encodes all genus-$g$ corrections to the bar
differential.  Its exponential is the full modular automorphy of~$\cA$:
the action of $\Gamma_g$ on the conformal blocks, i.e., the modular
functor.
\end{enumerate}

\noindent
The hierarchy is controlled by a single principle: \emph{each level is
the leading term of the next}.  The scalar characteristic~$\kappa$ is
the trace of the spectral data~$\Delta$; the spectral data is the
eigenvalue decomposition of the full data~$\Theta$.  In the language of
complex analysis: $\kappa$ is the winding number of the logarithm,
$\Delta$ is the set of periods, and $\Theta$ is the full multi-valued
function.

\subsection{Visible theorems}
\label{subsec:visible-theorems-log}
\index{categorical logarithm!predicted theorems}

The exp/log identification makes certain results visible that would
otherwise have no natural formulation:

\begin{enumerate}[label=(\arabic*)]
\item \textbf{Categorical Torelli.}  If the monodromy of the
categorical logarithm (the higher-genus corrections to the bar
differential, encoded in~$\Theta_\cA$) determines~$\cA$ up to
quasi-isomorphism, this is a \emph{categorical Torelli theorem}.  It
follows from bar-cobar inversion (Theorem~\ref{thm:bar-cobar-inversion-qi})
plus the fact that the curvature determines the flat connection up to
gauge equivalence on a simply connected base: since the Teichm\"uller
space is contractible, the flat bar family is determined by its
curvature, i.e., by~$\kappa(\cA)$.

\item \textbf{Functional equation.}  The complementarity theorem
(Theorem~\ref{thm:quantum-complementarity-main}) is the \emph{functional equation} of the categorical
logarithm: it relates $\log^{\mathrm{cat}}(\cA)$ and
$\log^{\mathrm{cat}}(\cA^!)$ through the total data
$H^*(\overline{\mathcal{M}}_g, Z_\cA)$, exactly as the functional
equation of the Riemann $\zeta$-function relates $\zeta(s)$ and
$\zeta(1-s)$ through the $\Gamma$-factor.

\item \textbf{Acyclicity of the trivial algebra.}  For the unit
(trivial) chiral algebra $\cA = k$, the bar complex $\barB(k)$ is
acyclic: $H^*(\barB(k)) = k$ concentrated in degree~$0$.  In the
exp/log framework this is $\log(1) = 0$: the logarithm of the
multiplicative identity is the additive identity.

\item \textbf{Convergence radius from OPE poles.}  The Koszul locus
should be characterized by the singularity structure of the OPE: the
convergence radius of the categorical logarithm equals the inverse of
the growth rate of multi-point OPE singularities, in analogy with the
Cauchy--Hadamard formula.

\item \textbf{Screening as infinitesimal generator.}  In the lattice
setting (Construction~\ref{constr:lattice:screening}), the screening charges $Q_i$
with $Q_i^2 = 0$ are the infinitesimal generators of the categorical
logarithm; the quantum group relations $e_i^N = 0$ at $\zeta =
e^{2\pi i/N}$ are their $N$-th powers; and the representation category
$\mathrm{Rep}(\mathfrak{u}_\zeta(\mathfrak{g}))$ carries the periodic
(exponential) structure.  The passage $Q^2 = 0 \to e^N = 0$ is
literally exponentiation: the root of unity
$\zeta = \exp(2\pi i/N)$ is the monodromy of the categorical
logarithm around the $N$-th root.
\end{enumerate}

\begin{conjecture}[Factorization finiteness criterion;
\ClaimStatusConjectured]\label{conj:factorization-finiteness-criterion}
\index{factorization bar-cobar!finiteness criterion}
\index{thick generation!bypassed by sectorwise finiteness}
Let $\cA$ be a chiral algebra whose bar complex
$\barB(\cA)$ admits a bigrading
$\barB^{p,q}(\cA)$ with $\dim \barB^{p,q} < \infty$ for all
$(p,q)$.  Suppose the associated spectral sequence
$E_r^{p,q} \Rightarrow H^{p+q}(\barB(\cA))$ satisfies:
\begin{enumerate}[label=\textup{(\roman*)}]
\item \textup{(Convergence)} $E_r^{p,q} = E_\infty^{p,q}$
  for each $(p,q)$ at some finite page~$r_0(p,q)$;
\item \textup{(Sub-exponential growth)} $\dim E_1^{0,p} < c^p$
  for every $c > 1$ and all sufficiently large~$p$.
\end{enumerate}
Then the factorization bar-cobar adjunction restricts to an
equivalence on the Koszul locus, and the factorization DK
equivalence holds for~$\cA$ without thick generation.

\emph{Instances.}
For lattice algebras, (i)--(ii) are satisfied by the
$\Lambda$-grading with $E_1$~degeneration
\textup{(}Theorem~\textup{\ref{thm:lattice:factorization-koszul})}.
For Yangians, the root weight/loop degree bigrading of
Proposition~\textup{\ref{prop:yangian-bar-loop-weight}} gives (i),
and (ii) is proved in
Proposition~\textup{\ref{prop:lqt-e1-subexponential-growth}}: the
growth rate is $\sim\exp(\pi\sqrt{rp/12})$, which is sub-exponential.
For $\mathcal{W}$-algebras, the DS filtration provides a candidate
bigrading whose (ii) would follow from Kac--Kazhdan character
estimates.
\end{conjecture}

\subsection{The monograph as construction of a categorical logarithm}
\label{subsec:monograph-as-log}
\index{categorical logarithm!monograph structure}

Taken at full depth, the four main theorems are the four fundamental
properties of \emph{any} logarithm:
\begin{itemize}
\item \textbf{Theorem~A} (existence): the categorical logarithm and
exponential exist, as a pair of adjoint functors between augmented
factorization algebras and conilpotent factorization coalgebras.
\item \textbf{Theorem~B} (invertibility): on the convergence domain
(the Koszul locus), $\exp \circ \log = \mathrm{id}$, and the
logarithm faithfully encodes its input.
\item \textbf{Theorem~C} (branch structure): the failure of
convergence off the Koszul locus decomposes into complementary
Lagrangians---the two branches of the multi-valued logarithm---and the
total logarithmic data is the sum of the two branches.
\item \textbf{Theorem~D} (leading coefficient): the leading
coefficient $\kappa(\cA)$ of the logarithmic expansion determines the
full scalar package, just as the leading coefficient of
$\log(1+z) = z + \cdots$ determines the local behavior.
\end{itemize}

Every open problem in the monograph is a question about this
logarithm.  The five master conjectures are five aspects of one
question: \emph{what is the categorical logarithm for factorization
algebras on algebraic curves, completely?}
\begin{itemize}
\item \textbf{MC1}: does the logarithm converge? (PBW concentration)
\item \textbf{MC3}: what is the convergence domain? (thick generation
in category~$\mathcal{O}$)
\item \textbf{MC4}: does the logarithm extend to infinite rank?
(stability of towers)
\item \textbf{MC5}: does the physical logarithm (BRST) equal the
mathematical one (bar)? (BV--BRST identification)
\item \textbf{Periodicity}: what is the monodromy group of the
categorical logarithm?
\end{itemize}
These are not five separate programmes.  They are five windows onto a
single object.

\subsection{From tree-level to modular homotopy theory}
\label{subsec:tree-to-modular}
\index{modular homotopy theory!three-level ladder}

The constructions of this chapter are \emph{tree-level}: the bar
differential is indexed by collisions on a single smooth curve, and
the operadic compositions are indexed by planar rooted trees.  Three
levels of homotopy-algebraic structure emerge as one opens the genus
filtration.

\begin{definition}[Tree-level operadic homotopy theory]
\label{def:tree-level-homotopy}
A \emph{tree-level operadic homotopy theory} for a cyclic operad~$P$
is an $\infty$-category~$\mathcal{H}$ with homotopy-coherent
$P$-action: operations $m_T$ indexed by $P$-decorated rooted
trees~$T$, satisfying the operadic master equation $\sum_T m_T = 0$
up to coherent homotopies.  The bar-cobar constructions of this
chapter produce such a theory for $P = \chirCom$ or $P = \chirAss$.
\end{definition}

\begin{definition}[Modular homotopy theory, abstract]
\label{def:modular-homotopy-abstract}
A \emph{modular homotopy theory} for a modular operad~$M$ is an
$\infty$-category~$\mathcal{H}$ with homotopy-coherent $M$-action:
operations $m_\Gamma$ indexed by $M$-decorated \emph{stable
graphs}~$\Gamma$ (graphs with genus labels $g_v$ at each vertex,
satisfying $\sum_v g_v + b_1(\Gamma) = g$), compatible with
separating and nonseparating clutching
(Definition~\ref{def:modular-homotopy-theory-intro}).
\end{definition}

\begin{remark}[The formal step is small; the existence problem is hard]
\label{rem:trees-to-graphs}
Replacing rooted trees by stable graphs is the
\emph{only} combinatorial change from tree-level to modular homotopy
theory.  The formal generalization is immediate.  What is not
immediate is \emph{existence}: constructing the operations $m_\Gamma$
for graphs with loops requires integration over
$\overline{\mathcal{M}}_{g,n}$, which produces curvature
($\dfib^{\,2} \neq 0$) and demands the full Maurer--Cartan package
$\Theta_\cA$.  This is the content of
Chapter~\ref{chap:higher-genus}.
\end{remark}

\medskip

In Chapter~\ref{chap:higher-genus}, the curve itself varies over
$\overline{\mathcal{M}}_g$, and the bar differential acquires
curvature: the genus-$g$ correction $m_0^{(g)}$ is a section of the
obstruction sheaf on $\overline{\mathcal{M}}_g$, controlled by a
single universal Maurer--Cartan class $\Theta_{\cA}$.  The passage
from this chapter to the next is the passage from operadic to
\emph{modular} Koszul duality --- from the combinatorics of
associahedra to the geometry of the Deligne--Mumford
compactification --- or, in the language of the categorical logarithm,
the passage from the single-valued to the multi-valued regime.
